\chapter{Evaluation}\label{chap:eval}
This chapter will evaluate CTaps against three benchmarks:
native UDP, native TCP and a Picoquic client. Native UDP/TCP uses OS-level sockets,
while the Picoquic client uses the Picoquic packet loop as described in \cref{sec:picoquic}.
Five experiments are used to compare CTaps with these benchmarks:
handshake time, UDP RTT, small file transfer, candidate racing and
connection migration.

In addition, we will compare a CTaps client with a TCP and a Picoquic
client in terms of lines of code, in order to
show the relative simplicity of a CTaps client given the
features it provides.

\section{Setup}
This section will describe the hardware, software and methodology used
when performing the experiments.

\subsection{Test Machine}
All experiments were performed on a single laptop which is described in~\cref{tab:machine-specs}.
Performance would not change significantly 
on a more performant system, as the network emulation described
in \cref{sec:softwaretools} makes all experiments, except UDP RTT, I/O-bound.

\begin{table}[htbp]
\centering
\begin{threeparttable}
\caption[Test machine specifications]{Specifications of the test machine used to run experiments.}
\label{tab:machine-specs}
\small
\begin{tabular}{ll}
\toprule
\rowcolor{gray!25}
\textbf{Component} & \textbf{Details} \\
\midrule
\multicolumn{1}{l||}{Model}      & ThinkPad X1 Carbon 6th Gen (2018) \\
\multicolumn{1}{l||}{CPU}        & Intel Core i5-8350U CPU @ 1.70GHz (4 cores, 8 threads) \\
\multicolumn{1}{l||}{RAM}        & \qty{16}{\giga\byte} LPDDR3-2133 \\
\multicolumn{1}{l||}{OS}         & Debian GNU/Linux 12 \\
\multicolumn{1}{l||}{Kernel}     & 6.1.0-44-amd64 \\
\bottomrule
\end{tabular}
\end{threeparttable}
\end{table}


\subsection{Network Emulation} \label{sec:softwaretools}
All five experiments were performed on localhost.\footnote{
The experiment servers and clients used can be found in the \mono{benchmark/src/server/}
and \mono{benchmark/src/client/}
folders in the CTaps repository.}
For all experiments except for UDP RTT, bandwidth emulation
and delay were added using the Linux utility
\mono{tc-netem}~\cite{tc-netem}. Jitter was added
to the candidate racing experiment, and is
described in \cref{sec:handshake-kde}.

The connection migration experiment used \mono{iptables}~\cite{iptables_man} to block the initial
path used and force migration. This block was performed based on
source IP, destination IP, protocol and destination port. The specific command
used can be seen in~\cref{lst:iptables-command}.
An equivalent command blocking
TCP data was also run, in order to illustrate the TCP connection failing after a path event.

\begin{listing}[htbp]
    \begin{minted}[fontsize=\small]{bash}
iptables -A OUTPUT -s 127.0.0.1 -d 127.0.0.1 -p udp --dport 8080 -j DROP 
\end{minted}
\caption[Command to block path using iptables]{
Command to block UDP data from Local Endpoint to file server on localhost using iptables.
}
\label{lst:iptables-command}
\end{listing}

This command adds a rule to the \mono{OUTPUT} chain of the \mono{filter} table via
\mono{iptables}~\cite{iptables_man}. When the outgoing packet is dropped, the Linux kernel
function \mono{nf_hook_slow()} returns \mono{-EPERM}~\cite{LinuxKernelNfHookSlow}.
This propagates back up through the networking stack and
\mono{sendmsg()} returns \mono{-1} with \mono{errno} set to
\mono{EPERM}.\footnote{
This simulates a synchronous failure. For asynchronous
failures, where the packet is lost somewhere along its path, \mono{sendmsg()}
would not return an error. The implications of this are discussed in
\cref{sec:connection-mig-benchmark-discssion}.}

For CTaps, this works well as it could
report any \mono{sendmsg()} failure directly to Picoquic.
Picoquic would then communicate back to CTaps that the path has become unavailable through
a callback and CTaps could start probing new paths.

For the Picoquic client, a small patch had to be applied in order to achieve
the same migration behavior. The Picoquic packet loop only
notifies the QUIC engine of a path failure if \mono{sendmsg()} fails with \mono{errno} set to a value implying 
that the remote is unreachable. In the Picoquic packet loop this means one of the following errors: \mono{EAFNOSUPPORT},
\mono{ECONNRESET}, \mono{EHOSTUNREACH}, \mono{ENETDOWN} or \mono{ENETUNREACH}
\cite{PicoquicImpliesUnreachable}.

For the migration experiment it was desirable for the Picoquic client
to also treat \mono{EPERM} as implying that the remote is unreachable. Picoquic was therefore forked
and \mono{EPERM} was added to the list of socket errors which imply that
the remote is
unreachable~\cite{PicobrigtImpliesUnreachable}. The forked version of Picoquic was only used
for the Picoquic experiment client. CTaps itself depends on the original version.

\subsection{Measurement Approach}
The UDP RTT experiment was performed through a bash script, while the remaining
experiments were performed through a dedicated Python script.
This Python script was invoked with a specific set of command-line arguments for each
experiment. Invocations of both the bash script and the Python script
were stored in dedicated per-experiment shell
scripts for reproducibility.\footnote{The associated shell scripts
can be found in the \mono{benchmark/scripts/} folder of the main CTaps repository,
together with a README on how to reproduce the experiments.}

The experiments were run sequentially while the test
machine was otherwise idle.
For I/O-dominated experiments without jitter, the
variation in performance was found to be so small
that increased sampling would not be visible in the results.
The handshake time, small file transfer and connection
migration experiments are therefore based on a single run
and serve as functional demonstrations rather than statistical performance characterizations.

\section{Experiments}
\subsection{Handshake Time Experiment}\label{sec:single-endpoint-handshake}
This experiment can be seen in~\cref{fig:racing-handshake}.
It demonstrates the handshake times of a CTaps client and various benchmark
clients when connecting to TCP and QUIC servers. It iterates
over delays of 25, 50, 75, and \qty{100}{\milli\second} to demonstrate the linear
relationship between network delay and handshake time.

CTaps does not support TLS, and the
TCP handshakes are therefore run without encryption.
The TCP 3-way handshake completes in one RTT, while 
QUIC mandates TLS 1.3 and requires two RTTs for a full handshake
when including the \mono{HANDSHAKE_DONE} frame.

The handshake client used for CTaps here is the same as
the client used for the candidate racing experiment in~\cref{sec:handshake-kde}.
This client has a preference
for QUIC and is given two Remote Endpoints. This is the reason for the
\qty{500}{\milli\second} stagger delay for the CTaps TCP
handshake time. CTaps tries two Remote Endpoints for QUIC before trying TCP
and the stagger delay is a static \qty{250}{\milli\second} for each attempt.
The high cost of this stagger delay is discussed further in~\cref{sec:discussion}.

\begin{figure}[H]
    %% Creator: Matplotlib, PGF backend
%%
%% To include the figure in your LaTeX document, write
%%   \input{<filename>.pgf}
%%
%% Make sure the required packages are loaded in your preamble
%%   \usepackage{pgf}
%%
%% Also ensure that all the required font packages are loaded; for instance,
%% the lmodern package is sometimes necessary when using math font.
%%   \usepackage{lmodern}
%%
%% Figures using additional raster images can only be included by \input if
%% they are in the same directory as the main LaTeX file. For loading figures
%% from other directories you can use the `import` package
%%   \usepackage{import}
%%
%% and then include the figures with
%%   \import{<path to file>}{<filename>.pgf}
%%
%% Matplotlib used the following preamble
%%   \def\mathdefault#1{#1}
%%   \everymath=\expandafter{\the\everymath\displaystyle}
%%   \IfFileExists{scrextend.sty}{
%%     \usepackage[fontsize=10.000000pt]{scrextend}
%%   }{
%%     \renewcommand{\normalsize}{\fontsize{10.000000}{12.000000}\selectfont}
%%     \normalsize
%%   }
%%   
%%   \makeatletter\@ifpackageloaded{underscore}{}{\usepackage[strings]{underscore}}\makeatother
%%
\begingroup%
\makeatletter%
\begin{pgfpicture}%
\pgfpathrectangle{\pgfpointorigin}{\pgfqpoint{5.369125in}{3.280176in}}%
\pgfusepath{use as bounding box, clip}%
\begin{pgfscope}%
\pgfsetbuttcap%
\pgfsetmiterjoin%
\definecolor{currentfill}{rgb}{1.000000,1.000000,1.000000}%
\pgfsetfillcolor{currentfill}%
\pgfsetlinewidth{0.000000pt}%
\definecolor{currentstroke}{rgb}{1.000000,1.000000,1.000000}%
\pgfsetstrokecolor{currentstroke}%
\pgfsetdash{}{0pt}%
\pgfpathmoveto{\pgfqpoint{-0.000000in}{0.000000in}}%
\pgfpathlineto{\pgfqpoint{5.369125in}{0.000000in}}%
\pgfpathlineto{\pgfqpoint{5.369125in}{3.280176in}}%
\pgfpathlineto{\pgfqpoint{-0.000000in}{3.280176in}}%
\pgfpathlineto{\pgfqpoint{-0.000000in}{0.000000in}}%
\pgfpathclose%
\pgfusepath{fill}%
\end{pgfscope}%
\begin{pgfscope}%
\pgfsetbuttcap%
\pgfsetmiterjoin%
\definecolor{currentfill}{rgb}{1.000000,1.000000,1.000000}%
\pgfsetfillcolor{currentfill}%
\pgfsetlinewidth{0.000000pt}%
\definecolor{currentstroke}{rgb}{0.000000,0.000000,0.000000}%
\pgfsetstrokecolor{currentstroke}%
\pgfsetstrokeopacity{0.000000}%
\pgfsetdash{}{0pt}%
\pgfpathmoveto{\pgfqpoint{0.613195in}{0.528318in}}%
\pgfpathlineto{\pgfqpoint{5.269125in}{0.528318in}}%
\pgfpathlineto{\pgfqpoint{5.269125in}{3.180176in}}%
\pgfpathlineto{\pgfqpoint{0.613195in}{3.180176in}}%
\pgfpathlineto{\pgfqpoint{0.613195in}{0.528318in}}%
\pgfpathclose%
\pgfusepath{fill}%
\end{pgfscope}%
\begin{pgfscope}%
\pgfpathrectangle{\pgfqpoint{0.613195in}{0.528318in}}{\pgfqpoint{4.655930in}{2.651859in}}%
\pgfusepath{clip}%
\pgfsetbuttcap%
\pgfsetroundjoin%
\pgfsetlinewidth{0.803000pt}%
\definecolor{currentstroke}{rgb}{0.690196,0.690196,0.690196}%
\pgfsetstrokecolor{currentstroke}%
\pgfsetstrokeopacity{0.400000}%
\pgfsetdash{{2.960000pt}{1.280000pt}}{0.000000pt}%
\pgfpathmoveto{\pgfqpoint{0.824828in}{0.528318in}}%
\pgfpathlineto{\pgfqpoint{0.824828in}{3.180176in}}%
\pgfusepath{stroke}%
\end{pgfscope}%
\begin{pgfscope}%
\pgfsetbuttcap%
\pgfsetroundjoin%
\definecolor{currentfill}{rgb}{0.000000,0.000000,0.000000}%
\pgfsetfillcolor{currentfill}%
\pgfsetlinewidth{0.803000pt}%
\definecolor{currentstroke}{rgb}{0.000000,0.000000,0.000000}%
\pgfsetstrokecolor{currentstroke}%
\pgfsetdash{}{0pt}%
\pgfsys@defobject{currentmarker}{\pgfqpoint{0.000000in}{-0.048611in}}{\pgfqpoint{0.000000in}{0.000000in}}{%
\pgfpathmoveto{\pgfqpoint{0.000000in}{0.000000in}}%
\pgfpathlineto{\pgfqpoint{0.000000in}{-0.048611in}}%
\pgfusepath{stroke,fill}%
}%
\begin{pgfscope}%
\pgfsys@transformshift{0.824828in}{0.528318in}%
\pgfsys@useobject{currentmarker}{}%
\end{pgfscope}%
\end{pgfscope}%
\begin{pgfscope}%
\definecolor{textcolor}{rgb}{0.000000,0.000000,0.000000}%
\pgfsetstrokecolor{textcolor}%
\pgfsetfillcolor{textcolor}%
\pgftext[x=0.824828in,y=0.431095in,,top]{\color{textcolor}{\rmfamily\fontsize{10.000000}{12.000000}\selectfont\catcode`\^=\active\def^{\ifmmode\sp\else\^{}\fi}\catcode`\%=\active\def%{\%}$\mathdefault{50}$}}%
\end{pgfscope}%
\begin{pgfscope}%
\pgfpathrectangle{\pgfqpoint{0.613195in}{0.528318in}}{\pgfqpoint{4.655930in}{2.651859in}}%
\pgfusepath{clip}%
\pgfsetbuttcap%
\pgfsetroundjoin%
\pgfsetlinewidth{0.803000pt}%
\definecolor{currentstroke}{rgb}{0.690196,0.690196,0.690196}%
\pgfsetstrokecolor{currentstroke}%
\pgfsetstrokeopacity{0.400000}%
\pgfsetdash{{2.960000pt}{1.280000pt}}{0.000000pt}%
\pgfpathmoveto{\pgfqpoint{2.235716in}{0.528318in}}%
\pgfpathlineto{\pgfqpoint{2.235716in}{3.180176in}}%
\pgfusepath{stroke}%
\end{pgfscope}%
\begin{pgfscope}%
\pgfsetbuttcap%
\pgfsetroundjoin%
\definecolor{currentfill}{rgb}{0.000000,0.000000,0.000000}%
\pgfsetfillcolor{currentfill}%
\pgfsetlinewidth{0.803000pt}%
\definecolor{currentstroke}{rgb}{0.000000,0.000000,0.000000}%
\pgfsetstrokecolor{currentstroke}%
\pgfsetdash{}{0pt}%
\pgfsys@defobject{currentmarker}{\pgfqpoint{0.000000in}{-0.048611in}}{\pgfqpoint{0.000000in}{0.000000in}}{%
\pgfpathmoveto{\pgfqpoint{0.000000in}{0.000000in}}%
\pgfpathlineto{\pgfqpoint{0.000000in}{-0.048611in}}%
\pgfusepath{stroke,fill}%
}%
\begin{pgfscope}%
\pgfsys@transformshift{2.235716in}{0.528318in}%
\pgfsys@useobject{currentmarker}{}%
\end{pgfscope}%
\end{pgfscope}%
\begin{pgfscope}%
\definecolor{textcolor}{rgb}{0.000000,0.000000,0.000000}%
\pgfsetstrokecolor{textcolor}%
\pgfsetfillcolor{textcolor}%
\pgftext[x=2.235716in,y=0.431095in,,top]{\color{textcolor}{\rmfamily\fontsize{10.000000}{12.000000}\selectfont\catcode`\^=\active\def^{\ifmmode\sp\else\^{}\fi}\catcode`\%=\active\def%{\%}$\mathdefault{100}$}}%
\end{pgfscope}%
\begin{pgfscope}%
\pgfpathrectangle{\pgfqpoint{0.613195in}{0.528318in}}{\pgfqpoint{4.655930in}{2.651859in}}%
\pgfusepath{clip}%
\pgfsetbuttcap%
\pgfsetroundjoin%
\pgfsetlinewidth{0.803000pt}%
\definecolor{currentstroke}{rgb}{0.690196,0.690196,0.690196}%
\pgfsetstrokecolor{currentstroke}%
\pgfsetstrokeopacity{0.400000}%
\pgfsetdash{{2.960000pt}{1.280000pt}}{0.000000pt}%
\pgfpathmoveto{\pgfqpoint{3.646604in}{0.528318in}}%
\pgfpathlineto{\pgfqpoint{3.646604in}{3.180176in}}%
\pgfusepath{stroke}%
\end{pgfscope}%
\begin{pgfscope}%
\pgfsetbuttcap%
\pgfsetroundjoin%
\definecolor{currentfill}{rgb}{0.000000,0.000000,0.000000}%
\pgfsetfillcolor{currentfill}%
\pgfsetlinewidth{0.803000pt}%
\definecolor{currentstroke}{rgb}{0.000000,0.000000,0.000000}%
\pgfsetstrokecolor{currentstroke}%
\pgfsetdash{}{0pt}%
\pgfsys@defobject{currentmarker}{\pgfqpoint{0.000000in}{-0.048611in}}{\pgfqpoint{0.000000in}{0.000000in}}{%
\pgfpathmoveto{\pgfqpoint{0.000000in}{0.000000in}}%
\pgfpathlineto{\pgfqpoint{0.000000in}{-0.048611in}}%
\pgfusepath{stroke,fill}%
}%
\begin{pgfscope}%
\pgfsys@transformshift{3.646604in}{0.528318in}%
\pgfsys@useobject{currentmarker}{}%
\end{pgfscope}%
\end{pgfscope}%
\begin{pgfscope}%
\definecolor{textcolor}{rgb}{0.000000,0.000000,0.000000}%
\pgfsetstrokecolor{textcolor}%
\pgfsetfillcolor{textcolor}%
\pgftext[x=3.646604in,y=0.431095in,,top]{\color{textcolor}{\rmfamily\fontsize{10.000000}{12.000000}\selectfont\catcode`\^=\active\def^{\ifmmode\sp\else\^{}\fi}\catcode`\%=\active\def%{\%}$\mathdefault{150}$}}%
\end{pgfscope}%
\begin{pgfscope}%
\pgfpathrectangle{\pgfqpoint{0.613195in}{0.528318in}}{\pgfqpoint{4.655930in}{2.651859in}}%
\pgfusepath{clip}%
\pgfsetbuttcap%
\pgfsetroundjoin%
\pgfsetlinewidth{0.803000pt}%
\definecolor{currentstroke}{rgb}{0.690196,0.690196,0.690196}%
\pgfsetstrokecolor{currentstroke}%
\pgfsetstrokeopacity{0.400000}%
\pgfsetdash{{2.960000pt}{1.280000pt}}{0.000000pt}%
\pgfpathmoveto{\pgfqpoint{5.057492in}{0.528318in}}%
\pgfpathlineto{\pgfqpoint{5.057492in}{3.180176in}}%
\pgfusepath{stroke}%
\end{pgfscope}%
\begin{pgfscope}%
\pgfsetbuttcap%
\pgfsetroundjoin%
\definecolor{currentfill}{rgb}{0.000000,0.000000,0.000000}%
\pgfsetfillcolor{currentfill}%
\pgfsetlinewidth{0.803000pt}%
\definecolor{currentstroke}{rgb}{0.000000,0.000000,0.000000}%
\pgfsetstrokecolor{currentstroke}%
\pgfsetdash{}{0pt}%
\pgfsys@defobject{currentmarker}{\pgfqpoint{0.000000in}{-0.048611in}}{\pgfqpoint{0.000000in}{0.000000in}}{%
\pgfpathmoveto{\pgfqpoint{0.000000in}{0.000000in}}%
\pgfpathlineto{\pgfqpoint{0.000000in}{-0.048611in}}%
\pgfusepath{stroke,fill}%
}%
\begin{pgfscope}%
\pgfsys@transformshift{5.057492in}{0.528318in}%
\pgfsys@useobject{currentmarker}{}%
\end{pgfscope}%
\end{pgfscope}%
\begin{pgfscope}%
\definecolor{textcolor}{rgb}{0.000000,0.000000,0.000000}%
\pgfsetstrokecolor{textcolor}%
\pgfsetfillcolor{textcolor}%
\pgftext[x=5.057492in,y=0.431095in,,top]{\color{textcolor}{\rmfamily\fontsize{10.000000}{12.000000}\selectfont\catcode`\^=\active\def^{\ifmmode\sp\else\^{}\fi}\catcode`\%=\active\def%{\%}$\mathdefault{200}$}}%
\end{pgfscope}%
\begin{pgfscope}%
\definecolor{textcolor}{rgb}{0.000000,0.000000,0.000000}%
\pgfsetstrokecolor{textcolor}%
\pgfsetfillcolor{textcolor}%
\pgftext[x=2.941160in,y=0.252083in,,top]{\color{textcolor}{\rmfamily\fontsize{11.000000}{13.200000}\selectfont\catcode`\^=\active\def^{\ifmmode\sp\else\^{}\fi}\catcode`\%=\active\def%{\%}RTT (ms)}}%
\end{pgfscope}%
\begin{pgfscope}%
\pgfpathrectangle{\pgfqpoint{0.613195in}{0.528318in}}{\pgfqpoint{4.655930in}{2.651859in}}%
\pgfusepath{clip}%
\pgfsetbuttcap%
\pgfsetroundjoin%
\pgfsetlinewidth{0.803000pt}%
\definecolor{currentstroke}{rgb}{0.690196,0.690196,0.690196}%
\pgfsetstrokecolor{currentstroke}%
\pgfsetstrokeopacity{0.400000}%
\pgfsetdash{{2.960000pt}{1.280000pt}}{0.000000pt}%
\pgfpathmoveto{\pgfqpoint{0.613195in}{0.830114in}}%
\pgfpathlineto{\pgfqpoint{5.269125in}{0.830114in}}%
\pgfusepath{stroke}%
\end{pgfscope}%
\begin{pgfscope}%
\pgfsetbuttcap%
\pgfsetroundjoin%
\definecolor{currentfill}{rgb}{0.000000,0.000000,0.000000}%
\pgfsetfillcolor{currentfill}%
\pgfsetlinewidth{0.803000pt}%
\definecolor{currentstroke}{rgb}{0.000000,0.000000,0.000000}%
\pgfsetstrokecolor{currentstroke}%
\pgfsetdash{}{0pt}%
\pgfsys@defobject{currentmarker}{\pgfqpoint{-0.048611in}{0.000000in}}{\pgfqpoint{-0.000000in}{0.000000in}}{%
\pgfpathmoveto{\pgfqpoint{-0.000000in}{0.000000in}}%
\pgfpathlineto{\pgfqpoint{-0.048611in}{0.000000in}}%
\pgfusepath{stroke,fill}%
}%
\begin{pgfscope}%
\pgfsys@transformshift{0.613195in}{0.830114in}%
\pgfsys@useobject{currentmarker}{}%
\end{pgfscope}%
\end{pgfscope}%
\begin{pgfscope}%
\definecolor{textcolor}{rgb}{0.000000,0.000000,0.000000}%
\pgfsetstrokecolor{textcolor}%
\pgfsetfillcolor{textcolor}%
\pgftext[x=0.307639in, y=0.781889in, left, base]{\color{textcolor}{\rmfamily\fontsize{10.000000}{12.000000}\selectfont\catcode`\^=\active\def^{\ifmmode\sp\else\^{}\fi}\catcode`\%=\active\def%{\%}$\mathdefault{100}$}}%
\end{pgfscope}%
\begin{pgfscope}%
\pgfpathrectangle{\pgfqpoint{0.613195in}{0.528318in}}{\pgfqpoint{4.655930in}{2.651859in}}%
\pgfusepath{clip}%
\pgfsetbuttcap%
\pgfsetroundjoin%
\pgfsetlinewidth{0.803000pt}%
\definecolor{currentstroke}{rgb}{0.690196,0.690196,0.690196}%
\pgfsetstrokecolor{currentstroke}%
\pgfsetstrokeopacity{0.400000}%
\pgfsetdash{{2.960000pt}{1.280000pt}}{0.000000pt}%
\pgfpathmoveto{\pgfqpoint{0.613195in}{1.200633in}}%
\pgfpathlineto{\pgfqpoint{5.269125in}{1.200633in}}%
\pgfusepath{stroke}%
\end{pgfscope}%
\begin{pgfscope}%
\pgfsetbuttcap%
\pgfsetroundjoin%
\definecolor{currentfill}{rgb}{0.000000,0.000000,0.000000}%
\pgfsetfillcolor{currentfill}%
\pgfsetlinewidth{0.803000pt}%
\definecolor{currentstroke}{rgb}{0.000000,0.000000,0.000000}%
\pgfsetstrokecolor{currentstroke}%
\pgfsetdash{}{0pt}%
\pgfsys@defobject{currentmarker}{\pgfqpoint{-0.048611in}{0.000000in}}{\pgfqpoint{-0.000000in}{0.000000in}}{%
\pgfpathmoveto{\pgfqpoint{-0.000000in}{0.000000in}}%
\pgfpathlineto{\pgfqpoint{-0.048611in}{0.000000in}}%
\pgfusepath{stroke,fill}%
}%
\begin{pgfscope}%
\pgfsys@transformshift{0.613195in}{1.200633in}%
\pgfsys@useobject{currentmarker}{}%
\end{pgfscope}%
\end{pgfscope}%
\begin{pgfscope}%
\definecolor{textcolor}{rgb}{0.000000,0.000000,0.000000}%
\pgfsetstrokecolor{textcolor}%
\pgfsetfillcolor{textcolor}%
\pgftext[x=0.307639in, y=1.152408in, left, base]{\color{textcolor}{\rmfamily\fontsize{10.000000}{12.000000}\selectfont\catcode`\^=\active\def^{\ifmmode\sp\else\^{}\fi}\catcode`\%=\active\def%{\%}$\mathdefault{200}$}}%
\end{pgfscope}%
\begin{pgfscope}%
\pgfpathrectangle{\pgfqpoint{0.613195in}{0.528318in}}{\pgfqpoint{4.655930in}{2.651859in}}%
\pgfusepath{clip}%
\pgfsetbuttcap%
\pgfsetroundjoin%
\pgfsetlinewidth{0.803000pt}%
\definecolor{currentstroke}{rgb}{0.690196,0.690196,0.690196}%
\pgfsetstrokecolor{currentstroke}%
\pgfsetstrokeopacity{0.400000}%
\pgfsetdash{{2.960000pt}{1.280000pt}}{0.000000pt}%
\pgfpathmoveto{\pgfqpoint{0.613195in}{1.571152in}}%
\pgfpathlineto{\pgfqpoint{5.269125in}{1.571152in}}%
\pgfusepath{stroke}%
\end{pgfscope}%
\begin{pgfscope}%
\pgfsetbuttcap%
\pgfsetroundjoin%
\definecolor{currentfill}{rgb}{0.000000,0.000000,0.000000}%
\pgfsetfillcolor{currentfill}%
\pgfsetlinewidth{0.803000pt}%
\definecolor{currentstroke}{rgb}{0.000000,0.000000,0.000000}%
\pgfsetstrokecolor{currentstroke}%
\pgfsetdash{}{0pt}%
\pgfsys@defobject{currentmarker}{\pgfqpoint{-0.048611in}{0.000000in}}{\pgfqpoint{-0.000000in}{0.000000in}}{%
\pgfpathmoveto{\pgfqpoint{-0.000000in}{0.000000in}}%
\pgfpathlineto{\pgfqpoint{-0.048611in}{0.000000in}}%
\pgfusepath{stroke,fill}%
}%
\begin{pgfscope}%
\pgfsys@transformshift{0.613195in}{1.571152in}%
\pgfsys@useobject{currentmarker}{}%
\end{pgfscope}%
\end{pgfscope}%
\begin{pgfscope}%
\definecolor{textcolor}{rgb}{0.000000,0.000000,0.000000}%
\pgfsetstrokecolor{textcolor}%
\pgfsetfillcolor{textcolor}%
\pgftext[x=0.307639in, y=1.522927in, left, base]{\color{textcolor}{\rmfamily\fontsize{10.000000}{12.000000}\selectfont\catcode`\^=\active\def^{\ifmmode\sp\else\^{}\fi}\catcode`\%=\active\def%{\%}$\mathdefault{300}$}}%
\end{pgfscope}%
\begin{pgfscope}%
\pgfpathrectangle{\pgfqpoint{0.613195in}{0.528318in}}{\pgfqpoint{4.655930in}{2.651859in}}%
\pgfusepath{clip}%
\pgfsetbuttcap%
\pgfsetroundjoin%
\pgfsetlinewidth{0.803000pt}%
\definecolor{currentstroke}{rgb}{0.690196,0.690196,0.690196}%
\pgfsetstrokecolor{currentstroke}%
\pgfsetstrokeopacity{0.400000}%
\pgfsetdash{{2.960000pt}{1.280000pt}}{0.000000pt}%
\pgfpathmoveto{\pgfqpoint{0.613195in}{1.941671in}}%
\pgfpathlineto{\pgfqpoint{5.269125in}{1.941671in}}%
\pgfusepath{stroke}%
\end{pgfscope}%
\begin{pgfscope}%
\pgfsetbuttcap%
\pgfsetroundjoin%
\definecolor{currentfill}{rgb}{0.000000,0.000000,0.000000}%
\pgfsetfillcolor{currentfill}%
\pgfsetlinewidth{0.803000pt}%
\definecolor{currentstroke}{rgb}{0.000000,0.000000,0.000000}%
\pgfsetstrokecolor{currentstroke}%
\pgfsetdash{}{0pt}%
\pgfsys@defobject{currentmarker}{\pgfqpoint{-0.048611in}{0.000000in}}{\pgfqpoint{-0.000000in}{0.000000in}}{%
\pgfpathmoveto{\pgfqpoint{-0.000000in}{0.000000in}}%
\pgfpathlineto{\pgfqpoint{-0.048611in}{0.000000in}}%
\pgfusepath{stroke,fill}%
}%
\begin{pgfscope}%
\pgfsys@transformshift{0.613195in}{1.941671in}%
\pgfsys@useobject{currentmarker}{}%
\end{pgfscope}%
\end{pgfscope}%
\begin{pgfscope}%
\definecolor{textcolor}{rgb}{0.000000,0.000000,0.000000}%
\pgfsetstrokecolor{textcolor}%
\pgfsetfillcolor{textcolor}%
\pgftext[x=0.307639in, y=1.893446in, left, base]{\color{textcolor}{\rmfamily\fontsize{10.000000}{12.000000}\selectfont\catcode`\^=\active\def^{\ifmmode\sp\else\^{}\fi}\catcode`\%=\active\def%{\%}$\mathdefault{400}$}}%
\end{pgfscope}%
\begin{pgfscope}%
\pgfpathrectangle{\pgfqpoint{0.613195in}{0.528318in}}{\pgfqpoint{4.655930in}{2.651859in}}%
\pgfusepath{clip}%
\pgfsetbuttcap%
\pgfsetroundjoin%
\pgfsetlinewidth{0.803000pt}%
\definecolor{currentstroke}{rgb}{0.690196,0.690196,0.690196}%
\pgfsetstrokecolor{currentstroke}%
\pgfsetstrokeopacity{0.400000}%
\pgfsetdash{{2.960000pt}{1.280000pt}}{0.000000pt}%
\pgfpathmoveto{\pgfqpoint{0.613195in}{2.312190in}}%
\pgfpathlineto{\pgfqpoint{5.269125in}{2.312190in}}%
\pgfusepath{stroke}%
\end{pgfscope}%
\begin{pgfscope}%
\pgfsetbuttcap%
\pgfsetroundjoin%
\definecolor{currentfill}{rgb}{0.000000,0.000000,0.000000}%
\pgfsetfillcolor{currentfill}%
\pgfsetlinewidth{0.803000pt}%
\definecolor{currentstroke}{rgb}{0.000000,0.000000,0.000000}%
\pgfsetstrokecolor{currentstroke}%
\pgfsetdash{}{0pt}%
\pgfsys@defobject{currentmarker}{\pgfqpoint{-0.048611in}{0.000000in}}{\pgfqpoint{-0.000000in}{0.000000in}}{%
\pgfpathmoveto{\pgfqpoint{-0.000000in}{0.000000in}}%
\pgfpathlineto{\pgfqpoint{-0.048611in}{0.000000in}}%
\pgfusepath{stroke,fill}%
}%
\begin{pgfscope}%
\pgfsys@transformshift{0.613195in}{2.312190in}%
\pgfsys@useobject{currentmarker}{}%
\end{pgfscope}%
\end{pgfscope}%
\begin{pgfscope}%
\definecolor{textcolor}{rgb}{0.000000,0.000000,0.000000}%
\pgfsetstrokecolor{textcolor}%
\pgfsetfillcolor{textcolor}%
\pgftext[x=0.307639in, y=2.263964in, left, base]{\color{textcolor}{\rmfamily\fontsize{10.000000}{12.000000}\selectfont\catcode`\^=\active\def^{\ifmmode\sp\else\^{}\fi}\catcode`\%=\active\def%{\%}$\mathdefault{500}$}}%
\end{pgfscope}%
\begin{pgfscope}%
\pgfpathrectangle{\pgfqpoint{0.613195in}{0.528318in}}{\pgfqpoint{4.655930in}{2.651859in}}%
\pgfusepath{clip}%
\pgfsetbuttcap%
\pgfsetroundjoin%
\pgfsetlinewidth{0.803000pt}%
\definecolor{currentstroke}{rgb}{0.690196,0.690196,0.690196}%
\pgfsetstrokecolor{currentstroke}%
\pgfsetstrokeopacity{0.400000}%
\pgfsetdash{{2.960000pt}{1.280000pt}}{0.000000pt}%
\pgfpathmoveto{\pgfqpoint{0.613195in}{2.682708in}}%
\pgfpathlineto{\pgfqpoint{5.269125in}{2.682708in}}%
\pgfusepath{stroke}%
\end{pgfscope}%
\begin{pgfscope}%
\pgfsetbuttcap%
\pgfsetroundjoin%
\definecolor{currentfill}{rgb}{0.000000,0.000000,0.000000}%
\pgfsetfillcolor{currentfill}%
\pgfsetlinewidth{0.803000pt}%
\definecolor{currentstroke}{rgb}{0.000000,0.000000,0.000000}%
\pgfsetstrokecolor{currentstroke}%
\pgfsetdash{}{0pt}%
\pgfsys@defobject{currentmarker}{\pgfqpoint{-0.048611in}{0.000000in}}{\pgfqpoint{-0.000000in}{0.000000in}}{%
\pgfpathmoveto{\pgfqpoint{-0.000000in}{0.000000in}}%
\pgfpathlineto{\pgfqpoint{-0.048611in}{0.000000in}}%
\pgfusepath{stroke,fill}%
}%
\begin{pgfscope}%
\pgfsys@transformshift{0.613195in}{2.682708in}%
\pgfsys@useobject{currentmarker}{}%
\end{pgfscope}%
\end{pgfscope}%
\begin{pgfscope}%
\definecolor{textcolor}{rgb}{0.000000,0.000000,0.000000}%
\pgfsetstrokecolor{textcolor}%
\pgfsetfillcolor{textcolor}%
\pgftext[x=0.307639in, y=2.634483in, left, base]{\color{textcolor}{\rmfamily\fontsize{10.000000}{12.000000}\selectfont\catcode`\^=\active\def^{\ifmmode\sp\else\^{}\fi}\catcode`\%=\active\def%{\%}$\mathdefault{600}$}}%
\end{pgfscope}%
\begin{pgfscope}%
\pgfpathrectangle{\pgfqpoint{0.613195in}{0.528318in}}{\pgfqpoint{4.655930in}{2.651859in}}%
\pgfusepath{clip}%
\pgfsetbuttcap%
\pgfsetroundjoin%
\pgfsetlinewidth{0.803000pt}%
\definecolor{currentstroke}{rgb}{0.690196,0.690196,0.690196}%
\pgfsetstrokecolor{currentstroke}%
\pgfsetstrokeopacity{0.400000}%
\pgfsetdash{{2.960000pt}{1.280000pt}}{0.000000pt}%
\pgfpathmoveto{\pgfqpoint{0.613195in}{3.053227in}}%
\pgfpathlineto{\pgfqpoint{5.269125in}{3.053227in}}%
\pgfusepath{stroke}%
\end{pgfscope}%
\begin{pgfscope}%
\pgfsetbuttcap%
\pgfsetroundjoin%
\definecolor{currentfill}{rgb}{0.000000,0.000000,0.000000}%
\pgfsetfillcolor{currentfill}%
\pgfsetlinewidth{0.803000pt}%
\definecolor{currentstroke}{rgb}{0.000000,0.000000,0.000000}%
\pgfsetstrokecolor{currentstroke}%
\pgfsetdash{}{0pt}%
\pgfsys@defobject{currentmarker}{\pgfqpoint{-0.048611in}{0.000000in}}{\pgfqpoint{-0.000000in}{0.000000in}}{%
\pgfpathmoveto{\pgfqpoint{-0.000000in}{0.000000in}}%
\pgfpathlineto{\pgfqpoint{-0.048611in}{0.000000in}}%
\pgfusepath{stroke,fill}%
}%
\begin{pgfscope}%
\pgfsys@transformshift{0.613195in}{3.053227in}%
\pgfsys@useobject{currentmarker}{}%
\end{pgfscope}%
\end{pgfscope}%
\begin{pgfscope}%
\definecolor{textcolor}{rgb}{0.000000,0.000000,0.000000}%
\pgfsetstrokecolor{textcolor}%
\pgfsetfillcolor{textcolor}%
\pgftext[x=0.307639in, y=3.005002in, left, base]{\color{textcolor}{\rmfamily\fontsize{10.000000}{12.000000}\selectfont\catcode`\^=\active\def^{\ifmmode\sp\else\^{}\fi}\catcode`\%=\active\def%{\%}$\mathdefault{700}$}}%
\end{pgfscope}%
\begin{pgfscope}%
\definecolor{textcolor}{rgb}{0.000000,0.000000,0.000000}%
\pgfsetstrokecolor{textcolor}%
\pgfsetfillcolor{textcolor}%
\pgftext[x=0.252083in,y=1.854247in,,bottom,rotate=90.000000]{\color{textcolor}{\rmfamily\fontsize{11.000000}{13.200000}\selectfont\catcode`\^=\active\def^{\ifmmode\sp\else\^{}\fi}\catcode`\%=\active\def%{\%}Handshake time (ms)}}%
\end{pgfscope}%
\begin{pgfscope}%
\pgfsetrectcap%
\pgfsetmiterjoin%
\pgfsetlinewidth{0.803000pt}%
\definecolor{currentstroke}{rgb}{0.000000,0.000000,0.000000}%
\pgfsetstrokecolor{currentstroke}%
\pgfsetdash{}{0pt}%
\pgfpathmoveto{\pgfqpoint{0.613195in}{0.528318in}}%
\pgfpathlineto{\pgfqpoint{0.613195in}{3.180176in}}%
\pgfusepath{stroke}%
\end{pgfscope}%
\begin{pgfscope}%
\pgfsetrectcap%
\pgfsetmiterjoin%
\pgfsetlinewidth{0.803000pt}%
\definecolor{currentstroke}{rgb}{0.000000,0.000000,0.000000}%
\pgfsetstrokecolor{currentstroke}%
\pgfsetdash{}{0pt}%
\pgfpathmoveto{\pgfqpoint{5.269125in}{0.528318in}}%
\pgfpathlineto{\pgfqpoint{5.269125in}{3.180176in}}%
\pgfusepath{stroke}%
\end{pgfscope}%
\begin{pgfscope}%
\pgfsetrectcap%
\pgfsetmiterjoin%
\pgfsetlinewidth{0.803000pt}%
\definecolor{currentstroke}{rgb}{0.000000,0.000000,0.000000}%
\pgfsetstrokecolor{currentstroke}%
\pgfsetdash{}{0pt}%
\pgfpathmoveto{\pgfqpoint{0.613195in}{0.528318in}}%
\pgfpathlineto{\pgfqpoint{5.269125in}{0.528318in}}%
\pgfusepath{stroke}%
\end{pgfscope}%
\begin{pgfscope}%
\pgfsetrectcap%
\pgfsetmiterjoin%
\pgfsetlinewidth{0.803000pt}%
\definecolor{currentstroke}{rgb}{0.000000,0.000000,0.000000}%
\pgfsetstrokecolor{currentstroke}%
\pgfsetdash{}{0pt}%
\pgfpathmoveto{\pgfqpoint{0.613195in}{3.180176in}}%
\pgfpathlineto{\pgfqpoint{5.269125in}{3.180176in}}%
\pgfusepath{stroke}%
\end{pgfscope}%
\begin{pgfscope}%
\pgfpathrectangle{\pgfqpoint{0.613195in}{0.528318in}}{\pgfqpoint{4.655930in}{2.651859in}}%
\pgfusepath{clip}%
\pgfsetrectcap%
\pgfsetroundjoin%
\pgfsetlinewidth{1.505625pt}%
\definecolor{currentstroke}{rgb}{0.082353,0.396078,0.752941}%
\pgfsetstrokecolor{currentstroke}%
\pgfsetdash{}{0pt}%
\pgfpathmoveto{\pgfqpoint{0.824828in}{0.648857in}}%
\pgfpathlineto{\pgfqpoint{2.235716in}{0.832819in}}%
\pgfpathlineto{\pgfqpoint{3.646604in}{1.018338in}}%
\pgfpathlineto{\pgfqpoint{5.057492in}{1.203968in}}%
\pgfusepath{stroke}%
\end{pgfscope}%
\begin{pgfscope}%
\pgfpathrectangle{\pgfqpoint{0.613195in}{0.528318in}}{\pgfqpoint{4.655930in}{2.651859in}}%
\pgfusepath{clip}%
\pgfsetbuttcap%
\pgfsetmiterjoin%
\definecolor{currentfill}{rgb}{0.082353,0.396078,0.752941}%
\pgfsetfillcolor{currentfill}%
\pgfsetlinewidth{1.003750pt}%
\definecolor{currentstroke}{rgb}{0.082353,0.396078,0.752941}%
\pgfsetstrokecolor{currentstroke}%
\pgfsetdash{}{0pt}%
\pgfsys@defobject{currentmarker}{\pgfqpoint{-0.027778in}{-0.027778in}}{\pgfqpoint{0.027778in}{0.027778in}}{%
\pgfpathmoveto{\pgfqpoint{-0.027778in}{-0.027778in}}%
\pgfpathlineto{\pgfqpoint{0.027778in}{-0.027778in}}%
\pgfpathlineto{\pgfqpoint{0.027778in}{0.027778in}}%
\pgfpathlineto{\pgfqpoint{-0.027778in}{0.027778in}}%
\pgfpathlineto{\pgfqpoint{-0.027778in}{-0.027778in}}%
\pgfpathclose%
\pgfusepath{stroke,fill}%
}%
\begin{pgfscope}%
\pgfsys@transformshift{0.824828in}{0.648857in}%
\pgfsys@useobject{currentmarker}{}%
\end{pgfscope}%
\begin{pgfscope}%
\pgfsys@transformshift{2.235716in}{0.832819in}%
\pgfsys@useobject{currentmarker}{}%
\end{pgfscope}%
\begin{pgfscope}%
\pgfsys@transformshift{3.646604in}{1.018338in}%
\pgfsys@useobject{currentmarker}{}%
\end{pgfscope}%
\begin{pgfscope}%
\pgfsys@transformshift{5.057492in}{1.203968in}%
\pgfsys@useobject{currentmarker}{}%
\end{pgfscope}%
\end{pgfscope}%
\begin{pgfscope}%
\pgfpathrectangle{\pgfqpoint{0.613195in}{0.528318in}}{\pgfqpoint{4.655930in}{2.651859in}}%
\pgfusepath{clip}%
\pgfsetrectcap%
\pgfsetroundjoin%
\pgfsetlinewidth{1.505625pt}%
\definecolor{currentstroke}{rgb}{0.180392,0.490196,0.196078}%
\pgfsetstrokecolor{currentstroke}%
\pgfsetdash{}{0pt}%
\pgfpathmoveto{\pgfqpoint{0.824828in}{0.872094in}}%
\pgfpathlineto{\pgfqpoint{2.235716in}{1.268549in}}%
\pgfpathlineto{\pgfqpoint{3.646604in}{1.617356in}}%
\pgfpathlineto{\pgfqpoint{5.057492in}{2.009846in}}%
\pgfusepath{stroke}%
\end{pgfscope}%
\begin{pgfscope}%
\pgfpathrectangle{\pgfqpoint{0.613195in}{0.528318in}}{\pgfqpoint{4.655930in}{2.651859in}}%
\pgfusepath{clip}%
\pgfsetbuttcap%
\pgfsetmiterjoin%
\definecolor{currentfill}{rgb}{0.180392,0.490196,0.196078}%
\pgfsetfillcolor{currentfill}%
\pgfsetlinewidth{1.003750pt}%
\definecolor{currentstroke}{rgb}{0.180392,0.490196,0.196078}%
\pgfsetstrokecolor{currentstroke}%
\pgfsetdash{}{0pt}%
\pgfsys@defobject{currentmarker}{\pgfqpoint{-0.027778in}{-0.027778in}}{\pgfqpoint{0.027778in}{0.027778in}}{%
\pgfpathmoveto{\pgfqpoint{-0.027778in}{-0.027778in}}%
\pgfpathlineto{\pgfqpoint{0.027778in}{-0.027778in}}%
\pgfpathlineto{\pgfqpoint{0.027778in}{0.027778in}}%
\pgfpathlineto{\pgfqpoint{-0.027778in}{0.027778in}}%
\pgfpathlineto{\pgfqpoint{-0.027778in}{-0.027778in}}%
\pgfpathclose%
\pgfusepath{stroke,fill}%
}%
\begin{pgfscope}%
\pgfsys@transformshift{0.824828in}{0.872094in}%
\pgfsys@useobject{currentmarker}{}%
\end{pgfscope}%
\begin{pgfscope}%
\pgfsys@transformshift{2.235716in}{1.268549in}%
\pgfsys@useobject{currentmarker}{}%
\end{pgfscope}%
\begin{pgfscope}%
\pgfsys@transformshift{3.646604in}{1.617356in}%
\pgfsys@useobject{currentmarker}{}%
\end{pgfscope}%
\begin{pgfscope}%
\pgfsys@transformshift{5.057492in}{2.009846in}%
\pgfsys@useobject{currentmarker}{}%
\end{pgfscope}%
\end{pgfscope}%
\begin{pgfscope}%
\pgfpathrectangle{\pgfqpoint{0.613195in}{0.528318in}}{\pgfqpoint{4.655930in}{2.651859in}}%
\pgfusepath{clip}%
\pgfsetbuttcap%
\pgfsetroundjoin%
\pgfsetlinewidth{1.505625pt}%
\definecolor{currentstroke}{rgb}{0.647059,0.839216,0.654902}%
\pgfsetstrokecolor{currentstroke}%
\pgfsetdash{{5.550000pt}{2.400000pt}}{0.000000pt}%
\pgfpathmoveto{\pgfqpoint{0.824828in}{0.952200in}}%
\pgfpathlineto{\pgfqpoint{2.235716in}{1.289780in}}%
\pgfpathlineto{\pgfqpoint{3.646604in}{1.667153in}}%
\pgfpathlineto{\pgfqpoint{5.057492in}{2.027187in}}%
\pgfusepath{stroke}%
\end{pgfscope}%
\begin{pgfscope}%
\pgfpathrectangle{\pgfqpoint{0.613195in}{0.528318in}}{\pgfqpoint{4.655930in}{2.651859in}}%
\pgfusepath{clip}%
\pgfsetbuttcap%
\pgfsetmiterjoin%
\definecolor{currentfill}{rgb}{0.647059,0.839216,0.654902}%
\pgfsetfillcolor{currentfill}%
\pgfsetlinewidth{1.003750pt}%
\definecolor{currentstroke}{rgb}{0.647059,0.839216,0.654902}%
\pgfsetstrokecolor{currentstroke}%
\pgfsetdash{}{0pt}%
\pgfsys@defobject{currentmarker}{\pgfqpoint{-0.027778in}{-0.027778in}}{\pgfqpoint{0.027778in}{0.027778in}}{%
\pgfpathmoveto{\pgfqpoint{-0.027778in}{-0.027778in}}%
\pgfpathlineto{\pgfqpoint{0.027778in}{-0.027778in}}%
\pgfpathlineto{\pgfqpoint{0.027778in}{0.027778in}}%
\pgfpathlineto{\pgfqpoint{-0.027778in}{0.027778in}}%
\pgfpathlineto{\pgfqpoint{-0.027778in}{-0.027778in}}%
\pgfpathclose%
\pgfusepath{stroke,fill}%
}%
\begin{pgfscope}%
\pgfsys@transformshift{0.824828in}{0.952200in}%
\pgfsys@useobject{currentmarker}{}%
\end{pgfscope}%
\begin{pgfscope}%
\pgfsys@transformshift{2.235716in}{1.289780in}%
\pgfsys@useobject{currentmarker}{}%
\end{pgfscope}%
\begin{pgfscope}%
\pgfsys@transformshift{3.646604in}{1.667153in}%
\pgfsys@useobject{currentmarker}{}%
\end{pgfscope}%
\begin{pgfscope}%
\pgfsys@transformshift{5.057492in}{2.027187in}%
\pgfsys@useobject{currentmarker}{}%
\end{pgfscope}%
\end{pgfscope}%
\begin{pgfscope}%
\pgfpathrectangle{\pgfqpoint{0.613195in}{0.528318in}}{\pgfqpoint{4.655930in}{2.651859in}}%
\pgfusepath{clip}%
\pgfsetbuttcap%
\pgfsetroundjoin%
\pgfsetlinewidth{1.505625pt}%
\definecolor{currentstroke}{rgb}{0.564706,0.792157,0.976471}%
\pgfsetstrokecolor{currentstroke}%
\pgfsetdash{{5.550000pt}{2.400000pt}}{0.000000pt}%
\pgfpathmoveto{\pgfqpoint{0.824828in}{2.513826in}}%
\pgfpathlineto{\pgfqpoint{2.235716in}{2.695565in}}%
\pgfpathlineto{\pgfqpoint{3.646604in}{2.882011in}}%
\pgfpathlineto{\pgfqpoint{5.057492in}{3.059637in}}%
\pgfusepath{stroke}%
\end{pgfscope}%
\begin{pgfscope}%
\pgfpathrectangle{\pgfqpoint{0.613195in}{0.528318in}}{\pgfqpoint{4.655930in}{2.651859in}}%
\pgfusepath{clip}%
\pgfsetbuttcap%
\pgfsetmiterjoin%
\definecolor{currentfill}{rgb}{0.564706,0.792157,0.976471}%
\pgfsetfillcolor{currentfill}%
\pgfsetlinewidth{1.003750pt}%
\definecolor{currentstroke}{rgb}{0.564706,0.792157,0.976471}%
\pgfsetstrokecolor{currentstroke}%
\pgfsetdash{}{0pt}%
\pgfsys@defobject{currentmarker}{\pgfqpoint{-0.027778in}{-0.027778in}}{\pgfqpoint{0.027778in}{0.027778in}}{%
\pgfpathmoveto{\pgfqpoint{-0.027778in}{-0.027778in}}%
\pgfpathlineto{\pgfqpoint{0.027778in}{-0.027778in}}%
\pgfpathlineto{\pgfqpoint{0.027778in}{0.027778in}}%
\pgfpathlineto{\pgfqpoint{-0.027778in}{0.027778in}}%
\pgfpathlineto{\pgfqpoint{-0.027778in}{-0.027778in}}%
\pgfpathclose%
\pgfusepath{stroke,fill}%
}%
\begin{pgfscope}%
\pgfsys@transformshift{0.824828in}{2.513826in}%
\pgfsys@useobject{currentmarker}{}%
\end{pgfscope}%
\begin{pgfscope}%
\pgfsys@transformshift{2.235716in}{2.695565in}%
\pgfsys@useobject{currentmarker}{}%
\end{pgfscope}%
\begin{pgfscope}%
\pgfsys@transformshift{3.646604in}{2.882011in}%
\pgfsys@useobject{currentmarker}{}%
\end{pgfscope}%
\begin{pgfscope}%
\pgfsys@transformshift{5.057492in}{3.059637in}%
\pgfsys@useobject{currentmarker}{}%
\end{pgfscope}%
\end{pgfscope}%
\begin{pgfscope}%
\pgfsetroundcap%
\pgfsetroundjoin%
\pgfsetlinewidth{1.003750pt}%
\definecolor{currentstroke}{rgb}{0.501961,0.501961,0.501961}%
\pgfsetstrokecolor{currentstroke}%
\pgfsetdash{}{0pt}%
\pgfpathmoveto{\pgfqpoint{1.220909in}{2.740872in}}%
\pgfpathquadraticcurveto{\pgfqpoint{1.034920in}{2.634257in}}{\pgfqpoint{0.862402in}{2.535365in}}%
\pgfusepath{stroke}%
\end{pgfscope}%
\begin{pgfscope}%
\pgfsetroundcap%
\pgfsetroundjoin%
\pgfsetlinewidth{1.003750pt}%
\definecolor{currentstroke}{rgb}{0.501961,0.501961,0.501961}%
\pgfsetstrokecolor{currentstroke}%
\pgfsetdash{}{0pt}%
\pgfpathmoveto{\pgfqpoint{0.912012in}{2.538188in}}%
\pgfpathlineto{\pgfqpoint{0.862402in}{2.535365in}}%
\pgfpathlineto{\pgfqpoint{0.889909in}{2.576747in}}%
\pgfusepath{stroke}%
\end{pgfscope}%
\begin{pgfscope}%
\definecolor{textcolor}{rgb}{0.501961,0.501961,0.501961}%
\pgfsetstrokecolor{textcolor}%
\pgfsetfillcolor{textcolor}%
\pgftext[x=1.107006in, y=2.931846in, left, base]{\color{textcolor}{\rmfamily\fontsize{8.000000}{9.600000}\selectfont\catcode`\^=\active\def^{\ifmmode\sp\else\^{}\fi}\catcode`\%=\active\def%{\%}$+500\,\mathrm{ms}$}}%
\end{pgfscope}%
\begin{pgfscope}%
\definecolor{textcolor}{rgb}{0.501961,0.501961,0.501961}%
\pgfsetstrokecolor{textcolor}%
\pgfsetfillcolor{textcolor}%
\pgftext[x=1.107006in, y=2.810241in, left, base]{\color{textcolor}{\rmfamily\fontsize{8.000000}{9.600000}\selectfont\catcode`\^=\active\def^{\ifmmode\sp\else\^{}\fi}\catcode`\%=\active\def%{\%}(stagger delay)}}%
\end{pgfscope}%
\begin{pgfscope}%
\pgfsetbuttcap%
\pgfsetmiterjoin%
\definecolor{currentfill}{rgb}{1.000000,1.000000,1.000000}%
\pgfsetfillcolor{currentfill}%
\pgfsetfillopacity{0.800000}%
\pgfsetlinewidth{1.003750pt}%
\definecolor{currentstroke}{rgb}{0.800000,0.800000,0.800000}%
\pgfsetstrokecolor{currentstroke}%
\pgfsetstrokeopacity{0.800000}%
\pgfsetdash{}{0pt}%
\pgfpathmoveto{\pgfqpoint{0.700695in}{1.460497in}}%
\pgfpathlineto{\pgfqpoint{2.058828in}{1.460497in}}%
\pgfpathquadraticcurveto{\pgfqpoint{2.083828in}{1.460497in}}{\pgfqpoint{2.083828in}{1.485497in}}%
\pgfpathlineto{\pgfqpoint{2.083828in}{2.222997in}}%
\pgfpathquadraticcurveto{\pgfqpoint{2.083828in}{2.247997in}}{\pgfqpoint{2.058828in}{2.247997in}}%
\pgfpathlineto{\pgfqpoint{0.700695in}{2.247997in}}%
\pgfpathquadraticcurveto{\pgfqpoint{0.675695in}{2.247997in}}{\pgfqpoint{0.675695in}{2.222997in}}%
\pgfpathlineto{\pgfqpoint{0.675695in}{1.485497in}}%
\pgfpathquadraticcurveto{\pgfqpoint{0.675695in}{1.460497in}}{\pgfqpoint{0.700695in}{1.460497in}}%
\pgfpathlineto{\pgfqpoint{0.700695in}{1.460497in}}%
\pgfpathclose%
\pgfusepath{stroke,fill}%
\end{pgfscope}%
\begin{pgfscope}%
\pgfsetrectcap%
\pgfsetroundjoin%
\pgfsetlinewidth{1.505625pt}%
\definecolor{currentstroke}{rgb}{0.082353,0.396078,0.752941}%
\pgfsetstrokecolor{currentstroke}%
\pgfsetdash{}{0pt}%
\pgfpathmoveto{\pgfqpoint{0.725695in}{2.147997in}}%
\pgfpathlineto{\pgfqpoint{0.850695in}{2.147997in}}%
\pgfpathlineto{\pgfqpoint{0.975695in}{2.147997in}}%
\pgfusepath{stroke}%
\end{pgfscope}%
\begin{pgfscope}%
\pgfsetbuttcap%
\pgfsetmiterjoin%
\definecolor{currentfill}{rgb}{0.082353,0.396078,0.752941}%
\pgfsetfillcolor{currentfill}%
\pgfsetlinewidth{1.003750pt}%
\definecolor{currentstroke}{rgb}{0.082353,0.396078,0.752941}%
\pgfsetstrokecolor{currentstroke}%
\pgfsetdash{}{0pt}%
\pgfsys@defobject{currentmarker}{\pgfqpoint{-0.027778in}{-0.027778in}}{\pgfqpoint{0.027778in}{0.027778in}}{%
\pgfpathmoveto{\pgfqpoint{-0.027778in}{-0.027778in}}%
\pgfpathlineto{\pgfqpoint{0.027778in}{-0.027778in}}%
\pgfpathlineto{\pgfqpoint{0.027778in}{0.027778in}}%
\pgfpathlineto{\pgfqpoint{-0.027778in}{0.027778in}}%
\pgfpathlineto{\pgfqpoint{-0.027778in}{-0.027778in}}%
\pgfpathclose%
\pgfusepath{stroke,fill}%
}%
\begin{pgfscope}%
\pgfsys@transformshift{0.850695in}{2.147997in}%
\pgfsys@useobject{currentmarker}{}%
\end{pgfscope}%
\end{pgfscope}%
\begin{pgfscope}%
\definecolor{textcolor}{rgb}{0.000000,0.000000,0.000000}%
\pgfsetstrokecolor{textcolor}%
\pgfsetfillcolor{textcolor}%
\pgftext[x=1.075695in,y=2.104247in,left,base]{\color{textcolor}{\rmfamily\fontsize{9.000000}{10.800000}\selectfont\catcode`\^=\active\def^{\ifmmode\sp\else\^{}\fi}\catcode`\%=\active\def%{\%}TCP (Native)}}%
\end{pgfscope}%
\begin{pgfscope}%
\pgfsetrectcap%
\pgfsetroundjoin%
\pgfsetlinewidth{1.505625pt}%
\definecolor{currentstroke}{rgb}{0.180392,0.490196,0.196078}%
\pgfsetstrokecolor{currentstroke}%
\pgfsetdash{}{0pt}%
\pgfpathmoveto{\pgfqpoint{0.725695in}{1.960497in}}%
\pgfpathlineto{\pgfqpoint{0.850695in}{1.960497in}}%
\pgfpathlineto{\pgfqpoint{0.975695in}{1.960497in}}%
\pgfusepath{stroke}%
\end{pgfscope}%
\begin{pgfscope}%
\pgfsetbuttcap%
\pgfsetmiterjoin%
\definecolor{currentfill}{rgb}{0.180392,0.490196,0.196078}%
\pgfsetfillcolor{currentfill}%
\pgfsetlinewidth{1.003750pt}%
\definecolor{currentstroke}{rgb}{0.180392,0.490196,0.196078}%
\pgfsetstrokecolor{currentstroke}%
\pgfsetdash{}{0pt}%
\pgfsys@defobject{currentmarker}{\pgfqpoint{-0.027778in}{-0.027778in}}{\pgfqpoint{0.027778in}{0.027778in}}{%
\pgfpathmoveto{\pgfqpoint{-0.027778in}{-0.027778in}}%
\pgfpathlineto{\pgfqpoint{0.027778in}{-0.027778in}}%
\pgfpathlineto{\pgfqpoint{0.027778in}{0.027778in}}%
\pgfpathlineto{\pgfqpoint{-0.027778in}{0.027778in}}%
\pgfpathlineto{\pgfqpoint{-0.027778in}{-0.027778in}}%
\pgfpathclose%
\pgfusepath{stroke,fill}%
}%
\begin{pgfscope}%
\pgfsys@transformshift{0.850695in}{1.960497in}%
\pgfsys@useobject{currentmarker}{}%
\end{pgfscope}%
\end{pgfscope}%
\begin{pgfscope}%
\definecolor{textcolor}{rgb}{0.000000,0.000000,0.000000}%
\pgfsetstrokecolor{textcolor}%
\pgfsetfillcolor{textcolor}%
\pgftext[x=1.075695in,y=1.916747in,left,base]{\color{textcolor}{\rmfamily\fontsize{9.000000}{10.800000}\selectfont\catcode`\^=\active\def^{\ifmmode\sp\else\^{}\fi}\catcode`\%=\active\def%{\%}QUIC (Picoquic)}}%
\end{pgfscope}%
\begin{pgfscope}%
\pgfsetbuttcap%
\pgfsetroundjoin%
\pgfsetlinewidth{1.505625pt}%
\definecolor{currentstroke}{rgb}{0.647059,0.839216,0.654902}%
\pgfsetstrokecolor{currentstroke}%
\pgfsetdash{{5.550000pt}{2.400000pt}}{0.000000pt}%
\pgfpathmoveto{\pgfqpoint{0.725695in}{1.772997in}}%
\pgfpathlineto{\pgfqpoint{0.850695in}{1.772997in}}%
\pgfpathlineto{\pgfqpoint{0.975695in}{1.772997in}}%
\pgfusepath{stroke}%
\end{pgfscope}%
\begin{pgfscope}%
\pgfsetbuttcap%
\pgfsetmiterjoin%
\definecolor{currentfill}{rgb}{0.647059,0.839216,0.654902}%
\pgfsetfillcolor{currentfill}%
\pgfsetlinewidth{1.003750pt}%
\definecolor{currentstroke}{rgb}{0.647059,0.839216,0.654902}%
\pgfsetstrokecolor{currentstroke}%
\pgfsetdash{}{0pt}%
\pgfsys@defobject{currentmarker}{\pgfqpoint{-0.027778in}{-0.027778in}}{\pgfqpoint{0.027778in}{0.027778in}}{%
\pgfpathmoveto{\pgfqpoint{-0.027778in}{-0.027778in}}%
\pgfpathlineto{\pgfqpoint{0.027778in}{-0.027778in}}%
\pgfpathlineto{\pgfqpoint{0.027778in}{0.027778in}}%
\pgfpathlineto{\pgfqpoint{-0.027778in}{0.027778in}}%
\pgfpathlineto{\pgfqpoint{-0.027778in}{-0.027778in}}%
\pgfpathclose%
\pgfusepath{stroke,fill}%
}%
\begin{pgfscope}%
\pgfsys@transformshift{0.850695in}{1.772997in}%
\pgfsys@useobject{currentmarker}{}%
\end{pgfscope}%
\end{pgfscope}%
\begin{pgfscope}%
\definecolor{textcolor}{rgb}{0.000000,0.000000,0.000000}%
\pgfsetstrokecolor{textcolor}%
\pgfsetfillcolor{textcolor}%
\pgftext[x=1.075695in,y=1.729247in,left,base]{\color{textcolor}{\rmfamily\fontsize{9.000000}{10.800000}\selectfont\catcode`\^=\active\def^{\ifmmode\sp\else\^{}\fi}\catcode`\%=\active\def%{\%}QUIC (CTaps)}}%
\end{pgfscope}%
\begin{pgfscope}%
\pgfsetbuttcap%
\pgfsetroundjoin%
\pgfsetlinewidth{1.505625pt}%
\definecolor{currentstroke}{rgb}{0.564706,0.792157,0.976471}%
\pgfsetstrokecolor{currentstroke}%
\pgfsetdash{{5.550000pt}{2.400000pt}}{0.000000pt}%
\pgfpathmoveto{\pgfqpoint{0.725695in}{1.585497in}}%
\pgfpathlineto{\pgfqpoint{0.850695in}{1.585497in}}%
\pgfpathlineto{\pgfqpoint{0.975695in}{1.585497in}}%
\pgfusepath{stroke}%
\end{pgfscope}%
\begin{pgfscope}%
\pgfsetbuttcap%
\pgfsetmiterjoin%
\definecolor{currentfill}{rgb}{0.564706,0.792157,0.976471}%
\pgfsetfillcolor{currentfill}%
\pgfsetlinewidth{1.003750pt}%
\definecolor{currentstroke}{rgb}{0.564706,0.792157,0.976471}%
\pgfsetstrokecolor{currentstroke}%
\pgfsetdash{}{0pt}%
\pgfsys@defobject{currentmarker}{\pgfqpoint{-0.027778in}{-0.027778in}}{\pgfqpoint{0.027778in}{0.027778in}}{%
\pgfpathmoveto{\pgfqpoint{-0.027778in}{-0.027778in}}%
\pgfpathlineto{\pgfqpoint{0.027778in}{-0.027778in}}%
\pgfpathlineto{\pgfqpoint{0.027778in}{0.027778in}}%
\pgfpathlineto{\pgfqpoint{-0.027778in}{0.027778in}}%
\pgfpathlineto{\pgfqpoint{-0.027778in}{-0.027778in}}%
\pgfpathclose%
\pgfusepath{stroke,fill}%
}%
\begin{pgfscope}%
\pgfsys@transformshift{0.850695in}{1.585497in}%
\pgfsys@useobject{currentmarker}{}%
\end{pgfscope}%
\end{pgfscope}%
\begin{pgfscope}%
\definecolor{textcolor}{rgb}{0.000000,0.000000,0.000000}%
\pgfsetstrokecolor{textcolor}%
\pgfsetfillcolor{textcolor}%
\pgftext[x=1.075695in,y=1.541747in,left,base]{\color{textcolor}{\rmfamily\fontsize{9.000000}{10.800000}\selectfont\catcode`\^=\active\def^{\ifmmode\sp\else\^{}\fi}\catcode`\%=\active\def%{\%}TCP (CTaps)}}%
\end{pgfscope}%
\end{pgfpicture}%
\makeatother%
\endgroup%

    \caption[Handshake time experiment]{Measured handshake times for TCP, Picoquic and CTaps clients, ordered by RTT.}
    \label{fig:racing-handshake}
\end{figure}

\subsection{UDP RTT Experiment}\label{sec:udp-rtt-ex}
After performing the experiment in \cref{sec:single-endpoint-handshake} it was of
interest to understand exactly how much overhead CTaps adds to RTT values in the ideal case.
The experiment in \cref{fig:udp-overhead} was performed using UDP over localhost without
simulating any network conditions. A small packet of \qty{64}{\Bytes} was transferred
between a UDP client and a UDP server, measuring the round-trip time for each echo.
The native UDP implementation shows the baseline for the Linux networking stack,
and any added overhead will be due to CTaps.

\begin{figure}[H]
    %% Creator: Matplotlib, PGF backend
%%
%% To include the figure in your LaTeX document, write
%%   \input{<filename>.pgf}
%%
%% Make sure the required packages are loaded in your preamble
%%   \usepackage{pgf}
%%
%% Also ensure that all the required font packages are loaded; for instance,
%% the lmodern package is sometimes necessary when using math font.
%%   \usepackage{lmodern}
%%
%% Figures using additional raster images can only be included by \input if
%% they are in the same directory as the main LaTeX file. For loading figures
%% from other directories you can use the `import` package
%%   \usepackage{import}
%%
%% and then include the figures with
%%   \import{<path to file>}{<filename>.pgf}
%%
%% Matplotlib used the following preamble
%%   \def\mathdefault#1{#1}
%%   \everymath=\expandafter{\the\everymath\displaystyle}
%%   \IfFileExists{scrextend.sty}{
%%     \usepackage[fontsize=10.000000pt]{scrextend}
%%   }{
%%     \renewcommand{\normalsize}{\fontsize{10.000000}{12.000000}\selectfont}
%%     \normalsize
%%   }
%%   \usepackage{siunitx}
%%   \makeatletter\@ifpackageloaded{underscore}{}{\usepackage[strings]{underscore}}\makeatother
%%
\begingroup%
\makeatletter%
\begin{pgfpicture}%
\pgfpathrectangle{\pgfpointorigin}{\pgfqpoint{5.369125in}{3.280176in}}%
\pgfusepath{use as bounding box, clip}%
\begin{pgfscope}%
\pgfsetbuttcap%
\pgfsetmiterjoin%
\definecolor{currentfill}{rgb}{1.000000,1.000000,1.000000}%
\pgfsetfillcolor{currentfill}%
\pgfsetlinewidth{0.000000pt}%
\definecolor{currentstroke}{rgb}{1.000000,1.000000,1.000000}%
\pgfsetstrokecolor{currentstroke}%
\pgfsetdash{}{0pt}%
\pgfpathmoveto{\pgfqpoint{0.000000in}{0.000000in}}%
\pgfpathlineto{\pgfqpoint{5.369125in}{0.000000in}}%
\pgfpathlineto{\pgfqpoint{5.369125in}{3.280176in}}%
\pgfpathlineto{\pgfqpoint{0.000000in}{3.280176in}}%
\pgfpathlineto{\pgfqpoint{0.000000in}{0.000000in}}%
\pgfpathclose%
\pgfusepath{fill}%
\end{pgfscope}%
\begin{pgfscope}%
\pgfsetbuttcap%
\pgfsetmiterjoin%
\definecolor{currentfill}{rgb}{1.000000,1.000000,1.000000}%
\pgfsetfillcolor{currentfill}%
\pgfsetlinewidth{0.000000pt}%
\definecolor{currentstroke}{rgb}{0.000000,0.000000,0.000000}%
\pgfsetstrokecolor{currentstroke}%
\pgfsetstrokeopacity{0.000000}%
\pgfsetdash{}{0pt}%
\pgfpathmoveto{\pgfqpoint{0.565433in}{0.528318in}}%
\pgfpathlineto{\pgfqpoint{5.269125in}{0.528318in}}%
\pgfpathlineto{\pgfqpoint{5.269125in}{3.131951in}}%
\pgfpathlineto{\pgfqpoint{0.565433in}{3.131951in}}%
\pgfpathlineto{\pgfqpoint{0.565433in}{0.528318in}}%
\pgfpathclose%
\pgfusepath{fill}%
\end{pgfscope}%
\begin{pgfscope}%
\pgfpathrectangle{\pgfqpoint{0.565433in}{0.528318in}}{\pgfqpoint{4.703693in}{2.603633in}}%
\pgfusepath{clip}%
\pgfsetbuttcap%
\pgfsetroundjoin%
\pgfsetlinewidth{0.803000pt}%
\definecolor{currentstroke}{rgb}{0.690196,0.690196,0.690196}%
\pgfsetstrokecolor{currentstroke}%
\pgfsetstrokeopacity{0.400000}%
\pgfsetdash{{2.960000pt}{1.280000pt}}{0.000000pt}%
\pgfpathmoveto{\pgfqpoint{1.053866in}{0.528318in}}%
\pgfpathlineto{\pgfqpoint{1.053866in}{3.131951in}}%
\pgfusepath{stroke}%
\end{pgfscope}%
\begin{pgfscope}%
\pgfsetbuttcap%
\pgfsetroundjoin%
\definecolor{currentfill}{rgb}{0.000000,0.000000,0.000000}%
\pgfsetfillcolor{currentfill}%
\pgfsetlinewidth{0.803000pt}%
\definecolor{currentstroke}{rgb}{0.000000,0.000000,0.000000}%
\pgfsetstrokecolor{currentstroke}%
\pgfsetdash{}{0pt}%
\pgfsys@defobject{currentmarker}{\pgfqpoint{0.000000in}{-0.048611in}}{\pgfqpoint{0.000000in}{0.000000in}}{%
\pgfpathmoveto{\pgfqpoint{0.000000in}{0.000000in}}%
\pgfpathlineto{\pgfqpoint{0.000000in}{-0.048611in}}%
\pgfusepath{stroke,fill}%
}%
\begin{pgfscope}%
\pgfsys@transformshift{1.053866in}{0.528318in}%
\pgfsys@useobject{currentmarker}{}%
\end{pgfscope}%
\end{pgfscope}%
\begin{pgfscope}%
\definecolor{textcolor}{rgb}{0.000000,0.000000,0.000000}%
\pgfsetstrokecolor{textcolor}%
\pgfsetfillcolor{textcolor}%
\pgftext[x=1.053866in,y=0.431095in,,top]{\color{textcolor}{\rmfamily\fontsize{10.000000}{12.000000}\selectfont\catcode`\^=\active\def^{\ifmmode\sp\else\^{}\fi}\catcode`\%=\active\def%{\%}$\mathdefault{20}$}}%
\end{pgfscope}%
\begin{pgfscope}%
\pgfpathrectangle{\pgfqpoint{0.565433in}{0.528318in}}{\pgfqpoint{4.703693in}{2.603633in}}%
\pgfusepath{clip}%
\pgfsetbuttcap%
\pgfsetroundjoin%
\pgfsetlinewidth{0.803000pt}%
\definecolor{currentstroke}{rgb}{0.690196,0.690196,0.690196}%
\pgfsetstrokecolor{currentstroke}%
\pgfsetstrokeopacity{0.400000}%
\pgfsetdash{{2.960000pt}{1.280000pt}}{0.000000pt}%
\pgfpathmoveto{\pgfqpoint{1.803687in}{0.528318in}}%
\pgfpathlineto{\pgfqpoint{1.803687in}{3.131951in}}%
\pgfusepath{stroke}%
\end{pgfscope}%
\begin{pgfscope}%
\pgfsetbuttcap%
\pgfsetroundjoin%
\definecolor{currentfill}{rgb}{0.000000,0.000000,0.000000}%
\pgfsetfillcolor{currentfill}%
\pgfsetlinewidth{0.803000pt}%
\definecolor{currentstroke}{rgb}{0.000000,0.000000,0.000000}%
\pgfsetstrokecolor{currentstroke}%
\pgfsetdash{}{0pt}%
\pgfsys@defobject{currentmarker}{\pgfqpoint{0.000000in}{-0.048611in}}{\pgfqpoint{0.000000in}{0.000000in}}{%
\pgfpathmoveto{\pgfqpoint{0.000000in}{0.000000in}}%
\pgfpathlineto{\pgfqpoint{0.000000in}{-0.048611in}}%
\pgfusepath{stroke,fill}%
}%
\begin{pgfscope}%
\pgfsys@transformshift{1.803687in}{0.528318in}%
\pgfsys@useobject{currentmarker}{}%
\end{pgfscope}%
\end{pgfscope}%
\begin{pgfscope}%
\definecolor{textcolor}{rgb}{0.000000,0.000000,0.000000}%
\pgfsetstrokecolor{textcolor}%
\pgfsetfillcolor{textcolor}%
\pgftext[x=1.803687in,y=0.431095in,,top]{\color{textcolor}{\rmfamily\fontsize{10.000000}{12.000000}\selectfont\catcode`\^=\active\def^{\ifmmode\sp\else\^{}\fi}\catcode`\%=\active\def%{\%}$\mathdefault{25}$}}%
\end{pgfscope}%
\begin{pgfscope}%
\pgfpathrectangle{\pgfqpoint{0.565433in}{0.528318in}}{\pgfqpoint{4.703693in}{2.603633in}}%
\pgfusepath{clip}%
\pgfsetbuttcap%
\pgfsetroundjoin%
\pgfsetlinewidth{0.803000pt}%
\definecolor{currentstroke}{rgb}{0.690196,0.690196,0.690196}%
\pgfsetstrokecolor{currentstroke}%
\pgfsetstrokeopacity{0.400000}%
\pgfsetdash{{2.960000pt}{1.280000pt}}{0.000000pt}%
\pgfpathmoveto{\pgfqpoint{2.553508in}{0.528318in}}%
\pgfpathlineto{\pgfqpoint{2.553508in}{3.131951in}}%
\pgfusepath{stroke}%
\end{pgfscope}%
\begin{pgfscope}%
\pgfsetbuttcap%
\pgfsetroundjoin%
\definecolor{currentfill}{rgb}{0.000000,0.000000,0.000000}%
\pgfsetfillcolor{currentfill}%
\pgfsetlinewidth{0.803000pt}%
\definecolor{currentstroke}{rgb}{0.000000,0.000000,0.000000}%
\pgfsetstrokecolor{currentstroke}%
\pgfsetdash{}{0pt}%
\pgfsys@defobject{currentmarker}{\pgfqpoint{0.000000in}{-0.048611in}}{\pgfqpoint{0.000000in}{0.000000in}}{%
\pgfpathmoveto{\pgfqpoint{0.000000in}{0.000000in}}%
\pgfpathlineto{\pgfqpoint{0.000000in}{-0.048611in}}%
\pgfusepath{stroke,fill}%
}%
\begin{pgfscope}%
\pgfsys@transformshift{2.553508in}{0.528318in}%
\pgfsys@useobject{currentmarker}{}%
\end{pgfscope}%
\end{pgfscope}%
\begin{pgfscope}%
\definecolor{textcolor}{rgb}{0.000000,0.000000,0.000000}%
\pgfsetstrokecolor{textcolor}%
\pgfsetfillcolor{textcolor}%
\pgftext[x=2.553508in,y=0.431095in,,top]{\color{textcolor}{\rmfamily\fontsize{10.000000}{12.000000}\selectfont\catcode`\^=\active\def^{\ifmmode\sp\else\^{}\fi}\catcode`\%=\active\def%{\%}$\mathdefault{30}$}}%
\end{pgfscope}%
\begin{pgfscope}%
\pgfpathrectangle{\pgfqpoint{0.565433in}{0.528318in}}{\pgfqpoint{4.703693in}{2.603633in}}%
\pgfusepath{clip}%
\pgfsetbuttcap%
\pgfsetroundjoin%
\pgfsetlinewidth{0.803000pt}%
\definecolor{currentstroke}{rgb}{0.690196,0.690196,0.690196}%
\pgfsetstrokecolor{currentstroke}%
\pgfsetstrokeopacity{0.400000}%
\pgfsetdash{{2.960000pt}{1.280000pt}}{0.000000pt}%
\pgfpathmoveto{\pgfqpoint{3.303330in}{0.528318in}}%
\pgfpathlineto{\pgfqpoint{3.303330in}{3.131951in}}%
\pgfusepath{stroke}%
\end{pgfscope}%
\begin{pgfscope}%
\pgfsetbuttcap%
\pgfsetroundjoin%
\definecolor{currentfill}{rgb}{0.000000,0.000000,0.000000}%
\pgfsetfillcolor{currentfill}%
\pgfsetlinewidth{0.803000pt}%
\definecolor{currentstroke}{rgb}{0.000000,0.000000,0.000000}%
\pgfsetstrokecolor{currentstroke}%
\pgfsetdash{}{0pt}%
\pgfsys@defobject{currentmarker}{\pgfqpoint{0.000000in}{-0.048611in}}{\pgfqpoint{0.000000in}{0.000000in}}{%
\pgfpathmoveto{\pgfqpoint{0.000000in}{0.000000in}}%
\pgfpathlineto{\pgfqpoint{0.000000in}{-0.048611in}}%
\pgfusepath{stroke,fill}%
}%
\begin{pgfscope}%
\pgfsys@transformshift{3.303330in}{0.528318in}%
\pgfsys@useobject{currentmarker}{}%
\end{pgfscope}%
\end{pgfscope}%
\begin{pgfscope}%
\definecolor{textcolor}{rgb}{0.000000,0.000000,0.000000}%
\pgfsetstrokecolor{textcolor}%
\pgfsetfillcolor{textcolor}%
\pgftext[x=3.303330in,y=0.431095in,,top]{\color{textcolor}{\rmfamily\fontsize{10.000000}{12.000000}\selectfont\catcode`\^=\active\def^{\ifmmode\sp\else\^{}\fi}\catcode`\%=\active\def%{\%}$\mathdefault{35}$}}%
\end{pgfscope}%
\begin{pgfscope}%
\pgfpathrectangle{\pgfqpoint{0.565433in}{0.528318in}}{\pgfqpoint{4.703693in}{2.603633in}}%
\pgfusepath{clip}%
\pgfsetbuttcap%
\pgfsetroundjoin%
\pgfsetlinewidth{0.803000pt}%
\definecolor{currentstroke}{rgb}{0.690196,0.690196,0.690196}%
\pgfsetstrokecolor{currentstroke}%
\pgfsetstrokeopacity{0.400000}%
\pgfsetdash{{2.960000pt}{1.280000pt}}{0.000000pt}%
\pgfpathmoveto{\pgfqpoint{4.053151in}{0.528318in}}%
\pgfpathlineto{\pgfqpoint{4.053151in}{3.131951in}}%
\pgfusepath{stroke}%
\end{pgfscope}%
\begin{pgfscope}%
\pgfsetbuttcap%
\pgfsetroundjoin%
\definecolor{currentfill}{rgb}{0.000000,0.000000,0.000000}%
\pgfsetfillcolor{currentfill}%
\pgfsetlinewidth{0.803000pt}%
\definecolor{currentstroke}{rgb}{0.000000,0.000000,0.000000}%
\pgfsetstrokecolor{currentstroke}%
\pgfsetdash{}{0pt}%
\pgfsys@defobject{currentmarker}{\pgfqpoint{0.000000in}{-0.048611in}}{\pgfqpoint{0.000000in}{0.000000in}}{%
\pgfpathmoveto{\pgfqpoint{0.000000in}{0.000000in}}%
\pgfpathlineto{\pgfqpoint{0.000000in}{-0.048611in}}%
\pgfusepath{stroke,fill}%
}%
\begin{pgfscope}%
\pgfsys@transformshift{4.053151in}{0.528318in}%
\pgfsys@useobject{currentmarker}{}%
\end{pgfscope}%
\end{pgfscope}%
\begin{pgfscope}%
\definecolor{textcolor}{rgb}{0.000000,0.000000,0.000000}%
\pgfsetstrokecolor{textcolor}%
\pgfsetfillcolor{textcolor}%
\pgftext[x=4.053151in,y=0.431095in,,top]{\color{textcolor}{\rmfamily\fontsize{10.000000}{12.000000}\selectfont\catcode`\^=\active\def^{\ifmmode\sp\else\^{}\fi}\catcode`\%=\active\def%{\%}$\mathdefault{40}$}}%
\end{pgfscope}%
\begin{pgfscope}%
\pgfpathrectangle{\pgfqpoint{0.565433in}{0.528318in}}{\pgfqpoint{4.703693in}{2.603633in}}%
\pgfusepath{clip}%
\pgfsetbuttcap%
\pgfsetroundjoin%
\pgfsetlinewidth{0.803000pt}%
\definecolor{currentstroke}{rgb}{0.690196,0.690196,0.690196}%
\pgfsetstrokecolor{currentstroke}%
\pgfsetstrokeopacity{0.400000}%
\pgfsetdash{{2.960000pt}{1.280000pt}}{0.000000pt}%
\pgfpathmoveto{\pgfqpoint{4.802972in}{0.528318in}}%
\pgfpathlineto{\pgfqpoint{4.802972in}{3.131951in}}%
\pgfusepath{stroke}%
\end{pgfscope}%
\begin{pgfscope}%
\pgfsetbuttcap%
\pgfsetroundjoin%
\definecolor{currentfill}{rgb}{0.000000,0.000000,0.000000}%
\pgfsetfillcolor{currentfill}%
\pgfsetlinewidth{0.803000pt}%
\definecolor{currentstroke}{rgb}{0.000000,0.000000,0.000000}%
\pgfsetstrokecolor{currentstroke}%
\pgfsetdash{}{0pt}%
\pgfsys@defobject{currentmarker}{\pgfqpoint{0.000000in}{-0.048611in}}{\pgfqpoint{0.000000in}{0.000000in}}{%
\pgfpathmoveto{\pgfqpoint{0.000000in}{0.000000in}}%
\pgfpathlineto{\pgfqpoint{0.000000in}{-0.048611in}}%
\pgfusepath{stroke,fill}%
}%
\begin{pgfscope}%
\pgfsys@transformshift{4.802972in}{0.528318in}%
\pgfsys@useobject{currentmarker}{}%
\end{pgfscope}%
\end{pgfscope}%
\begin{pgfscope}%
\definecolor{textcolor}{rgb}{0.000000,0.000000,0.000000}%
\pgfsetstrokecolor{textcolor}%
\pgfsetfillcolor{textcolor}%
\pgftext[x=4.802972in,y=0.431095in,,top]{\color{textcolor}{\rmfamily\fontsize{10.000000}{12.000000}\selectfont\catcode`\^=\active\def^{\ifmmode\sp\else\^{}\fi}\catcode`\%=\active\def%{\%}$\mathdefault{45}$}}%
\end{pgfscope}%
\begin{pgfscope}%
\definecolor{textcolor}{rgb}{0.000000,0.000000,0.000000}%
\pgfsetstrokecolor{textcolor}%
\pgfsetfillcolor{textcolor}%
\pgftext[x=2.917279in,y=0.252083in,,top]{\color{textcolor}{\rmfamily\fontsize{11.000000}{13.200000}\selectfont\catcode`\^=\active\def^{\ifmmode\sp\else\^{}\fi}\catcode`\%=\active\def%{\%}RTT (µs)}}%
\end{pgfscope}%
\begin{pgfscope}%
\pgfpathrectangle{\pgfqpoint{0.565433in}{0.528318in}}{\pgfqpoint{4.703693in}{2.603633in}}%
\pgfusepath{clip}%
\pgfsetbuttcap%
\pgfsetroundjoin%
\pgfsetlinewidth{0.803000pt}%
\definecolor{currentstroke}{rgb}{0.690196,0.690196,0.690196}%
\pgfsetstrokecolor{currentstroke}%
\pgfsetstrokeopacity{0.400000}%
\pgfsetdash{{2.960000pt}{1.280000pt}}{0.000000pt}%
\pgfpathmoveto{\pgfqpoint{0.565433in}{0.528318in}}%
\pgfpathlineto{\pgfqpoint{5.269125in}{0.528318in}}%
\pgfusepath{stroke}%
\end{pgfscope}%
\begin{pgfscope}%
\pgfsetbuttcap%
\pgfsetroundjoin%
\definecolor{currentfill}{rgb}{0.000000,0.000000,0.000000}%
\pgfsetfillcolor{currentfill}%
\pgfsetlinewidth{0.803000pt}%
\definecolor{currentstroke}{rgb}{0.000000,0.000000,0.000000}%
\pgfsetstrokecolor{currentstroke}%
\pgfsetdash{}{0pt}%
\pgfsys@defobject{currentmarker}{\pgfqpoint{-0.048611in}{0.000000in}}{\pgfqpoint{-0.000000in}{0.000000in}}{%
\pgfpathmoveto{\pgfqpoint{-0.000000in}{0.000000in}}%
\pgfpathlineto{\pgfqpoint{-0.048611in}{0.000000in}}%
\pgfusepath{stroke,fill}%
}%
\begin{pgfscope}%
\pgfsys@transformshift{0.565433in}{0.528318in}%
\pgfsys@useobject{currentmarker}{}%
\end{pgfscope}%
\end{pgfscope}%
\begin{pgfscope}%
\definecolor{textcolor}{rgb}{0.000000,0.000000,0.000000}%
\pgfsetstrokecolor{textcolor}%
\pgfsetfillcolor{textcolor}%
\pgftext[x=0.290741in, y=0.480092in, left, base]{\color{textcolor}{\rmfamily\fontsize{10.000000}{12.000000}\selectfont\catcode`\^=\active\def^{\ifmmode\sp\else\^{}\fi}\catcode`\%=\active\def%{\%}$\mathdefault{0.0}$}}%
\end{pgfscope}%
\begin{pgfscope}%
\pgfpathrectangle{\pgfqpoint{0.565433in}{0.528318in}}{\pgfqpoint{4.703693in}{2.603633in}}%
\pgfusepath{clip}%
\pgfsetbuttcap%
\pgfsetroundjoin%
\pgfsetlinewidth{0.803000pt}%
\definecolor{currentstroke}{rgb}{0.690196,0.690196,0.690196}%
\pgfsetstrokecolor{currentstroke}%
\pgfsetstrokeopacity{0.400000}%
\pgfsetdash{{2.960000pt}{1.280000pt}}{0.000000pt}%
\pgfpathmoveto{\pgfqpoint{0.565433in}{1.049044in}}%
\pgfpathlineto{\pgfqpoint{5.269125in}{1.049044in}}%
\pgfusepath{stroke}%
\end{pgfscope}%
\begin{pgfscope}%
\pgfsetbuttcap%
\pgfsetroundjoin%
\definecolor{currentfill}{rgb}{0.000000,0.000000,0.000000}%
\pgfsetfillcolor{currentfill}%
\pgfsetlinewidth{0.803000pt}%
\definecolor{currentstroke}{rgb}{0.000000,0.000000,0.000000}%
\pgfsetstrokecolor{currentstroke}%
\pgfsetdash{}{0pt}%
\pgfsys@defobject{currentmarker}{\pgfqpoint{-0.048611in}{0.000000in}}{\pgfqpoint{-0.000000in}{0.000000in}}{%
\pgfpathmoveto{\pgfqpoint{-0.000000in}{0.000000in}}%
\pgfpathlineto{\pgfqpoint{-0.048611in}{0.000000in}}%
\pgfusepath{stroke,fill}%
}%
\begin{pgfscope}%
\pgfsys@transformshift{0.565433in}{1.049044in}%
\pgfsys@useobject{currentmarker}{}%
\end{pgfscope}%
\end{pgfscope}%
\begin{pgfscope}%
\definecolor{textcolor}{rgb}{0.000000,0.000000,0.000000}%
\pgfsetstrokecolor{textcolor}%
\pgfsetfillcolor{textcolor}%
\pgftext[x=0.290741in, y=1.000819in, left, base]{\color{textcolor}{\rmfamily\fontsize{10.000000}{12.000000}\selectfont\catcode`\^=\active\def^{\ifmmode\sp\else\^{}\fi}\catcode`\%=\active\def%{\%}$\mathdefault{0.2}$}}%
\end{pgfscope}%
\begin{pgfscope}%
\pgfpathrectangle{\pgfqpoint{0.565433in}{0.528318in}}{\pgfqpoint{4.703693in}{2.603633in}}%
\pgfusepath{clip}%
\pgfsetbuttcap%
\pgfsetroundjoin%
\pgfsetlinewidth{0.803000pt}%
\definecolor{currentstroke}{rgb}{0.690196,0.690196,0.690196}%
\pgfsetstrokecolor{currentstroke}%
\pgfsetstrokeopacity{0.400000}%
\pgfsetdash{{2.960000pt}{1.280000pt}}{0.000000pt}%
\pgfpathmoveto{\pgfqpoint{0.565433in}{1.569771in}}%
\pgfpathlineto{\pgfqpoint{5.269125in}{1.569771in}}%
\pgfusepath{stroke}%
\end{pgfscope}%
\begin{pgfscope}%
\pgfsetbuttcap%
\pgfsetroundjoin%
\definecolor{currentfill}{rgb}{0.000000,0.000000,0.000000}%
\pgfsetfillcolor{currentfill}%
\pgfsetlinewidth{0.803000pt}%
\definecolor{currentstroke}{rgb}{0.000000,0.000000,0.000000}%
\pgfsetstrokecolor{currentstroke}%
\pgfsetdash{}{0pt}%
\pgfsys@defobject{currentmarker}{\pgfqpoint{-0.048611in}{0.000000in}}{\pgfqpoint{-0.000000in}{0.000000in}}{%
\pgfpathmoveto{\pgfqpoint{-0.000000in}{0.000000in}}%
\pgfpathlineto{\pgfqpoint{-0.048611in}{0.000000in}}%
\pgfusepath{stroke,fill}%
}%
\begin{pgfscope}%
\pgfsys@transformshift{0.565433in}{1.569771in}%
\pgfsys@useobject{currentmarker}{}%
\end{pgfscope}%
\end{pgfscope}%
\begin{pgfscope}%
\definecolor{textcolor}{rgb}{0.000000,0.000000,0.000000}%
\pgfsetstrokecolor{textcolor}%
\pgfsetfillcolor{textcolor}%
\pgftext[x=0.290741in, y=1.521546in, left, base]{\color{textcolor}{\rmfamily\fontsize{10.000000}{12.000000}\selectfont\catcode`\^=\active\def^{\ifmmode\sp\else\^{}\fi}\catcode`\%=\active\def%{\%}$\mathdefault{0.4}$}}%
\end{pgfscope}%
\begin{pgfscope}%
\pgfpathrectangle{\pgfqpoint{0.565433in}{0.528318in}}{\pgfqpoint{4.703693in}{2.603633in}}%
\pgfusepath{clip}%
\pgfsetbuttcap%
\pgfsetroundjoin%
\pgfsetlinewidth{0.803000pt}%
\definecolor{currentstroke}{rgb}{0.690196,0.690196,0.690196}%
\pgfsetstrokecolor{currentstroke}%
\pgfsetstrokeopacity{0.400000}%
\pgfsetdash{{2.960000pt}{1.280000pt}}{0.000000pt}%
\pgfpathmoveto{\pgfqpoint{0.565433in}{2.090498in}}%
\pgfpathlineto{\pgfqpoint{5.269125in}{2.090498in}}%
\pgfusepath{stroke}%
\end{pgfscope}%
\begin{pgfscope}%
\pgfsetbuttcap%
\pgfsetroundjoin%
\definecolor{currentfill}{rgb}{0.000000,0.000000,0.000000}%
\pgfsetfillcolor{currentfill}%
\pgfsetlinewidth{0.803000pt}%
\definecolor{currentstroke}{rgb}{0.000000,0.000000,0.000000}%
\pgfsetstrokecolor{currentstroke}%
\pgfsetdash{}{0pt}%
\pgfsys@defobject{currentmarker}{\pgfqpoint{-0.048611in}{0.000000in}}{\pgfqpoint{-0.000000in}{0.000000in}}{%
\pgfpathmoveto{\pgfqpoint{-0.000000in}{0.000000in}}%
\pgfpathlineto{\pgfqpoint{-0.048611in}{0.000000in}}%
\pgfusepath{stroke,fill}%
}%
\begin{pgfscope}%
\pgfsys@transformshift{0.565433in}{2.090498in}%
\pgfsys@useobject{currentmarker}{}%
\end{pgfscope}%
\end{pgfscope}%
\begin{pgfscope}%
\definecolor{textcolor}{rgb}{0.000000,0.000000,0.000000}%
\pgfsetstrokecolor{textcolor}%
\pgfsetfillcolor{textcolor}%
\pgftext[x=0.290741in, y=2.042272in, left, base]{\color{textcolor}{\rmfamily\fontsize{10.000000}{12.000000}\selectfont\catcode`\^=\active\def^{\ifmmode\sp\else\^{}\fi}\catcode`\%=\active\def%{\%}$\mathdefault{0.6}$}}%
\end{pgfscope}%
\begin{pgfscope}%
\pgfpathrectangle{\pgfqpoint{0.565433in}{0.528318in}}{\pgfqpoint{4.703693in}{2.603633in}}%
\pgfusepath{clip}%
\pgfsetbuttcap%
\pgfsetroundjoin%
\pgfsetlinewidth{0.803000pt}%
\definecolor{currentstroke}{rgb}{0.690196,0.690196,0.690196}%
\pgfsetstrokecolor{currentstroke}%
\pgfsetstrokeopacity{0.400000}%
\pgfsetdash{{2.960000pt}{1.280000pt}}{0.000000pt}%
\pgfpathmoveto{\pgfqpoint{0.565433in}{2.611224in}}%
\pgfpathlineto{\pgfqpoint{5.269125in}{2.611224in}}%
\pgfusepath{stroke}%
\end{pgfscope}%
\begin{pgfscope}%
\pgfsetbuttcap%
\pgfsetroundjoin%
\definecolor{currentfill}{rgb}{0.000000,0.000000,0.000000}%
\pgfsetfillcolor{currentfill}%
\pgfsetlinewidth{0.803000pt}%
\definecolor{currentstroke}{rgb}{0.000000,0.000000,0.000000}%
\pgfsetstrokecolor{currentstroke}%
\pgfsetdash{}{0pt}%
\pgfsys@defobject{currentmarker}{\pgfqpoint{-0.048611in}{0.000000in}}{\pgfqpoint{-0.000000in}{0.000000in}}{%
\pgfpathmoveto{\pgfqpoint{-0.000000in}{0.000000in}}%
\pgfpathlineto{\pgfqpoint{-0.048611in}{0.000000in}}%
\pgfusepath{stroke,fill}%
}%
\begin{pgfscope}%
\pgfsys@transformshift{0.565433in}{2.611224in}%
\pgfsys@useobject{currentmarker}{}%
\end{pgfscope}%
\end{pgfscope}%
\begin{pgfscope}%
\definecolor{textcolor}{rgb}{0.000000,0.000000,0.000000}%
\pgfsetstrokecolor{textcolor}%
\pgfsetfillcolor{textcolor}%
\pgftext[x=0.290741in, y=2.562999in, left, base]{\color{textcolor}{\rmfamily\fontsize{10.000000}{12.000000}\selectfont\catcode`\^=\active\def^{\ifmmode\sp\else\^{}\fi}\catcode`\%=\active\def%{\%}$\mathdefault{0.8}$}}%
\end{pgfscope}%
\begin{pgfscope}%
\pgfpathrectangle{\pgfqpoint{0.565433in}{0.528318in}}{\pgfqpoint{4.703693in}{2.603633in}}%
\pgfusepath{clip}%
\pgfsetbuttcap%
\pgfsetroundjoin%
\pgfsetlinewidth{0.803000pt}%
\definecolor{currentstroke}{rgb}{0.690196,0.690196,0.690196}%
\pgfsetstrokecolor{currentstroke}%
\pgfsetstrokeopacity{0.400000}%
\pgfsetdash{{2.960000pt}{1.280000pt}}{0.000000pt}%
\pgfpathmoveto{\pgfqpoint{0.565433in}{3.131951in}}%
\pgfpathlineto{\pgfqpoint{5.269125in}{3.131951in}}%
\pgfusepath{stroke}%
\end{pgfscope}%
\begin{pgfscope}%
\pgfsetbuttcap%
\pgfsetroundjoin%
\definecolor{currentfill}{rgb}{0.000000,0.000000,0.000000}%
\pgfsetfillcolor{currentfill}%
\pgfsetlinewidth{0.803000pt}%
\definecolor{currentstroke}{rgb}{0.000000,0.000000,0.000000}%
\pgfsetstrokecolor{currentstroke}%
\pgfsetdash{}{0pt}%
\pgfsys@defobject{currentmarker}{\pgfqpoint{-0.048611in}{0.000000in}}{\pgfqpoint{-0.000000in}{0.000000in}}{%
\pgfpathmoveto{\pgfqpoint{-0.000000in}{0.000000in}}%
\pgfpathlineto{\pgfqpoint{-0.048611in}{0.000000in}}%
\pgfusepath{stroke,fill}%
}%
\begin{pgfscope}%
\pgfsys@transformshift{0.565433in}{3.131951in}%
\pgfsys@useobject{currentmarker}{}%
\end{pgfscope}%
\end{pgfscope}%
\begin{pgfscope}%
\definecolor{textcolor}{rgb}{0.000000,0.000000,0.000000}%
\pgfsetstrokecolor{textcolor}%
\pgfsetfillcolor{textcolor}%
\pgftext[x=0.290741in, y=3.083726in, left, base]{\color{textcolor}{\rmfamily\fontsize{10.000000}{12.000000}\selectfont\catcode`\^=\active\def^{\ifmmode\sp\else\^{}\fi}\catcode`\%=\active\def%{\%}$\mathdefault{1.0}$}}%
\end{pgfscope}%
\begin{pgfscope}%
\definecolor{textcolor}{rgb}{0.000000,0.000000,0.000000}%
\pgfsetstrokecolor{textcolor}%
\pgfsetfillcolor{textcolor}%
\pgftext[x=0.235185in,y=1.830134in,,bottom,rotate=90.000000]{\color{textcolor}{\rmfamily\fontsize{11.000000}{13.200000}\selectfont\catcode`\^=\active\def^{\ifmmode\sp\else\^{}\fi}\catcode`\%=\active\def%{\%}Percentile}}%
\end{pgfscope}%
\begin{pgfscope}%
\pgfpathrectangle{\pgfqpoint{0.565433in}{0.528318in}}{\pgfqpoint{4.703693in}{2.603633in}}%
\pgfusepath{clip}%
\pgfsetrectcap%
\pgfsetroundjoin%
\pgfsetlinewidth{1.806750pt}%
\definecolor{currentstroke}{rgb}{0.082353,0.396078,0.752941}%
\pgfsetstrokecolor{currentstroke}%
\pgfsetdash{}{0pt}%
\pgfpathmoveto{\pgfqpoint{0.565433in}{0.528344in}}%
\pgfpathlineto{\pgfqpoint{0.587627in}{0.529151in}}%
\pgfpathlineto{\pgfqpoint{0.595875in}{0.530609in}}%
\pgfpathlineto{\pgfqpoint{0.603674in}{0.532692in}}%
\pgfpathlineto{\pgfqpoint{0.604273in}{0.532926in}}%
\pgfpathlineto{\pgfqpoint{0.609222in}{0.534514in}}%
\pgfpathlineto{\pgfqpoint{0.609822in}{0.534853in}}%
\pgfpathlineto{\pgfqpoint{0.614321in}{0.536363in}}%
\pgfpathlineto{\pgfqpoint{0.614921in}{0.536623in}}%
\pgfpathlineto{\pgfqpoint{0.619120in}{0.538081in}}%
\pgfpathlineto{\pgfqpoint{0.619720in}{0.538237in}}%
\pgfpathlineto{\pgfqpoint{0.623169in}{0.539279in}}%
\pgfpathlineto{\pgfqpoint{0.627818in}{0.540815in}}%
\pgfpathlineto{\pgfqpoint{0.628418in}{0.540971in}}%
\pgfpathlineto{\pgfqpoint{0.649863in}{0.543992in}}%
\pgfpathlineto{\pgfqpoint{0.681205in}{0.546543in}}%
\pgfpathlineto{\pgfqpoint{0.687953in}{0.548444in}}%
\pgfpathlineto{\pgfqpoint{0.688553in}{0.548548in}}%
\pgfpathlineto{\pgfqpoint{0.695152in}{0.550422in}}%
\pgfpathlineto{\pgfqpoint{0.695752in}{0.550761in}}%
\pgfpathlineto{\pgfqpoint{0.698301in}{0.551933in}}%
\pgfpathlineto{\pgfqpoint{0.698901in}{0.552219in}}%
\pgfpathlineto{\pgfqpoint{0.700850in}{0.553286in}}%
\pgfpathlineto{\pgfqpoint{0.701450in}{0.553547in}}%
\pgfpathlineto{\pgfqpoint{0.703100in}{0.554562in}}%
\pgfpathlineto{\pgfqpoint{0.703550in}{0.555057in}}%
\pgfpathlineto{\pgfqpoint{0.704749in}{0.555890in}}%
\pgfpathlineto{\pgfqpoint{0.705199in}{0.556124in}}%
\pgfpathlineto{\pgfqpoint{0.706699in}{0.556984in}}%
\pgfpathlineto{\pgfqpoint{0.707299in}{0.557244in}}%
\pgfpathlineto{\pgfqpoint{0.709248in}{0.558311in}}%
\pgfpathlineto{\pgfqpoint{0.709848in}{0.558884in}}%
\pgfpathlineto{\pgfqpoint{0.713747in}{0.560290in}}%
\pgfpathlineto{\pgfqpoint{0.714347in}{0.560707in}}%
\pgfpathlineto{\pgfqpoint{0.717946in}{0.561592in}}%
\pgfpathlineto{\pgfqpoint{0.722895in}{0.563180in}}%
\pgfpathlineto{\pgfqpoint{0.723495in}{0.563441in}}%
\pgfpathlineto{\pgfqpoint{0.726944in}{0.564378in}}%
\pgfpathlineto{\pgfqpoint{0.727394in}{0.564534in}}%
\pgfpathlineto{\pgfqpoint{0.730993in}{0.565576in}}%
\pgfpathlineto{\pgfqpoint{0.731593in}{0.565810in}}%
\pgfpathlineto{\pgfqpoint{0.737442in}{0.567060in}}%
\pgfpathlineto{\pgfqpoint{0.738041in}{0.567216in}}%
\pgfpathlineto{\pgfqpoint{0.748089in}{0.569689in}}%
\pgfpathlineto{\pgfqpoint{0.748689in}{0.569924in}}%
\pgfpathlineto{\pgfqpoint{0.751988in}{0.571226in}}%
\pgfpathlineto{\pgfqpoint{0.752588in}{0.571460in}}%
\pgfpathlineto{\pgfqpoint{0.755587in}{0.572710in}}%
\pgfpathlineto{\pgfqpoint{0.756187in}{0.573048in}}%
\pgfpathlineto{\pgfqpoint{0.757837in}{0.574063in}}%
\pgfpathlineto{\pgfqpoint{0.758437in}{0.574350in}}%
\pgfpathlineto{\pgfqpoint{0.760836in}{0.575495in}}%
\pgfpathlineto{\pgfqpoint{0.761436in}{0.575964in}}%
\pgfpathlineto{\pgfqpoint{0.763985in}{0.577136in}}%
\pgfpathlineto{\pgfqpoint{0.764585in}{0.577578in}}%
\pgfpathlineto{\pgfqpoint{0.765935in}{0.578542in}}%
\pgfpathlineto{\pgfqpoint{0.766535in}{0.578906in}}%
\pgfpathlineto{\pgfqpoint{0.767284in}{0.579765in}}%
\pgfpathlineto{\pgfqpoint{0.767884in}{0.580156in}}%
\pgfpathlineto{\pgfqpoint{0.769834in}{0.581223in}}%
\pgfpathlineto{\pgfqpoint{0.770434in}{0.581718in}}%
\pgfpathlineto{\pgfqpoint{0.771783in}{0.582682in}}%
\pgfpathlineto{\pgfqpoint{0.772383in}{0.583228in}}%
\pgfpathlineto{\pgfqpoint{0.773583in}{0.584166in}}%
\pgfpathlineto{\pgfqpoint{0.774183in}{0.584686in}}%
\pgfpathlineto{\pgfqpoint{0.775533in}{0.585650in}}%
\pgfpathlineto{\pgfqpoint{0.776132in}{0.586457in}}%
\pgfpathlineto{\pgfqpoint{0.777332in}{0.587394in}}%
\pgfpathlineto{\pgfqpoint{0.777932in}{0.587967in}}%
\pgfpathlineto{\pgfqpoint{0.779881in}{0.589034in}}%
\pgfpathlineto{\pgfqpoint{0.780481in}{0.589503in}}%
\pgfpathlineto{\pgfqpoint{0.782581in}{0.590597in}}%
\pgfpathlineto{\pgfqpoint{0.783181in}{0.590857in}}%
\pgfpathlineto{\pgfqpoint{0.785130in}{0.591924in}}%
\pgfpathlineto{\pgfqpoint{0.785730in}{0.592289in}}%
\pgfpathlineto{\pgfqpoint{0.787380in}{0.593304in}}%
\pgfpathlineto{\pgfqpoint{0.787980in}{0.593643in}}%
\pgfpathlineto{\pgfqpoint{0.790229in}{0.594762in}}%
\pgfpathlineto{\pgfqpoint{0.790529in}{0.594945in}}%
\pgfpathlineto{\pgfqpoint{0.792778in}{0.595804in}}%
\pgfpathlineto{\pgfqpoint{0.796528in}{0.597184in}}%
\pgfpathlineto{\pgfqpoint{0.797127in}{0.597574in}}%
\pgfpathlineto{\pgfqpoint{0.801026in}{0.598980in}}%
\pgfpathlineto{\pgfqpoint{0.801626in}{0.599293in}}%
\pgfpathlineto{\pgfqpoint{0.804626in}{0.600542in}}%
\pgfpathlineto{\pgfqpoint{0.805225in}{0.600699in}}%
\pgfpathlineto{\pgfqpoint{0.808225in}{0.601688in}}%
\pgfpathlineto{\pgfqpoint{0.808825in}{0.602026in}}%
\pgfpathlineto{\pgfqpoint{0.811224in}{0.602964in}}%
\pgfpathlineto{\pgfqpoint{0.811824in}{0.603302in}}%
\pgfpathlineto{\pgfqpoint{0.813473in}{0.604188in}}%
\pgfpathlineto{\pgfqpoint{0.814073in}{0.604552in}}%
\pgfpathlineto{\pgfqpoint{0.815423in}{0.605515in}}%
\pgfpathlineto{\pgfqpoint{0.816023in}{0.606062in}}%
\pgfpathlineto{\pgfqpoint{0.817972in}{0.607130in}}%
\pgfpathlineto{\pgfqpoint{0.818572in}{0.607494in}}%
\pgfpathlineto{\pgfqpoint{0.819472in}{0.608379in}}%
\pgfpathlineto{\pgfqpoint{0.820072in}{0.608848in}}%
\pgfpathlineto{\pgfqpoint{0.821871in}{0.609889in}}%
\pgfpathlineto{\pgfqpoint{0.822471in}{0.610749in}}%
\pgfpathlineto{\pgfqpoint{0.823521in}{0.611660in}}%
\pgfpathlineto{\pgfqpoint{0.824121in}{0.612050in}}%
\pgfpathlineto{\pgfqpoint{0.825771in}{0.613066in}}%
\pgfpathlineto{\pgfqpoint{0.826370in}{0.613587in}}%
\pgfpathlineto{\pgfqpoint{0.827420in}{0.614498in}}%
\pgfpathlineto{\pgfqpoint{0.828020in}{0.615071in}}%
\pgfpathlineto{\pgfqpoint{0.829520in}{0.616060in}}%
\pgfpathlineto{\pgfqpoint{0.830120in}{0.616737in}}%
\pgfpathlineto{\pgfqpoint{0.831019in}{0.617622in}}%
\pgfpathlineto{\pgfqpoint{0.831619in}{0.618247in}}%
\pgfpathlineto{\pgfqpoint{0.832819in}{0.619184in}}%
\pgfpathlineto{\pgfqpoint{0.833419in}{0.619861in}}%
\pgfpathlineto{\pgfqpoint{0.834169in}{0.620721in}}%
\pgfpathlineto{\pgfqpoint{0.834768in}{0.621189in}}%
\pgfpathlineto{\pgfqpoint{0.835368in}{0.622022in}}%
\pgfpathlineto{\pgfqpoint{0.835968in}{0.622986in}}%
\pgfpathlineto{\pgfqpoint{0.836718in}{0.623845in}}%
\pgfpathlineto{\pgfqpoint{0.837318in}{0.624965in}}%
\pgfpathlineto{\pgfqpoint{0.838068in}{0.625824in}}%
\pgfpathlineto{\pgfqpoint{0.838667in}{0.627178in}}%
\pgfpathlineto{\pgfqpoint{0.839267in}{0.628011in}}%
\pgfpathlineto{\pgfqpoint{0.839867in}{0.629000in}}%
\pgfpathlineto{\pgfqpoint{0.840467in}{0.629833in}}%
\pgfpathlineto{\pgfqpoint{0.841067in}{0.630562in}}%
\pgfpathlineto{\pgfqpoint{0.841817in}{0.631422in}}%
\pgfpathlineto{\pgfqpoint{0.842417in}{0.632515in}}%
\pgfpathlineto{\pgfqpoint{0.843016in}{0.633348in}}%
\pgfpathlineto{\pgfqpoint{0.843616in}{0.634416in}}%
\pgfpathlineto{\pgfqpoint{0.844216in}{0.635249in}}%
\pgfpathlineto{\pgfqpoint{0.844816in}{0.636551in}}%
\pgfpathlineto{\pgfqpoint{0.845266in}{0.637358in}}%
\pgfpathlineto{\pgfqpoint{0.845866in}{0.638686in}}%
\pgfpathlineto{\pgfqpoint{0.846466in}{0.639519in}}%
\pgfpathlineto{\pgfqpoint{0.847065in}{0.640951in}}%
\pgfpathlineto{\pgfqpoint{0.847665in}{0.641784in}}%
\pgfpathlineto{\pgfqpoint{0.848265in}{0.643320in}}%
\pgfpathlineto{\pgfqpoint{0.848865in}{0.644153in}}%
\pgfpathlineto{\pgfqpoint{0.849465in}{0.645872in}}%
\pgfpathlineto{\pgfqpoint{0.849765in}{0.646653in}}%
\pgfpathlineto{\pgfqpoint{0.850365in}{0.647824in}}%
\pgfpathlineto{\pgfqpoint{0.850815in}{0.648632in}}%
\pgfpathlineto{\pgfqpoint{0.851414in}{0.649829in}}%
\pgfpathlineto{\pgfqpoint{0.851864in}{0.650636in}}%
\pgfpathlineto{\pgfqpoint{0.852464in}{0.652485in}}%
\pgfpathlineto{\pgfqpoint{0.852764in}{0.653266in}}%
\pgfpathlineto{\pgfqpoint{0.853364in}{0.654568in}}%
\pgfpathlineto{\pgfqpoint{0.853814in}{0.655375in}}%
\pgfpathlineto{\pgfqpoint{0.854414in}{0.656989in}}%
\pgfpathlineto{\pgfqpoint{0.854864in}{0.657796in}}%
\pgfpathlineto{\pgfqpoint{0.855463in}{0.660114in}}%
\pgfpathlineto{\pgfqpoint{0.855763in}{0.660895in}}%
\pgfpathlineto{\pgfqpoint{0.856363in}{0.662353in}}%
\pgfpathlineto{\pgfqpoint{0.856663in}{0.663134in}}%
\pgfpathlineto{\pgfqpoint{0.857263in}{0.665008in}}%
\pgfpathlineto{\pgfqpoint{0.857713in}{0.665816in}}%
\pgfpathlineto{\pgfqpoint{0.858313in}{0.667716in}}%
\pgfpathlineto{\pgfqpoint{0.858763in}{0.668523in}}%
\pgfpathlineto{\pgfqpoint{0.859363in}{0.669903in}}%
\pgfpathlineto{\pgfqpoint{0.859812in}{0.670710in}}%
\pgfpathlineto{\pgfqpoint{0.860412in}{0.672689in}}%
\pgfpathlineto{\pgfqpoint{0.860712in}{0.673470in}}%
\pgfpathlineto{\pgfqpoint{0.861312in}{0.675267in}}%
\pgfpathlineto{\pgfqpoint{0.861762in}{0.676074in}}%
\pgfpathlineto{\pgfqpoint{0.862362in}{0.678495in}}%
\pgfpathlineto{\pgfqpoint{0.862662in}{0.679276in}}%
\pgfpathlineto{\pgfqpoint{0.863262in}{0.680838in}}%
\pgfpathlineto{\pgfqpoint{0.863562in}{0.681620in}}%
\pgfpathlineto{\pgfqpoint{0.864161in}{0.683520in}}%
\pgfpathlineto{\pgfqpoint{0.864461in}{0.684301in}}%
\pgfpathlineto{\pgfqpoint{0.865061in}{0.685994in}}%
\pgfpathlineto{\pgfqpoint{0.865511in}{0.686801in}}%
\pgfpathlineto{\pgfqpoint{0.866111in}{0.688832in}}%
\pgfpathlineto{\pgfqpoint{0.866561in}{0.689639in}}%
\pgfpathlineto{\pgfqpoint{0.867161in}{0.691045in}}%
\pgfpathlineto{\pgfqpoint{0.867461in}{0.691826in}}%
\pgfpathlineto{\pgfqpoint{0.868060in}{0.693831in}}%
\pgfpathlineto{\pgfqpoint{0.868510in}{0.694638in}}%
\pgfpathlineto{\pgfqpoint{0.869110in}{0.696226in}}%
\pgfpathlineto{\pgfqpoint{0.869560in}{0.697033in}}%
\pgfpathlineto{\pgfqpoint{0.870160in}{0.698621in}}%
\pgfpathlineto{\pgfqpoint{0.870610in}{0.699428in}}%
\pgfpathlineto{\pgfqpoint{0.871210in}{0.701017in}}%
\pgfpathlineto{\pgfqpoint{0.871660in}{0.701824in}}%
\pgfpathlineto{\pgfqpoint{0.872259in}{0.703542in}}%
\pgfpathlineto{\pgfqpoint{0.872709in}{0.704349in}}%
\pgfpathlineto{\pgfqpoint{0.873309in}{0.706380in}}%
\pgfpathlineto{\pgfqpoint{0.873609in}{0.707161in}}%
\pgfpathlineto{\pgfqpoint{0.874209in}{0.708645in}}%
\pgfpathlineto{\pgfqpoint{0.874659in}{0.709452in}}%
\pgfpathlineto{\pgfqpoint{0.875259in}{0.711353in}}%
\pgfpathlineto{\pgfqpoint{0.875559in}{0.712134in}}%
\pgfpathlineto{\pgfqpoint{0.876159in}{0.713618in}}%
\pgfpathlineto{\pgfqpoint{0.876608in}{0.714425in}}%
\pgfpathlineto{\pgfqpoint{0.877208in}{0.715753in}}%
\pgfpathlineto{\pgfqpoint{0.877658in}{0.716560in}}%
\pgfpathlineto{\pgfqpoint{0.878258in}{0.717940in}}%
\pgfpathlineto{\pgfqpoint{0.878708in}{0.718747in}}%
\pgfpathlineto{\pgfqpoint{0.879308in}{0.720101in}}%
\pgfpathlineto{\pgfqpoint{0.879608in}{0.720882in}}%
\pgfpathlineto{\pgfqpoint{0.880208in}{0.722132in}}%
\pgfpathlineto{\pgfqpoint{0.880807in}{0.722965in}}%
\pgfpathlineto{\pgfqpoint{0.881407in}{0.724267in}}%
\pgfpathlineto{\pgfqpoint{0.882007in}{0.725100in}}%
\pgfpathlineto{\pgfqpoint{0.882607in}{0.726246in}}%
\pgfpathlineto{\pgfqpoint{0.883057in}{0.727053in}}%
\pgfpathlineto{\pgfqpoint{0.883657in}{0.728329in}}%
\pgfpathlineto{\pgfqpoint{0.884107in}{0.729136in}}%
\pgfpathlineto{\pgfqpoint{0.884706in}{0.730151in}}%
\pgfpathlineto{\pgfqpoint{0.885006in}{0.730932in}}%
\pgfpathlineto{\pgfqpoint{0.885606in}{0.731870in}}%
\pgfpathlineto{\pgfqpoint{0.886206in}{0.732703in}}%
\pgfpathlineto{\pgfqpoint{0.886806in}{0.733796in}}%
\pgfpathlineto{\pgfqpoint{0.887256in}{0.734604in}}%
\pgfpathlineto{\pgfqpoint{0.887856in}{0.735749in}}%
\pgfpathlineto{\pgfqpoint{0.888456in}{0.736582in}}%
\pgfpathlineto{\pgfqpoint{0.889055in}{0.738301in}}%
\pgfpathlineto{\pgfqpoint{0.889505in}{0.739108in}}%
\pgfpathlineto{\pgfqpoint{0.890105in}{0.740071in}}%
\pgfpathlineto{\pgfqpoint{0.890555in}{0.740878in}}%
\pgfpathlineto{\pgfqpoint{0.891155in}{0.742310in}}%
\pgfpathlineto{\pgfqpoint{0.891605in}{0.743117in}}%
\pgfpathlineto{\pgfqpoint{0.892205in}{0.744523in}}%
\pgfpathlineto{\pgfqpoint{0.892655in}{0.745330in}}%
\pgfpathlineto{\pgfqpoint{0.893254in}{0.746971in}}%
\pgfpathlineto{\pgfqpoint{0.893704in}{0.747778in}}%
\pgfpathlineto{\pgfqpoint{0.894304in}{0.748976in}}%
\pgfpathlineto{\pgfqpoint{0.894754in}{0.749783in}}%
\pgfpathlineto{\pgfqpoint{0.895354in}{0.751058in}}%
\pgfpathlineto{\pgfqpoint{0.895954in}{0.751892in}}%
\pgfpathlineto{\pgfqpoint{0.896554in}{0.753324in}}%
\pgfpathlineto{\pgfqpoint{0.897154in}{0.754157in}}%
\pgfpathlineto{\pgfqpoint{0.897753in}{0.755693in}}%
\pgfpathlineto{\pgfqpoint{0.898203in}{0.756500in}}%
\pgfpathlineto{\pgfqpoint{0.898803in}{0.758010in}}%
\pgfpathlineto{\pgfqpoint{0.899253in}{0.758817in}}%
\pgfpathlineto{\pgfqpoint{0.899853in}{0.760484in}}%
\pgfpathlineto{\pgfqpoint{0.900303in}{0.761291in}}%
\pgfpathlineto{\pgfqpoint{0.900903in}{0.762488in}}%
\pgfpathlineto{\pgfqpoint{0.901353in}{0.763296in}}%
\pgfpathlineto{\pgfqpoint{0.901952in}{0.765014in}}%
\pgfpathlineto{\pgfqpoint{0.902402in}{0.765821in}}%
\pgfpathlineto{\pgfqpoint{0.903002in}{0.767748in}}%
\pgfpathlineto{\pgfqpoint{0.903452in}{0.768555in}}%
\pgfpathlineto{\pgfqpoint{0.904052in}{0.770456in}}%
\pgfpathlineto{\pgfqpoint{0.904502in}{0.771263in}}%
\pgfpathlineto{\pgfqpoint{0.905102in}{0.773137in}}%
\pgfpathlineto{\pgfqpoint{0.905552in}{0.773944in}}%
\pgfpathlineto{\pgfqpoint{0.906151in}{0.775663in}}%
\pgfpathlineto{\pgfqpoint{0.906601in}{0.776470in}}%
\pgfpathlineto{\pgfqpoint{0.907201in}{0.778683in}}%
\pgfpathlineto{\pgfqpoint{0.907651in}{0.779490in}}%
\pgfpathlineto{\pgfqpoint{0.908251in}{0.780922in}}%
\pgfpathlineto{\pgfqpoint{0.908551in}{0.781703in}}%
\pgfpathlineto{\pgfqpoint{0.909151in}{0.783526in}}%
\pgfpathlineto{\pgfqpoint{0.909451in}{0.784307in}}%
\pgfpathlineto{\pgfqpoint{0.910050in}{0.785999in}}%
\pgfpathlineto{\pgfqpoint{0.910350in}{0.786780in}}%
\pgfpathlineto{\pgfqpoint{0.910950in}{0.788785in}}%
\pgfpathlineto{\pgfqpoint{0.911250in}{0.789566in}}%
\pgfpathlineto{\pgfqpoint{0.911850in}{0.791597in}}%
\pgfpathlineto{\pgfqpoint{0.912150in}{0.792378in}}%
\pgfpathlineto{\pgfqpoint{0.912750in}{0.794383in}}%
\pgfpathlineto{\pgfqpoint{0.913050in}{0.795164in}}%
\pgfpathlineto{\pgfqpoint{0.913650in}{0.797013in}}%
\pgfpathlineto{\pgfqpoint{0.913950in}{0.797794in}}%
\pgfpathlineto{\pgfqpoint{0.914549in}{0.800007in}}%
\pgfpathlineto{\pgfqpoint{0.914849in}{0.800788in}}%
\pgfpathlineto{\pgfqpoint{0.915449in}{0.803131in}}%
\pgfpathlineto{\pgfqpoint{0.915749in}{0.803912in}}%
\pgfpathlineto{\pgfqpoint{0.916349in}{0.806073in}}%
\pgfpathlineto{\pgfqpoint{0.916649in}{0.806854in}}%
\pgfpathlineto{\pgfqpoint{0.917249in}{0.809510in}}%
\pgfpathlineto{\pgfqpoint{0.917549in}{0.810291in}}%
\pgfpathlineto{\pgfqpoint{0.918149in}{0.812608in}}%
\pgfpathlineto{\pgfqpoint{0.918448in}{0.813389in}}%
\pgfpathlineto{\pgfqpoint{0.919048in}{0.816149in}}%
\pgfpathlineto{\pgfqpoint{0.919348in}{0.816930in}}%
\pgfpathlineto{\pgfqpoint{0.919948in}{0.819664in}}%
\pgfpathlineto{\pgfqpoint{0.920248in}{0.820445in}}%
\pgfpathlineto{\pgfqpoint{0.920848in}{0.823023in}}%
\pgfpathlineto{\pgfqpoint{0.921148in}{0.823804in}}%
\pgfpathlineto{\pgfqpoint{0.921748in}{0.826694in}}%
\pgfpathlineto{\pgfqpoint{0.922048in}{0.827475in}}%
\pgfpathlineto{\pgfqpoint{0.922647in}{0.830157in}}%
\pgfpathlineto{\pgfqpoint{0.922947in}{0.830938in}}%
\pgfpathlineto{\pgfqpoint{0.923547in}{0.833906in}}%
\pgfpathlineto{\pgfqpoint{0.923847in}{0.834687in}}%
\pgfpathlineto{\pgfqpoint{0.924447in}{0.838176in}}%
\pgfpathlineto{\pgfqpoint{0.924747in}{0.838957in}}%
\pgfpathlineto{\pgfqpoint{0.925347in}{0.841535in}}%
\pgfpathlineto{\pgfqpoint{0.925647in}{0.842316in}}%
\pgfpathlineto{\pgfqpoint{0.926247in}{0.845831in}}%
\pgfpathlineto{\pgfqpoint{0.926547in}{0.846612in}}%
\pgfpathlineto{\pgfqpoint{0.927146in}{0.850335in}}%
\pgfpathlineto{\pgfqpoint{0.927296in}{0.851090in}}%
\pgfpathlineto{\pgfqpoint{0.927896in}{0.854371in}}%
\pgfpathlineto{\pgfqpoint{0.928196in}{0.855152in}}%
\pgfpathlineto{\pgfqpoint{0.928796in}{0.858589in}}%
\pgfpathlineto{\pgfqpoint{0.928946in}{0.859344in}}%
\pgfpathlineto{\pgfqpoint{0.929546in}{0.862077in}}%
\pgfpathlineto{\pgfqpoint{0.929696in}{0.862832in}}%
\pgfpathlineto{\pgfqpoint{0.930296in}{0.866582in}}%
\pgfpathlineto{\pgfqpoint{0.930596in}{0.867363in}}%
\pgfpathlineto{\pgfqpoint{0.931195in}{0.871477in}}%
\pgfpathlineto{\pgfqpoint{0.931495in}{0.872258in}}%
\pgfpathlineto{\pgfqpoint{0.932095in}{0.876137in}}%
\pgfpathlineto{\pgfqpoint{0.932245in}{0.876892in}}%
\pgfpathlineto{\pgfqpoint{0.932845in}{0.880850in}}%
\pgfpathlineto{\pgfqpoint{0.933145in}{0.881631in}}%
\pgfpathlineto{\pgfqpoint{0.933745in}{0.885849in}}%
\pgfpathlineto{\pgfqpoint{0.934045in}{0.886630in}}%
\pgfpathlineto{\pgfqpoint{0.934645in}{0.890743in}}%
\pgfpathlineto{\pgfqpoint{0.934795in}{0.891498in}}%
\pgfpathlineto{\pgfqpoint{0.935394in}{0.895222in}}%
\pgfpathlineto{\pgfqpoint{0.935544in}{0.895977in}}%
\pgfpathlineto{\pgfqpoint{0.936144in}{0.899283in}}%
\pgfpathlineto{\pgfqpoint{0.936294in}{0.900038in}}%
\pgfpathlineto{\pgfqpoint{0.936894in}{0.903501in}}%
\pgfpathlineto{\pgfqpoint{0.937044in}{0.904256in}}%
\pgfpathlineto{\pgfqpoint{0.937644in}{0.908214in}}%
\pgfpathlineto{\pgfqpoint{0.937794in}{0.908969in}}%
\pgfpathlineto{\pgfqpoint{0.938394in}{0.913187in}}%
\pgfpathlineto{\pgfqpoint{0.938544in}{0.913942in}}%
\pgfpathlineto{\pgfqpoint{0.939144in}{0.918029in}}%
\pgfpathlineto{\pgfqpoint{0.939293in}{0.918785in}}%
\pgfpathlineto{\pgfqpoint{0.939893in}{0.923185in}}%
\pgfpathlineto{\pgfqpoint{0.940043in}{0.923940in}}%
\pgfpathlineto{\pgfqpoint{0.940643in}{0.928835in}}%
\pgfpathlineto{\pgfqpoint{0.940793in}{0.929590in}}%
\pgfpathlineto{\pgfqpoint{0.941393in}{0.934146in}}%
\pgfpathlineto{\pgfqpoint{0.941543in}{0.934901in}}%
\pgfpathlineto{\pgfqpoint{0.942143in}{0.939119in}}%
\pgfpathlineto{\pgfqpoint{0.942293in}{0.939874in}}%
\pgfpathlineto{\pgfqpoint{0.942893in}{0.944691in}}%
\pgfpathlineto{\pgfqpoint{0.943043in}{0.945446in}}%
\pgfpathlineto{\pgfqpoint{0.943642in}{0.950471in}}%
\pgfpathlineto{\pgfqpoint{0.943792in}{0.951226in}}%
\pgfpathlineto{\pgfqpoint{0.944392in}{0.956277in}}%
\pgfpathlineto{\pgfqpoint{0.944542in}{0.957032in}}%
\pgfpathlineto{\pgfqpoint{0.945142in}{0.961692in}}%
\pgfpathlineto{\pgfqpoint{0.945292in}{0.962447in}}%
\pgfpathlineto{\pgfqpoint{0.945892in}{0.966926in}}%
\pgfpathlineto{\pgfqpoint{0.946042in}{0.967681in}}%
\pgfpathlineto{\pgfqpoint{0.946642in}{0.973174in}}%
\pgfpathlineto{\pgfqpoint{0.946792in}{0.973929in}}%
\pgfpathlineto{\pgfqpoint{0.947392in}{0.978460in}}%
\pgfpathlineto{\pgfqpoint{0.947541in}{0.979215in}}%
\pgfpathlineto{\pgfqpoint{0.948141in}{0.984006in}}%
\pgfpathlineto{\pgfqpoint{0.948291in}{0.984761in}}%
\pgfpathlineto{\pgfqpoint{0.948891in}{0.990306in}}%
\pgfpathlineto{\pgfqpoint{0.949041in}{0.991061in}}%
\pgfpathlineto{\pgfqpoint{0.949641in}{0.995982in}}%
\pgfpathlineto{\pgfqpoint{0.949791in}{0.996737in}}%
\pgfpathlineto{\pgfqpoint{0.950391in}{1.001111in}}%
\pgfpathlineto{\pgfqpoint{0.950541in}{1.001866in}}%
\pgfpathlineto{\pgfqpoint{0.951141in}{1.007074in}}%
\pgfpathlineto{\pgfqpoint{0.951291in}{1.007829in}}%
\pgfpathlineto{\pgfqpoint{0.951890in}{1.012724in}}%
\pgfpathlineto{\pgfqpoint{0.952040in}{1.013479in}}%
\pgfpathlineto{\pgfqpoint{0.952640in}{1.018894in}}%
\pgfpathlineto{\pgfqpoint{0.952790in}{1.019649in}}%
\pgfpathlineto{\pgfqpoint{0.953390in}{1.025013in}}%
\pgfpathlineto{\pgfqpoint{0.953540in}{1.025768in}}%
\pgfpathlineto{\pgfqpoint{0.954140in}{1.030897in}}%
\pgfpathlineto{\pgfqpoint{0.954290in}{1.031652in}}%
\pgfpathlineto{\pgfqpoint{0.954890in}{1.037016in}}%
\pgfpathlineto{\pgfqpoint{0.955040in}{1.037771in}}%
\pgfpathlineto{\pgfqpoint{0.955640in}{1.043160in}}%
\pgfpathlineto{\pgfqpoint{0.955790in}{1.043915in}}%
\pgfpathlineto{\pgfqpoint{0.956389in}{1.049461in}}%
\pgfpathlineto{\pgfqpoint{0.956539in}{1.050216in}}%
\pgfpathlineto{\pgfqpoint{0.957139in}{1.056751in}}%
\pgfpathlineto{\pgfqpoint{0.957289in}{1.057506in}}%
\pgfpathlineto{\pgfqpoint{0.957889in}{1.062948in}}%
\pgfpathlineto{\pgfqpoint{0.958039in}{1.063703in}}%
\pgfpathlineto{\pgfqpoint{0.958639in}{1.068389in}}%
\pgfpathlineto{\pgfqpoint{0.958789in}{1.069144in}}%
\pgfpathlineto{\pgfqpoint{0.959389in}{1.074326in}}%
\pgfpathlineto{\pgfqpoint{0.959539in}{1.075081in}}%
\pgfpathlineto{\pgfqpoint{0.960138in}{1.080522in}}%
\pgfpathlineto{\pgfqpoint{0.960288in}{1.081277in}}%
\pgfpathlineto{\pgfqpoint{0.960888in}{1.087109in}}%
\pgfpathlineto{\pgfqpoint{0.961038in}{1.087864in}}%
\pgfpathlineto{\pgfqpoint{0.961638in}{1.092447in}}%
\pgfpathlineto{\pgfqpoint{0.961788in}{1.093202in}}%
\pgfpathlineto{\pgfqpoint{0.962388in}{1.098800in}}%
\pgfpathlineto{\pgfqpoint{0.962538in}{1.099555in}}%
\pgfpathlineto{\pgfqpoint{0.963138in}{1.104840in}}%
\pgfpathlineto{\pgfqpoint{0.963288in}{1.105595in}}%
\pgfpathlineto{\pgfqpoint{0.963888in}{1.110698in}}%
\pgfpathlineto{\pgfqpoint{0.964038in}{1.111453in}}%
\pgfpathlineto{\pgfqpoint{0.964637in}{1.116374in}}%
\pgfpathlineto{\pgfqpoint{0.964787in}{1.117129in}}%
\pgfpathlineto{\pgfqpoint{0.965387in}{1.122545in}}%
\pgfpathlineto{\pgfqpoint{0.965537in}{1.123300in}}%
\pgfpathlineto{\pgfqpoint{0.966137in}{1.128169in}}%
\pgfpathlineto{\pgfqpoint{0.966287in}{1.128924in}}%
\pgfpathlineto{\pgfqpoint{0.966887in}{1.134105in}}%
\pgfpathlineto{\pgfqpoint{0.967037in}{1.134860in}}%
\pgfpathlineto{\pgfqpoint{0.967637in}{1.139468in}}%
\pgfpathlineto{\pgfqpoint{0.967787in}{1.140224in}}%
\pgfpathlineto{\pgfqpoint{0.968387in}{1.145275in}}%
\pgfpathlineto{\pgfqpoint{0.968536in}{1.146030in}}%
\pgfpathlineto{\pgfqpoint{0.969136in}{1.150977in}}%
\pgfpathlineto{\pgfqpoint{0.969286in}{1.151732in}}%
\pgfpathlineto{\pgfqpoint{0.969886in}{1.155949in}}%
\pgfpathlineto{\pgfqpoint{0.970036in}{1.156705in}}%
\pgfpathlineto{\pgfqpoint{0.970636in}{1.160975in}}%
\pgfpathlineto{\pgfqpoint{0.970786in}{1.161730in}}%
\pgfpathlineto{\pgfqpoint{0.971386in}{1.165479in}}%
\pgfpathlineto{\pgfqpoint{0.971536in}{1.166234in}}%
\pgfpathlineto{\pgfqpoint{0.972136in}{1.170426in}}%
\pgfpathlineto{\pgfqpoint{0.972286in}{1.171181in}}%
\pgfpathlineto{\pgfqpoint{0.972885in}{1.175294in}}%
\pgfpathlineto{\pgfqpoint{0.973035in}{1.176050in}}%
\pgfpathlineto{\pgfqpoint{0.973635in}{1.179955in}}%
\pgfpathlineto{\pgfqpoint{0.973785in}{1.180710in}}%
\pgfpathlineto{\pgfqpoint{0.974385in}{1.184980in}}%
\pgfpathlineto{\pgfqpoint{0.974685in}{1.185761in}}%
\pgfpathlineto{\pgfqpoint{0.975285in}{1.189901in}}%
\pgfpathlineto{\pgfqpoint{0.975585in}{1.190682in}}%
\pgfpathlineto{\pgfqpoint{0.976185in}{1.194535in}}%
\pgfpathlineto{\pgfqpoint{0.976485in}{1.195316in}}%
\pgfpathlineto{\pgfqpoint{0.977084in}{1.199300in}}%
\pgfpathlineto{\pgfqpoint{0.977384in}{1.200081in}}%
\pgfpathlineto{\pgfqpoint{0.977984in}{1.204403in}}%
\pgfpathlineto{\pgfqpoint{0.978284in}{1.205184in}}%
\pgfpathlineto{\pgfqpoint{0.978884in}{1.210157in}}%
\pgfpathlineto{\pgfqpoint{0.979184in}{1.210938in}}%
\pgfpathlineto{\pgfqpoint{0.979784in}{1.214740in}}%
\pgfpathlineto{\pgfqpoint{0.979934in}{1.215495in}}%
\pgfpathlineto{\pgfqpoint{0.980534in}{1.219166in}}%
\pgfpathlineto{\pgfqpoint{0.980834in}{1.219947in}}%
\pgfpathlineto{\pgfqpoint{0.981433in}{1.223279in}}%
\pgfpathlineto{\pgfqpoint{0.981733in}{1.224061in}}%
\pgfpathlineto{\pgfqpoint{0.982333in}{1.227914in}}%
\pgfpathlineto{\pgfqpoint{0.982483in}{1.228669in}}%
\pgfpathlineto{\pgfqpoint{0.983083in}{1.231585in}}%
\pgfpathlineto{\pgfqpoint{0.983383in}{1.232366in}}%
\pgfpathlineto{\pgfqpoint{0.983983in}{1.235881in}}%
\pgfpathlineto{\pgfqpoint{0.984283in}{1.236662in}}%
\pgfpathlineto{\pgfqpoint{0.984883in}{1.240099in}}%
\pgfpathlineto{\pgfqpoint{0.985033in}{1.240854in}}%
\pgfpathlineto{\pgfqpoint{0.985632in}{1.244239in}}%
\pgfpathlineto{\pgfqpoint{0.985782in}{1.244994in}}%
\pgfpathlineto{\pgfqpoint{0.986382in}{1.247910in}}%
\pgfpathlineto{\pgfqpoint{0.986532in}{1.248665in}}%
\pgfpathlineto{\pgfqpoint{0.987132in}{1.251347in}}%
\pgfpathlineto{\pgfqpoint{0.987282in}{1.252102in}}%
\pgfpathlineto{\pgfqpoint{0.987882in}{1.255200in}}%
\pgfpathlineto{\pgfqpoint{0.988032in}{1.255955in}}%
\pgfpathlineto{\pgfqpoint{0.988632in}{1.259158in}}%
\pgfpathlineto{\pgfqpoint{0.988932in}{1.259939in}}%
\pgfpathlineto{\pgfqpoint{0.989531in}{1.263219in}}%
\pgfpathlineto{\pgfqpoint{0.989831in}{1.264000in}}%
\pgfpathlineto{\pgfqpoint{0.990431in}{1.267359in}}%
\pgfpathlineto{\pgfqpoint{0.990731in}{1.268140in}}%
\pgfpathlineto{\pgfqpoint{0.991331in}{1.271499in}}%
\pgfpathlineto{\pgfqpoint{0.991631in}{1.272280in}}%
\pgfpathlineto{\pgfqpoint{0.992231in}{1.274883in}}%
\pgfpathlineto{\pgfqpoint{0.992531in}{1.275665in}}%
\pgfpathlineto{\pgfqpoint{0.993131in}{1.279127in}}%
\pgfpathlineto{\pgfqpoint{0.993281in}{1.279882in}}%
\pgfpathlineto{\pgfqpoint{0.993880in}{1.282460in}}%
\pgfpathlineto{\pgfqpoint{0.994180in}{1.283241in}}%
\pgfpathlineto{\pgfqpoint{0.994780in}{1.286964in}}%
\pgfpathlineto{\pgfqpoint{0.995080in}{1.287745in}}%
\pgfpathlineto{\pgfqpoint{0.995680in}{1.291391in}}%
\pgfpathlineto{\pgfqpoint{0.995980in}{1.292172in}}%
\pgfpathlineto{\pgfqpoint{0.996580in}{1.295660in}}%
\pgfpathlineto{\pgfqpoint{0.996880in}{1.296442in}}%
\pgfpathlineto{\pgfqpoint{0.997480in}{1.299227in}}%
\pgfpathlineto{\pgfqpoint{0.997630in}{1.299982in}}%
\pgfpathlineto{\pgfqpoint{0.998229in}{1.302873in}}%
\pgfpathlineto{\pgfqpoint{0.998529in}{1.303654in}}%
\pgfpathlineto{\pgfqpoint{0.999129in}{1.307403in}}%
\pgfpathlineto{\pgfqpoint{0.999429in}{1.308184in}}%
\pgfpathlineto{\pgfqpoint{1.000029in}{1.311412in}}%
\pgfpathlineto{\pgfqpoint{1.000179in}{1.312167in}}%
\pgfpathlineto{\pgfqpoint{1.000779in}{1.315240in}}%
\pgfpathlineto{\pgfqpoint{1.001079in}{1.316021in}}%
\pgfpathlineto{\pgfqpoint{1.001679in}{1.319380in}}%
\pgfpathlineto{\pgfqpoint{1.001829in}{1.320135in}}%
\pgfpathlineto{\pgfqpoint{1.002428in}{1.322686in}}%
\pgfpathlineto{\pgfqpoint{1.002728in}{1.323467in}}%
\pgfpathlineto{\pgfqpoint{1.003328in}{1.326670in}}%
\pgfpathlineto{\pgfqpoint{1.003628in}{1.327451in}}%
\pgfpathlineto{\pgfqpoint{1.004228in}{1.330810in}}%
\pgfpathlineto{\pgfqpoint{1.004528in}{1.331591in}}%
\pgfpathlineto{\pgfqpoint{1.005128in}{1.335027in}}%
\pgfpathlineto{\pgfqpoint{1.005278in}{1.335782in}}%
\pgfpathlineto{\pgfqpoint{1.005878in}{1.338855in}}%
\pgfpathlineto{\pgfqpoint{1.006178in}{1.339636in}}%
\pgfpathlineto{\pgfqpoint{1.006777in}{1.343645in}}%
\pgfpathlineto{\pgfqpoint{1.006927in}{1.344400in}}%
\pgfpathlineto{\pgfqpoint{1.007527in}{1.347264in}}%
\pgfpathlineto{\pgfqpoint{1.007827in}{1.348046in}}%
\pgfpathlineto{\pgfqpoint{1.008427in}{1.351560in}}%
\pgfpathlineto{\pgfqpoint{1.008727in}{1.352342in}}%
\pgfpathlineto{\pgfqpoint{1.009327in}{1.356455in}}%
\pgfpathlineto{\pgfqpoint{1.009477in}{1.357210in}}%
\pgfpathlineto{\pgfqpoint{1.010077in}{1.360491in}}%
\pgfpathlineto{\pgfqpoint{1.010227in}{1.361246in}}%
\pgfpathlineto{\pgfqpoint{1.010826in}{1.364579in}}%
\pgfpathlineto{\pgfqpoint{1.011126in}{1.365360in}}%
\pgfpathlineto{\pgfqpoint{1.011726in}{1.369734in}}%
\pgfpathlineto{\pgfqpoint{1.012026in}{1.370515in}}%
\pgfpathlineto{\pgfqpoint{1.012626in}{1.374472in}}%
\pgfpathlineto{\pgfqpoint{1.012776in}{1.375228in}}%
\pgfpathlineto{\pgfqpoint{1.013376in}{1.378899in}}%
\pgfpathlineto{\pgfqpoint{1.013526in}{1.379654in}}%
\pgfpathlineto{\pgfqpoint{1.014126in}{1.383038in}}%
\pgfpathlineto{\pgfqpoint{1.014426in}{1.383819in}}%
\pgfpathlineto{\pgfqpoint{1.015025in}{1.387855in}}%
\pgfpathlineto{\pgfqpoint{1.015325in}{1.388636in}}%
\pgfpathlineto{\pgfqpoint{1.015925in}{1.392620in}}%
\pgfpathlineto{\pgfqpoint{1.016075in}{1.393375in}}%
\pgfpathlineto{\pgfqpoint{1.016675in}{1.397046in}}%
\pgfpathlineto{\pgfqpoint{1.016825in}{1.397801in}}%
\pgfpathlineto{\pgfqpoint{1.017425in}{1.401550in}}%
\pgfpathlineto{\pgfqpoint{1.017575in}{1.402305in}}%
\pgfpathlineto{\pgfqpoint{1.018175in}{1.405898in}}%
\pgfpathlineto{\pgfqpoint{1.018325in}{1.406653in}}%
\pgfpathlineto{\pgfqpoint{1.018924in}{1.409986in}}%
\pgfpathlineto{\pgfqpoint{1.019224in}{1.410767in}}%
\pgfpathlineto{\pgfqpoint{1.019824in}{1.415610in}}%
\pgfpathlineto{\pgfqpoint{1.020124in}{1.416391in}}%
\pgfpathlineto{\pgfqpoint{1.020724in}{1.420531in}}%
\pgfpathlineto{\pgfqpoint{1.021024in}{1.421312in}}%
\pgfpathlineto{\pgfqpoint{1.021624in}{1.426129in}}%
\pgfpathlineto{\pgfqpoint{1.021924in}{1.426910in}}%
\pgfpathlineto{\pgfqpoint{1.022524in}{1.431101in}}%
\pgfpathlineto{\pgfqpoint{1.022824in}{1.431883in}}%
\pgfpathlineto{\pgfqpoint{1.023423in}{1.436153in}}%
\pgfpathlineto{\pgfqpoint{1.023573in}{1.436908in}}%
\pgfpathlineto{\pgfqpoint{1.024173in}{1.440995in}}%
\pgfpathlineto{\pgfqpoint{1.024323in}{1.441750in}}%
\pgfpathlineto{\pgfqpoint{1.024923in}{1.445421in}}%
\pgfpathlineto{\pgfqpoint{1.025073in}{1.446177in}}%
\pgfpathlineto{\pgfqpoint{1.025673in}{1.450889in}}%
\pgfpathlineto{\pgfqpoint{1.025823in}{1.451644in}}%
\pgfpathlineto{\pgfqpoint{1.026423in}{1.455419in}}%
\pgfpathlineto{\pgfqpoint{1.026573in}{1.456174in}}%
\pgfpathlineto{\pgfqpoint{1.027173in}{1.460288in}}%
\pgfpathlineto{\pgfqpoint{1.027322in}{1.461043in}}%
\pgfpathlineto{\pgfqpoint{1.027922in}{1.465730in}}%
\pgfpathlineto{\pgfqpoint{1.028072in}{1.466485in}}%
\pgfpathlineto{\pgfqpoint{1.028672in}{1.471223in}}%
\pgfpathlineto{\pgfqpoint{1.028822in}{1.471979in}}%
\pgfpathlineto{\pgfqpoint{1.029422in}{1.476873in}}%
\pgfpathlineto{\pgfqpoint{1.029572in}{1.477628in}}%
\pgfpathlineto{\pgfqpoint{1.030172in}{1.482367in}}%
\pgfpathlineto{\pgfqpoint{1.030322in}{1.483122in}}%
\pgfpathlineto{\pgfqpoint{1.030922in}{1.487496in}}%
\pgfpathlineto{\pgfqpoint{1.031072in}{1.488251in}}%
\pgfpathlineto{\pgfqpoint{1.031671in}{1.492703in}}%
\pgfpathlineto{\pgfqpoint{1.031821in}{1.493458in}}%
\pgfpathlineto{\pgfqpoint{1.032421in}{1.497937in}}%
\pgfpathlineto{\pgfqpoint{1.032571in}{1.498692in}}%
\pgfpathlineto{\pgfqpoint{1.033171in}{1.503014in}}%
\pgfpathlineto{\pgfqpoint{1.033321in}{1.503769in}}%
\pgfpathlineto{\pgfqpoint{1.033921in}{1.508560in}}%
\pgfpathlineto{\pgfqpoint{1.034071in}{1.509315in}}%
\pgfpathlineto{\pgfqpoint{1.034671in}{1.514079in}}%
\pgfpathlineto{\pgfqpoint{1.034821in}{1.514834in}}%
\pgfpathlineto{\pgfqpoint{1.035421in}{1.519521in}}%
\pgfpathlineto{\pgfqpoint{1.035571in}{1.520276in}}%
\pgfpathlineto{\pgfqpoint{1.036170in}{1.524936in}}%
\pgfpathlineto{\pgfqpoint{1.036320in}{1.525691in}}%
\pgfpathlineto{\pgfqpoint{1.036920in}{1.530378in}}%
\pgfpathlineto{\pgfqpoint{1.037070in}{1.531133in}}%
\pgfpathlineto{\pgfqpoint{1.037670in}{1.536080in}}%
\pgfpathlineto{\pgfqpoint{1.037820in}{1.536835in}}%
\pgfpathlineto{\pgfqpoint{1.038420in}{1.541834in}}%
\pgfpathlineto{\pgfqpoint{1.038570in}{1.542589in}}%
\pgfpathlineto{\pgfqpoint{1.039170in}{1.547250in}}%
\pgfpathlineto{\pgfqpoint{1.039320in}{1.548005in}}%
\pgfpathlineto{\pgfqpoint{1.039919in}{1.552873in}}%
\pgfpathlineto{\pgfqpoint{1.040069in}{1.553628in}}%
\pgfpathlineto{\pgfqpoint{1.040669in}{1.558836in}}%
\pgfpathlineto{\pgfqpoint{1.040819in}{1.559591in}}%
\pgfpathlineto{\pgfqpoint{1.041419in}{1.564694in}}%
\pgfpathlineto{\pgfqpoint{1.041569in}{1.565449in}}%
\pgfpathlineto{\pgfqpoint{1.042169in}{1.570995in}}%
\pgfpathlineto{\pgfqpoint{1.042319in}{1.571750in}}%
\pgfpathlineto{\pgfqpoint{1.042919in}{1.576358in}}%
\pgfpathlineto{\pgfqpoint{1.043069in}{1.577113in}}%
\pgfpathlineto{\pgfqpoint{1.043669in}{1.582711in}}%
\pgfpathlineto{\pgfqpoint{1.043819in}{1.583466in}}%
\pgfpathlineto{\pgfqpoint{1.044418in}{1.588543in}}%
\pgfpathlineto{\pgfqpoint{1.044568in}{1.589298in}}%
\pgfpathlineto{\pgfqpoint{1.045168in}{1.595156in}}%
\pgfpathlineto{\pgfqpoint{1.045318in}{1.595911in}}%
\pgfpathlineto{\pgfqpoint{1.045918in}{1.601249in}}%
\pgfpathlineto{\pgfqpoint{1.046068in}{1.602004in}}%
\pgfpathlineto{\pgfqpoint{1.046668in}{1.608018in}}%
\pgfpathlineto{\pgfqpoint{1.046818in}{1.608773in}}%
\pgfpathlineto{\pgfqpoint{1.047418in}{1.613642in}}%
\pgfpathlineto{\pgfqpoint{1.047568in}{1.614397in}}%
\pgfpathlineto{\pgfqpoint{1.048167in}{1.619214in}}%
\pgfpathlineto{\pgfqpoint{1.048317in}{1.619969in}}%
\pgfpathlineto{\pgfqpoint{1.048917in}{1.624395in}}%
\pgfpathlineto{\pgfqpoint{1.049067in}{1.625150in}}%
\pgfpathlineto{\pgfqpoint{1.049667in}{1.630384in}}%
\pgfpathlineto{\pgfqpoint{1.049817in}{1.631139in}}%
\pgfpathlineto{\pgfqpoint{1.050417in}{1.635825in}}%
\pgfpathlineto{\pgfqpoint{1.050567in}{1.636580in}}%
\pgfpathlineto{\pgfqpoint{1.051167in}{1.642543in}}%
\pgfpathlineto{\pgfqpoint{1.051317in}{1.643298in}}%
\pgfpathlineto{\pgfqpoint{1.051917in}{1.648739in}}%
\pgfpathlineto{\pgfqpoint{1.052067in}{1.649494in}}%
\pgfpathlineto{\pgfqpoint{1.052666in}{1.655326in}}%
\pgfpathlineto{\pgfqpoint{1.052816in}{1.656081in}}%
\pgfpathlineto{\pgfqpoint{1.053416in}{1.661497in}}%
\pgfpathlineto{\pgfqpoint{1.053566in}{1.662252in}}%
\pgfpathlineto{\pgfqpoint{1.054166in}{1.667589in}}%
\pgfpathlineto{\pgfqpoint{1.054316in}{1.668345in}}%
\pgfpathlineto{\pgfqpoint{1.054916in}{1.672667in}}%
\pgfpathlineto{\pgfqpoint{1.055066in}{1.673422in}}%
\pgfpathlineto{\pgfqpoint{1.055666in}{1.678186in}}%
\pgfpathlineto{\pgfqpoint{1.055816in}{1.678941in}}%
\pgfpathlineto{\pgfqpoint{1.056416in}{1.684956in}}%
\pgfpathlineto{\pgfqpoint{1.056565in}{1.685711in}}%
\pgfpathlineto{\pgfqpoint{1.057165in}{1.690892in}}%
\pgfpathlineto{\pgfqpoint{1.057315in}{1.691647in}}%
\pgfpathlineto{\pgfqpoint{1.057915in}{1.697037in}}%
\pgfpathlineto{\pgfqpoint{1.058065in}{1.697792in}}%
\pgfpathlineto{\pgfqpoint{1.058665in}{1.702452in}}%
\pgfpathlineto{\pgfqpoint{1.058815in}{1.703207in}}%
\pgfpathlineto{\pgfqpoint{1.059415in}{1.708050in}}%
\pgfpathlineto{\pgfqpoint{1.059565in}{1.708805in}}%
\pgfpathlineto{\pgfqpoint{1.060165in}{1.713986in}}%
\pgfpathlineto{\pgfqpoint{1.060315in}{1.714741in}}%
\pgfpathlineto{\pgfqpoint{1.060914in}{1.719844in}}%
\pgfpathlineto{\pgfqpoint{1.061064in}{1.720599in}}%
\pgfpathlineto{\pgfqpoint{1.061664in}{1.725833in}}%
\pgfpathlineto{\pgfqpoint{1.061814in}{1.726588in}}%
\pgfpathlineto{\pgfqpoint{1.062414in}{1.731066in}}%
\pgfpathlineto{\pgfqpoint{1.062564in}{1.731821in}}%
\pgfpathlineto{\pgfqpoint{1.063164in}{1.736404in}}%
\pgfpathlineto{\pgfqpoint{1.063314in}{1.737159in}}%
\pgfpathlineto{\pgfqpoint{1.063914in}{1.741897in}}%
\pgfpathlineto{\pgfqpoint{1.064064in}{1.742652in}}%
\pgfpathlineto{\pgfqpoint{1.064664in}{1.747130in}}%
\pgfpathlineto{\pgfqpoint{1.064814in}{1.747886in}}%
\pgfpathlineto{\pgfqpoint{1.065413in}{1.752572in}}%
\pgfpathlineto{\pgfqpoint{1.065563in}{1.753327in}}%
\pgfpathlineto{\pgfqpoint{1.066163in}{1.757883in}}%
\pgfpathlineto{\pgfqpoint{1.066313in}{1.758639in}}%
\pgfpathlineto{\pgfqpoint{1.066913in}{1.763611in}}%
\pgfpathlineto{\pgfqpoint{1.067063in}{1.764367in}}%
\pgfpathlineto{\pgfqpoint{1.067663in}{1.768584in}}%
\pgfpathlineto{\pgfqpoint{1.067813in}{1.769339in}}%
\pgfpathlineto{\pgfqpoint{1.068413in}{1.774156in}}%
\pgfpathlineto{\pgfqpoint{1.068563in}{1.774911in}}%
\pgfpathlineto{\pgfqpoint{1.069162in}{1.779468in}}%
\pgfpathlineto{\pgfqpoint{1.069312in}{1.780223in}}%
\pgfpathlineto{\pgfqpoint{1.069912in}{1.785378in}}%
\pgfpathlineto{\pgfqpoint{1.070062in}{1.786133in}}%
\pgfpathlineto{\pgfqpoint{1.070662in}{1.790507in}}%
\pgfpathlineto{\pgfqpoint{1.070812in}{1.791262in}}%
\pgfpathlineto{\pgfqpoint{1.071412in}{1.795063in}}%
\pgfpathlineto{\pgfqpoint{1.071562in}{1.795818in}}%
\pgfpathlineto{\pgfqpoint{1.072162in}{1.799646in}}%
\pgfpathlineto{\pgfqpoint{1.072312in}{1.800401in}}%
\pgfpathlineto{\pgfqpoint{1.072912in}{1.804801in}}%
\pgfpathlineto{\pgfqpoint{1.073212in}{1.805582in}}%
\pgfpathlineto{\pgfqpoint{1.073811in}{1.810164in}}%
\pgfpathlineto{\pgfqpoint{1.073961in}{1.810920in}}%
\pgfpathlineto{\pgfqpoint{1.074561in}{1.815398in}}%
\pgfpathlineto{\pgfqpoint{1.074711in}{1.816153in}}%
\pgfpathlineto{\pgfqpoint{1.075311in}{1.820683in}}%
\pgfpathlineto{\pgfqpoint{1.075461in}{1.821438in}}%
\pgfpathlineto{\pgfqpoint{1.076061in}{1.825552in}}%
\pgfpathlineto{\pgfqpoint{1.076211in}{1.826307in}}%
\pgfpathlineto{\pgfqpoint{1.076811in}{1.830421in}}%
\pgfpathlineto{\pgfqpoint{1.076961in}{1.831176in}}%
\pgfpathlineto{\pgfqpoint{1.077560in}{1.834925in}}%
\pgfpathlineto{\pgfqpoint{1.077710in}{1.835680in}}%
\pgfpathlineto{\pgfqpoint{1.078310in}{1.839794in}}%
\pgfpathlineto{\pgfqpoint{1.078460in}{1.840549in}}%
\pgfpathlineto{\pgfqpoint{1.079060in}{1.844142in}}%
\pgfpathlineto{\pgfqpoint{1.079210in}{1.844897in}}%
\pgfpathlineto{\pgfqpoint{1.079810in}{1.848907in}}%
\pgfpathlineto{\pgfqpoint{1.079960in}{1.849662in}}%
\pgfpathlineto{\pgfqpoint{1.080560in}{1.854036in}}%
\pgfpathlineto{\pgfqpoint{1.080710in}{1.854791in}}%
\pgfpathlineto{\pgfqpoint{1.081310in}{1.858878in}}%
\pgfpathlineto{\pgfqpoint{1.081460in}{1.859633in}}%
\pgfpathlineto{\pgfqpoint{1.082059in}{1.863279in}}%
\pgfpathlineto{\pgfqpoint{1.082209in}{1.864034in}}%
\pgfpathlineto{\pgfqpoint{1.082809in}{1.867575in}}%
\pgfpathlineto{\pgfqpoint{1.082959in}{1.868330in}}%
\pgfpathlineto{\pgfqpoint{1.083559in}{1.872365in}}%
\pgfpathlineto{\pgfqpoint{1.083709in}{1.873120in}}%
\pgfpathlineto{\pgfqpoint{1.084309in}{1.876765in}}%
\pgfpathlineto{\pgfqpoint{1.084609in}{1.877546in}}%
\pgfpathlineto{\pgfqpoint{1.085209in}{1.881556in}}%
\pgfpathlineto{\pgfqpoint{1.085359in}{1.882311in}}%
\pgfpathlineto{\pgfqpoint{1.085958in}{1.886659in}}%
\pgfpathlineto{\pgfqpoint{1.086108in}{1.887414in}}%
\pgfpathlineto{\pgfqpoint{1.086708in}{1.891632in}}%
\pgfpathlineto{\pgfqpoint{1.086858in}{1.892387in}}%
\pgfpathlineto{\pgfqpoint{1.087458in}{1.895616in}}%
\pgfpathlineto{\pgfqpoint{1.087608in}{1.896371in}}%
\pgfpathlineto{\pgfqpoint{1.088208in}{1.899599in}}%
\pgfpathlineto{\pgfqpoint{1.088358in}{1.900354in}}%
\pgfpathlineto{\pgfqpoint{1.088958in}{1.903401in}}%
\pgfpathlineto{\pgfqpoint{1.089108in}{1.904156in}}%
\pgfpathlineto{\pgfqpoint{1.089708in}{1.907384in}}%
\pgfpathlineto{\pgfqpoint{1.090008in}{1.908165in}}%
\pgfpathlineto{\pgfqpoint{1.090607in}{1.911836in}}%
\pgfpathlineto{\pgfqpoint{1.090757in}{1.912591in}}%
\pgfpathlineto{\pgfqpoint{1.091357in}{1.915664in}}%
\pgfpathlineto{\pgfqpoint{1.091657in}{1.916445in}}%
\pgfpathlineto{\pgfqpoint{1.092257in}{1.920480in}}%
\pgfpathlineto{\pgfqpoint{1.092407in}{1.921235in}}%
\pgfpathlineto{\pgfqpoint{1.093007in}{1.924750in}}%
\pgfpathlineto{\pgfqpoint{1.093157in}{1.925505in}}%
\pgfpathlineto{\pgfqpoint{1.093757in}{1.929072in}}%
\pgfpathlineto{\pgfqpoint{1.093907in}{1.929827in}}%
\pgfpathlineto{\pgfqpoint{1.094506in}{1.933473in}}%
\pgfpathlineto{\pgfqpoint{1.094806in}{1.934254in}}%
\pgfpathlineto{\pgfqpoint{1.095406in}{1.938758in}}%
\pgfpathlineto{\pgfqpoint{1.095556in}{1.939513in}}%
\pgfpathlineto{\pgfqpoint{1.096156in}{1.942689in}}%
\pgfpathlineto{\pgfqpoint{1.096306in}{1.943444in}}%
\pgfpathlineto{\pgfqpoint{1.096906in}{1.947142in}}%
\pgfpathlineto{\pgfqpoint{1.097056in}{1.947897in}}%
\pgfpathlineto{\pgfqpoint{1.097656in}{1.951047in}}%
\pgfpathlineto{\pgfqpoint{1.097806in}{1.951802in}}%
\pgfpathlineto{\pgfqpoint{1.098406in}{1.955265in}}%
\pgfpathlineto{\pgfqpoint{1.098705in}{1.956046in}}%
\pgfpathlineto{\pgfqpoint{1.099305in}{1.959457in}}%
\pgfpathlineto{\pgfqpoint{1.099455in}{1.960212in}}%
\pgfpathlineto{\pgfqpoint{1.100055in}{1.963571in}}%
\pgfpathlineto{\pgfqpoint{1.100205in}{1.964326in}}%
\pgfpathlineto{\pgfqpoint{1.100805in}{1.967814in}}%
\pgfpathlineto{\pgfqpoint{1.100955in}{1.968570in}}%
\pgfpathlineto{\pgfqpoint{1.101555in}{1.971381in}}%
\pgfpathlineto{\pgfqpoint{1.101855in}{1.972163in}}%
\pgfpathlineto{\pgfqpoint{1.102455in}{1.977266in}}%
\pgfpathlineto{\pgfqpoint{1.102605in}{1.978021in}}%
\pgfpathlineto{\pgfqpoint{1.103204in}{1.981353in}}%
\pgfpathlineto{\pgfqpoint{1.103354in}{1.982108in}}%
\pgfpathlineto{\pgfqpoint{1.103954in}{1.985675in}}%
\pgfpathlineto{\pgfqpoint{1.104104in}{1.986430in}}%
\pgfpathlineto{\pgfqpoint{1.104704in}{1.990102in}}%
\pgfpathlineto{\pgfqpoint{1.105004in}{1.990883in}}%
\pgfpathlineto{\pgfqpoint{1.105604in}{1.994372in}}%
\pgfpathlineto{\pgfqpoint{1.105904in}{1.995153in}}%
\pgfpathlineto{\pgfqpoint{1.106504in}{1.999006in}}%
\pgfpathlineto{\pgfqpoint{1.106804in}{1.999787in}}%
\pgfpathlineto{\pgfqpoint{1.107403in}{2.003823in}}%
\pgfpathlineto{\pgfqpoint{1.107553in}{2.004578in}}%
\pgfpathlineto{\pgfqpoint{1.108153in}{2.007832in}}%
\pgfpathlineto{\pgfqpoint{1.108303in}{2.008587in}}%
\pgfpathlineto{\pgfqpoint{1.108903in}{2.011764in}}%
\pgfpathlineto{\pgfqpoint{1.109203in}{2.012545in}}%
\pgfpathlineto{\pgfqpoint{1.109803in}{2.016190in}}%
\pgfpathlineto{\pgfqpoint{1.110103in}{2.016971in}}%
\pgfpathlineto{\pgfqpoint{1.110703in}{2.021163in}}%
\pgfpathlineto{\pgfqpoint{1.111003in}{2.021944in}}%
\pgfpathlineto{\pgfqpoint{1.111602in}{2.025745in}}%
\pgfpathlineto{\pgfqpoint{1.111752in}{2.026500in}}%
\pgfpathlineto{\pgfqpoint{1.112352in}{2.029833in}}%
\pgfpathlineto{\pgfqpoint{1.112652in}{2.030614in}}%
\pgfpathlineto{\pgfqpoint{1.113252in}{2.034676in}}%
\pgfpathlineto{\pgfqpoint{1.113552in}{2.035457in}}%
\pgfpathlineto{\pgfqpoint{1.114152in}{2.039414in}}%
\pgfpathlineto{\pgfqpoint{1.114452in}{2.040195in}}%
\pgfpathlineto{\pgfqpoint{1.115052in}{2.043242in}}%
\pgfpathlineto{\pgfqpoint{1.115351in}{2.044023in}}%
\pgfpathlineto{\pgfqpoint{1.115951in}{2.047251in}}%
\pgfpathlineto{\pgfqpoint{1.116101in}{2.048006in}}%
\pgfpathlineto{\pgfqpoint{1.116701in}{2.050870in}}%
\pgfpathlineto{\pgfqpoint{1.116851in}{2.051625in}}%
\pgfpathlineto{\pgfqpoint{1.117451in}{2.055010in}}%
\pgfpathlineto{\pgfqpoint{1.117751in}{2.055791in}}%
\pgfpathlineto{\pgfqpoint{1.118351in}{2.058837in}}%
\pgfpathlineto{\pgfqpoint{1.118651in}{2.059619in}}%
\pgfpathlineto{\pgfqpoint{1.119251in}{2.063758in}}%
\pgfpathlineto{\pgfqpoint{1.119401in}{2.064513in}}%
\pgfpathlineto{\pgfqpoint{1.120000in}{2.067586in}}%
\pgfpathlineto{\pgfqpoint{1.120300in}{2.068367in}}%
\pgfpathlineto{\pgfqpoint{1.120900in}{2.071673in}}%
\pgfpathlineto{\pgfqpoint{1.121200in}{2.072454in}}%
\pgfpathlineto{\pgfqpoint{1.121800in}{2.076152in}}%
\pgfpathlineto{\pgfqpoint{1.122100in}{2.076933in}}%
\pgfpathlineto{\pgfqpoint{1.122700in}{2.080265in}}%
\pgfpathlineto{\pgfqpoint{1.122850in}{2.081020in}}%
\pgfpathlineto{\pgfqpoint{1.123450in}{2.083702in}}%
\pgfpathlineto{\pgfqpoint{1.123749in}{2.084483in}}%
\pgfpathlineto{\pgfqpoint{1.124349in}{2.087738in}}%
\pgfpathlineto{\pgfqpoint{1.124649in}{2.088519in}}%
\pgfpathlineto{\pgfqpoint{1.125249in}{2.092086in}}%
\pgfpathlineto{\pgfqpoint{1.125549in}{2.092867in}}%
\pgfpathlineto{\pgfqpoint{1.126149in}{2.096538in}}%
\pgfpathlineto{\pgfqpoint{1.126449in}{2.097319in}}%
\pgfpathlineto{\pgfqpoint{1.127049in}{2.100834in}}%
\pgfpathlineto{\pgfqpoint{1.127349in}{2.101615in}}%
\pgfpathlineto{\pgfqpoint{1.127948in}{2.104844in}}%
\pgfpathlineto{\pgfqpoint{1.128098in}{2.105599in}}%
\pgfpathlineto{\pgfqpoint{1.128698in}{2.108489in}}%
\pgfpathlineto{\pgfqpoint{1.128998in}{2.109270in}}%
\pgfpathlineto{\pgfqpoint{1.129598in}{2.112837in}}%
\pgfpathlineto{\pgfqpoint{1.129748in}{2.113592in}}%
\pgfpathlineto{\pgfqpoint{1.130348in}{2.116586in}}%
\pgfpathlineto{\pgfqpoint{1.130498in}{2.117341in}}%
\pgfpathlineto{\pgfqpoint{1.131098in}{2.119919in}}%
\pgfpathlineto{\pgfqpoint{1.131398in}{2.120700in}}%
\pgfpathlineto{\pgfqpoint{1.131998in}{2.123538in}}%
\pgfpathlineto{\pgfqpoint{1.132297in}{2.124319in}}%
\pgfpathlineto{\pgfqpoint{1.132897in}{2.127495in}}%
\pgfpathlineto{\pgfqpoint{1.133197in}{2.128276in}}%
\pgfpathlineto{\pgfqpoint{1.133797in}{2.131661in}}%
\pgfpathlineto{\pgfqpoint{1.134097in}{2.132442in}}%
\pgfpathlineto{\pgfqpoint{1.134697in}{2.135202in}}%
\pgfpathlineto{\pgfqpoint{1.134997in}{2.135983in}}%
\pgfpathlineto{\pgfqpoint{1.135597in}{2.138665in}}%
\pgfpathlineto{\pgfqpoint{1.135897in}{2.139446in}}%
\pgfpathlineto{\pgfqpoint{1.136496in}{2.143013in}}%
\pgfpathlineto{\pgfqpoint{1.136796in}{2.143794in}}%
\pgfpathlineto{\pgfqpoint{1.137396in}{2.147153in}}%
\pgfpathlineto{\pgfqpoint{1.137546in}{2.147908in}}%
\pgfpathlineto{\pgfqpoint{1.138146in}{2.150720in}}%
\pgfpathlineto{\pgfqpoint{1.138446in}{2.151501in}}%
\pgfpathlineto{\pgfqpoint{1.139046in}{2.154990in}}%
\pgfpathlineto{\pgfqpoint{1.139196in}{2.155745in}}%
\pgfpathlineto{\pgfqpoint{1.139796in}{2.158791in}}%
\pgfpathlineto{\pgfqpoint{1.139946in}{2.159546in}}%
\pgfpathlineto{\pgfqpoint{1.140545in}{2.162228in}}%
\pgfpathlineto{\pgfqpoint{1.140845in}{2.163009in}}%
\pgfpathlineto{\pgfqpoint{1.141445in}{2.165482in}}%
\pgfpathlineto{\pgfqpoint{1.141595in}{2.166237in}}%
\pgfpathlineto{\pgfqpoint{1.142195in}{2.168242in}}%
\pgfpathlineto{\pgfqpoint{1.142495in}{2.169023in}}%
\pgfpathlineto{\pgfqpoint{1.143095in}{2.172877in}}%
\pgfpathlineto{\pgfqpoint{1.143395in}{2.173658in}}%
\pgfpathlineto{\pgfqpoint{1.143995in}{2.176730in}}%
\pgfpathlineto{\pgfqpoint{1.144295in}{2.177511in}}%
\pgfpathlineto{\pgfqpoint{1.144894in}{2.180427in}}%
\pgfpathlineto{\pgfqpoint{1.145194in}{2.181208in}}%
\pgfpathlineto{\pgfqpoint{1.145794in}{2.184385in}}%
\pgfpathlineto{\pgfqpoint{1.146094in}{2.185166in}}%
\pgfpathlineto{\pgfqpoint{1.146694in}{2.188603in}}%
\pgfpathlineto{\pgfqpoint{1.146994in}{2.189384in}}%
\pgfpathlineto{\pgfqpoint{1.147594in}{2.192690in}}%
\pgfpathlineto{\pgfqpoint{1.147894in}{2.193471in}}%
\pgfpathlineto{\pgfqpoint{1.148494in}{2.196622in}}%
\pgfpathlineto{\pgfqpoint{1.148793in}{2.197403in}}%
\pgfpathlineto{\pgfqpoint{1.149393in}{2.200579in}}%
\pgfpathlineto{\pgfqpoint{1.149693in}{2.201360in}}%
\pgfpathlineto{\pgfqpoint{1.150293in}{2.204589in}}%
\pgfpathlineto{\pgfqpoint{1.150593in}{2.205370in}}%
\pgfpathlineto{\pgfqpoint{1.151193in}{2.208520in}}%
\pgfpathlineto{\pgfqpoint{1.151493in}{2.209301in}}%
\pgfpathlineto{\pgfqpoint{1.152093in}{2.212504in}}%
\pgfpathlineto{\pgfqpoint{1.152393in}{2.213285in}}%
\pgfpathlineto{\pgfqpoint{1.152992in}{2.216670in}}%
\pgfpathlineto{\pgfqpoint{1.153142in}{2.217425in}}%
\pgfpathlineto{\pgfqpoint{1.153742in}{2.220185in}}%
\pgfpathlineto{\pgfqpoint{1.154042in}{2.220966in}}%
\pgfpathlineto{\pgfqpoint{1.154642in}{2.224168in}}%
\pgfpathlineto{\pgfqpoint{1.154942in}{2.224949in}}%
\pgfpathlineto{\pgfqpoint{1.155542in}{2.229089in}}%
\pgfpathlineto{\pgfqpoint{1.155842in}{2.229870in}}%
\pgfpathlineto{\pgfqpoint{1.156442in}{2.232344in}}%
\pgfpathlineto{\pgfqpoint{1.156742in}{2.233125in}}%
\pgfpathlineto{\pgfqpoint{1.157341in}{2.235728in}}%
\pgfpathlineto{\pgfqpoint{1.157641in}{2.236509in}}%
\pgfpathlineto{\pgfqpoint{1.158241in}{2.239478in}}%
\pgfpathlineto{\pgfqpoint{1.158391in}{2.240233in}}%
\pgfpathlineto{\pgfqpoint{1.158991in}{2.243331in}}%
\pgfpathlineto{\pgfqpoint{1.159291in}{2.244112in}}%
\pgfpathlineto{\pgfqpoint{1.159891in}{2.247679in}}%
\pgfpathlineto{\pgfqpoint{1.160191in}{2.248460in}}%
\pgfpathlineto{\pgfqpoint{1.160791in}{2.252209in}}%
\pgfpathlineto{\pgfqpoint{1.160941in}{2.252964in}}%
\pgfpathlineto{\pgfqpoint{1.161540in}{2.255646in}}%
\pgfpathlineto{\pgfqpoint{1.161690in}{2.256401in}}%
\pgfpathlineto{\pgfqpoint{1.162290in}{2.259291in}}%
\pgfpathlineto{\pgfqpoint{1.162590in}{2.260072in}}%
\pgfpathlineto{\pgfqpoint{1.163190in}{2.262884in}}%
\pgfpathlineto{\pgfqpoint{1.163490in}{2.263665in}}%
\pgfpathlineto{\pgfqpoint{1.164090in}{2.266061in}}%
\pgfpathlineto{\pgfqpoint{1.164390in}{2.266842in}}%
\pgfpathlineto{\pgfqpoint{1.164990in}{2.269758in}}%
\pgfpathlineto{\pgfqpoint{1.165290in}{2.270539in}}%
\pgfpathlineto{\pgfqpoint{1.165889in}{2.273273in}}%
\pgfpathlineto{\pgfqpoint{1.166189in}{2.274054in}}%
\pgfpathlineto{\pgfqpoint{1.166789in}{2.277621in}}%
\pgfpathlineto{\pgfqpoint{1.167089in}{2.278402in}}%
\pgfpathlineto{\pgfqpoint{1.167689in}{2.281917in}}%
\pgfpathlineto{\pgfqpoint{1.167989in}{2.282698in}}%
\pgfpathlineto{\pgfqpoint{1.168589in}{2.285926in}}%
\pgfpathlineto{\pgfqpoint{1.168889in}{2.286707in}}%
\pgfpathlineto{\pgfqpoint{1.169489in}{2.289728in}}%
\pgfpathlineto{\pgfqpoint{1.169788in}{2.290509in}}%
\pgfpathlineto{\pgfqpoint{1.170388in}{2.294258in}}%
\pgfpathlineto{\pgfqpoint{1.170538in}{2.295013in}}%
\pgfpathlineto{\pgfqpoint{1.171138in}{2.297565in}}%
\pgfpathlineto{\pgfqpoint{1.171288in}{2.298320in}}%
\pgfpathlineto{\pgfqpoint{1.171888in}{2.301392in}}%
\pgfpathlineto{\pgfqpoint{1.172188in}{2.302173in}}%
\pgfpathlineto{\pgfqpoint{1.172788in}{2.305766in}}%
\pgfpathlineto{\pgfqpoint{1.173088in}{2.306547in}}%
\pgfpathlineto{\pgfqpoint{1.173688in}{2.309984in}}%
\pgfpathlineto{\pgfqpoint{1.173987in}{2.310765in}}%
\pgfpathlineto{\pgfqpoint{1.174587in}{2.314332in}}%
\pgfpathlineto{\pgfqpoint{1.174887in}{2.315113in}}%
\pgfpathlineto{\pgfqpoint{1.175487in}{2.318550in}}%
\pgfpathlineto{\pgfqpoint{1.175637in}{2.319305in}}%
\pgfpathlineto{\pgfqpoint{1.176237in}{2.322091in}}%
\pgfpathlineto{\pgfqpoint{1.176387in}{2.322846in}}%
\pgfpathlineto{\pgfqpoint{1.176987in}{2.325944in}}%
\pgfpathlineto{\pgfqpoint{1.177287in}{2.326725in}}%
\pgfpathlineto{\pgfqpoint{1.177887in}{2.330110in}}%
\pgfpathlineto{\pgfqpoint{1.178186in}{2.330891in}}%
\pgfpathlineto{\pgfqpoint{1.178786in}{2.334120in}}%
\pgfpathlineto{\pgfqpoint{1.179086in}{2.334901in}}%
\pgfpathlineto{\pgfqpoint{1.179686in}{2.338520in}}%
\pgfpathlineto{\pgfqpoint{1.179986in}{2.339301in}}%
\pgfpathlineto{\pgfqpoint{1.180586in}{2.342634in}}%
\pgfpathlineto{\pgfqpoint{1.180886in}{2.343415in}}%
\pgfpathlineto{\pgfqpoint{1.181486in}{2.346851in}}%
\pgfpathlineto{\pgfqpoint{1.181786in}{2.347632in}}%
\pgfpathlineto{\pgfqpoint{1.182385in}{2.351069in}}%
\pgfpathlineto{\pgfqpoint{1.182535in}{2.351824in}}%
\pgfpathlineto{\pgfqpoint{1.183135in}{2.354532in}}%
\pgfpathlineto{\pgfqpoint{1.183435in}{2.355313in}}%
\pgfpathlineto{\pgfqpoint{1.184035in}{2.359115in}}%
\pgfpathlineto{\pgfqpoint{1.184335in}{2.359896in}}%
\pgfpathlineto{\pgfqpoint{1.184935in}{2.363124in}}%
\pgfpathlineto{\pgfqpoint{1.185235in}{2.363905in}}%
\pgfpathlineto{\pgfqpoint{1.185835in}{2.367472in}}%
\pgfpathlineto{\pgfqpoint{1.186135in}{2.368253in}}%
\pgfpathlineto{\pgfqpoint{1.186734in}{2.371300in}}%
\pgfpathlineto{\pgfqpoint{1.187034in}{2.372081in}}%
\pgfpathlineto{\pgfqpoint{1.187634in}{2.375205in}}%
\pgfpathlineto{\pgfqpoint{1.187784in}{2.375960in}}%
\pgfpathlineto{\pgfqpoint{1.188384in}{2.378538in}}%
\pgfpathlineto{\pgfqpoint{1.188684in}{2.379319in}}%
\pgfpathlineto{\pgfqpoint{1.189284in}{2.381922in}}%
\pgfpathlineto{\pgfqpoint{1.189584in}{2.382703in}}%
\pgfpathlineto{\pgfqpoint{1.190184in}{2.385750in}}%
\pgfpathlineto{\pgfqpoint{1.190484in}{2.386531in}}%
\pgfpathlineto{\pgfqpoint{1.191083in}{2.389942in}}%
\pgfpathlineto{\pgfqpoint{1.191383in}{2.390723in}}%
\pgfpathlineto{\pgfqpoint{1.191983in}{2.393899in}}%
\pgfpathlineto{\pgfqpoint{1.192283in}{2.394680in}}%
\pgfpathlineto{\pgfqpoint{1.192883in}{2.398143in}}%
\pgfpathlineto{\pgfqpoint{1.193033in}{2.398898in}}%
\pgfpathlineto{\pgfqpoint{1.193633in}{2.401319in}}%
\pgfpathlineto{\pgfqpoint{1.193933in}{2.402100in}}%
\pgfpathlineto{\pgfqpoint{1.194533in}{2.405459in}}%
\pgfpathlineto{\pgfqpoint{1.194833in}{2.406240in}}%
\pgfpathlineto{\pgfqpoint{1.195432in}{2.408948in}}%
\pgfpathlineto{\pgfqpoint{1.195732in}{2.409729in}}%
\pgfpathlineto{\pgfqpoint{1.196332in}{2.413036in}}%
\pgfpathlineto{\pgfqpoint{1.196482in}{2.413791in}}%
\pgfpathlineto{\pgfqpoint{1.197082in}{2.416551in}}%
\pgfpathlineto{\pgfqpoint{1.197382in}{2.417332in}}%
\pgfpathlineto{\pgfqpoint{1.197982in}{2.420534in}}%
\pgfpathlineto{\pgfqpoint{1.198282in}{2.421315in}}%
\pgfpathlineto{\pgfqpoint{1.198882in}{2.425039in}}%
\pgfpathlineto{\pgfqpoint{1.199181in}{2.425820in}}%
\pgfpathlineto{\pgfqpoint{1.199781in}{2.429178in}}%
\pgfpathlineto{\pgfqpoint{1.200081in}{2.429959in}}%
\pgfpathlineto{\pgfqpoint{1.200681in}{2.433006in}}%
\pgfpathlineto{\pgfqpoint{1.200981in}{2.433787in}}%
\pgfpathlineto{\pgfqpoint{1.201581in}{2.436911in}}%
\pgfpathlineto{\pgfqpoint{1.201881in}{2.437692in}}%
\pgfpathlineto{\pgfqpoint{1.202481in}{2.440712in}}%
\pgfpathlineto{\pgfqpoint{1.202631in}{2.441467in}}%
\pgfpathlineto{\pgfqpoint{1.203231in}{2.443759in}}%
\pgfpathlineto{\pgfqpoint{1.203530in}{2.444540in}}%
\pgfpathlineto{\pgfqpoint{1.204130in}{2.446935in}}%
\pgfpathlineto{\pgfqpoint{1.204430in}{2.447716in}}%
\pgfpathlineto{\pgfqpoint{1.205030in}{2.450762in}}%
\pgfpathlineto{\pgfqpoint{1.205330in}{2.451543in}}%
\pgfpathlineto{\pgfqpoint{1.205930in}{2.454407in}}%
\pgfpathlineto{\pgfqpoint{1.206080in}{2.455163in}}%
\pgfpathlineto{\pgfqpoint{1.206680in}{2.457506in}}%
\pgfpathlineto{\pgfqpoint{1.206830in}{2.458261in}}%
\pgfpathlineto{\pgfqpoint{1.207430in}{2.460656in}}%
\pgfpathlineto{\pgfqpoint{1.207579in}{2.461411in}}%
\pgfpathlineto{\pgfqpoint{1.208179in}{2.464093in}}%
\pgfpathlineto{\pgfqpoint{1.208329in}{2.464848in}}%
\pgfpathlineto{\pgfqpoint{1.208929in}{2.467400in}}%
\pgfpathlineto{\pgfqpoint{1.209229in}{2.468181in}}%
\pgfpathlineto{\pgfqpoint{1.209829in}{2.471409in}}%
\pgfpathlineto{\pgfqpoint{1.210129in}{2.472190in}}%
\pgfpathlineto{\pgfqpoint{1.210729in}{2.474638in}}%
\pgfpathlineto{\pgfqpoint{1.211029in}{2.475419in}}%
\pgfpathlineto{\pgfqpoint{1.211629in}{2.478153in}}%
\pgfpathlineto{\pgfqpoint{1.211928in}{2.478934in}}%
\pgfpathlineto{\pgfqpoint{1.212528in}{2.481511in}}%
\pgfpathlineto{\pgfqpoint{1.212828in}{2.482292in}}%
\pgfpathlineto{\pgfqpoint{1.213428in}{2.485261in}}%
\pgfpathlineto{\pgfqpoint{1.213728in}{2.486042in}}%
\pgfpathlineto{\pgfqpoint{1.214328in}{2.488723in}}%
\pgfpathlineto{\pgfqpoint{1.214628in}{2.489504in}}%
\pgfpathlineto{\pgfqpoint{1.215228in}{2.492056in}}%
\pgfpathlineto{\pgfqpoint{1.215528in}{2.492837in}}%
\pgfpathlineto{\pgfqpoint{1.216127in}{2.495389in}}%
\pgfpathlineto{\pgfqpoint{1.216427in}{2.496170in}}%
\pgfpathlineto{\pgfqpoint{1.217027in}{2.498539in}}%
\pgfpathlineto{\pgfqpoint{1.217327in}{2.499320in}}%
\pgfpathlineto{\pgfqpoint{1.217927in}{2.501716in}}%
\pgfpathlineto{\pgfqpoint{1.218227in}{2.502497in}}%
\pgfpathlineto{\pgfqpoint{1.218827in}{2.504840in}}%
\pgfpathlineto{\pgfqpoint{1.219127in}{2.505621in}}%
\pgfpathlineto{\pgfqpoint{1.219727in}{2.507990in}}%
\pgfpathlineto{\pgfqpoint{1.220027in}{2.508771in}}%
\pgfpathlineto{\pgfqpoint{1.220626in}{2.511141in}}%
\pgfpathlineto{\pgfqpoint{1.220926in}{2.511922in}}%
\pgfpathlineto{\pgfqpoint{1.221526in}{2.514161in}}%
\pgfpathlineto{\pgfqpoint{1.221826in}{2.514942in}}%
\pgfpathlineto{\pgfqpoint{1.222426in}{2.517337in}}%
\pgfpathlineto{\pgfqpoint{1.222876in}{2.518144in}}%
\pgfpathlineto{\pgfqpoint{1.223476in}{2.520227in}}%
\pgfpathlineto{\pgfqpoint{1.223776in}{2.521008in}}%
\pgfpathlineto{\pgfqpoint{1.224375in}{2.523222in}}%
\pgfpathlineto{\pgfqpoint{1.224675in}{2.524003in}}%
\pgfpathlineto{\pgfqpoint{1.225275in}{2.526294in}}%
\pgfpathlineto{\pgfqpoint{1.225575in}{2.527075in}}%
\pgfpathlineto{\pgfqpoint{1.226175in}{2.529418in}}%
\pgfpathlineto{\pgfqpoint{1.226475in}{2.530199in}}%
\pgfpathlineto{\pgfqpoint{1.227075in}{2.532152in}}%
\pgfpathlineto{\pgfqpoint{1.227375in}{2.532933in}}%
\pgfpathlineto{\pgfqpoint{1.227975in}{2.535563in}}%
\pgfpathlineto{\pgfqpoint{1.228275in}{2.536344in}}%
\pgfpathlineto{\pgfqpoint{1.228874in}{2.538713in}}%
\pgfpathlineto{\pgfqpoint{1.229174in}{2.539494in}}%
\pgfpathlineto{\pgfqpoint{1.229774in}{2.541343in}}%
\pgfpathlineto{\pgfqpoint{1.230074in}{2.542124in}}%
\pgfpathlineto{\pgfqpoint{1.230674in}{2.544051in}}%
\pgfpathlineto{\pgfqpoint{1.230974in}{2.544832in}}%
\pgfpathlineto{\pgfqpoint{1.231574in}{2.547592in}}%
\pgfpathlineto{\pgfqpoint{1.231874in}{2.548373in}}%
\pgfpathlineto{\pgfqpoint{1.232474in}{2.551002in}}%
\pgfpathlineto{\pgfqpoint{1.232773in}{2.551783in}}%
\pgfpathlineto{\pgfqpoint{1.233373in}{2.553866in}}%
\pgfpathlineto{\pgfqpoint{1.233673in}{2.554647in}}%
\pgfpathlineto{\pgfqpoint{1.234273in}{2.556522in}}%
\pgfpathlineto{\pgfqpoint{1.234573in}{2.557303in}}%
\pgfpathlineto{\pgfqpoint{1.235173in}{2.559803in}}%
\pgfpathlineto{\pgfqpoint{1.235473in}{2.560584in}}%
\pgfpathlineto{\pgfqpoint{1.236073in}{2.562458in}}%
\pgfpathlineto{\pgfqpoint{1.236373in}{2.563239in}}%
\pgfpathlineto{\pgfqpoint{1.236972in}{2.564932in}}%
\pgfpathlineto{\pgfqpoint{1.237272in}{2.565713in}}%
\pgfpathlineto{\pgfqpoint{1.237872in}{2.567744in}}%
\pgfpathlineto{\pgfqpoint{1.238172in}{2.568525in}}%
\pgfpathlineto{\pgfqpoint{1.238772in}{2.570712in}}%
\pgfpathlineto{\pgfqpoint{1.239222in}{2.571519in}}%
\pgfpathlineto{\pgfqpoint{1.239822in}{2.573914in}}%
\pgfpathlineto{\pgfqpoint{1.240122in}{2.574695in}}%
\pgfpathlineto{\pgfqpoint{1.240722in}{2.577039in}}%
\pgfpathlineto{\pgfqpoint{1.241021in}{2.577820in}}%
\pgfpathlineto{\pgfqpoint{1.241621in}{2.579642in}}%
\pgfpathlineto{\pgfqpoint{1.241921in}{2.580423in}}%
\pgfpathlineto{\pgfqpoint{1.242521in}{2.582116in}}%
\pgfpathlineto{\pgfqpoint{1.242821in}{2.582897in}}%
\pgfpathlineto{\pgfqpoint{1.243421in}{2.584641in}}%
\pgfpathlineto{\pgfqpoint{1.243871in}{2.585448in}}%
\pgfpathlineto{\pgfqpoint{1.244471in}{2.587427in}}%
\pgfpathlineto{\pgfqpoint{1.244771in}{2.588208in}}%
\pgfpathlineto{\pgfqpoint{1.245370in}{2.590291in}}%
\pgfpathlineto{\pgfqpoint{1.245670in}{2.591072in}}%
\pgfpathlineto{\pgfqpoint{1.246270in}{2.593207in}}%
\pgfpathlineto{\pgfqpoint{1.246570in}{2.593988in}}%
\pgfpathlineto{\pgfqpoint{1.247170in}{2.595472in}}%
\pgfpathlineto{\pgfqpoint{1.247620in}{2.596279in}}%
\pgfpathlineto{\pgfqpoint{1.248220in}{2.598597in}}%
\pgfpathlineto{\pgfqpoint{1.248670in}{2.599404in}}%
\pgfpathlineto{\pgfqpoint{1.249270in}{2.601643in}}%
\pgfpathlineto{\pgfqpoint{1.249569in}{2.602424in}}%
\pgfpathlineto{\pgfqpoint{1.250169in}{2.604559in}}%
\pgfpathlineto{\pgfqpoint{1.250469in}{2.605340in}}%
\pgfpathlineto{\pgfqpoint{1.251069in}{2.606928in}}%
\pgfpathlineto{\pgfqpoint{1.251369in}{2.607709in}}%
\pgfpathlineto{\pgfqpoint{1.251969in}{2.609506in}}%
\pgfpathlineto{\pgfqpoint{1.252269in}{2.610287in}}%
\pgfpathlineto{\pgfqpoint{1.252869in}{2.612188in}}%
\pgfpathlineto{\pgfqpoint{1.253319in}{2.612995in}}%
\pgfpathlineto{\pgfqpoint{1.253918in}{2.615130in}}%
\pgfpathlineto{\pgfqpoint{1.254218in}{2.615911in}}%
\pgfpathlineto{\pgfqpoint{1.254818in}{2.617473in}}%
\pgfpathlineto{\pgfqpoint{1.255268in}{2.618280in}}%
\pgfpathlineto{\pgfqpoint{1.255868in}{2.620285in}}%
\pgfpathlineto{\pgfqpoint{1.256168in}{2.621066in}}%
\pgfpathlineto{\pgfqpoint{1.256768in}{2.622628in}}%
\pgfpathlineto{\pgfqpoint{1.257218in}{2.623435in}}%
\pgfpathlineto{\pgfqpoint{1.257817in}{2.625440in}}%
\pgfpathlineto{\pgfqpoint{1.258267in}{2.626247in}}%
\pgfpathlineto{\pgfqpoint{1.258867in}{2.628148in}}%
\pgfpathlineto{\pgfqpoint{1.259317in}{2.628955in}}%
\pgfpathlineto{\pgfqpoint{1.259917in}{2.631064in}}%
\pgfpathlineto{\pgfqpoint{1.260367in}{2.631871in}}%
\pgfpathlineto{\pgfqpoint{1.260967in}{2.633616in}}%
\pgfpathlineto{\pgfqpoint{1.261267in}{2.634397in}}%
\pgfpathlineto{\pgfqpoint{1.261867in}{2.636245in}}%
\pgfpathlineto{\pgfqpoint{1.262316in}{2.637052in}}%
\pgfpathlineto{\pgfqpoint{1.262916in}{2.639239in}}%
\pgfpathlineto{\pgfqpoint{1.263366in}{2.640047in}}%
\pgfpathlineto{\pgfqpoint{1.263966in}{2.641921in}}%
\pgfpathlineto{\pgfqpoint{1.264266in}{2.642702in}}%
\pgfpathlineto{\pgfqpoint{1.264866in}{2.644395in}}%
\pgfpathlineto{\pgfqpoint{1.265316in}{2.645202in}}%
\pgfpathlineto{\pgfqpoint{1.265916in}{2.647181in}}%
\pgfpathlineto{\pgfqpoint{1.266365in}{2.647988in}}%
\pgfpathlineto{\pgfqpoint{1.266965in}{2.649784in}}%
\pgfpathlineto{\pgfqpoint{1.267415in}{2.650591in}}%
\pgfpathlineto{\pgfqpoint{1.268015in}{2.652700in}}%
\pgfpathlineto{\pgfqpoint{1.268465in}{2.653507in}}%
\pgfpathlineto{\pgfqpoint{1.269065in}{2.655070in}}%
\pgfpathlineto{\pgfqpoint{1.269515in}{2.655877in}}%
\pgfpathlineto{\pgfqpoint{1.270115in}{2.657725in}}%
\pgfpathlineto{\pgfqpoint{1.270414in}{2.658506in}}%
\pgfpathlineto{\pgfqpoint{1.271014in}{2.659964in}}%
\pgfpathlineto{\pgfqpoint{1.271464in}{2.660771in}}%
\pgfpathlineto{\pgfqpoint{1.272064in}{2.662542in}}%
\pgfpathlineto{\pgfqpoint{1.272364in}{2.663323in}}%
\pgfpathlineto{\pgfqpoint{1.272964in}{2.664651in}}%
\pgfpathlineto{\pgfqpoint{1.273264in}{2.665432in}}%
\pgfpathlineto{\pgfqpoint{1.273864in}{2.666994in}}%
\pgfpathlineto{\pgfqpoint{1.274164in}{2.667775in}}%
\pgfpathlineto{\pgfqpoint{1.274763in}{2.669363in}}%
\pgfpathlineto{\pgfqpoint{1.275063in}{2.670145in}}%
\pgfpathlineto{\pgfqpoint{1.275663in}{2.671551in}}%
\pgfpathlineto{\pgfqpoint{1.276113in}{2.672358in}}%
\pgfpathlineto{\pgfqpoint{1.276713in}{2.674024in}}%
\pgfpathlineto{\pgfqpoint{1.277163in}{2.674831in}}%
\pgfpathlineto{\pgfqpoint{1.277763in}{2.676706in}}%
\pgfpathlineto{\pgfqpoint{1.278213in}{2.677513in}}%
\pgfpathlineto{\pgfqpoint{1.278812in}{2.679778in}}%
\pgfpathlineto{\pgfqpoint{1.279112in}{2.680559in}}%
\pgfpathlineto{\pgfqpoint{1.279712in}{2.682382in}}%
\pgfpathlineto{\pgfqpoint{1.280162in}{2.683189in}}%
\pgfpathlineto{\pgfqpoint{1.280762in}{2.684803in}}%
\pgfpathlineto{\pgfqpoint{1.281212in}{2.685610in}}%
\pgfpathlineto{\pgfqpoint{1.281812in}{2.687381in}}%
\pgfpathlineto{\pgfqpoint{1.282112in}{2.688162in}}%
\pgfpathlineto{\pgfqpoint{1.282712in}{2.689776in}}%
\pgfpathlineto{\pgfqpoint{1.283161in}{2.690583in}}%
\pgfpathlineto{\pgfqpoint{1.283761in}{2.692145in}}%
\pgfpathlineto{\pgfqpoint{1.284361in}{2.692978in}}%
\pgfpathlineto{\pgfqpoint{1.284961in}{2.694671in}}%
\pgfpathlineto{\pgfqpoint{1.285411in}{2.695478in}}%
\pgfpathlineto{\pgfqpoint{1.286011in}{2.697248in}}%
\pgfpathlineto{\pgfqpoint{1.286461in}{2.698056in}}%
\pgfpathlineto{\pgfqpoint{1.287061in}{2.699904in}}%
\pgfpathlineto{\pgfqpoint{1.287510in}{2.700711in}}%
\pgfpathlineto{\pgfqpoint{1.288110in}{2.702247in}}%
\pgfpathlineto{\pgfqpoint{1.288410in}{2.703028in}}%
\pgfpathlineto{\pgfqpoint{1.289010in}{2.704122in}}%
\pgfpathlineto{\pgfqpoint{1.289460in}{2.704929in}}%
\pgfpathlineto{\pgfqpoint{1.290060in}{2.706595in}}%
\pgfpathlineto{\pgfqpoint{1.290360in}{2.707377in}}%
\pgfpathlineto{\pgfqpoint{1.290960in}{2.708470in}}%
\pgfpathlineto{\pgfqpoint{1.291559in}{2.709303in}}%
\pgfpathlineto{\pgfqpoint{1.292159in}{2.711412in}}%
\pgfpathlineto{\pgfqpoint{1.292609in}{2.712219in}}%
\pgfpathlineto{\pgfqpoint{1.293209in}{2.713834in}}%
\pgfpathlineto{\pgfqpoint{1.293659in}{2.714641in}}%
\pgfpathlineto{\pgfqpoint{1.294259in}{2.716463in}}%
\pgfpathlineto{\pgfqpoint{1.294709in}{2.717270in}}%
\pgfpathlineto{\pgfqpoint{1.295309in}{2.718494in}}%
\pgfpathlineto{\pgfqpoint{1.295758in}{2.719301in}}%
\pgfpathlineto{\pgfqpoint{1.296358in}{2.720967in}}%
\pgfpathlineto{\pgfqpoint{1.296658in}{2.721749in}}%
\pgfpathlineto{\pgfqpoint{1.297258in}{2.723050in}}%
\pgfpathlineto{\pgfqpoint{1.297558in}{2.723831in}}%
\pgfpathlineto{\pgfqpoint{1.298158in}{2.725524in}}%
\pgfpathlineto{\pgfqpoint{1.298608in}{2.726331in}}%
\pgfpathlineto{\pgfqpoint{1.299208in}{2.727711in}}%
\pgfpathlineto{\pgfqpoint{1.299508in}{2.728492in}}%
\pgfpathlineto{\pgfqpoint{1.300107in}{2.730341in}}%
\pgfpathlineto{\pgfqpoint{1.300557in}{2.731148in}}%
\pgfpathlineto{\pgfqpoint{1.301157in}{2.732814in}}%
\pgfpathlineto{\pgfqpoint{1.301607in}{2.733621in}}%
\pgfpathlineto{\pgfqpoint{1.302207in}{2.735261in}}%
\pgfpathlineto{\pgfqpoint{1.302507in}{2.736043in}}%
\pgfpathlineto{\pgfqpoint{1.303107in}{2.737891in}}%
\pgfpathlineto{\pgfqpoint{1.303707in}{2.738724in}}%
\pgfpathlineto{\pgfqpoint{1.304306in}{2.740339in}}%
\pgfpathlineto{\pgfqpoint{1.304756in}{2.741146in}}%
\pgfpathlineto{\pgfqpoint{1.305356in}{2.742500in}}%
\pgfpathlineto{\pgfqpoint{1.305806in}{2.743307in}}%
\pgfpathlineto{\pgfqpoint{1.306406in}{2.745155in}}%
\pgfpathlineto{\pgfqpoint{1.306856in}{2.745962in}}%
\pgfpathlineto{\pgfqpoint{1.307456in}{2.747030in}}%
\pgfpathlineto{\pgfqpoint{1.307906in}{2.747837in}}%
\pgfpathlineto{\pgfqpoint{1.308505in}{2.749191in}}%
\pgfpathlineto{\pgfqpoint{1.308955in}{2.749998in}}%
\pgfpathlineto{\pgfqpoint{1.309555in}{2.751144in}}%
\pgfpathlineto{\pgfqpoint{1.310005in}{2.751951in}}%
\pgfpathlineto{\pgfqpoint{1.310605in}{2.753383in}}%
\pgfpathlineto{\pgfqpoint{1.311055in}{2.754190in}}%
\pgfpathlineto{\pgfqpoint{1.311655in}{2.755752in}}%
\pgfpathlineto{\pgfqpoint{1.312255in}{2.756585in}}%
\pgfpathlineto{\pgfqpoint{1.312854in}{2.758225in}}%
\pgfpathlineto{\pgfqpoint{1.313304in}{2.759033in}}%
\pgfpathlineto{\pgfqpoint{1.313904in}{2.760334in}}%
\pgfpathlineto{\pgfqpoint{1.314504in}{2.761168in}}%
\pgfpathlineto{\pgfqpoint{1.315104in}{2.762834in}}%
\pgfpathlineto{\pgfqpoint{1.315554in}{2.763641in}}%
\pgfpathlineto{\pgfqpoint{1.316154in}{2.765411in}}%
\pgfpathlineto{\pgfqpoint{1.316603in}{2.766219in}}%
\pgfpathlineto{\pgfqpoint{1.317203in}{2.767989in}}%
\pgfpathlineto{\pgfqpoint{1.317653in}{2.768796in}}%
\pgfpathlineto{\pgfqpoint{1.318253in}{2.770202in}}%
\pgfpathlineto{\pgfqpoint{1.318703in}{2.771009in}}%
\pgfpathlineto{\pgfqpoint{1.319303in}{2.772129in}}%
\pgfpathlineto{\pgfqpoint{1.319903in}{2.772962in}}%
\pgfpathlineto{\pgfqpoint{1.320503in}{2.774290in}}%
\pgfpathlineto{\pgfqpoint{1.320952in}{2.775097in}}%
\pgfpathlineto{\pgfqpoint{1.321552in}{2.776373in}}%
\pgfpathlineto{\pgfqpoint{1.322302in}{2.777232in}}%
\pgfpathlineto{\pgfqpoint{1.322902in}{2.778794in}}%
\pgfpathlineto{\pgfqpoint{1.323352in}{2.779601in}}%
\pgfpathlineto{\pgfqpoint{1.323952in}{2.780955in}}%
\pgfpathlineto{\pgfqpoint{1.324402in}{2.781762in}}%
\pgfpathlineto{\pgfqpoint{1.325001in}{2.783168in}}%
\pgfpathlineto{\pgfqpoint{1.325451in}{2.783975in}}%
\pgfpathlineto{\pgfqpoint{1.326051in}{2.785173in}}%
\pgfpathlineto{\pgfqpoint{1.326501in}{2.785980in}}%
\pgfpathlineto{\pgfqpoint{1.327101in}{2.787334in}}%
\pgfpathlineto{\pgfqpoint{1.327551in}{2.788141in}}%
\pgfpathlineto{\pgfqpoint{1.328151in}{2.789339in}}%
\pgfpathlineto{\pgfqpoint{1.328751in}{2.790172in}}%
\pgfpathlineto{\pgfqpoint{1.329350in}{2.791474in}}%
\pgfpathlineto{\pgfqpoint{1.329950in}{2.792307in}}%
\pgfpathlineto{\pgfqpoint{1.330550in}{2.793479in}}%
\pgfpathlineto{\pgfqpoint{1.331150in}{2.794312in}}%
\pgfpathlineto{\pgfqpoint{1.331750in}{2.795301in}}%
\pgfpathlineto{\pgfqpoint{1.332200in}{2.796108in}}%
\pgfpathlineto{\pgfqpoint{1.332800in}{2.797046in}}%
\pgfpathlineto{\pgfqpoint{1.333399in}{2.797879in}}%
\pgfpathlineto{\pgfqpoint{1.333999in}{2.799129in}}%
\pgfpathlineto{\pgfqpoint{1.334599in}{2.799962in}}%
\pgfpathlineto{\pgfqpoint{1.335199in}{2.801081in}}%
\pgfpathlineto{\pgfqpoint{1.335799in}{2.801914in}}%
\pgfpathlineto{\pgfqpoint{1.336399in}{2.802956in}}%
\pgfpathlineto{\pgfqpoint{1.336999in}{2.803789in}}%
\pgfpathlineto{\pgfqpoint{1.337598in}{2.805013in}}%
\pgfpathlineto{\pgfqpoint{1.338348in}{2.805872in}}%
\pgfpathlineto{\pgfqpoint{1.338948in}{2.807148in}}%
\pgfpathlineto{\pgfqpoint{1.339548in}{2.807981in}}%
\pgfpathlineto{\pgfqpoint{1.340148in}{2.809283in}}%
\pgfpathlineto{\pgfqpoint{1.340598in}{2.810090in}}%
\pgfpathlineto{\pgfqpoint{1.341198in}{2.811470in}}%
\pgfpathlineto{\pgfqpoint{1.341797in}{2.812303in}}%
\pgfpathlineto{\pgfqpoint{1.342397in}{2.813631in}}%
\pgfpathlineto{\pgfqpoint{1.342997in}{2.814464in}}%
\pgfpathlineto{\pgfqpoint{1.343597in}{2.815505in}}%
\pgfpathlineto{\pgfqpoint{1.344197in}{2.816339in}}%
\pgfpathlineto{\pgfqpoint{1.344797in}{2.817094in}}%
\pgfpathlineto{\pgfqpoint{1.345547in}{2.817953in}}%
\pgfpathlineto{\pgfqpoint{1.346146in}{2.818760in}}%
\pgfpathlineto{\pgfqpoint{1.347046in}{2.819645in}}%
\pgfpathlineto{\pgfqpoint{1.347646in}{2.820582in}}%
\pgfpathlineto{\pgfqpoint{1.348396in}{2.821442in}}%
\pgfpathlineto{\pgfqpoint{1.348996in}{2.822353in}}%
\pgfpathlineto{\pgfqpoint{1.349746in}{2.823212in}}%
\pgfpathlineto{\pgfqpoint{1.350345in}{2.824149in}}%
\pgfpathlineto{\pgfqpoint{1.351095in}{2.825009in}}%
\pgfpathlineto{\pgfqpoint{1.351695in}{2.825998in}}%
\pgfpathlineto{\pgfqpoint{1.352295in}{2.826831in}}%
\pgfpathlineto{\pgfqpoint{1.352895in}{2.827690in}}%
\pgfpathlineto{\pgfqpoint{1.353795in}{2.828576in}}%
\pgfpathlineto{\pgfqpoint{1.354394in}{2.829617in}}%
\pgfpathlineto{\pgfqpoint{1.354994in}{2.830450in}}%
\pgfpathlineto{\pgfqpoint{1.355594in}{2.831153in}}%
\pgfpathlineto{\pgfqpoint{1.356644in}{2.832065in}}%
\pgfpathlineto{\pgfqpoint{1.357244in}{2.833132in}}%
\pgfpathlineto{\pgfqpoint{1.357994in}{2.833991in}}%
\pgfpathlineto{\pgfqpoint{1.358593in}{2.835423in}}%
\pgfpathlineto{\pgfqpoint{1.359193in}{2.836256in}}%
\pgfpathlineto{\pgfqpoint{1.359793in}{2.837090in}}%
\pgfpathlineto{\pgfqpoint{1.360543in}{2.837949in}}%
\pgfpathlineto{\pgfqpoint{1.361143in}{2.838600in}}%
\pgfpathlineto{\pgfqpoint{1.361743in}{2.839433in}}%
\pgfpathlineto{\pgfqpoint{1.362343in}{2.840266in}}%
\pgfpathlineto{\pgfqpoint{1.363092in}{2.841125in}}%
\pgfpathlineto{\pgfqpoint{1.363692in}{2.841906in}}%
\pgfpathlineto{\pgfqpoint{1.364292in}{2.842739in}}%
\pgfpathlineto{\pgfqpoint{1.364892in}{2.843468in}}%
\pgfpathlineto{\pgfqpoint{1.365792in}{2.844354in}}%
\pgfpathlineto{\pgfqpoint{1.366392in}{2.845629in}}%
\pgfpathlineto{\pgfqpoint{1.366991in}{2.846463in}}%
\pgfpathlineto{\pgfqpoint{1.367591in}{2.847192in}}%
\pgfpathlineto{\pgfqpoint{1.368641in}{2.848103in}}%
\pgfpathlineto{\pgfqpoint{1.369241in}{2.849144in}}%
\pgfpathlineto{\pgfqpoint{1.370141in}{2.850030in}}%
\pgfpathlineto{\pgfqpoint{1.370741in}{2.850837in}}%
\pgfpathlineto{\pgfqpoint{1.371490in}{2.851696in}}%
\pgfpathlineto{\pgfqpoint{1.372090in}{2.852165in}}%
\pgfpathlineto{\pgfqpoint{1.372990in}{2.853050in}}%
\pgfpathlineto{\pgfqpoint{1.373590in}{2.854169in}}%
\pgfpathlineto{\pgfqpoint{1.374490in}{2.855055in}}%
\pgfpathlineto{\pgfqpoint{1.375090in}{2.855966in}}%
\pgfpathlineto{\pgfqpoint{1.375689in}{2.856799in}}%
\pgfpathlineto{\pgfqpoint{1.376289in}{2.857710in}}%
\pgfpathlineto{\pgfqpoint{1.377189in}{2.858596in}}%
\pgfpathlineto{\pgfqpoint{1.377789in}{2.859299in}}%
\pgfpathlineto{\pgfqpoint{1.378689in}{2.860184in}}%
\pgfpathlineto{\pgfqpoint{1.379289in}{2.861017in}}%
\pgfpathlineto{\pgfqpoint{1.380188in}{2.861902in}}%
\pgfpathlineto{\pgfqpoint{1.380788in}{2.862683in}}%
\pgfpathlineto{\pgfqpoint{1.381988in}{2.863621in}}%
\pgfpathlineto{\pgfqpoint{1.382588in}{2.864219in}}%
\pgfpathlineto{\pgfqpoint{1.383188in}{2.865053in}}%
\pgfpathlineto{\pgfqpoint{1.383787in}{2.865782in}}%
\pgfpathlineto{\pgfqpoint{1.384687in}{2.866667in}}%
\pgfpathlineto{\pgfqpoint{1.385287in}{2.867604in}}%
\pgfpathlineto{\pgfqpoint{1.386337in}{2.868515in}}%
\pgfpathlineto{\pgfqpoint{1.386937in}{2.869349in}}%
\pgfpathlineto{\pgfqpoint{1.387986in}{2.870260in}}%
\pgfpathlineto{\pgfqpoint{1.388586in}{2.871457in}}%
\pgfpathlineto{\pgfqpoint{1.389336in}{2.872317in}}%
\pgfpathlineto{\pgfqpoint{1.389936in}{2.873020in}}%
\pgfpathlineto{\pgfqpoint{1.390836in}{2.873905in}}%
\pgfpathlineto{\pgfqpoint{1.391436in}{2.874452in}}%
\pgfpathlineto{\pgfqpoint{1.392185in}{2.875311in}}%
\pgfpathlineto{\pgfqpoint{1.392785in}{2.876222in}}%
\pgfpathlineto{\pgfqpoint{1.393835in}{2.877133in}}%
\pgfpathlineto{\pgfqpoint{1.394435in}{2.877784in}}%
\pgfpathlineto{\pgfqpoint{1.395035in}{2.878617in}}%
\pgfpathlineto{\pgfqpoint{1.395635in}{2.879216in}}%
\pgfpathlineto{\pgfqpoint{1.396684in}{2.880128in}}%
\pgfpathlineto{\pgfqpoint{1.397284in}{2.880805in}}%
\pgfpathlineto{\pgfqpoint{1.398184in}{2.881690in}}%
\pgfpathlineto{\pgfqpoint{1.398784in}{2.882184in}}%
\pgfpathlineto{\pgfqpoint{1.399534in}{2.883044in}}%
\pgfpathlineto{\pgfqpoint{1.400134in}{2.883564in}}%
\pgfpathlineto{\pgfqpoint{1.400883in}{2.884424in}}%
\pgfpathlineto{\pgfqpoint{1.401483in}{2.884970in}}%
\pgfpathlineto{\pgfqpoint{1.402233in}{2.885830in}}%
\pgfpathlineto{\pgfqpoint{1.402833in}{2.886715in}}%
\pgfpathlineto{\pgfqpoint{1.403583in}{2.887574in}}%
\pgfpathlineto{\pgfqpoint{1.404183in}{2.888121in}}%
\pgfpathlineto{\pgfqpoint{1.405382in}{2.889058in}}%
\pgfpathlineto{\pgfqpoint{1.405982in}{2.889891in}}%
\pgfpathlineto{\pgfqpoint{1.406882in}{2.890776in}}%
\pgfpathlineto{\pgfqpoint{1.407482in}{2.891427in}}%
\pgfpathlineto{\pgfqpoint{1.408532in}{2.892339in}}%
\pgfpathlineto{\pgfqpoint{1.409131in}{2.893172in}}%
\pgfpathlineto{\pgfqpoint{1.410031in}{2.894057in}}%
\pgfpathlineto{\pgfqpoint{1.410631in}{2.894526in}}%
\pgfpathlineto{\pgfqpoint{1.411681in}{2.895437in}}%
\pgfpathlineto{\pgfqpoint{1.412281in}{2.895984in}}%
\pgfpathlineto{\pgfqpoint{1.413180in}{2.896869in}}%
\pgfpathlineto{\pgfqpoint{1.413780in}{2.897546in}}%
\pgfpathlineto{\pgfqpoint{1.414980in}{2.898483in}}%
\pgfpathlineto{\pgfqpoint{1.415580in}{2.899056in}}%
\pgfpathlineto{\pgfqpoint{1.416480in}{2.899941in}}%
\pgfpathlineto{\pgfqpoint{1.417080in}{2.900514in}}%
\pgfpathlineto{\pgfqpoint{1.417829in}{2.901373in}}%
\pgfpathlineto{\pgfqpoint{1.418429in}{2.901868in}}%
\pgfpathlineto{\pgfqpoint{1.419629in}{2.902805in}}%
\pgfpathlineto{\pgfqpoint{1.420229in}{2.903534in}}%
\pgfpathlineto{\pgfqpoint{1.421279in}{2.904446in}}%
\pgfpathlineto{\pgfqpoint{1.421878in}{2.905070in}}%
\pgfpathlineto{\pgfqpoint{1.423528in}{2.906086in}}%
\pgfpathlineto{\pgfqpoint{1.424128in}{2.906554in}}%
\pgfpathlineto{\pgfqpoint{1.424878in}{2.907414in}}%
\pgfpathlineto{\pgfqpoint{1.425478in}{2.907986in}}%
\pgfpathlineto{\pgfqpoint{1.426827in}{2.908950in}}%
\pgfpathlineto{\pgfqpoint{1.427427in}{2.909705in}}%
\pgfpathlineto{\pgfqpoint{1.428777in}{2.910668in}}%
\pgfpathlineto{\pgfqpoint{1.429377in}{2.911111in}}%
\pgfpathlineto{\pgfqpoint{1.430426in}{2.912022in}}%
\pgfpathlineto{\pgfqpoint{1.431026in}{2.912647in}}%
\pgfpathlineto{\pgfqpoint{1.432226in}{2.913584in}}%
\pgfpathlineto{\pgfqpoint{1.432826in}{2.914053in}}%
\pgfpathlineto{\pgfqpoint{1.433726in}{2.914938in}}%
\pgfpathlineto{\pgfqpoint{1.434325in}{2.915511in}}%
\pgfpathlineto{\pgfqpoint{1.435825in}{2.916500in}}%
\pgfpathlineto{\pgfqpoint{1.436425in}{2.916969in}}%
\pgfpathlineto{\pgfqpoint{1.437475in}{2.917880in}}%
\pgfpathlineto{\pgfqpoint{1.438074in}{2.918401in}}%
\pgfpathlineto{\pgfqpoint{1.439274in}{2.919338in}}%
\pgfpathlineto{\pgfqpoint{1.439874in}{2.919885in}}%
\pgfpathlineto{\pgfqpoint{1.441524in}{2.920900in}}%
\pgfpathlineto{\pgfqpoint{1.442124in}{2.921630in}}%
\pgfpathlineto{\pgfqpoint{1.443923in}{2.922671in}}%
\pgfpathlineto{\pgfqpoint{1.444523in}{2.923270in}}%
\pgfpathlineto{\pgfqpoint{1.446023in}{2.924259in}}%
\pgfpathlineto{\pgfqpoint{1.446622in}{2.925014in}}%
\pgfpathlineto{\pgfqpoint{1.448272in}{2.926030in}}%
\pgfpathlineto{\pgfqpoint{1.448872in}{2.926394in}}%
\pgfpathlineto{\pgfqpoint{1.450671in}{2.927436in}}%
\pgfpathlineto{\pgfqpoint{1.451271in}{2.927982in}}%
\pgfpathlineto{\pgfqpoint{1.452471in}{2.928920in}}%
\pgfpathlineto{\pgfqpoint{1.453071in}{2.929388in}}%
\pgfpathlineto{\pgfqpoint{1.454271in}{2.930326in}}%
\pgfpathlineto{\pgfqpoint{1.454870in}{2.931081in}}%
\pgfpathlineto{\pgfqpoint{1.456970in}{2.932174in}}%
\pgfpathlineto{\pgfqpoint{1.457570in}{2.932825in}}%
\pgfpathlineto{\pgfqpoint{1.458770in}{2.933762in}}%
\pgfpathlineto{\pgfqpoint{1.459369in}{2.934101in}}%
\pgfpathlineto{\pgfqpoint{1.461619in}{2.935220in}}%
\pgfpathlineto{\pgfqpoint{1.462219in}{2.935741in}}%
\pgfpathlineto{\pgfqpoint{1.463418in}{2.936679in}}%
\pgfpathlineto{\pgfqpoint{1.464018in}{2.936965in}}%
\pgfpathlineto{\pgfqpoint{1.466268in}{2.938084in}}%
\pgfpathlineto{\pgfqpoint{1.466868in}{2.938683in}}%
\pgfpathlineto{\pgfqpoint{1.468667in}{2.939725in}}%
\pgfpathlineto{\pgfqpoint{1.469267in}{2.940219in}}%
\pgfpathlineto{\pgfqpoint{1.471367in}{2.941313in}}%
\pgfpathlineto{\pgfqpoint{1.471966in}{2.941625in}}%
\pgfpathlineto{\pgfqpoint{1.473916in}{2.942693in}}%
\pgfpathlineto{\pgfqpoint{1.474516in}{2.943083in}}%
\pgfpathlineto{\pgfqpoint{1.476165in}{2.944099in}}%
\pgfpathlineto{\pgfqpoint{1.476765in}{2.944594in}}%
\pgfpathlineto{\pgfqpoint{1.478565in}{2.945635in}}%
\pgfpathlineto{\pgfqpoint{1.479165in}{2.946130in}}%
\pgfpathlineto{\pgfqpoint{1.480964in}{2.947171in}}%
\pgfpathlineto{\pgfqpoint{1.481564in}{2.947718in}}%
\pgfpathlineto{\pgfqpoint{1.483814in}{2.948837in}}%
\pgfpathlineto{\pgfqpoint{1.484413in}{2.949280in}}%
\pgfpathlineto{\pgfqpoint{1.486063in}{2.950296in}}%
\pgfpathlineto{\pgfqpoint{1.486663in}{2.950790in}}%
\pgfpathlineto{\pgfqpoint{1.488612in}{2.951858in}}%
\pgfpathlineto{\pgfqpoint{1.489212in}{2.952222in}}%
\pgfpathlineto{\pgfqpoint{1.491462in}{2.953342in}}%
\pgfpathlineto{\pgfqpoint{1.492062in}{2.953889in}}%
\pgfpathlineto{\pgfqpoint{1.495361in}{2.955190in}}%
\pgfpathlineto{\pgfqpoint{1.495961in}{2.955607in}}%
\pgfpathlineto{\pgfqpoint{1.497760in}{2.956648in}}%
\pgfpathlineto{\pgfqpoint{1.498360in}{2.956961in}}%
\pgfpathlineto{\pgfqpoint{1.499860in}{2.957950in}}%
\pgfpathlineto{\pgfqpoint{1.500460in}{2.958158in}}%
\pgfpathlineto{\pgfqpoint{1.502859in}{2.959304in}}%
\pgfpathlineto{\pgfqpoint{1.503459in}{2.959590in}}%
\pgfpathlineto{\pgfqpoint{1.505558in}{2.960684in}}%
\pgfpathlineto{\pgfqpoint{1.506158in}{2.961022in}}%
\pgfpathlineto{\pgfqpoint{1.509158in}{2.962272in}}%
\pgfpathlineto{\pgfqpoint{1.509607in}{2.962481in}}%
\pgfpathlineto{\pgfqpoint{1.511857in}{2.963418in}}%
\pgfpathlineto{\pgfqpoint{1.513806in}{2.964485in}}%
\pgfpathlineto{\pgfqpoint{1.514406in}{2.965084in}}%
\pgfpathlineto{\pgfqpoint{1.516806in}{2.966230in}}%
\pgfpathlineto{\pgfqpoint{1.517406in}{2.966464in}}%
\pgfpathlineto{\pgfqpoint{1.519505in}{2.967558in}}%
\pgfpathlineto{\pgfqpoint{1.520105in}{2.967922in}}%
\pgfpathlineto{\pgfqpoint{1.522504in}{2.969068in}}%
\pgfpathlineto{\pgfqpoint{1.523104in}{2.969432in}}%
\pgfpathlineto{\pgfqpoint{1.525354in}{2.970552in}}%
\pgfpathlineto{\pgfqpoint{1.525954in}{2.970708in}}%
\pgfpathlineto{\pgfqpoint{1.529103in}{2.971984in}}%
\pgfpathlineto{\pgfqpoint{1.529703in}{2.972296in}}%
\pgfpathlineto{\pgfqpoint{1.533452in}{2.973676in}}%
\pgfpathlineto{\pgfqpoint{1.534052in}{2.973806in}}%
\pgfpathlineto{\pgfqpoint{1.535851in}{2.974848in}}%
\pgfpathlineto{\pgfqpoint{1.536451in}{2.975186in}}%
\pgfpathlineto{\pgfqpoint{1.539750in}{2.976488in}}%
\pgfpathlineto{\pgfqpoint{1.540350in}{2.976748in}}%
\pgfpathlineto{\pgfqpoint{1.542600in}{2.977868in}}%
\pgfpathlineto{\pgfqpoint{1.543199in}{2.978232in}}%
\pgfpathlineto{\pgfqpoint{1.545299in}{2.979326in}}%
\pgfpathlineto{\pgfqpoint{1.545899in}{2.979612in}}%
\pgfpathlineto{\pgfqpoint{1.551447in}{2.981305in}}%
\pgfpathlineto{\pgfqpoint{1.552047in}{2.981643in}}%
\pgfpathlineto{\pgfqpoint{1.555197in}{2.982919in}}%
\pgfpathlineto{\pgfqpoint{1.555796in}{2.983258in}}%
\pgfpathlineto{\pgfqpoint{1.559845in}{2.984690in}}%
\pgfpathlineto{\pgfqpoint{1.560295in}{2.984846in}}%
\pgfpathlineto{\pgfqpoint{1.563445in}{2.985861in}}%
\pgfpathlineto{\pgfqpoint{1.564044in}{2.986174in}}%
\pgfpathlineto{\pgfqpoint{1.569593in}{2.987866in}}%
\pgfpathlineto{\pgfqpoint{1.570043in}{2.988100in}}%
\pgfpathlineto{\pgfqpoint{1.572892in}{2.989090in}}%
\pgfpathlineto{\pgfqpoint{1.573492in}{2.989350in}}%
\pgfpathlineto{\pgfqpoint{1.579191in}{2.991068in}}%
\pgfpathlineto{\pgfqpoint{1.579791in}{2.991277in}}%
\pgfpathlineto{\pgfqpoint{1.584590in}{2.992839in}}%
\pgfpathlineto{\pgfqpoint{1.585189in}{2.993229in}}%
\pgfpathlineto{\pgfqpoint{1.589838in}{2.994375in}}%
\pgfpathlineto{\pgfqpoint{1.590438in}{2.994531in}}%
\pgfpathlineto{\pgfqpoint{1.595237in}{2.996093in}}%
\pgfpathlineto{\pgfqpoint{1.595837in}{2.996406in}}%
\pgfpathlineto{\pgfqpoint{1.600186in}{2.997525in}}%
\pgfpathlineto{\pgfqpoint{1.600636in}{2.997630in}}%
\pgfpathlineto{\pgfqpoint{1.604235in}{2.998671in}}%
\pgfpathlineto{\pgfqpoint{1.604685in}{2.998931in}}%
\pgfpathlineto{\pgfqpoint{1.608284in}{2.999869in}}%
\pgfpathlineto{\pgfqpoint{1.608884in}{3.000051in}}%
\pgfpathlineto{\pgfqpoint{1.619081in}{3.002550in}}%
\pgfpathlineto{\pgfqpoint{1.619531in}{3.002733in}}%
\pgfpathlineto{\pgfqpoint{1.624030in}{3.003852in}}%
\pgfpathlineto{\pgfqpoint{1.624630in}{3.004034in}}%
\pgfpathlineto{\pgfqpoint{1.634378in}{3.006456in}}%
\pgfpathlineto{\pgfqpoint{1.634978in}{3.006690in}}%
\pgfpathlineto{\pgfqpoint{1.649524in}{3.009945in}}%
\pgfpathlineto{\pgfqpoint{1.649974in}{3.010049in}}%
\pgfpathlineto{\pgfqpoint{1.657022in}{3.011403in}}%
\pgfpathlineto{\pgfqpoint{1.657472in}{3.011481in}}%
\pgfpathlineto{\pgfqpoint{1.661521in}{3.012444in}}%
\pgfpathlineto{\pgfqpoint{1.662121in}{3.012548in}}%
\pgfpathlineto{\pgfqpoint{1.689265in}{3.017990in}}%
\pgfpathlineto{\pgfqpoint{1.689864in}{3.018120in}}%
\pgfpathlineto{\pgfqpoint{1.734404in}{3.025202in}}%
\pgfpathlineto{\pgfqpoint{1.742652in}{3.026217in}}%
\pgfpathlineto{\pgfqpoint{1.757498in}{3.028248in}}%
\pgfpathlineto{\pgfqpoint{1.757648in}{3.028326in}}%
\pgfpathlineto{\pgfqpoint{1.765146in}{3.029238in}}%
\pgfpathlineto{\pgfqpoint{1.776694in}{3.030982in}}%
\pgfpathlineto{\pgfqpoint{1.785392in}{3.031893in}}%
\pgfpathlineto{\pgfqpoint{1.791840in}{3.032753in}}%
\pgfpathlineto{\pgfqpoint{1.809386in}{3.035096in}}%
\pgfpathlineto{\pgfqpoint{1.823633in}{3.036320in}}%
\pgfpathlineto{\pgfqpoint{1.830981in}{3.036892in}}%
\pgfpathlineto{\pgfqpoint{1.845977in}{3.038272in}}%
\pgfpathlineto{\pgfqpoint{1.846427in}{3.038376in}}%
\pgfpathlineto{\pgfqpoint{1.858274in}{3.039522in}}%
\pgfpathlineto{\pgfqpoint{1.943454in}{3.045172in}}%
\pgfpathlineto{\pgfqpoint{1.944054in}{3.045224in}}%
\pgfpathlineto{\pgfqpoint{1.964449in}{3.046344in}}%
\pgfpathlineto{\pgfqpoint{1.983494in}{3.047333in}}%
\pgfpathlineto{\pgfqpoint{2.007939in}{3.048791in}}%
\pgfpathlineto{\pgfqpoint{2.024285in}{3.049754in}}%
\pgfpathlineto{\pgfqpoint{2.176498in}{3.057852in}}%
\pgfpathlineto{\pgfqpoint{2.209790in}{3.058997in}}%
\pgfpathlineto{\pgfqpoint{2.221938in}{3.059362in}}%
\pgfpathlineto{\pgfqpoint{2.242033in}{3.060377in}}%
\pgfpathlineto{\pgfqpoint{2.247581in}{3.060611in}}%
\pgfpathlineto{\pgfqpoint{2.359455in}{3.064569in}}%
\pgfpathlineto{\pgfqpoint{2.359455in}{3.064595in}}%
\pgfpathlineto{\pgfqpoint{2.382399in}{3.065532in}}%
\pgfpathlineto{\pgfqpoint{2.417941in}{3.067042in}}%
\pgfpathlineto{\pgfqpoint{2.432637in}{3.067719in}}%
\pgfpathlineto{\pgfqpoint{2.460681in}{3.068657in}}%
\pgfpathlineto{\pgfqpoint{2.490973in}{3.069698in}}%
\pgfpathlineto{\pgfqpoint{2.649785in}{3.073187in}}%
\pgfpathlineto{\pgfqpoint{2.764508in}{3.077223in}}%
\pgfpathlineto{\pgfqpoint{2.764808in}{3.077275in}}%
\pgfpathlineto{\pgfqpoint{2.846089in}{3.081519in}}%
\pgfpathlineto{\pgfqpoint{2.846689in}{3.081571in}}%
\pgfpathlineto{\pgfqpoint{2.887479in}{3.084070in}}%
\pgfpathlineto{\pgfqpoint{2.911023in}{3.085372in}}%
\pgfpathlineto{\pgfqpoint{3.216350in}{3.097974in}}%
\pgfpathlineto{\pgfqpoint{3.244994in}{3.098989in}}%
\pgfpathlineto{\pgfqpoint{3.287883in}{3.100759in}}%
\pgfpathlineto{\pgfqpoint{3.600859in}{3.107841in}}%
\pgfpathlineto{\pgfqpoint{3.660694in}{3.108987in}}%
\pgfpathlineto{\pgfqpoint{3.671492in}{3.109195in}}%
\pgfpathlineto{\pgfqpoint{3.710482in}{3.109898in}}%
\pgfpathlineto{\pgfqpoint{4.227409in}{3.113179in}}%
\pgfpathlineto{\pgfqpoint{4.962984in}{3.114897in}}%
\pgfpathlineto{\pgfqpoint{5.275792in}{3.115437in}}%
\pgfpathlineto{\pgfqpoint{5.275792in}{3.115437in}}%
\pgfusepath{stroke}%
\end{pgfscope}%
\begin{pgfscope}%
\pgfpathrectangle{\pgfqpoint{0.565433in}{0.528318in}}{\pgfqpoint{4.703693in}{2.603633in}}%
\pgfusepath{clip}%
\pgfsetrectcap%
\pgfsetroundjoin%
\pgfsetlinewidth{1.806750pt}%
\definecolor{currentstroke}{rgb}{0.180392,0.490196,0.196078}%
\pgfsetstrokecolor{currentstroke}%
\pgfsetdash{}{0pt}%
\pgfpathmoveto{\pgfqpoint{1.608884in}{0.528344in}}%
\pgfpathlineto{\pgfqpoint{1.744601in}{0.531025in}}%
\pgfpathlineto{\pgfqpoint{1.750300in}{0.531338in}}%
\pgfpathlineto{\pgfqpoint{1.816734in}{0.534957in}}%
\pgfpathlineto{\pgfqpoint{1.833380in}{0.537144in}}%
\pgfpathlineto{\pgfqpoint{1.833380in}{0.537170in}}%
\pgfpathlineto{\pgfqpoint{1.842978in}{0.538237in}}%
\pgfpathlineto{\pgfqpoint{1.946603in}{0.556958in}}%
\pgfpathlineto{\pgfqpoint{1.947203in}{0.557166in}}%
\pgfpathlineto{\pgfqpoint{1.954101in}{0.559093in}}%
\pgfpathlineto{\pgfqpoint{1.954551in}{0.559275in}}%
\pgfpathlineto{\pgfqpoint{1.959950in}{0.560473in}}%
\pgfpathlineto{\pgfqpoint{1.960550in}{0.560759in}}%
\pgfpathlineto{\pgfqpoint{1.965799in}{0.562399in}}%
\pgfpathlineto{\pgfqpoint{1.966398in}{0.562712in}}%
\pgfpathlineto{\pgfqpoint{1.970597in}{0.564170in}}%
\pgfpathlineto{\pgfqpoint{1.971197in}{0.564404in}}%
\pgfpathlineto{\pgfqpoint{1.978096in}{0.566331in}}%
\pgfpathlineto{\pgfqpoint{1.978546in}{0.566487in}}%
\pgfpathlineto{\pgfqpoint{1.981545in}{0.567476in}}%
\pgfpathlineto{\pgfqpoint{1.982145in}{0.567763in}}%
\pgfpathlineto{\pgfqpoint{1.987843in}{0.569481in}}%
\pgfpathlineto{\pgfqpoint{1.988443in}{0.569715in}}%
\pgfpathlineto{\pgfqpoint{1.992792in}{0.570835in}}%
\pgfpathlineto{\pgfqpoint{1.993392in}{0.571043in}}%
\pgfpathlineto{\pgfqpoint{2.015737in}{0.575652in}}%
\pgfpathlineto{\pgfqpoint{2.016337in}{0.575808in}}%
\pgfpathlineto{\pgfqpoint{2.041830in}{0.580963in}}%
\pgfpathlineto{\pgfqpoint{2.042430in}{0.581093in}}%
\pgfpathlineto{\pgfqpoint{2.049479in}{0.583046in}}%
\pgfpathlineto{\pgfqpoint{2.050079in}{0.583228in}}%
\pgfpathlineto{\pgfqpoint{2.055027in}{0.584816in}}%
\pgfpathlineto{\pgfqpoint{2.055627in}{0.585103in}}%
\pgfpathlineto{\pgfqpoint{2.059976in}{0.586587in}}%
\pgfpathlineto{\pgfqpoint{2.060576in}{0.586743in}}%
\pgfpathlineto{\pgfqpoint{2.063275in}{0.587941in}}%
\pgfpathlineto{\pgfqpoint{2.063875in}{0.588462in}}%
\pgfpathlineto{\pgfqpoint{2.067774in}{0.589868in}}%
\pgfpathlineto{\pgfqpoint{2.068374in}{0.590206in}}%
\pgfpathlineto{\pgfqpoint{2.071523in}{0.591482in}}%
\pgfpathlineto{\pgfqpoint{2.072123in}{0.591768in}}%
\pgfpathlineto{\pgfqpoint{2.076172in}{0.593200in}}%
\pgfpathlineto{\pgfqpoint{2.076772in}{0.593565in}}%
\pgfpathlineto{\pgfqpoint{2.078572in}{0.594606in}}%
\pgfpathlineto{\pgfqpoint{2.079172in}{0.594919in}}%
\pgfpathlineto{\pgfqpoint{2.082171in}{0.596168in}}%
\pgfpathlineto{\pgfqpoint{2.082771in}{0.596455in}}%
\pgfpathlineto{\pgfqpoint{2.084420in}{0.597470in}}%
\pgfpathlineto{\pgfqpoint{2.085020in}{0.597783in}}%
\pgfpathlineto{\pgfqpoint{2.086520in}{0.598772in}}%
\pgfpathlineto{\pgfqpoint{2.087120in}{0.599136in}}%
\pgfpathlineto{\pgfqpoint{2.088469in}{0.600100in}}%
\pgfpathlineto{\pgfqpoint{2.089069in}{0.600438in}}%
\pgfpathlineto{\pgfqpoint{2.091019in}{0.601506in}}%
\pgfpathlineto{\pgfqpoint{2.091619in}{0.602000in}}%
\pgfpathlineto{\pgfqpoint{2.093718in}{0.603094in}}%
\pgfpathlineto{\pgfqpoint{2.094318in}{0.603641in}}%
\pgfpathlineto{\pgfqpoint{2.095668in}{0.604604in}}%
\pgfpathlineto{\pgfqpoint{2.096268in}{0.605047in}}%
\pgfpathlineto{\pgfqpoint{2.098067in}{0.606088in}}%
\pgfpathlineto{\pgfqpoint{2.098667in}{0.606531in}}%
\pgfpathlineto{\pgfqpoint{2.100017in}{0.607494in}}%
\pgfpathlineto{\pgfqpoint{2.100616in}{0.607859in}}%
\pgfpathlineto{\pgfqpoint{2.102116in}{0.608848in}}%
\pgfpathlineto{\pgfqpoint{2.102716in}{0.609421in}}%
\pgfpathlineto{\pgfqpoint{2.103766in}{0.610332in}}%
\pgfpathlineto{\pgfqpoint{2.104366in}{0.610853in}}%
\pgfpathlineto{\pgfqpoint{2.106015in}{0.611868in}}%
\pgfpathlineto{\pgfqpoint{2.106615in}{0.612441in}}%
\pgfpathlineto{\pgfqpoint{2.108265in}{0.613456in}}%
\pgfpathlineto{\pgfqpoint{2.108865in}{0.613873in}}%
\pgfpathlineto{\pgfqpoint{2.110664in}{0.614914in}}%
\pgfpathlineto{\pgfqpoint{2.111264in}{0.615383in}}%
\pgfpathlineto{\pgfqpoint{2.112464in}{0.616320in}}%
\pgfpathlineto{\pgfqpoint{2.113064in}{0.616711in}}%
\pgfpathlineto{\pgfqpoint{2.114113in}{0.617622in}}%
\pgfpathlineto{\pgfqpoint{2.114713in}{0.618065in}}%
\pgfpathlineto{\pgfqpoint{2.116213in}{0.619054in}}%
\pgfpathlineto{\pgfqpoint{2.116813in}{0.619783in}}%
\pgfpathlineto{\pgfqpoint{2.118012in}{0.620721in}}%
\pgfpathlineto{\pgfqpoint{2.118612in}{0.621241in}}%
\pgfpathlineto{\pgfqpoint{2.120262in}{0.622257in}}%
\pgfpathlineto{\pgfqpoint{2.120862in}{0.622830in}}%
\pgfpathlineto{\pgfqpoint{2.121761in}{0.623715in}}%
\pgfpathlineto{\pgfqpoint{2.122361in}{0.624262in}}%
\pgfpathlineto{\pgfqpoint{2.123861in}{0.625251in}}%
\pgfpathlineto{\pgfqpoint{2.124461in}{0.625876in}}%
\pgfpathlineto{\pgfqpoint{2.125660in}{0.626813in}}%
\pgfpathlineto{\pgfqpoint{2.126260in}{0.627672in}}%
\pgfpathlineto{\pgfqpoint{2.127460in}{0.628610in}}%
\pgfpathlineto{\pgfqpoint{2.128060in}{0.629156in}}%
\pgfpathlineto{\pgfqpoint{2.129710in}{0.630172in}}%
\pgfpathlineto{\pgfqpoint{2.130309in}{0.630771in}}%
\pgfpathlineto{\pgfqpoint{2.131809in}{0.631760in}}%
\pgfpathlineto{\pgfqpoint{2.132409in}{0.632437in}}%
\pgfpathlineto{\pgfqpoint{2.134058in}{0.633452in}}%
\pgfpathlineto{\pgfqpoint{2.134658in}{0.633921in}}%
\pgfpathlineto{\pgfqpoint{2.135558in}{0.634806in}}%
\pgfpathlineto{\pgfqpoint{2.136158in}{0.635223in}}%
\pgfpathlineto{\pgfqpoint{2.137508in}{0.636186in}}%
\pgfpathlineto{\pgfqpoint{2.138108in}{0.636811in}}%
\pgfpathlineto{\pgfqpoint{2.139157in}{0.637722in}}%
\pgfpathlineto{\pgfqpoint{2.139757in}{0.638399in}}%
\pgfpathlineto{\pgfqpoint{2.140657in}{0.639284in}}%
\pgfpathlineto{\pgfqpoint{2.141257in}{0.639701in}}%
\pgfpathlineto{\pgfqpoint{2.142307in}{0.640612in}}%
\pgfpathlineto{\pgfqpoint{2.142906in}{0.641159in}}%
\pgfpathlineto{\pgfqpoint{2.143956in}{0.642070in}}%
\pgfpathlineto{\pgfqpoint{2.144556in}{0.642878in}}%
\pgfpathlineto{\pgfqpoint{2.145756in}{0.643815in}}%
\pgfpathlineto{\pgfqpoint{2.146356in}{0.644336in}}%
\pgfpathlineto{\pgfqpoint{2.147705in}{0.645299in}}%
\pgfpathlineto{\pgfqpoint{2.148305in}{0.646158in}}%
\pgfpathlineto{\pgfqpoint{2.149505in}{0.647095in}}%
\pgfpathlineto{\pgfqpoint{2.150105in}{0.647590in}}%
\pgfpathlineto{\pgfqpoint{2.151154in}{0.648501in}}%
\pgfpathlineto{\pgfqpoint{2.151754in}{0.649022in}}%
\pgfpathlineto{\pgfqpoint{2.152654in}{0.649907in}}%
\pgfpathlineto{\pgfqpoint{2.153254in}{0.650767in}}%
\pgfpathlineto{\pgfqpoint{2.154454in}{0.651704in}}%
\pgfpathlineto{\pgfqpoint{2.155053in}{0.652433in}}%
\pgfpathlineto{\pgfqpoint{2.155953in}{0.653318in}}%
\pgfpathlineto{\pgfqpoint{2.156553in}{0.654073in}}%
\pgfpathlineto{\pgfqpoint{2.157153in}{0.654906in}}%
\pgfpathlineto{\pgfqpoint{2.157753in}{0.655818in}}%
\pgfpathlineto{\pgfqpoint{2.158653in}{0.656703in}}%
\pgfpathlineto{\pgfqpoint{2.159252in}{0.657458in}}%
\pgfpathlineto{\pgfqpoint{2.160602in}{0.658421in}}%
\pgfpathlineto{\pgfqpoint{2.161202in}{0.659150in}}%
\pgfpathlineto{\pgfqpoint{2.161802in}{0.659983in}}%
\pgfpathlineto{\pgfqpoint{2.162402in}{0.661051in}}%
\pgfpathlineto{\pgfqpoint{2.163302in}{0.661936in}}%
\pgfpathlineto{\pgfqpoint{2.163901in}{0.662639in}}%
\pgfpathlineto{\pgfqpoint{2.164801in}{0.663524in}}%
\pgfpathlineto{\pgfqpoint{2.165401in}{0.664305in}}%
\pgfpathlineto{\pgfqpoint{2.166301in}{0.665191in}}%
\pgfpathlineto{\pgfqpoint{2.166901in}{0.665894in}}%
\pgfpathlineto{\pgfqpoint{2.167950in}{0.666805in}}%
\pgfpathlineto{\pgfqpoint{2.168550in}{0.667742in}}%
\pgfpathlineto{\pgfqpoint{2.169150in}{0.668575in}}%
\pgfpathlineto{\pgfqpoint{2.169750in}{0.669070in}}%
\pgfpathlineto{\pgfqpoint{2.170500in}{0.669929in}}%
\pgfpathlineto{\pgfqpoint{2.171100in}{0.671075in}}%
\pgfpathlineto{\pgfqpoint{2.171999in}{0.671960in}}%
\pgfpathlineto{\pgfqpoint{2.172599in}{0.672923in}}%
\pgfpathlineto{\pgfqpoint{2.173499in}{0.673809in}}%
\pgfpathlineto{\pgfqpoint{2.174099in}{0.674486in}}%
\pgfpathlineto{\pgfqpoint{2.175149in}{0.675397in}}%
\pgfpathlineto{\pgfqpoint{2.175749in}{0.676204in}}%
\pgfpathlineto{\pgfqpoint{2.176498in}{0.677063in}}%
\pgfpathlineto{\pgfqpoint{2.177098in}{0.678027in}}%
\pgfpathlineto{\pgfqpoint{2.177698in}{0.678860in}}%
\pgfpathlineto{\pgfqpoint{2.178298in}{0.679797in}}%
\pgfpathlineto{\pgfqpoint{2.179048in}{0.680656in}}%
\pgfpathlineto{\pgfqpoint{2.179648in}{0.681307in}}%
\pgfpathlineto{\pgfqpoint{2.180098in}{0.682114in}}%
\pgfpathlineto{\pgfqpoint{2.180697in}{0.683104in}}%
\pgfpathlineto{\pgfqpoint{2.181297in}{0.683937in}}%
\pgfpathlineto{\pgfqpoint{2.181897in}{0.684822in}}%
\pgfpathlineto{\pgfqpoint{2.182497in}{0.685655in}}%
\pgfpathlineto{\pgfqpoint{2.183097in}{0.686957in}}%
\pgfpathlineto{\pgfqpoint{2.183697in}{0.687790in}}%
\pgfpathlineto{\pgfqpoint{2.184297in}{0.688754in}}%
\pgfpathlineto{\pgfqpoint{2.184896in}{0.689587in}}%
\pgfpathlineto{\pgfqpoint{2.185496in}{0.690889in}}%
\pgfpathlineto{\pgfqpoint{2.185946in}{0.691696in}}%
\pgfpathlineto{\pgfqpoint{2.186546in}{0.692789in}}%
\pgfpathlineto{\pgfqpoint{2.187146in}{0.693622in}}%
\pgfpathlineto{\pgfqpoint{2.187746in}{0.694664in}}%
\pgfpathlineto{\pgfqpoint{2.188346in}{0.695497in}}%
\pgfpathlineto{\pgfqpoint{2.188945in}{0.696669in}}%
\pgfpathlineto{\pgfqpoint{2.189695in}{0.697528in}}%
\pgfpathlineto{\pgfqpoint{2.190295in}{0.698621in}}%
\pgfpathlineto{\pgfqpoint{2.190895in}{0.699454in}}%
\pgfpathlineto{\pgfqpoint{2.191495in}{0.700756in}}%
\pgfpathlineto{\pgfqpoint{2.192095in}{0.701589in}}%
\pgfpathlineto{\pgfqpoint{2.192695in}{0.703152in}}%
\pgfpathlineto{\pgfqpoint{2.193294in}{0.703985in}}%
\pgfpathlineto{\pgfqpoint{2.193894in}{0.705052in}}%
\pgfpathlineto{\pgfqpoint{2.194344in}{0.705859in}}%
\pgfpathlineto{\pgfqpoint{2.194944in}{0.706927in}}%
\pgfpathlineto{\pgfqpoint{2.195544in}{0.707760in}}%
\pgfpathlineto{\pgfqpoint{2.196144in}{0.708958in}}%
\pgfpathlineto{\pgfqpoint{2.196744in}{0.709791in}}%
\pgfpathlineto{\pgfqpoint{2.197343in}{0.711171in}}%
\pgfpathlineto{\pgfqpoint{2.197793in}{0.711978in}}%
\pgfpathlineto{\pgfqpoint{2.198393in}{0.713384in}}%
\pgfpathlineto{\pgfqpoint{2.198993in}{0.714217in}}%
\pgfpathlineto{\pgfqpoint{2.199593in}{0.715701in}}%
\pgfpathlineto{\pgfqpoint{2.200193in}{0.716534in}}%
\pgfpathlineto{\pgfqpoint{2.200793in}{0.717680in}}%
\pgfpathlineto{\pgfqpoint{2.201242in}{0.718487in}}%
\pgfpathlineto{\pgfqpoint{2.201842in}{0.719528in}}%
\pgfpathlineto{\pgfqpoint{2.202292in}{0.720336in}}%
\pgfpathlineto{\pgfqpoint{2.202892in}{0.721768in}}%
\pgfpathlineto{\pgfqpoint{2.203342in}{0.722575in}}%
\pgfpathlineto{\pgfqpoint{2.203942in}{0.723720in}}%
\pgfpathlineto{\pgfqpoint{2.204542in}{0.724553in}}%
\pgfpathlineto{\pgfqpoint{2.205142in}{0.726090in}}%
\pgfpathlineto{\pgfqpoint{2.205591in}{0.726897in}}%
\pgfpathlineto{\pgfqpoint{2.206191in}{0.728459in}}%
\pgfpathlineto{\pgfqpoint{2.206641in}{0.729266in}}%
\pgfpathlineto{\pgfqpoint{2.207241in}{0.730516in}}%
\pgfpathlineto{\pgfqpoint{2.207841in}{0.731349in}}%
\pgfpathlineto{\pgfqpoint{2.208441in}{0.733067in}}%
\pgfpathlineto{\pgfqpoint{2.208741in}{0.733848in}}%
\pgfpathlineto{\pgfqpoint{2.209341in}{0.735359in}}%
\pgfpathlineto{\pgfqpoint{2.209790in}{0.736166in}}%
\pgfpathlineto{\pgfqpoint{2.210390in}{0.737884in}}%
\pgfpathlineto{\pgfqpoint{2.210840in}{0.738691in}}%
\pgfpathlineto{\pgfqpoint{2.211440in}{0.740175in}}%
\pgfpathlineto{\pgfqpoint{2.211890in}{0.740982in}}%
\pgfpathlineto{\pgfqpoint{2.212490in}{0.742466in}}%
\pgfpathlineto{\pgfqpoint{2.212940in}{0.743274in}}%
\pgfpathlineto{\pgfqpoint{2.213540in}{0.744732in}}%
\pgfpathlineto{\pgfqpoint{2.213989in}{0.745539in}}%
\pgfpathlineto{\pgfqpoint{2.214589in}{0.747491in}}%
\pgfpathlineto{\pgfqpoint{2.214889in}{0.748273in}}%
\pgfpathlineto{\pgfqpoint{2.215489in}{0.749887in}}%
\pgfpathlineto{\pgfqpoint{2.215789in}{0.750668in}}%
\pgfpathlineto{\pgfqpoint{2.216389in}{0.752282in}}%
\pgfpathlineto{\pgfqpoint{2.216839in}{0.753089in}}%
\pgfpathlineto{\pgfqpoint{2.217439in}{0.755068in}}%
\pgfpathlineto{\pgfqpoint{2.217888in}{0.755875in}}%
\pgfpathlineto{\pgfqpoint{2.218488in}{0.757724in}}%
\pgfpathlineto{\pgfqpoint{2.218788in}{0.758505in}}%
\pgfpathlineto{\pgfqpoint{2.219388in}{0.760249in}}%
\pgfpathlineto{\pgfqpoint{2.219838in}{0.761056in}}%
\pgfpathlineto{\pgfqpoint{2.220438in}{0.762957in}}%
\pgfpathlineto{\pgfqpoint{2.220738in}{0.763738in}}%
\pgfpathlineto{\pgfqpoint{2.221338in}{0.765665in}}%
\pgfpathlineto{\pgfqpoint{2.221788in}{0.766472in}}%
\pgfpathlineto{\pgfqpoint{2.222387in}{0.768190in}}%
\pgfpathlineto{\pgfqpoint{2.222837in}{0.768998in}}%
\pgfpathlineto{\pgfqpoint{2.223437in}{0.770820in}}%
\pgfpathlineto{\pgfqpoint{2.223737in}{0.771601in}}%
\pgfpathlineto{\pgfqpoint{2.224337in}{0.773476in}}%
\pgfpathlineto{\pgfqpoint{2.224637in}{0.774257in}}%
\pgfpathlineto{\pgfqpoint{2.225237in}{0.776079in}}%
\pgfpathlineto{\pgfqpoint{2.225537in}{0.776860in}}%
\pgfpathlineto{\pgfqpoint{2.226137in}{0.778865in}}%
\pgfpathlineto{\pgfqpoint{2.226436in}{0.779646in}}%
\pgfpathlineto{\pgfqpoint{2.227036in}{0.781651in}}%
\pgfpathlineto{\pgfqpoint{2.227336in}{0.782432in}}%
\pgfpathlineto{\pgfqpoint{2.227936in}{0.784099in}}%
\pgfpathlineto{\pgfqpoint{2.228236in}{0.784880in}}%
\pgfpathlineto{\pgfqpoint{2.228836in}{0.786702in}}%
\pgfpathlineto{\pgfqpoint{2.229136in}{0.787483in}}%
\pgfpathlineto{\pgfqpoint{2.229736in}{0.789514in}}%
\pgfpathlineto{\pgfqpoint{2.230036in}{0.790295in}}%
\pgfpathlineto{\pgfqpoint{2.230635in}{0.792196in}}%
\pgfpathlineto{\pgfqpoint{2.230935in}{0.792977in}}%
\pgfpathlineto{\pgfqpoint{2.231535in}{0.794878in}}%
\pgfpathlineto{\pgfqpoint{2.231985in}{0.795685in}}%
\pgfpathlineto{\pgfqpoint{2.232585in}{0.797716in}}%
\pgfpathlineto{\pgfqpoint{2.232885in}{0.798497in}}%
\pgfpathlineto{\pgfqpoint{2.233485in}{0.800085in}}%
\pgfpathlineto{\pgfqpoint{2.233785in}{0.800866in}}%
\pgfpathlineto{\pgfqpoint{2.234385in}{0.802689in}}%
\pgfpathlineto{\pgfqpoint{2.234834in}{0.803496in}}%
\pgfpathlineto{\pgfqpoint{2.235434in}{0.805891in}}%
\pgfpathlineto{\pgfqpoint{2.235734in}{0.806672in}}%
\pgfpathlineto{\pgfqpoint{2.236334in}{0.808885in}}%
\pgfpathlineto{\pgfqpoint{2.236784in}{0.809692in}}%
\pgfpathlineto{\pgfqpoint{2.237384in}{0.811853in}}%
\pgfpathlineto{\pgfqpoint{2.237684in}{0.812634in}}%
\pgfpathlineto{\pgfqpoint{2.238284in}{0.814639in}}%
\pgfpathlineto{\pgfqpoint{2.238584in}{0.815420in}}%
\pgfpathlineto{\pgfqpoint{2.239183in}{0.817894in}}%
\pgfpathlineto{\pgfqpoint{2.239483in}{0.818675in}}%
\pgfpathlineto{\pgfqpoint{2.240083in}{0.820628in}}%
\pgfpathlineto{\pgfqpoint{2.240383in}{0.821409in}}%
\pgfpathlineto{\pgfqpoint{2.240983in}{0.824090in}}%
\pgfpathlineto{\pgfqpoint{2.241283in}{0.824871in}}%
\pgfpathlineto{\pgfqpoint{2.241883in}{0.827189in}}%
\pgfpathlineto{\pgfqpoint{2.242183in}{0.827970in}}%
\pgfpathlineto{\pgfqpoint{2.242783in}{0.830313in}}%
\pgfpathlineto{\pgfqpoint{2.243082in}{0.831094in}}%
\pgfpathlineto{\pgfqpoint{2.243682in}{0.833333in}}%
\pgfpathlineto{\pgfqpoint{2.243982in}{0.834114in}}%
\pgfpathlineto{\pgfqpoint{2.244582in}{0.836432in}}%
\pgfpathlineto{\pgfqpoint{2.244882in}{0.837213in}}%
\pgfpathlineto{\pgfqpoint{2.245482in}{0.839868in}}%
\pgfpathlineto{\pgfqpoint{2.245932in}{0.840676in}}%
\pgfpathlineto{\pgfqpoint{2.246532in}{0.843331in}}%
\pgfpathlineto{\pgfqpoint{2.246832in}{0.844112in}}%
\pgfpathlineto{\pgfqpoint{2.247431in}{0.846872in}}%
\pgfpathlineto{\pgfqpoint{2.247731in}{0.847653in}}%
\pgfpathlineto{\pgfqpoint{2.248331in}{0.850257in}}%
\pgfpathlineto{\pgfqpoint{2.248631in}{0.851038in}}%
\pgfpathlineto{\pgfqpoint{2.249231in}{0.853381in}}%
\pgfpathlineto{\pgfqpoint{2.249531in}{0.854162in}}%
\pgfpathlineto{\pgfqpoint{2.250131in}{0.856662in}}%
\pgfpathlineto{\pgfqpoint{2.250431in}{0.857443in}}%
\pgfpathlineto{\pgfqpoint{2.251031in}{0.859578in}}%
\pgfpathlineto{\pgfqpoint{2.251331in}{0.860359in}}%
\pgfpathlineto{\pgfqpoint{2.251930in}{0.862676in}}%
\pgfpathlineto{\pgfqpoint{2.252230in}{0.863457in}}%
\pgfpathlineto{\pgfqpoint{2.252830in}{0.866139in}}%
\pgfpathlineto{\pgfqpoint{2.253130in}{0.866920in}}%
\pgfpathlineto{\pgfqpoint{2.253730in}{0.869029in}}%
\pgfpathlineto{\pgfqpoint{2.254030in}{0.869810in}}%
\pgfpathlineto{\pgfqpoint{2.254630in}{0.872440in}}%
\pgfpathlineto{\pgfqpoint{2.254930in}{0.873221in}}%
\pgfpathlineto{\pgfqpoint{2.255530in}{0.875720in}}%
\pgfpathlineto{\pgfqpoint{2.255829in}{0.876502in}}%
\pgfpathlineto{\pgfqpoint{2.256429in}{0.879053in}}%
\pgfpathlineto{\pgfqpoint{2.256729in}{0.879834in}}%
\pgfpathlineto{\pgfqpoint{2.257329in}{0.882386in}}%
\pgfpathlineto{\pgfqpoint{2.257629in}{0.883167in}}%
\pgfpathlineto{\pgfqpoint{2.258229in}{0.885458in}}%
\pgfpathlineto{\pgfqpoint{2.258529in}{0.886239in}}%
\pgfpathlineto{\pgfqpoint{2.259129in}{0.888791in}}%
\pgfpathlineto{\pgfqpoint{2.259429in}{0.889572in}}%
\pgfpathlineto{\pgfqpoint{2.260028in}{0.892045in}}%
\pgfpathlineto{\pgfqpoint{2.260328in}{0.892826in}}%
\pgfpathlineto{\pgfqpoint{2.260928in}{0.895352in}}%
\pgfpathlineto{\pgfqpoint{2.261078in}{0.896107in}}%
\pgfpathlineto{\pgfqpoint{2.261678in}{0.898320in}}%
\pgfpathlineto{\pgfqpoint{2.261978in}{0.899101in}}%
\pgfpathlineto{\pgfqpoint{2.262578in}{0.901627in}}%
\pgfpathlineto{\pgfqpoint{2.262878in}{0.902408in}}%
\pgfpathlineto{\pgfqpoint{2.263478in}{0.904543in}}%
\pgfpathlineto{\pgfqpoint{2.263778in}{0.905324in}}%
\pgfpathlineto{\pgfqpoint{2.264377in}{0.908084in}}%
\pgfpathlineto{\pgfqpoint{2.264677in}{0.908865in}}%
\pgfpathlineto{\pgfqpoint{2.265277in}{0.911651in}}%
\pgfpathlineto{\pgfqpoint{2.265577in}{0.912432in}}%
\pgfpathlineto{\pgfqpoint{2.266177in}{0.914983in}}%
\pgfpathlineto{\pgfqpoint{2.266477in}{0.915764in}}%
\pgfpathlineto{\pgfqpoint{2.267077in}{0.918342in}}%
\pgfpathlineto{\pgfqpoint{2.267377in}{0.919123in}}%
\pgfpathlineto{\pgfqpoint{2.267977in}{0.922013in}}%
\pgfpathlineto{\pgfqpoint{2.268276in}{0.922794in}}%
\pgfpathlineto{\pgfqpoint{2.268876in}{0.925450in}}%
\pgfpathlineto{\pgfqpoint{2.269176in}{0.926231in}}%
\pgfpathlineto{\pgfqpoint{2.269776in}{0.928939in}}%
\pgfpathlineto{\pgfqpoint{2.270076in}{0.929720in}}%
\pgfpathlineto{\pgfqpoint{2.270676in}{0.932454in}}%
\pgfpathlineto{\pgfqpoint{2.270976in}{0.933235in}}%
\pgfpathlineto{\pgfqpoint{2.271576in}{0.935604in}}%
\pgfpathlineto{\pgfqpoint{2.272026in}{0.936411in}}%
\pgfpathlineto{\pgfqpoint{2.272625in}{0.939171in}}%
\pgfpathlineto{\pgfqpoint{2.272925in}{0.939952in}}%
\pgfpathlineto{\pgfqpoint{2.273525in}{0.942868in}}%
\pgfpathlineto{\pgfqpoint{2.273825in}{0.943649in}}%
\pgfpathlineto{\pgfqpoint{2.274425in}{0.946123in}}%
\pgfpathlineto{\pgfqpoint{2.274725in}{0.946904in}}%
\pgfpathlineto{\pgfqpoint{2.275325in}{0.949247in}}%
\pgfpathlineto{\pgfqpoint{2.275625in}{0.950028in}}%
\pgfpathlineto{\pgfqpoint{2.276225in}{0.952528in}}%
\pgfpathlineto{\pgfqpoint{2.276525in}{0.953309in}}%
\pgfpathlineto{\pgfqpoint{2.277124in}{0.956251in}}%
\pgfpathlineto{\pgfqpoint{2.277424in}{0.957032in}}%
\pgfpathlineto{\pgfqpoint{2.278024in}{0.960026in}}%
\pgfpathlineto{\pgfqpoint{2.278324in}{0.960807in}}%
\pgfpathlineto{\pgfqpoint{2.278924in}{0.963515in}}%
\pgfpathlineto{\pgfqpoint{2.279224in}{0.964296in}}%
\pgfpathlineto{\pgfqpoint{2.279824in}{0.967186in}}%
\pgfpathlineto{\pgfqpoint{2.280124in}{0.967967in}}%
\pgfpathlineto{\pgfqpoint{2.280724in}{0.970753in}}%
\pgfpathlineto{\pgfqpoint{2.281023in}{0.971534in}}%
\pgfpathlineto{\pgfqpoint{2.281623in}{0.974242in}}%
\pgfpathlineto{\pgfqpoint{2.281923in}{0.975023in}}%
\pgfpathlineto{\pgfqpoint{2.282523in}{0.977887in}}%
\pgfpathlineto{\pgfqpoint{2.282673in}{0.978642in}}%
\pgfpathlineto{\pgfqpoint{2.283273in}{0.980881in}}%
\pgfpathlineto{\pgfqpoint{2.283573in}{0.981662in}}%
\pgfpathlineto{\pgfqpoint{2.284173in}{0.984630in}}%
\pgfpathlineto{\pgfqpoint{2.284473in}{0.985412in}}%
\pgfpathlineto{\pgfqpoint{2.285072in}{0.988093in}}%
\pgfpathlineto{\pgfqpoint{2.285372in}{0.988874in}}%
\pgfpathlineto{\pgfqpoint{2.285972in}{0.991895in}}%
\pgfpathlineto{\pgfqpoint{2.286272in}{0.992676in}}%
\pgfpathlineto{\pgfqpoint{2.286872in}{0.995149in}}%
\pgfpathlineto{\pgfqpoint{2.287172in}{0.995930in}}%
\pgfpathlineto{\pgfqpoint{2.287772in}{0.998846in}}%
\pgfpathlineto{\pgfqpoint{2.288072in}{0.999627in}}%
\pgfpathlineto{\pgfqpoint{2.288672in}{1.002309in}}%
\pgfpathlineto{\pgfqpoint{2.288972in}{1.003090in}}%
\pgfpathlineto{\pgfqpoint{2.289571in}{1.005694in}}%
\pgfpathlineto{\pgfqpoint{2.289871in}{1.006475in}}%
\pgfpathlineto{\pgfqpoint{2.290471in}{1.009703in}}%
\pgfpathlineto{\pgfqpoint{2.290771in}{1.010485in}}%
\pgfpathlineto{\pgfqpoint{2.291371in}{1.013557in}}%
\pgfpathlineto{\pgfqpoint{2.291671in}{1.014338in}}%
\pgfpathlineto{\pgfqpoint{2.292271in}{1.017358in}}%
\pgfpathlineto{\pgfqpoint{2.292571in}{1.018139in}}%
\pgfpathlineto{\pgfqpoint{2.293171in}{1.020769in}}%
\pgfpathlineto{\pgfqpoint{2.293470in}{1.021550in}}%
\pgfpathlineto{\pgfqpoint{2.294070in}{1.023971in}}%
\pgfpathlineto{\pgfqpoint{2.294370in}{1.024752in}}%
\pgfpathlineto{\pgfqpoint{2.294970in}{1.027590in}}%
\pgfpathlineto{\pgfqpoint{2.295270in}{1.028371in}}%
\pgfpathlineto{\pgfqpoint{2.295870in}{1.030975in}}%
\pgfpathlineto{\pgfqpoint{2.296170in}{1.031756in}}%
\pgfpathlineto{\pgfqpoint{2.296770in}{1.034282in}}%
\pgfpathlineto{\pgfqpoint{2.297070in}{1.035063in}}%
\pgfpathlineto{\pgfqpoint{2.297669in}{1.037302in}}%
\pgfpathlineto{\pgfqpoint{2.297969in}{1.038083in}}%
\pgfpathlineto{\pgfqpoint{2.298569in}{1.041129in}}%
\pgfpathlineto{\pgfqpoint{2.298869in}{1.041910in}}%
\pgfpathlineto{\pgfqpoint{2.299469in}{1.044983in}}%
\pgfpathlineto{\pgfqpoint{2.299769in}{1.045764in}}%
\pgfpathlineto{\pgfqpoint{2.300369in}{1.048628in}}%
\pgfpathlineto{\pgfqpoint{2.300669in}{1.049409in}}%
\pgfpathlineto{\pgfqpoint{2.301269in}{1.051986in}}%
\pgfpathlineto{\pgfqpoint{2.301569in}{1.052768in}}%
\pgfpathlineto{\pgfqpoint{2.302168in}{1.055527in}}%
\pgfpathlineto{\pgfqpoint{2.302468in}{1.056308in}}%
\pgfpathlineto{\pgfqpoint{2.303068in}{1.059146in}}%
\pgfpathlineto{\pgfqpoint{2.303368in}{1.059927in}}%
\pgfpathlineto{\pgfqpoint{2.303968in}{1.063182in}}%
\pgfpathlineto{\pgfqpoint{2.304268in}{1.063963in}}%
\pgfpathlineto{\pgfqpoint{2.304868in}{1.066957in}}%
\pgfpathlineto{\pgfqpoint{2.305168in}{1.067738in}}%
\pgfpathlineto{\pgfqpoint{2.305768in}{1.070837in}}%
\pgfpathlineto{\pgfqpoint{2.305918in}{1.071592in}}%
\pgfpathlineto{\pgfqpoint{2.306517in}{1.074352in}}%
\pgfpathlineto{\pgfqpoint{2.306817in}{1.075133in}}%
\pgfpathlineto{\pgfqpoint{2.307417in}{1.078517in}}%
\pgfpathlineto{\pgfqpoint{2.307567in}{1.079272in}}%
\pgfpathlineto{\pgfqpoint{2.308167in}{1.082032in}}%
\pgfpathlineto{\pgfqpoint{2.308467in}{1.082813in}}%
\pgfpathlineto{\pgfqpoint{2.309067in}{1.085443in}}%
\pgfpathlineto{\pgfqpoint{2.309367in}{1.086224in}}%
\pgfpathlineto{\pgfqpoint{2.309967in}{1.089427in}}%
\pgfpathlineto{\pgfqpoint{2.310266in}{1.090208in}}%
\pgfpathlineto{\pgfqpoint{2.310866in}{1.093254in}}%
\pgfpathlineto{\pgfqpoint{2.311166in}{1.094035in}}%
\pgfpathlineto{\pgfqpoint{2.311766in}{1.097446in}}%
\pgfpathlineto{\pgfqpoint{2.311916in}{1.098201in}}%
\pgfpathlineto{\pgfqpoint{2.312516in}{1.100935in}}%
\pgfpathlineto{\pgfqpoint{2.312816in}{1.101716in}}%
\pgfpathlineto{\pgfqpoint{2.313416in}{1.104632in}}%
\pgfpathlineto{\pgfqpoint{2.313716in}{1.105413in}}%
\pgfpathlineto{\pgfqpoint{2.314315in}{1.108537in}}%
\pgfpathlineto{\pgfqpoint{2.314465in}{1.109292in}}%
\pgfpathlineto{\pgfqpoint{2.315065in}{1.111610in}}%
\pgfpathlineto{\pgfqpoint{2.315365in}{1.112391in}}%
\pgfpathlineto{\pgfqpoint{2.315965in}{1.115515in}}%
\pgfpathlineto{\pgfqpoint{2.316265in}{1.116296in}}%
\pgfpathlineto{\pgfqpoint{2.316865in}{1.119655in}}%
\pgfpathlineto{\pgfqpoint{2.317165in}{1.120436in}}%
\pgfpathlineto{\pgfqpoint{2.317765in}{1.123456in}}%
\pgfpathlineto{\pgfqpoint{2.318065in}{1.124237in}}%
\pgfpathlineto{\pgfqpoint{2.318664in}{1.127830in}}%
\pgfpathlineto{\pgfqpoint{2.318964in}{1.128611in}}%
\pgfpathlineto{\pgfqpoint{2.319564in}{1.131658in}}%
\pgfpathlineto{\pgfqpoint{2.319714in}{1.132413in}}%
\pgfpathlineto{\pgfqpoint{2.320314in}{1.135277in}}%
\pgfpathlineto{\pgfqpoint{2.320614in}{1.136058in}}%
\pgfpathlineto{\pgfqpoint{2.321214in}{1.139234in}}%
\pgfpathlineto{\pgfqpoint{2.321364in}{1.139989in}}%
\pgfpathlineto{\pgfqpoint{2.321964in}{1.142697in}}%
\pgfpathlineto{\pgfqpoint{2.322264in}{1.143478in}}%
\pgfpathlineto{\pgfqpoint{2.322863in}{1.146785in}}%
\pgfpathlineto{\pgfqpoint{2.323163in}{1.147566in}}%
\pgfpathlineto{\pgfqpoint{2.323763in}{1.151055in}}%
\pgfpathlineto{\pgfqpoint{2.324063in}{1.151836in}}%
\pgfpathlineto{\pgfqpoint{2.324663in}{1.154960in}}%
\pgfpathlineto{\pgfqpoint{2.324963in}{1.155741in}}%
\pgfpathlineto{\pgfqpoint{2.325563in}{1.159334in}}%
\pgfpathlineto{\pgfqpoint{2.326013in}{1.160141in}}%
\pgfpathlineto{\pgfqpoint{2.326613in}{1.163943in}}%
\pgfpathlineto{\pgfqpoint{2.326763in}{1.164698in}}%
\pgfpathlineto{\pgfqpoint{2.327362in}{1.167562in}}%
\pgfpathlineto{\pgfqpoint{2.327662in}{1.168343in}}%
\pgfpathlineto{\pgfqpoint{2.328262in}{1.171311in}}%
\pgfpathlineto{\pgfqpoint{2.328562in}{1.172092in}}%
\pgfpathlineto{\pgfqpoint{2.329162in}{1.175841in}}%
\pgfpathlineto{\pgfqpoint{2.329312in}{1.176596in}}%
\pgfpathlineto{\pgfqpoint{2.329912in}{1.179799in}}%
\pgfpathlineto{\pgfqpoint{2.330212in}{1.180580in}}%
\pgfpathlineto{\pgfqpoint{2.330812in}{1.184095in}}%
\pgfpathlineto{\pgfqpoint{2.330962in}{1.184850in}}%
\pgfpathlineto{\pgfqpoint{2.331561in}{1.187870in}}%
\pgfpathlineto{\pgfqpoint{2.331861in}{1.188651in}}%
\pgfpathlineto{\pgfqpoint{2.332461in}{1.191984in}}%
\pgfpathlineto{\pgfqpoint{2.332761in}{1.192765in}}%
\pgfpathlineto{\pgfqpoint{2.333361in}{1.196306in}}%
\pgfpathlineto{\pgfqpoint{2.333661in}{1.197087in}}%
\pgfpathlineto{\pgfqpoint{2.334261in}{1.199821in}}%
\pgfpathlineto{\pgfqpoint{2.334561in}{1.200602in}}%
\pgfpathlineto{\pgfqpoint{2.335161in}{1.203726in}}%
\pgfpathlineto{\pgfqpoint{2.335310in}{1.204481in}}%
\pgfpathlineto{\pgfqpoint{2.335910in}{1.207866in}}%
\pgfpathlineto{\pgfqpoint{2.336060in}{1.208621in}}%
\pgfpathlineto{\pgfqpoint{2.336660in}{1.211042in}}%
\pgfpathlineto{\pgfqpoint{2.336960in}{1.211823in}}%
\pgfpathlineto{\pgfqpoint{2.337560in}{1.215364in}}%
\pgfpathlineto{\pgfqpoint{2.337710in}{1.216119in}}%
\pgfpathlineto{\pgfqpoint{2.338310in}{1.219140in}}%
\pgfpathlineto{\pgfqpoint{2.338610in}{1.219921in}}%
\pgfpathlineto{\pgfqpoint{2.339210in}{1.223800in}}%
\pgfpathlineto{\pgfqpoint{2.339360in}{1.224555in}}%
\pgfpathlineto{\pgfqpoint{2.339959in}{1.227654in}}%
\pgfpathlineto{\pgfqpoint{2.340259in}{1.228435in}}%
\pgfpathlineto{\pgfqpoint{2.340859in}{1.232236in}}%
\pgfpathlineto{\pgfqpoint{2.341159in}{1.233017in}}%
\pgfpathlineto{\pgfqpoint{2.341759in}{1.237079in}}%
\pgfpathlineto{\pgfqpoint{2.342059in}{1.237860in}}%
\pgfpathlineto{\pgfqpoint{2.342659in}{1.242182in}}%
\pgfpathlineto{\pgfqpoint{2.342959in}{1.242963in}}%
\pgfpathlineto{\pgfqpoint{2.343559in}{1.246816in}}%
\pgfpathlineto{\pgfqpoint{2.343708in}{1.247571in}}%
\pgfpathlineto{\pgfqpoint{2.344308in}{1.251321in}}%
\pgfpathlineto{\pgfqpoint{2.344458in}{1.252076in}}%
\pgfpathlineto{\pgfqpoint{2.345058in}{1.254627in}}%
\pgfpathlineto{\pgfqpoint{2.345208in}{1.255382in}}%
\pgfpathlineto{\pgfqpoint{2.345808in}{1.258481in}}%
\pgfpathlineto{\pgfqpoint{2.345958in}{1.259236in}}%
\pgfpathlineto{\pgfqpoint{2.346558in}{1.262438in}}%
\pgfpathlineto{\pgfqpoint{2.346858in}{1.263219in}}%
\pgfpathlineto{\pgfqpoint{2.347458in}{1.267489in}}%
\pgfpathlineto{\pgfqpoint{2.347758in}{1.268270in}}%
\pgfpathlineto{\pgfqpoint{2.348357in}{1.271655in}}%
\pgfpathlineto{\pgfqpoint{2.348507in}{1.272410in}}%
\pgfpathlineto{\pgfqpoint{2.349107in}{1.275534in}}%
\pgfpathlineto{\pgfqpoint{2.349257in}{1.276289in}}%
\pgfpathlineto{\pgfqpoint{2.349857in}{1.279987in}}%
\pgfpathlineto{\pgfqpoint{2.350157in}{1.280768in}}%
\pgfpathlineto{\pgfqpoint{2.350757in}{1.284517in}}%
\pgfpathlineto{\pgfqpoint{2.350907in}{1.285272in}}%
\pgfpathlineto{\pgfqpoint{2.351507in}{1.288735in}}%
\pgfpathlineto{\pgfqpoint{2.351807in}{1.289516in}}%
\pgfpathlineto{\pgfqpoint{2.352406in}{1.293525in}}%
\pgfpathlineto{\pgfqpoint{2.352556in}{1.294281in}}%
\pgfpathlineto{\pgfqpoint{2.353156in}{1.297405in}}%
\pgfpathlineto{\pgfqpoint{2.353306in}{1.298160in}}%
\pgfpathlineto{\pgfqpoint{2.353906in}{1.301883in}}%
\pgfpathlineto{\pgfqpoint{2.354056in}{1.302638in}}%
\pgfpathlineto{\pgfqpoint{2.354656in}{1.305684in}}%
\pgfpathlineto{\pgfqpoint{2.354806in}{1.306440in}}%
\pgfpathlineto{\pgfqpoint{2.355406in}{1.310553in}}%
\pgfpathlineto{\pgfqpoint{2.355556in}{1.311308in}}%
\pgfpathlineto{\pgfqpoint{2.356156in}{1.314667in}}%
\pgfpathlineto{\pgfqpoint{2.356455in}{1.315448in}}%
\pgfpathlineto{\pgfqpoint{2.357055in}{1.319770in}}%
\pgfpathlineto{\pgfqpoint{2.357205in}{1.320525in}}%
\pgfpathlineto{\pgfqpoint{2.357805in}{1.324248in}}%
\pgfpathlineto{\pgfqpoint{2.357955in}{1.325003in}}%
\pgfpathlineto{\pgfqpoint{2.358555in}{1.328701in}}%
\pgfpathlineto{\pgfqpoint{2.358855in}{1.329482in}}%
\pgfpathlineto{\pgfqpoint{2.359455in}{1.333674in}}%
\pgfpathlineto{\pgfqpoint{2.359605in}{1.334429in}}%
\pgfpathlineto{\pgfqpoint{2.360205in}{1.337865in}}%
\pgfpathlineto{\pgfqpoint{2.360355in}{1.338620in}}%
\pgfpathlineto{\pgfqpoint{2.360954in}{1.342083in}}%
\pgfpathlineto{\pgfqpoint{2.361104in}{1.342838in}}%
\pgfpathlineto{\pgfqpoint{2.361704in}{1.346379in}}%
\pgfpathlineto{\pgfqpoint{2.362004in}{1.347160in}}%
\pgfpathlineto{\pgfqpoint{2.362604in}{1.351300in}}%
\pgfpathlineto{\pgfqpoint{2.362754in}{1.352055in}}%
\pgfpathlineto{\pgfqpoint{2.363354in}{1.355883in}}%
\pgfpathlineto{\pgfqpoint{2.363504in}{1.356638in}}%
\pgfpathlineto{\pgfqpoint{2.364104in}{1.360491in}}%
\pgfpathlineto{\pgfqpoint{2.364254in}{1.361246in}}%
\pgfpathlineto{\pgfqpoint{2.364853in}{1.365047in}}%
\pgfpathlineto{\pgfqpoint{2.365003in}{1.365802in}}%
\pgfpathlineto{\pgfqpoint{2.365603in}{1.369473in}}%
\pgfpathlineto{\pgfqpoint{2.365753in}{1.370229in}}%
\pgfpathlineto{\pgfqpoint{2.366353in}{1.373639in}}%
\pgfpathlineto{\pgfqpoint{2.366503in}{1.374394in}}%
\pgfpathlineto{\pgfqpoint{2.367103in}{1.377805in}}%
\pgfpathlineto{\pgfqpoint{2.367253in}{1.378560in}}%
\pgfpathlineto{\pgfqpoint{2.367853in}{1.382544in}}%
\pgfpathlineto{\pgfqpoint{2.368003in}{1.383299in}}%
\pgfpathlineto{\pgfqpoint{2.368603in}{1.387829in}}%
\pgfpathlineto{\pgfqpoint{2.368753in}{1.388584in}}%
\pgfpathlineto{\pgfqpoint{2.369352in}{1.393219in}}%
\pgfpathlineto{\pgfqpoint{2.369652in}{1.394000in}}%
\pgfpathlineto{\pgfqpoint{2.370252in}{1.398973in}}%
\pgfpathlineto{\pgfqpoint{2.370552in}{1.399754in}}%
\pgfpathlineto{\pgfqpoint{2.371152in}{1.404805in}}%
\pgfpathlineto{\pgfqpoint{2.371302in}{1.405560in}}%
\pgfpathlineto{\pgfqpoint{2.371902in}{1.409257in}}%
\pgfpathlineto{\pgfqpoint{2.372052in}{1.410012in}}%
\pgfpathlineto{\pgfqpoint{2.372652in}{1.414022in}}%
\pgfpathlineto{\pgfqpoint{2.372802in}{1.414777in}}%
\pgfpathlineto{\pgfqpoint{2.373401in}{1.418448in}}%
\pgfpathlineto{\pgfqpoint{2.373551in}{1.419203in}}%
\pgfpathlineto{\pgfqpoint{2.374151in}{1.422926in}}%
\pgfpathlineto{\pgfqpoint{2.374301in}{1.423681in}}%
\pgfpathlineto{\pgfqpoint{2.374901in}{1.427847in}}%
\pgfpathlineto{\pgfqpoint{2.375051in}{1.428602in}}%
\pgfpathlineto{\pgfqpoint{2.375651in}{1.432455in}}%
\pgfpathlineto{\pgfqpoint{2.375951in}{1.433236in}}%
\pgfpathlineto{\pgfqpoint{2.376551in}{1.438001in}}%
\pgfpathlineto{\pgfqpoint{2.376701in}{1.438756in}}%
\pgfpathlineto{\pgfqpoint{2.377300in}{1.442375in}}%
\pgfpathlineto{\pgfqpoint{2.377450in}{1.443130in}}%
\pgfpathlineto{\pgfqpoint{2.378050in}{1.447244in}}%
\pgfpathlineto{\pgfqpoint{2.378200in}{1.447999in}}%
\pgfpathlineto{\pgfqpoint{2.378800in}{1.452425in}}%
\pgfpathlineto{\pgfqpoint{2.378950in}{1.453180in}}%
\pgfpathlineto{\pgfqpoint{2.379550in}{1.456982in}}%
\pgfpathlineto{\pgfqpoint{2.379700in}{1.457737in}}%
\pgfpathlineto{\pgfqpoint{2.380300in}{1.461928in}}%
\pgfpathlineto{\pgfqpoint{2.380450in}{1.462684in}}%
\pgfpathlineto{\pgfqpoint{2.381050in}{1.467240in}}%
\pgfpathlineto{\pgfqpoint{2.381200in}{1.467995in}}%
\pgfpathlineto{\pgfqpoint{2.381799in}{1.472083in}}%
\pgfpathlineto{\pgfqpoint{2.381949in}{1.472838in}}%
\pgfpathlineto{\pgfqpoint{2.382549in}{1.476795in}}%
\pgfpathlineto{\pgfqpoint{2.382699in}{1.477550in}}%
\pgfpathlineto{\pgfqpoint{2.383299in}{1.482523in}}%
\pgfpathlineto{\pgfqpoint{2.383449in}{1.483278in}}%
\pgfpathlineto{\pgfqpoint{2.384049in}{1.487809in}}%
\pgfpathlineto{\pgfqpoint{2.384199in}{1.488564in}}%
\pgfpathlineto{\pgfqpoint{2.384799in}{1.493432in}}%
\pgfpathlineto{\pgfqpoint{2.384949in}{1.494188in}}%
\pgfpathlineto{\pgfqpoint{2.385549in}{1.498301in}}%
\pgfpathlineto{\pgfqpoint{2.385698in}{1.499056in}}%
\pgfpathlineto{\pgfqpoint{2.386298in}{1.503482in}}%
\pgfpathlineto{\pgfqpoint{2.386448in}{1.504238in}}%
\pgfpathlineto{\pgfqpoint{2.387048in}{1.508325in}}%
\pgfpathlineto{\pgfqpoint{2.387198in}{1.509080in}}%
\pgfpathlineto{\pgfqpoint{2.387798in}{1.514209in}}%
\pgfpathlineto{\pgfqpoint{2.387948in}{1.514965in}}%
\pgfpathlineto{\pgfqpoint{2.388548in}{1.519547in}}%
\pgfpathlineto{\pgfqpoint{2.388698in}{1.520302in}}%
\pgfpathlineto{\pgfqpoint{2.389298in}{1.524286in}}%
\pgfpathlineto{\pgfqpoint{2.389448in}{1.525041in}}%
\pgfpathlineto{\pgfqpoint{2.390047in}{1.529675in}}%
\pgfpathlineto{\pgfqpoint{2.390197in}{1.530430in}}%
\pgfpathlineto{\pgfqpoint{2.390797in}{1.535507in}}%
\pgfpathlineto{\pgfqpoint{2.390947in}{1.536262in}}%
\pgfpathlineto{\pgfqpoint{2.391547in}{1.541001in}}%
\pgfpathlineto{\pgfqpoint{2.391697in}{1.541756in}}%
\pgfpathlineto{\pgfqpoint{2.392297in}{1.546859in}}%
\pgfpathlineto{\pgfqpoint{2.392447in}{1.547614in}}%
\pgfpathlineto{\pgfqpoint{2.393047in}{1.552509in}}%
\pgfpathlineto{\pgfqpoint{2.393197in}{1.553264in}}%
\pgfpathlineto{\pgfqpoint{2.393797in}{1.558367in}}%
\pgfpathlineto{\pgfqpoint{2.393947in}{1.559122in}}%
\pgfpathlineto{\pgfqpoint{2.394546in}{1.563132in}}%
\pgfpathlineto{\pgfqpoint{2.394696in}{1.563887in}}%
\pgfpathlineto{\pgfqpoint{2.395296in}{1.568651in}}%
\pgfpathlineto{\pgfqpoint{2.395446in}{1.569406in}}%
\pgfpathlineto{\pgfqpoint{2.396046in}{1.573781in}}%
\pgfpathlineto{\pgfqpoint{2.396196in}{1.574536in}}%
\pgfpathlineto{\pgfqpoint{2.396796in}{1.579066in}}%
\pgfpathlineto{\pgfqpoint{2.396946in}{1.579821in}}%
\pgfpathlineto{\pgfqpoint{2.397546in}{1.584534in}}%
\pgfpathlineto{\pgfqpoint{2.397696in}{1.585289in}}%
\pgfpathlineto{\pgfqpoint{2.398295in}{1.589585in}}%
\pgfpathlineto{\pgfqpoint{2.398445in}{1.590340in}}%
\pgfpathlineto{\pgfqpoint{2.399045in}{1.594766in}}%
\pgfpathlineto{\pgfqpoint{2.399195in}{1.595521in}}%
\pgfpathlineto{\pgfqpoint{2.399795in}{1.599999in}}%
\pgfpathlineto{\pgfqpoint{2.399945in}{1.600754in}}%
\pgfpathlineto{\pgfqpoint{2.400545in}{1.605831in}}%
\pgfpathlineto{\pgfqpoint{2.400695in}{1.606586in}}%
\pgfpathlineto{\pgfqpoint{2.401295in}{1.611716in}}%
\pgfpathlineto{\pgfqpoint{2.401445in}{1.612471in}}%
\pgfpathlineto{\pgfqpoint{2.402045in}{1.617131in}}%
\pgfpathlineto{\pgfqpoint{2.402195in}{1.617886in}}%
\pgfpathlineto{\pgfqpoint{2.402794in}{1.622573in}}%
\pgfpathlineto{\pgfqpoint{2.402944in}{1.623328in}}%
\pgfpathlineto{\pgfqpoint{2.403544in}{1.628144in}}%
\pgfpathlineto{\pgfqpoint{2.403694in}{1.628899in}}%
\pgfpathlineto{\pgfqpoint{2.404294in}{1.633690in}}%
\pgfpathlineto{\pgfqpoint{2.404444in}{1.634445in}}%
\pgfpathlineto{\pgfqpoint{2.405044in}{1.638715in}}%
\pgfpathlineto{\pgfqpoint{2.405194in}{1.639470in}}%
\pgfpathlineto{\pgfqpoint{2.405794in}{1.644105in}}%
\pgfpathlineto{\pgfqpoint{2.405944in}{1.644860in}}%
\pgfpathlineto{\pgfqpoint{2.406544in}{1.649833in}}%
\pgfpathlineto{\pgfqpoint{2.406693in}{1.650588in}}%
\pgfpathlineto{\pgfqpoint{2.407293in}{1.655352in}}%
\pgfpathlineto{\pgfqpoint{2.407443in}{1.656107in}}%
\pgfpathlineto{\pgfqpoint{2.408043in}{1.660976in}}%
\pgfpathlineto{\pgfqpoint{2.408193in}{1.661731in}}%
\pgfpathlineto{\pgfqpoint{2.408793in}{1.666157in}}%
\pgfpathlineto{\pgfqpoint{2.408943in}{1.666913in}}%
\pgfpathlineto{\pgfqpoint{2.409543in}{1.671625in}}%
\pgfpathlineto{\pgfqpoint{2.409843in}{1.672406in}}%
\pgfpathlineto{\pgfqpoint{2.410443in}{1.677692in}}%
\pgfpathlineto{\pgfqpoint{2.410593in}{1.678447in}}%
\pgfpathlineto{\pgfqpoint{2.411192in}{1.683289in}}%
\pgfpathlineto{\pgfqpoint{2.411342in}{1.684044in}}%
\pgfpathlineto{\pgfqpoint{2.411942in}{1.688835in}}%
\pgfpathlineto{\pgfqpoint{2.412092in}{1.689590in}}%
\pgfpathlineto{\pgfqpoint{2.412692in}{1.693912in}}%
\pgfpathlineto{\pgfqpoint{2.412842in}{1.694667in}}%
\pgfpathlineto{\pgfqpoint{2.413442in}{1.699901in}}%
\pgfpathlineto{\pgfqpoint{2.413592in}{1.700656in}}%
\pgfpathlineto{\pgfqpoint{2.414192in}{1.705342in}}%
\pgfpathlineto{\pgfqpoint{2.414342in}{1.706097in}}%
\pgfpathlineto{\pgfqpoint{2.414941in}{1.710758in}}%
\pgfpathlineto{\pgfqpoint{2.415091in}{1.711513in}}%
\pgfpathlineto{\pgfqpoint{2.415691in}{1.716095in}}%
\pgfpathlineto{\pgfqpoint{2.415841in}{1.716850in}}%
\pgfpathlineto{\pgfqpoint{2.416441in}{1.722110in}}%
\pgfpathlineto{\pgfqpoint{2.416741in}{1.722891in}}%
\pgfpathlineto{\pgfqpoint{2.417341in}{1.727499in}}%
\pgfpathlineto{\pgfqpoint{2.417491in}{1.728254in}}%
\pgfpathlineto{\pgfqpoint{2.418091in}{1.732732in}}%
\pgfpathlineto{\pgfqpoint{2.418241in}{1.733487in}}%
\pgfpathlineto{\pgfqpoint{2.418841in}{1.738304in}}%
\pgfpathlineto{\pgfqpoint{2.419140in}{1.739085in}}%
\pgfpathlineto{\pgfqpoint{2.419740in}{1.744293in}}%
\pgfpathlineto{\pgfqpoint{2.419890in}{1.745048in}}%
\pgfpathlineto{\pgfqpoint{2.420490in}{1.749682in}}%
\pgfpathlineto{\pgfqpoint{2.420640in}{1.750437in}}%
\pgfpathlineto{\pgfqpoint{2.421240in}{1.754785in}}%
\pgfpathlineto{\pgfqpoint{2.421390in}{1.755540in}}%
\pgfpathlineto{\pgfqpoint{2.421990in}{1.760149in}}%
\pgfpathlineto{\pgfqpoint{2.422140in}{1.760904in}}%
\pgfpathlineto{\pgfqpoint{2.422740in}{1.765226in}}%
\pgfpathlineto{\pgfqpoint{2.422890in}{1.765981in}}%
\pgfpathlineto{\pgfqpoint{2.423489in}{1.770563in}}%
\pgfpathlineto{\pgfqpoint{2.423639in}{1.771318in}}%
\pgfpathlineto{\pgfqpoint{2.424239in}{1.775536in}}%
\pgfpathlineto{\pgfqpoint{2.424389in}{1.776291in}}%
\pgfpathlineto{\pgfqpoint{2.424989in}{1.780509in}}%
\pgfpathlineto{\pgfqpoint{2.425139in}{1.781264in}}%
\pgfpathlineto{\pgfqpoint{2.425739in}{1.785560in}}%
\pgfpathlineto{\pgfqpoint{2.425889in}{1.786315in}}%
\pgfpathlineto{\pgfqpoint{2.426489in}{1.791002in}}%
\pgfpathlineto{\pgfqpoint{2.426639in}{1.791757in}}%
\pgfpathlineto{\pgfqpoint{2.427239in}{1.796339in}}%
\pgfpathlineto{\pgfqpoint{2.427389in}{1.797094in}}%
\pgfpathlineto{\pgfqpoint{2.427988in}{1.801546in}}%
\pgfpathlineto{\pgfqpoint{2.428138in}{1.802301in}}%
\pgfpathlineto{\pgfqpoint{2.428738in}{1.806676in}}%
\pgfpathlineto{\pgfqpoint{2.428888in}{1.807431in}}%
\pgfpathlineto{\pgfqpoint{2.429488in}{1.811831in}}%
\pgfpathlineto{\pgfqpoint{2.429638in}{1.812586in}}%
\pgfpathlineto{\pgfqpoint{2.430238in}{1.816648in}}%
\pgfpathlineto{\pgfqpoint{2.430388in}{1.817403in}}%
\pgfpathlineto{\pgfqpoint{2.430988in}{1.821126in}}%
\pgfpathlineto{\pgfqpoint{2.431138in}{1.821881in}}%
\pgfpathlineto{\pgfqpoint{2.431737in}{1.825734in}}%
\pgfpathlineto{\pgfqpoint{2.431887in}{1.826489in}}%
\pgfpathlineto{\pgfqpoint{2.432487in}{1.830525in}}%
\pgfpathlineto{\pgfqpoint{2.432637in}{1.831280in}}%
\pgfpathlineto{\pgfqpoint{2.433237in}{1.835576in}}%
\pgfpathlineto{\pgfqpoint{2.433387in}{1.836331in}}%
\pgfpathlineto{\pgfqpoint{2.433987in}{1.841044in}}%
\pgfpathlineto{\pgfqpoint{2.434137in}{1.841799in}}%
\pgfpathlineto{\pgfqpoint{2.434737in}{1.845808in}}%
\pgfpathlineto{\pgfqpoint{2.434887in}{1.846563in}}%
\pgfpathlineto{\pgfqpoint{2.435487in}{1.851406in}}%
\pgfpathlineto{\pgfqpoint{2.435637in}{1.852161in}}%
\pgfpathlineto{\pgfqpoint{2.436236in}{1.856067in}}%
\pgfpathlineto{\pgfqpoint{2.436386in}{1.856822in}}%
\pgfpathlineto{\pgfqpoint{2.436986in}{1.861768in}}%
\pgfpathlineto{\pgfqpoint{2.437136in}{1.862524in}}%
\pgfpathlineto{\pgfqpoint{2.437736in}{1.866611in}}%
\pgfpathlineto{\pgfqpoint{2.437886in}{1.867366in}}%
\pgfpathlineto{\pgfqpoint{2.438486in}{1.870881in}}%
\pgfpathlineto{\pgfqpoint{2.438636in}{1.871636in}}%
\pgfpathlineto{\pgfqpoint{2.439236in}{1.876219in}}%
\pgfpathlineto{\pgfqpoint{2.439386in}{1.876974in}}%
\pgfpathlineto{\pgfqpoint{2.439986in}{1.881218in}}%
\pgfpathlineto{\pgfqpoint{2.440135in}{1.881973in}}%
\pgfpathlineto{\pgfqpoint{2.440735in}{1.885774in}}%
\pgfpathlineto{\pgfqpoint{2.440885in}{1.886529in}}%
\pgfpathlineto{\pgfqpoint{2.441485in}{1.890565in}}%
\pgfpathlineto{\pgfqpoint{2.441635in}{1.891320in}}%
\pgfpathlineto{\pgfqpoint{2.442235in}{1.895538in}}%
\pgfpathlineto{\pgfqpoint{2.442385in}{1.896293in}}%
\pgfpathlineto{\pgfqpoint{2.442985in}{1.900250in}}%
\pgfpathlineto{\pgfqpoint{2.443135in}{1.901005in}}%
\pgfpathlineto{\pgfqpoint{2.443735in}{1.904885in}}%
\pgfpathlineto{\pgfqpoint{2.443885in}{1.905640in}}%
\pgfpathlineto{\pgfqpoint{2.444484in}{1.909727in}}%
\pgfpathlineto{\pgfqpoint{2.444634in}{1.910482in}}%
\pgfpathlineto{\pgfqpoint{2.445234in}{1.914518in}}%
\pgfpathlineto{\pgfqpoint{2.445534in}{1.915299in}}%
\pgfpathlineto{\pgfqpoint{2.446134in}{1.918918in}}%
\pgfpathlineto{\pgfqpoint{2.446284in}{1.919673in}}%
\pgfpathlineto{\pgfqpoint{2.446884in}{1.923813in}}%
\pgfpathlineto{\pgfqpoint{2.447034in}{1.924568in}}%
\pgfpathlineto{\pgfqpoint{2.447634in}{1.928656in}}%
\pgfpathlineto{\pgfqpoint{2.447784in}{1.929411in}}%
\pgfpathlineto{\pgfqpoint{2.448384in}{1.933941in}}%
\pgfpathlineto{\pgfqpoint{2.448533in}{1.934696in}}%
\pgfpathlineto{\pgfqpoint{2.449133in}{1.938628in}}%
\pgfpathlineto{\pgfqpoint{2.449283in}{1.939383in}}%
\pgfpathlineto{\pgfqpoint{2.449883in}{1.943705in}}%
\pgfpathlineto{\pgfqpoint{2.450033in}{1.944460in}}%
\pgfpathlineto{\pgfqpoint{2.450633in}{1.947923in}}%
\pgfpathlineto{\pgfqpoint{2.450783in}{1.948678in}}%
\pgfpathlineto{\pgfqpoint{2.451383in}{1.952167in}}%
\pgfpathlineto{\pgfqpoint{2.451533in}{1.952922in}}%
\pgfpathlineto{\pgfqpoint{2.452133in}{1.956645in}}%
\pgfpathlineto{\pgfqpoint{2.452283in}{1.957400in}}%
\pgfpathlineto{\pgfqpoint{2.452882in}{1.961514in}}%
\pgfpathlineto{\pgfqpoint{2.453032in}{1.962269in}}%
\pgfpathlineto{\pgfqpoint{2.453632in}{1.965862in}}%
\pgfpathlineto{\pgfqpoint{2.453782in}{1.966617in}}%
\pgfpathlineto{\pgfqpoint{2.454382in}{1.970496in}}%
\pgfpathlineto{\pgfqpoint{2.454682in}{1.971277in}}%
\pgfpathlineto{\pgfqpoint{2.455282in}{1.975235in}}%
\pgfpathlineto{\pgfqpoint{2.455582in}{1.976016in}}%
\pgfpathlineto{\pgfqpoint{2.456182in}{1.980312in}}%
\pgfpathlineto{\pgfqpoint{2.456482in}{1.981093in}}%
\pgfpathlineto{\pgfqpoint{2.457081in}{1.985051in}}%
\pgfpathlineto{\pgfqpoint{2.457381in}{1.985832in}}%
\pgfpathlineto{\pgfqpoint{2.457981in}{1.989503in}}%
\pgfpathlineto{\pgfqpoint{2.458131in}{1.990258in}}%
\pgfpathlineto{\pgfqpoint{2.458731in}{1.993252in}}%
\pgfpathlineto{\pgfqpoint{2.459031in}{1.994033in}}%
\pgfpathlineto{\pgfqpoint{2.459631in}{1.998355in}}%
\pgfpathlineto{\pgfqpoint{2.459781in}{1.999110in}}%
\pgfpathlineto{\pgfqpoint{2.460381in}{2.002573in}}%
\pgfpathlineto{\pgfqpoint{2.460531in}{2.003328in}}%
\pgfpathlineto{\pgfqpoint{2.461130in}{2.006270in}}%
\pgfpathlineto{\pgfqpoint{2.461280in}{2.007025in}}%
\pgfpathlineto{\pgfqpoint{2.461880in}{2.009837in}}%
\pgfpathlineto{\pgfqpoint{2.462030in}{2.010592in}}%
\pgfpathlineto{\pgfqpoint{2.462630in}{2.013430in}}%
\pgfpathlineto{\pgfqpoint{2.462780in}{2.014185in}}%
\pgfpathlineto{\pgfqpoint{2.463380in}{2.017856in}}%
\pgfpathlineto{\pgfqpoint{2.463530in}{2.018611in}}%
\pgfpathlineto{\pgfqpoint{2.464130in}{2.022100in}}%
\pgfpathlineto{\pgfqpoint{2.464280in}{2.022855in}}%
\pgfpathlineto{\pgfqpoint{2.464880in}{2.026110in}}%
\pgfpathlineto{\pgfqpoint{2.465030in}{2.026865in}}%
\pgfpathlineto{\pgfqpoint{2.465629in}{2.030145in}}%
\pgfpathlineto{\pgfqpoint{2.465929in}{2.030927in}}%
\pgfpathlineto{\pgfqpoint{2.466529in}{2.034572in}}%
\pgfpathlineto{\pgfqpoint{2.466829in}{2.035353in}}%
\pgfpathlineto{\pgfqpoint{2.467429in}{2.038972in}}%
\pgfpathlineto{\pgfqpoint{2.467729in}{2.039753in}}%
\pgfpathlineto{\pgfqpoint{2.468329in}{2.043346in}}%
\pgfpathlineto{\pgfqpoint{2.468479in}{2.044101in}}%
\pgfpathlineto{\pgfqpoint{2.469079in}{2.047121in}}%
\pgfpathlineto{\pgfqpoint{2.469229in}{2.047876in}}%
\pgfpathlineto{\pgfqpoint{2.469828in}{2.051443in}}%
\pgfpathlineto{\pgfqpoint{2.469978in}{2.052198in}}%
\pgfpathlineto{\pgfqpoint{2.470578in}{2.055401in}}%
\pgfpathlineto{\pgfqpoint{2.470728in}{2.056156in}}%
\pgfpathlineto{\pgfqpoint{2.471328in}{2.059124in}}%
\pgfpathlineto{\pgfqpoint{2.471628in}{2.059905in}}%
\pgfpathlineto{\pgfqpoint{2.472228in}{2.063550in}}%
\pgfpathlineto{\pgfqpoint{2.472528in}{2.064331in}}%
\pgfpathlineto{\pgfqpoint{2.473128in}{2.067742in}}%
\pgfpathlineto{\pgfqpoint{2.473428in}{2.068523in}}%
\pgfpathlineto{\pgfqpoint{2.474027in}{2.071595in}}%
\pgfpathlineto{\pgfqpoint{2.474327in}{2.072376in}}%
\pgfpathlineto{\pgfqpoint{2.474927in}{2.075735in}}%
\pgfpathlineto{\pgfqpoint{2.475227in}{2.076516in}}%
\pgfpathlineto{\pgfqpoint{2.475827in}{2.079667in}}%
\pgfpathlineto{\pgfqpoint{2.475977in}{2.080422in}}%
\pgfpathlineto{\pgfqpoint{2.476577in}{2.083286in}}%
\pgfpathlineto{\pgfqpoint{2.476877in}{2.084067in}}%
\pgfpathlineto{\pgfqpoint{2.477477in}{2.087347in}}%
\pgfpathlineto{\pgfqpoint{2.477777in}{2.088128in}}%
\pgfpathlineto{\pgfqpoint{2.478376in}{2.091878in}}%
\pgfpathlineto{\pgfqpoint{2.478526in}{2.092633in}}%
\pgfpathlineto{\pgfqpoint{2.479126in}{2.095366in}}%
\pgfpathlineto{\pgfqpoint{2.479426in}{2.096148in}}%
\pgfpathlineto{\pgfqpoint{2.480026in}{2.099480in}}%
\pgfpathlineto{\pgfqpoint{2.480326in}{2.100261in}}%
\pgfpathlineto{\pgfqpoint{2.480926in}{2.103386in}}%
\pgfpathlineto{\pgfqpoint{2.481076in}{2.104141in}}%
\pgfpathlineto{\pgfqpoint{2.481676in}{2.106796in}}%
\pgfpathlineto{\pgfqpoint{2.481976in}{2.107577in}}%
\pgfpathlineto{\pgfqpoint{2.482575in}{2.111118in}}%
\pgfpathlineto{\pgfqpoint{2.482875in}{2.111900in}}%
\pgfpathlineto{\pgfqpoint{2.483475in}{2.114633in}}%
\pgfpathlineto{\pgfqpoint{2.483775in}{2.115414in}}%
\pgfpathlineto{\pgfqpoint{2.484375in}{2.118435in}}%
\pgfpathlineto{\pgfqpoint{2.484675in}{2.119216in}}%
\pgfpathlineto{\pgfqpoint{2.485275in}{2.122236in}}%
\pgfpathlineto{\pgfqpoint{2.485575in}{2.123017in}}%
\pgfpathlineto{\pgfqpoint{2.486175in}{2.126141in}}%
\pgfpathlineto{\pgfqpoint{2.486324in}{2.126896in}}%
\pgfpathlineto{\pgfqpoint{2.486924in}{2.129578in}}%
\pgfpathlineto{\pgfqpoint{2.487224in}{2.130359in}}%
\pgfpathlineto{\pgfqpoint{2.487824in}{2.133536in}}%
\pgfpathlineto{\pgfqpoint{2.488124in}{2.134317in}}%
\pgfpathlineto{\pgfqpoint{2.488724in}{2.136764in}}%
\pgfpathlineto{\pgfqpoint{2.488874in}{2.137519in}}%
\pgfpathlineto{\pgfqpoint{2.489474in}{2.140149in}}%
\pgfpathlineto{\pgfqpoint{2.489624in}{2.140904in}}%
\pgfpathlineto{\pgfqpoint{2.490224in}{2.143742in}}%
\pgfpathlineto{\pgfqpoint{2.490374in}{2.144497in}}%
\pgfpathlineto{\pgfqpoint{2.490973in}{2.147387in}}%
\pgfpathlineto{\pgfqpoint{2.491273in}{2.148168in}}%
\pgfpathlineto{\pgfqpoint{2.491873in}{2.151527in}}%
\pgfpathlineto{\pgfqpoint{2.492173in}{2.152308in}}%
\pgfpathlineto{\pgfqpoint{2.492773in}{2.155588in}}%
\pgfpathlineto{\pgfqpoint{2.493073in}{2.156370in}}%
\pgfpathlineto{\pgfqpoint{2.493673in}{2.159312in}}%
\pgfpathlineto{\pgfqpoint{2.493973in}{2.160093in}}%
\pgfpathlineto{\pgfqpoint{2.494573in}{2.162879in}}%
\pgfpathlineto{\pgfqpoint{2.494872in}{2.163660in}}%
\pgfpathlineto{\pgfqpoint{2.495472in}{2.166966in}}%
\pgfpathlineto{\pgfqpoint{2.495772in}{2.167747in}}%
\pgfpathlineto{\pgfqpoint{2.496372in}{2.171054in}}%
\pgfpathlineto{\pgfqpoint{2.496672in}{2.171835in}}%
\pgfpathlineto{\pgfqpoint{2.497272in}{2.174543in}}%
\pgfpathlineto{\pgfqpoint{2.497572in}{2.175324in}}%
\pgfpathlineto{\pgfqpoint{2.498172in}{2.178500in}}%
\pgfpathlineto{\pgfqpoint{2.498472in}{2.179282in}}%
\pgfpathlineto{\pgfqpoint{2.499071in}{2.182146in}}%
\pgfpathlineto{\pgfqpoint{2.499371in}{2.182927in}}%
\pgfpathlineto{\pgfqpoint{2.499971in}{2.185999in}}%
\pgfpathlineto{\pgfqpoint{2.500271in}{2.186780in}}%
\pgfpathlineto{\pgfqpoint{2.500871in}{2.189956in}}%
\pgfpathlineto{\pgfqpoint{2.501171in}{2.190738in}}%
\pgfpathlineto{\pgfqpoint{2.501771in}{2.193628in}}%
\pgfpathlineto{\pgfqpoint{2.502071in}{2.194409in}}%
\pgfpathlineto{\pgfqpoint{2.502671in}{2.197064in}}%
\pgfpathlineto{\pgfqpoint{2.502971in}{2.197845in}}%
\pgfpathlineto{\pgfqpoint{2.503570in}{2.200449in}}%
\pgfpathlineto{\pgfqpoint{2.503870in}{2.201230in}}%
\pgfpathlineto{\pgfqpoint{2.504470in}{2.203990in}}%
\pgfpathlineto{\pgfqpoint{2.504770in}{2.204771in}}%
\pgfpathlineto{\pgfqpoint{2.505370in}{2.207791in}}%
\pgfpathlineto{\pgfqpoint{2.505670in}{2.208572in}}%
\pgfpathlineto{\pgfqpoint{2.506270in}{2.211567in}}%
\pgfpathlineto{\pgfqpoint{2.506570in}{2.212348in}}%
\pgfpathlineto{\pgfqpoint{2.507170in}{2.215134in}}%
\pgfpathlineto{\pgfqpoint{2.507469in}{2.215915in}}%
\pgfpathlineto{\pgfqpoint{2.508069in}{2.218857in}}%
\pgfpathlineto{\pgfqpoint{2.508219in}{2.219612in}}%
\pgfpathlineto{\pgfqpoint{2.508819in}{2.222762in}}%
\pgfpathlineto{\pgfqpoint{2.509119in}{2.223543in}}%
\pgfpathlineto{\pgfqpoint{2.509719in}{2.226616in}}%
\pgfpathlineto{\pgfqpoint{2.510019in}{2.227397in}}%
\pgfpathlineto{\pgfqpoint{2.510619in}{2.229662in}}%
\pgfpathlineto{\pgfqpoint{2.510919in}{2.230443in}}%
\pgfpathlineto{\pgfqpoint{2.511518in}{2.233359in}}%
\pgfpathlineto{\pgfqpoint{2.511818in}{2.234140in}}%
\pgfpathlineto{\pgfqpoint{2.512418in}{2.237212in}}%
\pgfpathlineto{\pgfqpoint{2.512718in}{2.237993in}}%
\pgfpathlineto{\pgfqpoint{2.513318in}{2.241326in}}%
\pgfpathlineto{\pgfqpoint{2.513618in}{2.242107in}}%
\pgfpathlineto{\pgfqpoint{2.514218in}{2.245127in}}%
\pgfpathlineto{\pgfqpoint{2.514518in}{2.245909in}}%
\pgfpathlineto{\pgfqpoint{2.515118in}{2.249189in}}%
\pgfpathlineto{\pgfqpoint{2.515418in}{2.249970in}}%
\pgfpathlineto{\pgfqpoint{2.516017in}{2.252756in}}%
\pgfpathlineto{\pgfqpoint{2.516317in}{2.253537in}}%
\pgfpathlineto{\pgfqpoint{2.516917in}{2.256453in}}%
\pgfpathlineto{\pgfqpoint{2.517217in}{2.257234in}}%
\pgfpathlineto{\pgfqpoint{2.517817in}{2.260124in}}%
\pgfpathlineto{\pgfqpoint{2.518117in}{2.260905in}}%
\pgfpathlineto{\pgfqpoint{2.518717in}{2.263665in}}%
\pgfpathlineto{\pgfqpoint{2.519017in}{2.264446in}}%
\pgfpathlineto{\pgfqpoint{2.519617in}{2.267649in}}%
\pgfpathlineto{\pgfqpoint{2.519916in}{2.268430in}}%
\pgfpathlineto{\pgfqpoint{2.520516in}{2.271554in}}%
\pgfpathlineto{\pgfqpoint{2.520816in}{2.272335in}}%
\pgfpathlineto{\pgfqpoint{2.521416in}{2.275668in}}%
\pgfpathlineto{\pgfqpoint{2.521716in}{2.276449in}}%
\pgfpathlineto{\pgfqpoint{2.522316in}{2.279521in}}%
\pgfpathlineto{\pgfqpoint{2.522466in}{2.280276in}}%
\pgfpathlineto{\pgfqpoint{2.523066in}{2.282958in}}%
\pgfpathlineto{\pgfqpoint{2.523366in}{2.283739in}}%
\pgfpathlineto{\pgfqpoint{2.523965in}{2.286343in}}%
\pgfpathlineto{\pgfqpoint{2.524265in}{2.287124in}}%
\pgfpathlineto{\pgfqpoint{2.524865in}{2.290222in}}%
\pgfpathlineto{\pgfqpoint{2.525165in}{2.291003in}}%
\pgfpathlineto{\pgfqpoint{2.525765in}{2.293451in}}%
\pgfpathlineto{\pgfqpoint{2.526065in}{2.294232in}}%
\pgfpathlineto{\pgfqpoint{2.526665in}{2.297330in}}%
\pgfpathlineto{\pgfqpoint{2.526965in}{2.298111in}}%
\pgfpathlineto{\pgfqpoint{2.527565in}{2.301106in}}%
\pgfpathlineto{\pgfqpoint{2.527865in}{2.301887in}}%
\pgfpathlineto{\pgfqpoint{2.528464in}{2.304647in}}%
\pgfpathlineto{\pgfqpoint{2.528764in}{2.305428in}}%
\pgfpathlineto{\pgfqpoint{2.529364in}{2.308213in}}%
\pgfpathlineto{\pgfqpoint{2.529664in}{2.308995in}}%
\pgfpathlineto{\pgfqpoint{2.530264in}{2.311806in}}%
\pgfpathlineto{\pgfqpoint{2.530564in}{2.312588in}}%
\pgfpathlineto{\pgfqpoint{2.531164in}{2.315113in}}%
\pgfpathlineto{\pgfqpoint{2.531464in}{2.315894in}}%
\pgfpathlineto{\pgfqpoint{2.532064in}{2.318810in}}%
\pgfpathlineto{\pgfqpoint{2.532363in}{2.319591in}}%
\pgfpathlineto{\pgfqpoint{2.532963in}{2.322950in}}%
\pgfpathlineto{\pgfqpoint{2.533263in}{2.323731in}}%
\pgfpathlineto{\pgfqpoint{2.533863in}{2.326803in}}%
\pgfpathlineto{\pgfqpoint{2.534163in}{2.327585in}}%
\pgfpathlineto{\pgfqpoint{2.534763in}{2.330449in}}%
\pgfpathlineto{\pgfqpoint{2.535063in}{2.331230in}}%
\pgfpathlineto{\pgfqpoint{2.535663in}{2.333859in}}%
\pgfpathlineto{\pgfqpoint{2.535963in}{2.334640in}}%
\pgfpathlineto{\pgfqpoint{2.536562in}{2.337426in}}%
\pgfpathlineto{\pgfqpoint{2.536862in}{2.338207in}}%
\pgfpathlineto{\pgfqpoint{2.537462in}{2.340837in}}%
\pgfpathlineto{\pgfqpoint{2.537612in}{2.341592in}}%
\pgfpathlineto{\pgfqpoint{2.538212in}{2.344144in}}%
\pgfpathlineto{\pgfqpoint{2.538512in}{2.344925in}}%
\pgfpathlineto{\pgfqpoint{2.539112in}{2.347737in}}%
\pgfpathlineto{\pgfqpoint{2.539412in}{2.348518in}}%
\pgfpathlineto{\pgfqpoint{2.540012in}{2.351252in}}%
\pgfpathlineto{\pgfqpoint{2.540312in}{2.352033in}}%
\pgfpathlineto{\pgfqpoint{2.540911in}{2.354845in}}%
\pgfpathlineto{\pgfqpoint{2.541211in}{2.355626in}}%
\pgfpathlineto{\pgfqpoint{2.541811in}{2.358646in}}%
\pgfpathlineto{\pgfqpoint{2.542111in}{2.359427in}}%
\pgfpathlineto{\pgfqpoint{2.542711in}{2.362265in}}%
\pgfpathlineto{\pgfqpoint{2.543011in}{2.363046in}}%
\pgfpathlineto{\pgfqpoint{2.543611in}{2.365754in}}%
\pgfpathlineto{\pgfqpoint{2.543911in}{2.366535in}}%
\pgfpathlineto{\pgfqpoint{2.544511in}{2.369060in}}%
\pgfpathlineto{\pgfqpoint{2.544811in}{2.369841in}}%
\pgfpathlineto{\pgfqpoint{2.545410in}{2.372523in}}%
\pgfpathlineto{\pgfqpoint{2.545710in}{2.373304in}}%
\pgfpathlineto{\pgfqpoint{2.546310in}{2.375335in}}%
\pgfpathlineto{\pgfqpoint{2.546610in}{2.376116in}}%
\pgfpathlineto{\pgfqpoint{2.547210in}{2.378798in}}%
\pgfpathlineto{\pgfqpoint{2.547510in}{2.379579in}}%
\pgfpathlineto{\pgfqpoint{2.548110in}{2.382000in}}%
\pgfpathlineto{\pgfqpoint{2.548410in}{2.382782in}}%
\pgfpathlineto{\pgfqpoint{2.549010in}{2.385541in}}%
\pgfpathlineto{\pgfqpoint{2.549309in}{2.386322in}}%
\pgfpathlineto{\pgfqpoint{2.549909in}{2.388562in}}%
\pgfpathlineto{\pgfqpoint{2.550209in}{2.389343in}}%
\pgfpathlineto{\pgfqpoint{2.550809in}{2.392441in}}%
\pgfpathlineto{\pgfqpoint{2.551109in}{2.393222in}}%
\pgfpathlineto{\pgfqpoint{2.551709in}{2.395904in}}%
\pgfpathlineto{\pgfqpoint{2.552009in}{2.396685in}}%
\pgfpathlineto{\pgfqpoint{2.552609in}{2.399705in}}%
\pgfpathlineto{\pgfqpoint{2.552909in}{2.400486in}}%
\pgfpathlineto{\pgfqpoint{2.553508in}{2.402777in}}%
\pgfpathlineto{\pgfqpoint{2.553808in}{2.403559in}}%
\pgfpathlineto{\pgfqpoint{2.554408in}{2.406318in}}%
\pgfpathlineto{\pgfqpoint{2.554708in}{2.407099in}}%
\pgfpathlineto{\pgfqpoint{2.555308in}{2.409677in}}%
\pgfpathlineto{\pgfqpoint{2.555608in}{2.410458in}}%
\pgfpathlineto{\pgfqpoint{2.556208in}{2.413244in}}%
\pgfpathlineto{\pgfqpoint{2.556358in}{2.413999in}}%
\pgfpathlineto{\pgfqpoint{2.556958in}{2.416811in}}%
\pgfpathlineto{\pgfqpoint{2.557258in}{2.417592in}}%
\pgfpathlineto{\pgfqpoint{2.557857in}{2.420040in}}%
\pgfpathlineto{\pgfqpoint{2.558157in}{2.420821in}}%
\pgfpathlineto{\pgfqpoint{2.558757in}{2.423867in}}%
\pgfpathlineto{\pgfqpoint{2.559057in}{2.424648in}}%
\pgfpathlineto{\pgfqpoint{2.559657in}{2.427304in}}%
\pgfpathlineto{\pgfqpoint{2.559957in}{2.428085in}}%
\pgfpathlineto{\pgfqpoint{2.560557in}{2.430819in}}%
\pgfpathlineto{\pgfqpoint{2.561007in}{2.431626in}}%
\pgfpathlineto{\pgfqpoint{2.561607in}{2.434386in}}%
\pgfpathlineto{\pgfqpoint{2.561906in}{2.435167in}}%
\pgfpathlineto{\pgfqpoint{2.562506in}{2.437692in}}%
\pgfpathlineto{\pgfqpoint{2.562806in}{2.438473in}}%
\pgfpathlineto{\pgfqpoint{2.563406in}{2.440712in}}%
\pgfpathlineto{\pgfqpoint{2.563706in}{2.441493in}}%
\pgfpathlineto{\pgfqpoint{2.564306in}{2.443837in}}%
\pgfpathlineto{\pgfqpoint{2.564606in}{2.444618in}}%
\pgfpathlineto{\pgfqpoint{2.565206in}{2.447091in}}%
\pgfpathlineto{\pgfqpoint{2.565506in}{2.447872in}}%
\pgfpathlineto{\pgfqpoint{2.566105in}{2.450762in}}%
\pgfpathlineto{\pgfqpoint{2.566405in}{2.451543in}}%
\pgfpathlineto{\pgfqpoint{2.567005in}{2.453757in}}%
\pgfpathlineto{\pgfqpoint{2.567305in}{2.454538in}}%
\pgfpathlineto{\pgfqpoint{2.567905in}{2.457193in}}%
\pgfpathlineto{\pgfqpoint{2.568055in}{2.457948in}}%
\pgfpathlineto{\pgfqpoint{2.568655in}{2.460266in}}%
\pgfpathlineto{\pgfqpoint{2.568955in}{2.461047in}}%
\pgfpathlineto{\pgfqpoint{2.569555in}{2.463260in}}%
\pgfpathlineto{\pgfqpoint{2.569855in}{2.464041in}}%
\pgfpathlineto{\pgfqpoint{2.570454in}{2.467035in}}%
\pgfpathlineto{\pgfqpoint{2.570754in}{2.467816in}}%
\pgfpathlineto{\pgfqpoint{2.571354in}{2.470472in}}%
\pgfpathlineto{\pgfqpoint{2.571654in}{2.471253in}}%
\pgfpathlineto{\pgfqpoint{2.572254in}{2.473779in}}%
\pgfpathlineto{\pgfqpoint{2.572404in}{2.474534in}}%
\pgfpathlineto{\pgfqpoint{2.573004in}{2.476616in}}%
\pgfpathlineto{\pgfqpoint{2.573454in}{2.477424in}}%
\pgfpathlineto{\pgfqpoint{2.574054in}{2.479975in}}%
\pgfpathlineto{\pgfqpoint{2.574353in}{2.480756in}}%
\pgfpathlineto{\pgfqpoint{2.574953in}{2.482813in}}%
\pgfpathlineto{\pgfqpoint{2.575253in}{2.483594in}}%
\pgfpathlineto{\pgfqpoint{2.575853in}{2.486250in}}%
\pgfpathlineto{\pgfqpoint{2.576303in}{2.487057in}}%
\pgfpathlineto{\pgfqpoint{2.576903in}{2.489609in}}%
\pgfpathlineto{\pgfqpoint{2.577203in}{2.490390in}}%
\pgfpathlineto{\pgfqpoint{2.577803in}{2.493228in}}%
\pgfpathlineto{\pgfqpoint{2.578253in}{2.494035in}}%
\pgfpathlineto{\pgfqpoint{2.578852in}{2.495883in}}%
\pgfpathlineto{\pgfqpoint{2.579152in}{2.496664in}}%
\pgfpathlineto{\pgfqpoint{2.579752in}{2.498669in}}%
\pgfpathlineto{\pgfqpoint{2.580052in}{2.499450in}}%
\pgfpathlineto{\pgfqpoint{2.580652in}{2.502080in}}%
\pgfpathlineto{\pgfqpoint{2.580952in}{2.502861in}}%
\pgfpathlineto{\pgfqpoint{2.581552in}{2.504918in}}%
\pgfpathlineto{\pgfqpoint{2.581852in}{2.505699in}}%
\pgfpathlineto{\pgfqpoint{2.582452in}{2.507704in}}%
\pgfpathlineto{\pgfqpoint{2.582751in}{2.508485in}}%
\pgfpathlineto{\pgfqpoint{2.583351in}{2.510360in}}%
\pgfpathlineto{\pgfqpoint{2.583651in}{2.511141in}}%
\pgfpathlineto{\pgfqpoint{2.584251in}{2.512885in}}%
\pgfpathlineto{\pgfqpoint{2.584551in}{2.513666in}}%
\pgfpathlineto{\pgfqpoint{2.585151in}{2.515905in}}%
\pgfpathlineto{\pgfqpoint{2.585451in}{2.516686in}}%
\pgfpathlineto{\pgfqpoint{2.586051in}{2.518847in}}%
\pgfpathlineto{\pgfqpoint{2.586501in}{2.519655in}}%
\pgfpathlineto{\pgfqpoint{2.587100in}{2.522258in}}%
\pgfpathlineto{\pgfqpoint{2.587400in}{2.523039in}}%
\pgfpathlineto{\pgfqpoint{2.588000in}{2.524784in}}%
\pgfpathlineto{\pgfqpoint{2.588450in}{2.525591in}}%
\pgfpathlineto{\pgfqpoint{2.589050in}{2.527361in}}%
\pgfpathlineto{\pgfqpoint{2.589350in}{2.528142in}}%
\pgfpathlineto{\pgfqpoint{2.589950in}{2.529809in}}%
\pgfpathlineto{\pgfqpoint{2.590250in}{2.530590in}}%
\pgfpathlineto{\pgfqpoint{2.590850in}{2.532803in}}%
\pgfpathlineto{\pgfqpoint{2.591299in}{2.533610in}}%
\pgfpathlineto{\pgfqpoint{2.591899in}{2.535693in}}%
\pgfpathlineto{\pgfqpoint{2.592199in}{2.536474in}}%
\pgfpathlineto{\pgfqpoint{2.592799in}{2.538244in}}%
\pgfpathlineto{\pgfqpoint{2.593099in}{2.539026in}}%
\pgfpathlineto{\pgfqpoint{2.593699in}{2.540978in}}%
\pgfpathlineto{\pgfqpoint{2.593999in}{2.541759in}}%
\pgfpathlineto{\pgfqpoint{2.594599in}{2.543478in}}%
\pgfpathlineto{\pgfqpoint{2.594899in}{2.544259in}}%
\pgfpathlineto{\pgfqpoint{2.595498in}{2.546420in}}%
\pgfpathlineto{\pgfqpoint{2.595798in}{2.547201in}}%
\pgfpathlineto{\pgfqpoint{2.596398in}{2.548841in}}%
\pgfpathlineto{\pgfqpoint{2.596698in}{2.549622in}}%
\pgfpathlineto{\pgfqpoint{2.597298in}{2.551341in}}%
\pgfpathlineto{\pgfqpoint{2.597598in}{2.552122in}}%
\pgfpathlineto{\pgfqpoint{2.598198in}{2.554101in}}%
\pgfpathlineto{\pgfqpoint{2.598648in}{2.554908in}}%
\pgfpathlineto{\pgfqpoint{2.599248in}{2.557095in}}%
\pgfpathlineto{\pgfqpoint{2.599547in}{2.557876in}}%
\pgfpathlineto{\pgfqpoint{2.600147in}{2.559907in}}%
\pgfpathlineto{\pgfqpoint{2.600597in}{2.560714in}}%
\pgfpathlineto{\pgfqpoint{2.601197in}{2.562953in}}%
\pgfpathlineto{\pgfqpoint{2.601647in}{2.563760in}}%
\pgfpathlineto{\pgfqpoint{2.602247in}{2.566077in}}%
\pgfpathlineto{\pgfqpoint{2.602547in}{2.566858in}}%
\pgfpathlineto{\pgfqpoint{2.603147in}{2.568447in}}%
\pgfpathlineto{\pgfqpoint{2.603597in}{2.569254in}}%
\pgfpathlineto{\pgfqpoint{2.604196in}{2.570842in}}%
\pgfpathlineto{\pgfqpoint{2.604496in}{2.571623in}}%
\pgfpathlineto{\pgfqpoint{2.605096in}{2.573550in}}%
\pgfpathlineto{\pgfqpoint{2.605546in}{2.574357in}}%
\pgfpathlineto{\pgfqpoint{2.606146in}{2.576205in}}%
\pgfpathlineto{\pgfqpoint{2.606596in}{2.577013in}}%
\pgfpathlineto{\pgfqpoint{2.607196in}{2.579460in}}%
\pgfpathlineto{\pgfqpoint{2.607496in}{2.580241in}}%
\pgfpathlineto{\pgfqpoint{2.608095in}{2.582168in}}%
\pgfpathlineto{\pgfqpoint{2.608395in}{2.582949in}}%
\pgfpathlineto{\pgfqpoint{2.608995in}{2.584407in}}%
\pgfpathlineto{\pgfqpoint{2.609445in}{2.585214in}}%
\pgfpathlineto{\pgfqpoint{2.610045in}{2.587037in}}%
\pgfpathlineto{\pgfqpoint{2.610495in}{2.587844in}}%
\pgfpathlineto{\pgfqpoint{2.611095in}{2.589458in}}%
\pgfpathlineto{\pgfqpoint{2.611545in}{2.590265in}}%
\pgfpathlineto{\pgfqpoint{2.612144in}{2.591697in}}%
\pgfpathlineto{\pgfqpoint{2.612594in}{2.592504in}}%
\pgfpathlineto{\pgfqpoint{2.613194in}{2.594040in}}%
\pgfpathlineto{\pgfqpoint{2.613494in}{2.594821in}}%
\pgfpathlineto{\pgfqpoint{2.614094in}{2.596514in}}%
\pgfpathlineto{\pgfqpoint{2.614544in}{2.597321in}}%
\pgfpathlineto{\pgfqpoint{2.615144in}{2.598987in}}%
\pgfpathlineto{\pgfqpoint{2.615594in}{2.599794in}}%
\pgfpathlineto{\pgfqpoint{2.616193in}{2.601435in}}%
\pgfpathlineto{\pgfqpoint{2.616493in}{2.602216in}}%
\pgfpathlineto{\pgfqpoint{2.617093in}{2.603570in}}%
\pgfpathlineto{\pgfqpoint{2.617543in}{2.604377in}}%
\pgfpathlineto{\pgfqpoint{2.618143in}{2.605627in}}%
\pgfpathlineto{\pgfqpoint{2.618593in}{2.606434in}}%
\pgfpathlineto{\pgfqpoint{2.619193in}{2.607996in}}%
\pgfpathlineto{\pgfqpoint{2.619493in}{2.608777in}}%
\pgfpathlineto{\pgfqpoint{2.620093in}{2.610105in}}%
\pgfpathlineto{\pgfqpoint{2.620542in}{2.610912in}}%
\pgfpathlineto{\pgfqpoint{2.621142in}{2.612292in}}%
\pgfpathlineto{\pgfqpoint{2.621592in}{2.613099in}}%
\pgfpathlineto{\pgfqpoint{2.622192in}{2.614661in}}%
\pgfpathlineto{\pgfqpoint{2.622642in}{2.615468in}}%
\pgfpathlineto{\pgfqpoint{2.623242in}{2.616614in}}%
\pgfpathlineto{\pgfqpoint{2.623692in}{2.617421in}}%
\pgfpathlineto{\pgfqpoint{2.624292in}{2.618879in}}%
\pgfpathlineto{\pgfqpoint{2.624741in}{2.619686in}}%
\pgfpathlineto{\pgfqpoint{2.625341in}{2.621300in}}%
\pgfpathlineto{\pgfqpoint{2.625791in}{2.622108in}}%
\pgfpathlineto{\pgfqpoint{2.626391in}{2.623201in}}%
\pgfpathlineto{\pgfqpoint{2.626991in}{2.624034in}}%
\pgfpathlineto{\pgfqpoint{2.627591in}{2.625518in}}%
\pgfpathlineto{\pgfqpoint{2.628041in}{2.626325in}}%
\pgfpathlineto{\pgfqpoint{2.628641in}{2.628148in}}%
\pgfpathlineto{\pgfqpoint{2.629090in}{2.628955in}}%
\pgfpathlineto{\pgfqpoint{2.629690in}{2.630647in}}%
\pgfpathlineto{\pgfqpoint{2.630140in}{2.631455in}}%
\pgfpathlineto{\pgfqpoint{2.630740in}{2.632574in}}%
\pgfpathlineto{\pgfqpoint{2.631190in}{2.633381in}}%
\pgfpathlineto{\pgfqpoint{2.631790in}{2.635048in}}%
\pgfpathlineto{\pgfqpoint{2.632240in}{2.635855in}}%
\pgfpathlineto{\pgfqpoint{2.632840in}{2.637443in}}%
\pgfpathlineto{\pgfqpoint{2.633289in}{2.638250in}}%
\pgfpathlineto{\pgfqpoint{2.633889in}{2.639552in}}%
\pgfpathlineto{\pgfqpoint{2.634489in}{2.640385in}}%
\pgfpathlineto{\pgfqpoint{2.635089in}{2.641687in}}%
\pgfpathlineto{\pgfqpoint{2.635389in}{2.642468in}}%
\pgfpathlineto{\pgfqpoint{2.635989in}{2.643822in}}%
\pgfpathlineto{\pgfqpoint{2.636589in}{2.644655in}}%
\pgfpathlineto{\pgfqpoint{2.637188in}{2.646113in}}%
\pgfpathlineto{\pgfqpoint{2.637788in}{2.646946in}}%
\pgfpathlineto{\pgfqpoint{2.638388in}{2.648456in}}%
\pgfpathlineto{\pgfqpoint{2.638838in}{2.649263in}}%
\pgfpathlineto{\pgfqpoint{2.639438in}{2.650305in}}%
\pgfpathlineto{\pgfqpoint{2.640038in}{2.651138in}}%
\pgfpathlineto{\pgfqpoint{2.640638in}{2.651945in}}%
\pgfpathlineto{\pgfqpoint{2.641238in}{2.652778in}}%
\pgfpathlineto{\pgfqpoint{2.641837in}{2.654054in}}%
\pgfpathlineto{\pgfqpoint{2.642287in}{2.654861in}}%
\pgfpathlineto{\pgfqpoint{2.642887in}{2.656267in}}%
\pgfpathlineto{\pgfqpoint{2.643337in}{2.657074in}}%
\pgfpathlineto{\pgfqpoint{2.643937in}{2.658689in}}%
\pgfpathlineto{\pgfqpoint{2.644537in}{2.659522in}}%
\pgfpathlineto{\pgfqpoint{2.645137in}{2.660954in}}%
\pgfpathlineto{\pgfqpoint{2.645437in}{2.661735in}}%
\pgfpathlineto{\pgfqpoint{2.646036in}{2.662620in}}%
\pgfpathlineto{\pgfqpoint{2.646486in}{2.663427in}}%
\pgfpathlineto{\pgfqpoint{2.647086in}{2.664625in}}%
\pgfpathlineto{\pgfqpoint{2.647536in}{2.665432in}}%
\pgfpathlineto{\pgfqpoint{2.648136in}{2.666473in}}%
\pgfpathlineto{\pgfqpoint{2.648586in}{2.667281in}}%
\pgfpathlineto{\pgfqpoint{2.649186in}{2.668426in}}%
\pgfpathlineto{\pgfqpoint{2.649636in}{2.669233in}}%
\pgfpathlineto{\pgfqpoint{2.650235in}{2.670327in}}%
\pgfpathlineto{\pgfqpoint{2.650985in}{2.671186in}}%
\pgfpathlineto{\pgfqpoint{2.651585in}{2.672566in}}%
\pgfpathlineto{\pgfqpoint{2.652185in}{2.673399in}}%
\pgfpathlineto{\pgfqpoint{2.652785in}{2.674935in}}%
\pgfpathlineto{\pgfqpoint{2.653235in}{2.675742in}}%
\pgfpathlineto{\pgfqpoint{2.653835in}{2.676784in}}%
\pgfpathlineto{\pgfqpoint{2.654284in}{2.677591in}}%
\pgfpathlineto{\pgfqpoint{2.654884in}{2.678737in}}%
\pgfpathlineto{\pgfqpoint{2.655484in}{2.679570in}}%
\pgfpathlineto{\pgfqpoint{2.656084in}{2.680507in}}%
\pgfpathlineto{\pgfqpoint{2.656984in}{2.681392in}}%
\pgfpathlineto{\pgfqpoint{2.657584in}{2.682564in}}%
\pgfpathlineto{\pgfqpoint{2.658183in}{2.683397in}}%
\pgfpathlineto{\pgfqpoint{2.658783in}{2.684959in}}%
\pgfpathlineto{\pgfqpoint{2.659533in}{2.685818in}}%
\pgfpathlineto{\pgfqpoint{2.660133in}{2.687172in}}%
\pgfpathlineto{\pgfqpoint{2.660583in}{2.687979in}}%
\pgfpathlineto{\pgfqpoint{2.661183in}{2.688891in}}%
\pgfpathlineto{\pgfqpoint{2.661783in}{2.689724in}}%
\pgfpathlineto{\pgfqpoint{2.662382in}{2.690583in}}%
\pgfpathlineto{\pgfqpoint{2.662982in}{2.691416in}}%
\pgfpathlineto{\pgfqpoint{2.663582in}{2.692536in}}%
\pgfpathlineto{\pgfqpoint{2.664032in}{2.693343in}}%
\pgfpathlineto{\pgfqpoint{2.664632in}{2.694358in}}%
\pgfpathlineto{\pgfqpoint{2.665382in}{2.695218in}}%
\pgfpathlineto{\pgfqpoint{2.665982in}{2.696259in}}%
\pgfpathlineto{\pgfqpoint{2.666731in}{2.697118in}}%
\pgfpathlineto{\pgfqpoint{2.667331in}{2.698186in}}%
\pgfpathlineto{\pgfqpoint{2.668081in}{2.699045in}}%
\pgfpathlineto{\pgfqpoint{2.668681in}{2.699982in}}%
\pgfpathlineto{\pgfqpoint{2.669281in}{2.700815in}}%
\pgfpathlineto{\pgfqpoint{2.669881in}{2.702039in}}%
\pgfpathlineto{\pgfqpoint{2.670481in}{2.702872in}}%
\pgfpathlineto{\pgfqpoint{2.671080in}{2.704070in}}%
\pgfpathlineto{\pgfqpoint{2.671830in}{2.704929in}}%
\pgfpathlineto{\pgfqpoint{2.672430in}{2.706231in}}%
\pgfpathlineto{\pgfqpoint{2.673180in}{2.707090in}}%
\pgfpathlineto{\pgfqpoint{2.673780in}{2.707897in}}%
\pgfpathlineto{\pgfqpoint{2.674380in}{2.708730in}}%
\pgfpathlineto{\pgfqpoint{2.674979in}{2.709642in}}%
\pgfpathlineto{\pgfqpoint{2.675579in}{2.710475in}}%
\pgfpathlineto{\pgfqpoint{2.676179in}{2.711464in}}%
\pgfpathlineto{\pgfqpoint{2.676929in}{2.712323in}}%
\pgfpathlineto{\pgfqpoint{2.677529in}{2.713183in}}%
\pgfpathlineto{\pgfqpoint{2.678129in}{2.714016in}}%
\pgfpathlineto{\pgfqpoint{2.678729in}{2.715005in}}%
\pgfpathlineto{\pgfqpoint{2.679328in}{2.715838in}}%
\pgfpathlineto{\pgfqpoint{2.679928in}{2.716671in}}%
\pgfpathlineto{\pgfqpoint{2.680528in}{2.717505in}}%
\pgfpathlineto{\pgfqpoint{2.681128in}{2.718286in}}%
\pgfpathlineto{\pgfqpoint{2.681878in}{2.719145in}}%
\pgfpathlineto{\pgfqpoint{2.682478in}{2.720108in}}%
\pgfpathlineto{\pgfqpoint{2.683228in}{2.720967in}}%
\pgfpathlineto{\pgfqpoint{2.683827in}{2.721879in}}%
\pgfpathlineto{\pgfqpoint{2.684427in}{2.722712in}}%
\pgfpathlineto{\pgfqpoint{2.685027in}{2.723363in}}%
\pgfpathlineto{\pgfqpoint{2.685627in}{2.724196in}}%
\pgfpathlineto{\pgfqpoint{2.686227in}{2.725159in}}%
\pgfpathlineto{\pgfqpoint{2.686977in}{2.726019in}}%
\pgfpathlineto{\pgfqpoint{2.687576in}{2.726904in}}%
\pgfpathlineto{\pgfqpoint{2.688176in}{2.727737in}}%
\pgfpathlineto{\pgfqpoint{2.688776in}{2.728518in}}%
\pgfpathlineto{\pgfqpoint{2.689526in}{2.729377in}}%
\pgfpathlineto{\pgfqpoint{2.690126in}{2.730341in}}%
\pgfpathlineto{\pgfqpoint{2.690876in}{2.731200in}}%
\pgfpathlineto{\pgfqpoint{2.691476in}{2.732215in}}%
\pgfpathlineto{\pgfqpoint{2.692225in}{2.733074in}}%
\pgfpathlineto{\pgfqpoint{2.692825in}{2.734090in}}%
\pgfpathlineto{\pgfqpoint{2.693575in}{2.734949in}}%
\pgfpathlineto{\pgfqpoint{2.694175in}{2.735782in}}%
\pgfpathlineto{\pgfqpoint{2.695225in}{2.736693in}}%
\pgfpathlineto{\pgfqpoint{2.695825in}{2.737579in}}%
\pgfpathlineto{\pgfqpoint{2.697024in}{2.738516in}}%
\pgfpathlineto{\pgfqpoint{2.697624in}{2.739531in}}%
\pgfpathlineto{\pgfqpoint{2.698374in}{2.740391in}}%
\pgfpathlineto{\pgfqpoint{2.698974in}{2.741041in}}%
\pgfpathlineto{\pgfqpoint{2.699574in}{2.741875in}}%
\pgfpathlineto{\pgfqpoint{2.700173in}{2.742421in}}%
\pgfpathlineto{\pgfqpoint{2.700773in}{2.743255in}}%
\pgfpathlineto{\pgfqpoint{2.701373in}{2.743905in}}%
\pgfpathlineto{\pgfqpoint{2.702123in}{2.744765in}}%
\pgfpathlineto{\pgfqpoint{2.702723in}{2.745390in}}%
\pgfpathlineto{\pgfqpoint{2.703473in}{2.746249in}}%
\pgfpathlineto{\pgfqpoint{2.704073in}{2.747212in}}%
\pgfpathlineto{\pgfqpoint{2.704822in}{2.748071in}}%
\pgfpathlineto{\pgfqpoint{2.705422in}{2.748670in}}%
\pgfpathlineto{\pgfqpoint{2.706322in}{2.749555in}}%
\pgfpathlineto{\pgfqpoint{2.706922in}{2.750571in}}%
\pgfpathlineto{\pgfqpoint{2.707672in}{2.751430in}}%
\pgfpathlineto{\pgfqpoint{2.708272in}{2.752237in}}%
\pgfpathlineto{\pgfqpoint{2.709171in}{2.753122in}}%
\pgfpathlineto{\pgfqpoint{2.709771in}{2.753669in}}%
\pgfpathlineto{\pgfqpoint{2.710221in}{2.754476in}}%
\pgfpathlineto{\pgfqpoint{2.710821in}{2.755257in}}%
\pgfpathlineto{\pgfqpoint{2.711571in}{2.756117in}}%
\pgfpathlineto{\pgfqpoint{2.712171in}{2.756820in}}%
\pgfpathlineto{\pgfqpoint{2.713070in}{2.757705in}}%
\pgfpathlineto{\pgfqpoint{2.713670in}{2.758252in}}%
\pgfpathlineto{\pgfqpoint{2.714270in}{2.759085in}}%
\pgfpathlineto{\pgfqpoint{2.714870in}{2.759996in}}%
\pgfpathlineto{\pgfqpoint{2.715470in}{2.760829in}}%
\pgfpathlineto{\pgfqpoint{2.716070in}{2.761871in}}%
\pgfpathlineto{\pgfqpoint{2.716819in}{2.762730in}}%
\pgfpathlineto{\pgfqpoint{2.717419in}{2.763277in}}%
\pgfpathlineto{\pgfqpoint{2.718169in}{2.764136in}}%
\pgfpathlineto{\pgfqpoint{2.718769in}{2.765021in}}%
\pgfpathlineto{\pgfqpoint{2.719669in}{2.765906in}}%
\pgfpathlineto{\pgfqpoint{2.720269in}{2.766739in}}%
\pgfpathlineto{\pgfqpoint{2.721318in}{2.767651in}}%
\pgfpathlineto{\pgfqpoint{2.721918in}{2.768302in}}%
\pgfpathlineto{\pgfqpoint{2.722518in}{2.769135in}}%
\pgfpathlineto{\pgfqpoint{2.723118in}{2.769786in}}%
\pgfpathlineto{\pgfqpoint{2.724018in}{2.770671in}}%
\pgfpathlineto{\pgfqpoint{2.724618in}{2.771738in}}%
\pgfpathlineto{\pgfqpoint{2.725367in}{2.772598in}}%
\pgfpathlineto{\pgfqpoint{2.725967in}{2.773222in}}%
\pgfpathlineto{\pgfqpoint{2.726867in}{2.774108in}}%
\pgfpathlineto{\pgfqpoint{2.727467in}{2.774733in}}%
\pgfpathlineto{\pgfqpoint{2.728367in}{2.775618in}}%
\pgfpathlineto{\pgfqpoint{2.728967in}{2.776425in}}%
\pgfpathlineto{\pgfqpoint{2.729716in}{2.777284in}}%
\pgfpathlineto{\pgfqpoint{2.730316in}{2.777961in}}%
\pgfpathlineto{\pgfqpoint{2.731366in}{2.778872in}}%
\pgfpathlineto{\pgfqpoint{2.731966in}{2.779705in}}%
\pgfpathlineto{\pgfqpoint{2.732866in}{2.780591in}}%
\pgfpathlineto{\pgfqpoint{2.733466in}{2.781111in}}%
\pgfpathlineto{\pgfqpoint{2.734365in}{2.781997in}}%
\pgfpathlineto{\pgfqpoint{2.734965in}{2.782595in}}%
\pgfpathlineto{\pgfqpoint{2.736165in}{2.783533in}}%
\pgfpathlineto{\pgfqpoint{2.736765in}{2.784210in}}%
\pgfpathlineto{\pgfqpoint{2.737365in}{2.785043in}}%
\pgfpathlineto{\pgfqpoint{2.737964in}{2.785538in}}%
\pgfpathlineto{\pgfqpoint{2.739014in}{2.786449in}}%
\pgfpathlineto{\pgfqpoint{2.739614in}{2.787152in}}%
\pgfpathlineto{\pgfqpoint{2.740514in}{2.788037in}}%
\pgfpathlineto{\pgfqpoint{2.741114in}{2.788714in}}%
\pgfpathlineto{\pgfqpoint{2.742013in}{2.789599in}}%
\pgfpathlineto{\pgfqpoint{2.742613in}{2.790276in}}%
\pgfpathlineto{\pgfqpoint{2.743813in}{2.791214in}}%
\pgfpathlineto{\pgfqpoint{2.744413in}{2.791864in}}%
\pgfpathlineto{\pgfqpoint{2.745463in}{2.792776in}}%
\pgfpathlineto{\pgfqpoint{2.746063in}{2.793505in}}%
\pgfpathlineto{\pgfqpoint{2.746962in}{2.794390in}}%
\pgfpathlineto{\pgfqpoint{2.747562in}{2.795093in}}%
\pgfpathlineto{\pgfqpoint{2.748762in}{2.796030in}}%
\pgfpathlineto{\pgfqpoint{2.749362in}{2.796681in}}%
\pgfpathlineto{\pgfqpoint{2.750112in}{2.797540in}}%
\pgfpathlineto{\pgfqpoint{2.750711in}{2.798347in}}%
\pgfpathlineto{\pgfqpoint{2.751461in}{2.799207in}}%
\pgfpathlineto{\pgfqpoint{2.752061in}{2.799753in}}%
\pgfpathlineto{\pgfqpoint{2.753111in}{2.800665in}}%
\pgfpathlineto{\pgfqpoint{2.753711in}{2.801446in}}%
\pgfpathlineto{\pgfqpoint{2.754461in}{2.802305in}}%
\pgfpathlineto{\pgfqpoint{2.755060in}{2.802617in}}%
\pgfpathlineto{\pgfqpoint{2.756410in}{2.803581in}}%
\pgfpathlineto{\pgfqpoint{2.757010in}{2.804206in}}%
\pgfpathlineto{\pgfqpoint{2.758809in}{2.805247in}}%
\pgfpathlineto{\pgfqpoint{2.759409in}{2.806132in}}%
\pgfpathlineto{\pgfqpoint{2.760459in}{2.807044in}}%
\pgfpathlineto{\pgfqpoint{2.761059in}{2.807668in}}%
\pgfpathlineto{\pgfqpoint{2.761959in}{2.808554in}}%
\pgfpathlineto{\pgfqpoint{2.762559in}{2.809127in}}%
\pgfpathlineto{\pgfqpoint{2.764058in}{2.810116in}}%
\pgfpathlineto{\pgfqpoint{2.764658in}{2.810611in}}%
\pgfpathlineto{\pgfqpoint{2.765408in}{2.811470in}}%
\pgfpathlineto{\pgfqpoint{2.766008in}{2.811860in}}%
\pgfpathlineto{\pgfqpoint{2.767207in}{2.812798in}}%
\pgfpathlineto{\pgfqpoint{2.767807in}{2.813266in}}%
\pgfpathlineto{\pgfqpoint{2.769157in}{2.814230in}}%
\pgfpathlineto{\pgfqpoint{2.769757in}{2.814724in}}%
\pgfpathlineto{\pgfqpoint{2.771107in}{2.815688in}}%
\pgfpathlineto{\pgfqpoint{2.771706in}{2.816391in}}%
\pgfpathlineto{\pgfqpoint{2.773356in}{2.817406in}}%
\pgfpathlineto{\pgfqpoint{2.773956in}{2.818187in}}%
\pgfpathlineto{\pgfqpoint{2.774856in}{2.819072in}}%
\pgfpathlineto{\pgfqpoint{2.775456in}{2.819619in}}%
\pgfpathlineto{\pgfqpoint{2.776505in}{2.820530in}}%
\pgfpathlineto{\pgfqpoint{2.777105in}{2.820999in}}%
\pgfpathlineto{\pgfqpoint{2.778155in}{2.821910in}}%
\pgfpathlineto{\pgfqpoint{2.778755in}{2.822405in}}%
\pgfpathlineto{\pgfqpoint{2.779804in}{2.823316in}}%
\pgfpathlineto{\pgfqpoint{2.780404in}{2.823681in}}%
\pgfpathlineto{\pgfqpoint{2.781754in}{2.824644in}}%
\pgfpathlineto{\pgfqpoint{2.782354in}{2.825321in}}%
\pgfpathlineto{\pgfqpoint{2.783404in}{2.826232in}}%
\pgfpathlineto{\pgfqpoint{2.784003in}{2.826883in}}%
\pgfpathlineto{\pgfqpoint{2.785503in}{2.827873in}}%
\pgfpathlineto{\pgfqpoint{2.786103in}{2.828341in}}%
\pgfpathlineto{\pgfqpoint{2.788202in}{2.829435in}}%
\pgfpathlineto{\pgfqpoint{2.788802in}{2.830190in}}%
\pgfpathlineto{\pgfqpoint{2.790002in}{2.831127in}}%
\pgfpathlineto{\pgfqpoint{2.790602in}{2.831596in}}%
\pgfpathlineto{\pgfqpoint{2.791652in}{2.832507in}}%
\pgfpathlineto{\pgfqpoint{2.792252in}{2.832872in}}%
\pgfpathlineto{\pgfqpoint{2.793601in}{2.833835in}}%
\pgfpathlineto{\pgfqpoint{2.794201in}{2.834304in}}%
\pgfpathlineto{\pgfqpoint{2.795551in}{2.835267in}}%
\pgfpathlineto{\pgfqpoint{2.796151in}{2.836048in}}%
\pgfpathlineto{\pgfqpoint{2.797200in}{2.836959in}}%
\pgfpathlineto{\pgfqpoint{2.797800in}{2.837454in}}%
\pgfpathlineto{\pgfqpoint{2.799300in}{2.838443in}}%
\pgfpathlineto{\pgfqpoint{2.799900in}{2.838990in}}%
\pgfpathlineto{\pgfqpoint{2.801549in}{2.840006in}}%
\pgfpathlineto{\pgfqpoint{2.802149in}{2.840657in}}%
\pgfpathlineto{\pgfqpoint{2.803649in}{2.841646in}}%
\pgfpathlineto{\pgfqpoint{2.804249in}{2.842323in}}%
\pgfpathlineto{\pgfqpoint{2.805598in}{2.843182in}}%
\pgfpathlineto{\pgfqpoint{2.806198in}{2.843547in}}%
\pgfpathlineto{\pgfqpoint{2.807698in}{2.844536in}}%
\pgfpathlineto{\pgfqpoint{2.808298in}{2.845135in}}%
\pgfpathlineto{\pgfqpoint{2.810247in}{2.846202in}}%
\pgfpathlineto{\pgfqpoint{2.810847in}{2.846593in}}%
\pgfpathlineto{\pgfqpoint{2.812047in}{2.847530in}}%
\pgfpathlineto{\pgfqpoint{2.812647in}{2.847947in}}%
\pgfpathlineto{\pgfqpoint{2.814296in}{2.848962in}}%
\pgfpathlineto{\pgfqpoint{2.814896in}{2.849405in}}%
\pgfpathlineto{\pgfqpoint{2.816396in}{2.850394in}}%
\pgfpathlineto{\pgfqpoint{2.816996in}{2.851019in}}%
\pgfpathlineto{\pgfqpoint{2.819395in}{2.852165in}}%
\pgfpathlineto{\pgfqpoint{2.819995in}{2.852607in}}%
\pgfpathlineto{\pgfqpoint{2.821495in}{2.853597in}}%
\pgfpathlineto{\pgfqpoint{2.822094in}{2.854013in}}%
\pgfpathlineto{\pgfqpoint{2.823444in}{2.854976in}}%
\pgfpathlineto{\pgfqpoint{2.824044in}{2.855497in}}%
\pgfpathlineto{\pgfqpoint{2.826143in}{2.856591in}}%
\pgfpathlineto{\pgfqpoint{2.826743in}{2.856981in}}%
\pgfpathlineto{\pgfqpoint{2.828393in}{2.857997in}}%
\pgfpathlineto{\pgfqpoint{2.828993in}{2.858413in}}%
\pgfpathlineto{\pgfqpoint{2.830642in}{2.859429in}}%
\pgfpathlineto{\pgfqpoint{2.831242in}{2.859975in}}%
\pgfpathlineto{\pgfqpoint{2.833042in}{2.861017in}}%
\pgfpathlineto{\pgfqpoint{2.833642in}{2.861381in}}%
\pgfpathlineto{\pgfqpoint{2.835591in}{2.862449in}}%
\pgfpathlineto{\pgfqpoint{2.836191in}{2.862709in}}%
\pgfpathlineto{\pgfqpoint{2.838291in}{2.863803in}}%
\pgfpathlineto{\pgfqpoint{2.838890in}{2.864193in}}%
\pgfpathlineto{\pgfqpoint{2.840990in}{2.865287in}}%
\pgfpathlineto{\pgfqpoint{2.841590in}{2.865808in}}%
\pgfpathlineto{\pgfqpoint{2.842490in}{2.866693in}}%
\pgfpathlineto{\pgfqpoint{2.843089in}{2.867135in}}%
\pgfpathlineto{\pgfqpoint{2.845039in}{2.868203in}}%
\pgfpathlineto{\pgfqpoint{2.845639in}{2.868515in}}%
\pgfpathlineto{\pgfqpoint{2.848488in}{2.869739in}}%
\pgfpathlineto{\pgfqpoint{2.849088in}{2.870364in}}%
\pgfpathlineto{\pgfqpoint{2.850738in}{2.871379in}}%
\pgfpathlineto{\pgfqpoint{2.851337in}{2.871718in}}%
\pgfpathlineto{\pgfqpoint{2.853437in}{2.872811in}}%
\pgfpathlineto{\pgfqpoint{2.854037in}{2.873150in}}%
\pgfpathlineto{\pgfqpoint{2.855986in}{2.874217in}}%
\pgfpathlineto{\pgfqpoint{2.856586in}{2.874426in}}%
\pgfpathlineto{\pgfqpoint{2.859585in}{2.875675in}}%
\pgfpathlineto{\pgfqpoint{2.860185in}{2.876092in}}%
\pgfpathlineto{\pgfqpoint{2.862585in}{2.877238in}}%
\pgfpathlineto{\pgfqpoint{2.863185in}{2.877680in}}%
\pgfpathlineto{\pgfqpoint{2.865734in}{2.878852in}}%
\pgfpathlineto{\pgfqpoint{2.866334in}{2.879294in}}%
\pgfpathlineto{\pgfqpoint{2.869183in}{2.880518in}}%
\pgfpathlineto{\pgfqpoint{2.869633in}{2.880726in}}%
\pgfpathlineto{\pgfqpoint{2.871733in}{2.881586in}}%
\pgfpathlineto{\pgfqpoint{2.872332in}{2.882210in}}%
\pgfpathlineto{\pgfqpoint{2.874882in}{2.883382in}}%
\pgfpathlineto{\pgfqpoint{2.875482in}{2.883877in}}%
\pgfpathlineto{\pgfqpoint{2.877431in}{2.884944in}}%
\pgfpathlineto{\pgfqpoint{2.878031in}{2.885309in}}%
\pgfpathlineto{\pgfqpoint{2.879831in}{2.886350in}}%
\pgfpathlineto{\pgfqpoint{2.880430in}{2.886689in}}%
\pgfpathlineto{\pgfqpoint{2.883430in}{2.887938in}}%
\pgfpathlineto{\pgfqpoint{2.884030in}{2.888173in}}%
\pgfpathlineto{\pgfqpoint{2.886879in}{2.889397in}}%
\pgfpathlineto{\pgfqpoint{2.887479in}{2.889631in}}%
\pgfpathlineto{\pgfqpoint{2.889728in}{2.890750in}}%
\pgfpathlineto{\pgfqpoint{2.890328in}{2.891115in}}%
\pgfpathlineto{\pgfqpoint{2.892578in}{2.892234in}}%
\pgfpathlineto{\pgfqpoint{2.893177in}{2.892651in}}%
\pgfpathlineto{\pgfqpoint{2.896777in}{2.894005in}}%
\pgfpathlineto{\pgfqpoint{2.897376in}{2.894291in}}%
\pgfpathlineto{\pgfqpoint{2.899926in}{2.895463in}}%
\pgfpathlineto{\pgfqpoint{2.900526in}{2.895723in}}%
\pgfpathlineto{\pgfqpoint{2.902775in}{2.896843in}}%
\pgfpathlineto{\pgfqpoint{2.903375in}{2.897259in}}%
\pgfpathlineto{\pgfqpoint{2.905924in}{2.898223in}}%
\pgfpathlineto{\pgfqpoint{2.909374in}{2.899551in}}%
\pgfpathlineto{\pgfqpoint{2.909973in}{2.899915in}}%
\pgfpathlineto{\pgfqpoint{2.912223in}{2.901035in}}%
\pgfpathlineto{\pgfqpoint{2.912823in}{2.901295in}}%
\pgfpathlineto{\pgfqpoint{2.917322in}{2.902805in}}%
\pgfpathlineto{\pgfqpoint{2.917922in}{2.903066in}}%
\pgfpathlineto{\pgfqpoint{2.921071in}{2.904341in}}%
\pgfpathlineto{\pgfqpoint{2.921671in}{2.904732in}}%
\pgfpathlineto{\pgfqpoint{2.925120in}{2.906060in}}%
\pgfpathlineto{\pgfqpoint{2.925720in}{2.906372in}}%
\pgfpathlineto{\pgfqpoint{2.928569in}{2.907596in}}%
\pgfpathlineto{\pgfqpoint{2.929169in}{2.908039in}}%
\pgfpathlineto{\pgfqpoint{2.933368in}{2.909497in}}%
\pgfpathlineto{\pgfqpoint{2.933668in}{2.909705in}}%
\pgfpathlineto{\pgfqpoint{2.936367in}{2.910538in}}%
\pgfpathlineto{\pgfqpoint{2.936967in}{2.910772in}}%
\pgfpathlineto{\pgfqpoint{2.940266in}{2.912074in}}%
\pgfpathlineto{\pgfqpoint{2.940866in}{2.912543in}}%
\pgfpathlineto{\pgfqpoint{2.943265in}{2.913688in}}%
\pgfpathlineto{\pgfqpoint{2.943865in}{2.913871in}}%
\pgfpathlineto{\pgfqpoint{2.948364in}{2.915381in}}%
\pgfpathlineto{\pgfqpoint{2.948814in}{2.915719in}}%
\pgfpathlineto{\pgfqpoint{2.953163in}{2.916839in}}%
\pgfpathlineto{\pgfqpoint{2.959462in}{2.918661in}}%
\pgfpathlineto{\pgfqpoint{2.960061in}{2.918844in}}%
\pgfpathlineto{\pgfqpoint{2.965610in}{2.920536in}}%
\pgfpathlineto{\pgfqpoint{2.966210in}{2.920822in}}%
\pgfpathlineto{\pgfqpoint{2.972658in}{2.922671in}}%
\pgfpathlineto{\pgfqpoint{2.973108in}{2.922853in}}%
\pgfpathlineto{\pgfqpoint{2.975358in}{2.923712in}}%
\pgfpathlineto{\pgfqpoint{2.983306in}{2.925821in}}%
\pgfpathlineto{\pgfqpoint{2.983906in}{2.926030in}}%
\pgfpathlineto{\pgfqpoint{2.987355in}{2.927357in}}%
\pgfpathlineto{\pgfqpoint{2.987955in}{2.927592in}}%
\pgfpathlineto{\pgfqpoint{2.991104in}{2.928607in}}%
\pgfpathlineto{\pgfqpoint{2.991704in}{2.928789in}}%
\pgfpathlineto{\pgfqpoint{2.995303in}{2.930143in}}%
\pgfpathlineto{\pgfqpoint{2.995903in}{2.930404in}}%
\pgfpathlineto{\pgfqpoint{3.001152in}{2.932044in}}%
\pgfpathlineto{\pgfqpoint{3.001752in}{2.932174in}}%
\pgfpathlineto{\pgfqpoint{3.006850in}{2.933788in}}%
\pgfpathlineto{\pgfqpoint{3.007450in}{2.934049in}}%
\pgfpathlineto{\pgfqpoint{3.015248in}{2.936132in}}%
\pgfpathlineto{\pgfqpoint{3.015848in}{2.936314in}}%
\pgfpathlineto{\pgfqpoint{3.024096in}{2.938475in}}%
\pgfpathlineto{\pgfqpoint{3.024696in}{2.938761in}}%
\pgfpathlineto{\pgfqpoint{3.030545in}{2.940506in}}%
\pgfpathlineto{\pgfqpoint{3.031145in}{2.940740in}}%
\pgfpathlineto{\pgfqpoint{3.036093in}{2.941912in}}%
\pgfpathlineto{\pgfqpoint{3.036693in}{2.942094in}}%
\pgfpathlineto{\pgfqpoint{3.040592in}{2.943500in}}%
\pgfpathlineto{\pgfqpoint{3.041192in}{2.943734in}}%
\pgfpathlineto{\pgfqpoint{3.052739in}{2.946468in}}%
\pgfpathlineto{\pgfqpoint{3.053339in}{2.946676in}}%
\pgfpathlineto{\pgfqpoint{3.066686in}{2.949723in}}%
\pgfpathlineto{\pgfqpoint{3.067136in}{2.949879in}}%
\pgfpathlineto{\pgfqpoint{3.072235in}{2.951051in}}%
\pgfpathlineto{\pgfqpoint{3.072835in}{2.951233in}}%
\pgfpathlineto{\pgfqpoint{3.083782in}{2.953862in}}%
\pgfpathlineto{\pgfqpoint{3.084382in}{2.954097in}}%
\pgfpathlineto{\pgfqpoint{3.094729in}{2.956622in}}%
\pgfpathlineto{\pgfqpoint{3.095329in}{2.956831in}}%
\pgfpathlineto{\pgfqpoint{3.100578in}{2.958028in}}%
\pgfpathlineto{\pgfqpoint{3.101178in}{2.958289in}}%
\pgfpathlineto{\pgfqpoint{3.141518in}{2.966021in}}%
\pgfpathlineto{\pgfqpoint{3.142118in}{2.966256in}}%
\pgfpathlineto{\pgfqpoint{3.161463in}{2.970343in}}%
\pgfpathlineto{\pgfqpoint{3.162063in}{2.970656in}}%
\pgfpathlineto{\pgfqpoint{3.179609in}{2.974431in}}%
\pgfpathlineto{\pgfqpoint{3.180209in}{2.974666in}}%
\pgfpathlineto{\pgfqpoint{3.186807in}{2.976540in}}%
\pgfpathlineto{\pgfqpoint{3.187407in}{2.976670in}}%
\pgfpathlineto{\pgfqpoint{3.190557in}{2.977686in}}%
\pgfpathlineto{\pgfqpoint{3.191156in}{2.977920in}}%
\pgfpathlineto{\pgfqpoint{3.200304in}{2.980237in}}%
\pgfpathlineto{\pgfqpoint{3.200754in}{2.980420in}}%
\pgfpathlineto{\pgfqpoint{3.206003in}{2.981617in}}%
\pgfpathlineto{\pgfqpoint{3.214701in}{2.983856in}}%
\pgfpathlineto{\pgfqpoint{3.215001in}{2.983987in}}%
\pgfpathlineto{\pgfqpoint{3.219500in}{2.984976in}}%
\pgfpathlineto{\pgfqpoint{3.232247in}{2.987918in}}%
\pgfpathlineto{\pgfqpoint{3.232846in}{2.988048in}}%
\pgfpathlineto{\pgfqpoint{3.235996in}{2.989324in}}%
\pgfpathlineto{\pgfqpoint{3.236296in}{2.989454in}}%
\pgfpathlineto{\pgfqpoint{3.241244in}{2.990391in}}%
\pgfpathlineto{\pgfqpoint{3.241844in}{2.990652in}}%
\pgfpathlineto{\pgfqpoint{3.248443in}{2.991954in}}%
\pgfpathlineto{\pgfqpoint{3.249043in}{2.992084in}}%
\pgfpathlineto{\pgfqpoint{3.256391in}{2.994089in}}%
\pgfpathlineto{\pgfqpoint{3.256991in}{2.994297in}}%
\pgfpathlineto{\pgfqpoint{3.278885in}{2.998827in}}%
\pgfpathlineto{\pgfqpoint{3.279485in}{2.999036in}}%
\pgfpathlineto{\pgfqpoint{3.290283in}{3.001639in}}%
\pgfpathlineto{\pgfqpoint{3.290733in}{3.001743in}}%
\pgfpathlineto{\pgfqpoint{3.295981in}{3.002602in}}%
\pgfpathlineto{\pgfqpoint{3.296581in}{3.002837in}}%
\pgfpathlineto{\pgfqpoint{3.330923in}{3.009528in}}%
\pgfpathlineto{\pgfqpoint{3.331523in}{3.009762in}}%
\pgfpathlineto{\pgfqpoint{3.407405in}{3.022234in}}%
\pgfpathlineto{\pgfqpoint{3.415353in}{3.023301in}}%
\pgfpathlineto{\pgfqpoint{3.415653in}{3.023353in}}%
\pgfpathlineto{\pgfqpoint{3.429899in}{3.024916in}}%
\pgfpathlineto{\pgfqpoint{3.445646in}{3.026920in}}%
\pgfpathlineto{\pgfqpoint{3.456143in}{3.028561in}}%
\pgfpathlineto{\pgfqpoint{3.456443in}{3.028639in}}%
\pgfpathlineto{\pgfqpoint{3.463341in}{3.029758in}}%
\pgfpathlineto{\pgfqpoint{3.463941in}{3.029863in}}%
\pgfpathlineto{\pgfqpoint{3.481937in}{3.032154in}}%
\pgfpathlineto{\pgfqpoint{3.482537in}{3.032284in}}%
\pgfpathlineto{\pgfqpoint{3.493784in}{3.034002in}}%
\pgfpathlineto{\pgfqpoint{3.505331in}{3.035278in}}%
\pgfpathlineto{\pgfqpoint{3.515829in}{3.036918in}}%
\pgfpathlineto{\pgfqpoint{3.516429in}{3.037049in}}%
\pgfpathlineto{\pgfqpoint{3.526926in}{3.038220in}}%
\pgfpathlineto{\pgfqpoint{3.536674in}{3.039288in}}%
\pgfpathlineto{\pgfqpoint{3.547771in}{3.040485in}}%
\pgfpathlineto{\pgfqpoint{3.671792in}{3.050509in}}%
\pgfpathlineto{\pgfqpoint{3.786214in}{3.056394in}}%
\pgfpathlineto{\pgfqpoint{3.999614in}{3.065064in}}%
\pgfpathlineto{\pgfqpoint{4.027807in}{3.065949in}}%
\pgfpathlineto{\pgfqpoint{4.033355in}{3.066209in}}%
\pgfpathlineto{\pgfqpoint{4.151077in}{3.069854in}}%
\pgfpathlineto{\pgfqpoint{4.156476in}{3.070063in}}%
\pgfpathlineto{\pgfqpoint{4.232358in}{3.072640in}}%
\pgfpathlineto{\pgfqpoint{4.269549in}{3.073838in}}%
\pgfpathlineto{\pgfqpoint{4.274348in}{3.074046in}}%
\pgfpathlineto{\pgfqpoint{4.310639in}{3.075348in}}%
\pgfpathlineto{\pgfqpoint{4.340332in}{3.076363in}}%
\pgfpathlineto{\pgfqpoint{4.367776in}{3.077691in}}%
\pgfpathlineto{\pgfqpoint{4.380373in}{3.078056in}}%
\pgfpathlineto{\pgfqpoint{4.413365in}{3.079123in}}%
\pgfpathlineto{\pgfqpoint{4.418614in}{3.079332in}}%
\pgfpathlineto{\pgfqpoint{4.506643in}{3.082456in}}%
\pgfpathlineto{\pgfqpoint{4.544284in}{3.083628in}}%
\pgfpathlineto{\pgfqpoint{4.579375in}{3.084487in}}%
\pgfpathlineto{\pgfqpoint{4.621665in}{3.085580in}}%
\pgfpathlineto{\pgfqpoint{4.706395in}{3.087299in}}%
\pgfpathlineto{\pgfqpoint{4.897899in}{3.090735in}}%
\pgfpathlineto{\pgfqpoint{4.967333in}{3.092141in}}%
\pgfpathlineto{\pgfqpoint{4.982029in}{3.092376in}}%
\pgfpathlineto{\pgfqpoint{5.040065in}{3.093365in}}%
\pgfpathlineto{\pgfqpoint{5.247466in}{3.096281in}}%
\pgfpathlineto{\pgfqpoint{5.275509in}{3.096437in}}%
\pgfpathlineto{\pgfqpoint{5.275509in}{3.096437in}}%
\pgfusepath{stroke}%
\end{pgfscope}%
\begin{pgfscope}%
\pgfsetrectcap%
\pgfsetmiterjoin%
\pgfsetlinewidth{0.803000pt}%
\definecolor{currentstroke}{rgb}{0.000000,0.000000,0.000000}%
\pgfsetstrokecolor{currentstroke}%
\pgfsetdash{}{0pt}%
\pgfpathmoveto{\pgfqpoint{0.565433in}{0.528318in}}%
\pgfpathlineto{\pgfqpoint{0.565433in}{3.131951in}}%
\pgfusepath{stroke}%
\end{pgfscope}%
\begin{pgfscope}%
\pgfsetrectcap%
\pgfsetmiterjoin%
\pgfsetlinewidth{0.803000pt}%
\definecolor{currentstroke}{rgb}{0.000000,0.000000,0.000000}%
\pgfsetstrokecolor{currentstroke}%
\pgfsetdash{}{0pt}%
\pgfpathmoveto{\pgfqpoint{5.269125in}{0.528318in}}%
\pgfpathlineto{\pgfqpoint{5.269125in}{3.131951in}}%
\pgfusepath{stroke}%
\end{pgfscope}%
\begin{pgfscope}%
\pgfsetrectcap%
\pgfsetmiterjoin%
\pgfsetlinewidth{0.803000pt}%
\definecolor{currentstroke}{rgb}{0.000000,0.000000,0.000000}%
\pgfsetstrokecolor{currentstroke}%
\pgfsetdash{}{0pt}%
\pgfpathmoveto{\pgfqpoint{0.565433in}{0.528318in}}%
\pgfpathlineto{\pgfqpoint{5.269125in}{0.528318in}}%
\pgfusepath{stroke}%
\end{pgfscope}%
\begin{pgfscope}%
\pgfsetrectcap%
\pgfsetmiterjoin%
\pgfsetlinewidth{0.803000pt}%
\definecolor{currentstroke}{rgb}{0.000000,0.000000,0.000000}%
\pgfsetstrokecolor{currentstroke}%
\pgfsetdash{}{0pt}%
\pgfpathmoveto{\pgfqpoint{0.565433in}{3.131951in}}%
\pgfpathlineto{\pgfqpoint{5.269125in}{3.131951in}}%
\pgfusepath{stroke}%
\end{pgfscope}%
\begin{pgfscope}%
\pgfpathrectangle{\pgfqpoint{0.565433in}{0.528318in}}{\pgfqpoint{4.703693in}{2.603633in}}%
\pgfusepath{clip}%
\pgfsetbuttcap%
\pgfsetroundjoin%
\definecolor{currentfill}{rgb}{0.082353,0.396078,0.752941}%
\pgfsetfillcolor{currentfill}%
\pgfsetlinewidth{1.003750pt}%
\definecolor{currentstroke}{rgb}{0.082353,0.396078,0.752941}%
\pgfsetstrokecolor{currentstroke}%
\pgfsetdash{}{0pt}%
\pgfsys@defobject{currentmarker}{\pgfqpoint{-0.031056in}{-0.031056in}}{\pgfqpoint{0.031056in}{0.031056in}}{%
\pgfpathmoveto{\pgfqpoint{0.000000in}{-0.031056in}}%
\pgfpathcurveto{\pgfqpoint{0.008236in}{-0.031056in}}{\pgfqpoint{0.016136in}{-0.027784in}}{\pgfqpoint{0.021960in}{-0.021960in}}%
\pgfpathcurveto{\pgfqpoint{0.027784in}{-0.016136in}}{\pgfqpoint{0.031056in}{-0.008236in}}{\pgfqpoint{0.031056in}{0.000000in}}%
\pgfpathcurveto{\pgfqpoint{0.031056in}{0.008236in}}{\pgfqpoint{0.027784in}{0.016136in}}{\pgfqpoint{0.021960in}{0.021960in}}%
\pgfpathcurveto{\pgfqpoint{0.016136in}{0.027784in}}{\pgfqpoint{0.008236in}{0.031056in}}{\pgfqpoint{0.000000in}{0.031056in}}%
\pgfpathcurveto{\pgfqpoint{-0.008236in}{0.031056in}}{\pgfqpoint{-0.016136in}{0.027784in}}{\pgfqpoint{-0.021960in}{0.021960in}}%
\pgfpathcurveto{\pgfqpoint{-0.027784in}{0.016136in}}{\pgfqpoint{-0.031056in}{0.008236in}}{\pgfqpoint{-0.031056in}{0.000000in}}%
\pgfpathcurveto{\pgfqpoint{-0.031056in}{-0.008236in}}{\pgfqpoint{-0.027784in}{-0.016136in}}{\pgfqpoint{-0.021960in}{-0.021960in}}%
\pgfpathcurveto{\pgfqpoint{-0.016136in}{-0.027784in}}{\pgfqpoint{-0.008236in}{-0.031056in}}{\pgfqpoint{0.000000in}{-0.031056in}}%
\pgfpathlineto{\pgfqpoint{0.000000in}{-0.031056in}}%
\pgfpathclose%
\pgfusepath{stroke,fill}%
}%
\begin{pgfscope}%
\pgfsys@transformshift{1.076811in}{1.830134in}%
\pgfsys@useobject{currentmarker}{}%
\end{pgfscope}%
\begin{pgfscope}%
\pgfsys@transformshift{1.388736in}{2.871588in}%
\pgfsys@useobject{currentmarker}{}%
\end{pgfscope}%
\begin{pgfscope}%
\pgfsys@transformshift{3.517045in}{3.105915in}%
\pgfsys@useobject{currentmarker}{}%
\end{pgfscope}%
\end{pgfscope}%
\begin{pgfscope}%
\pgfpathrectangle{\pgfqpoint{0.565433in}{0.528318in}}{\pgfqpoint{4.703693in}{2.603633in}}%
\pgfusepath{clip}%
\pgfsetbuttcap%
\pgfsetroundjoin%
\definecolor{currentfill}{rgb}{0.180392,0.490196,0.196078}%
\pgfsetfillcolor{currentfill}%
\pgfsetlinewidth{1.003750pt}%
\definecolor{currentstroke}{rgb}{0.180392,0.490196,0.196078}%
\pgfsetstrokecolor{currentstroke}%
\pgfsetdash{}{0pt}%
\pgfsys@defobject{currentmarker}{\pgfqpoint{-0.031056in}{-0.031056in}}{\pgfqpoint{0.031056in}{0.031056in}}{%
\pgfpathmoveto{\pgfqpoint{-0.031056in}{-0.031056in}}%
\pgfpathlineto{\pgfqpoint{0.031056in}{-0.031056in}}%
\pgfpathlineto{\pgfqpoint{0.031056in}{0.031056in}}%
\pgfpathlineto{\pgfqpoint{-0.031056in}{0.031056in}}%
\pgfpathlineto{\pgfqpoint{-0.031056in}{-0.031056in}}%
\pgfpathclose%
\pgfusepath{stroke,fill}%
}%
\begin{pgfscope}%
\pgfsys@transformshift{2.432487in}{1.830134in}%
\pgfsys@useobject{currentmarker}{}%
\end{pgfscope}%
\begin{pgfscope}%
\pgfsys@transformshift{2.851037in}{2.871588in}%
\pgfsys@useobject{currentmarker}{}%
\end{pgfscope}%
\end{pgfscope}%
\begin{pgfscope}%
\pgfsetroundcap%
\pgfsetroundjoin%
\pgfsetlinewidth{1.003750pt}%
\definecolor{currentstroke}{rgb}{0.000000,0.000000,0.000000}%
\pgfsetstrokecolor{currentstroke}%
\pgfsetdash{}{0pt}%
\pgfpathmoveto{\pgfqpoint{4.844246in}{2.802606in}}%
\pgfpathquadraticcurveto{\pgfqpoint{5.045392in}{2.946198in}}{\pgfqpoint{5.233900in}{3.080768in}}%
\pgfusepath{stroke}%
\end{pgfscope}%
\begin{pgfscope}%
\pgfsetroundcap%
\pgfsetroundjoin%
\pgfsetlinewidth{1.003750pt}%
\definecolor{currentstroke}{rgb}{0.000000,0.000000,0.000000}%
\pgfsetstrokecolor{currentstroke}%
\pgfsetdash{}{0pt}%
\pgfpathmoveto{\pgfqpoint{5.184815in}{3.073032in}}%
\pgfpathlineto{\pgfqpoint{5.233900in}{3.080768in}}%
\pgfpathlineto{\pgfqpoint{5.210638in}{3.036859in}}%
\pgfusepath{stroke}%
\end{pgfscope}%
\begin{pgfscope}%
\definecolor{textcolor}{rgb}{0.000000,0.000000,0.000000}%
\pgfsetstrokecolor{textcolor}%
\pgfsetfillcolor{textcolor}%
\pgftext[x=4.048436in, y=2.681532in, left, base]{\color{textcolor}{\rmfamily\fontsize{8.000000}{9.600000}\selectfont\catcode`\^=\active\def^{\ifmmode\sp\else\^{}\fi}\catcode`\%=\active\def%{\%}CTaps p99 $\approx$ 62.93 $\mu$s}}%
\end{pgfscope}%
\begin{pgfscope}%
\definecolor{textcolor}{rgb}{0.000000,0.000000,0.000000}%
\pgfsetstrokecolor{textcolor}%
\pgfsetfillcolor{textcolor}%
\pgftext[x=4.316778in, y=2.559927in, left, base]{\color{textcolor}{\rmfamily\fontsize{8.000000}{9.600000}\selectfont\catcode`\^=\active\def^{\ifmmode\sp\else\^{}\fi}\catcode`\%=\active\def%{\%}(outside x-range)}}%
\end{pgfscope}%
\begin{pgfscope}%
\pgfsetbuttcap%
\pgfsetmiterjoin%
\definecolor{currentfill}{rgb}{1.000000,1.000000,1.000000}%
\pgfsetfillcolor{currentfill}%
\pgfsetfillopacity{0.800000}%
\pgfsetlinewidth{1.003750pt}%
\definecolor{currentstroke}{rgb}{0.501961,0.501961,0.501961}%
\pgfsetstrokecolor{currentstroke}%
\pgfsetstrokeopacity{0.800000}%
\pgfsetdash{}{0pt}%
\pgfpathmoveto{\pgfqpoint{3.827708in}{1.024739in}}%
\pgfpathlineto{\pgfqpoint{5.175051in}{1.024739in}}%
\pgfpathquadraticcurveto{\pgfqpoint{5.199357in}{1.024739in}}{\pgfqpoint{5.199357in}{1.049044in}}%
\pgfpathlineto{\pgfqpoint{5.199357in}{1.515710in}}%
\pgfpathquadraticcurveto{\pgfqpoint{5.199357in}{1.540015in}}{\pgfqpoint{5.175051in}{1.540015in}}%
\pgfpathlineto{\pgfqpoint{3.827708in}{1.540015in}}%
\pgfpathquadraticcurveto{\pgfqpoint{3.803403in}{1.540015in}}{\pgfqpoint{3.803403in}{1.515710in}}%
\pgfpathlineto{\pgfqpoint{3.803403in}{1.049044in}}%
\pgfpathquadraticcurveto{\pgfqpoint{3.803403in}{1.024739in}}{\pgfqpoint{3.827708in}{1.024739in}}%
\pgfpathlineto{\pgfqpoint{3.827708in}{1.024739in}}%
\pgfpathclose%
\pgfusepath{stroke,fill}%
\end{pgfscope}%
\begin{pgfscope}%
\definecolor{textcolor}{rgb}{0.000000,0.000000,0.000000}%
\pgfsetstrokecolor{textcolor}%
\pgfsetfillcolor{textcolor}%
\pgftext[x=5.175051in,y=1.049044in,right,bottom]{\color{textcolor}{\rmfamily\fontsize{7.000000}{8.400000}\selectfont\catcode`\^=\active\def^{\ifmmode\sp\else\^{}\fi}\catcode`\%=\active\def%{\%}\setlength{\tabcolsep}{4pt}\begin{tabular}{lrrr}      & \multicolumn{3}{c}{/ \si{\micro\second}} \\      & p50 & p90 & p99 \\Base  & 20.2 & 22.2 & 36.4 \\CTaps & 29.2 & 32.0 & 62.9\end{tabular}}}%
\end{pgfscope}%
\begin{pgfscope}%
\pgfsetbuttcap%
\pgfsetmiterjoin%
\definecolor{currentfill}{rgb}{1.000000,1.000000,1.000000}%
\pgfsetfillcolor{currentfill}%
\pgfsetfillopacity{0.800000}%
\pgfsetlinewidth{1.003750pt}%
\definecolor{currentstroke}{rgb}{0.800000,0.800000,0.800000}%
\pgfsetstrokecolor{currentstroke}%
\pgfsetstrokeopacity{0.800000}%
\pgfsetdash{}{0pt}%
\pgfpathmoveto{\pgfqpoint{4.112325in}{0.583873in}}%
\pgfpathlineto{\pgfqpoint{5.191347in}{0.583873in}}%
\pgfpathquadraticcurveto{\pgfqpoint{5.213570in}{0.583873in}}{\pgfqpoint{5.213570in}{0.606095in}}%
\pgfpathlineto{\pgfqpoint{5.213570in}{0.904861in}}%
\pgfpathquadraticcurveto{\pgfqpoint{5.213570in}{0.927083in}}{\pgfqpoint{5.191347in}{0.927083in}}%
\pgfpathlineto{\pgfqpoint{4.112325in}{0.927083in}}%
\pgfpathquadraticcurveto{\pgfqpoint{4.090102in}{0.927083in}}{\pgfqpoint{4.090102in}{0.904861in}}%
\pgfpathlineto{\pgfqpoint{4.090102in}{0.606095in}}%
\pgfpathquadraticcurveto{\pgfqpoint{4.090102in}{0.583873in}}{\pgfqpoint{4.112325in}{0.583873in}}%
\pgfpathlineto{\pgfqpoint{4.112325in}{0.583873in}}%
\pgfpathclose%
\pgfusepath{stroke,fill}%
\end{pgfscope}%
\begin{pgfscope}%
\pgfsetrectcap%
\pgfsetroundjoin%
\pgfsetlinewidth{1.806750pt}%
\definecolor{currentstroke}{rgb}{0.082353,0.396078,0.752941}%
\pgfsetstrokecolor{currentstroke}%
\pgfsetdash{}{0pt}%
\pgfpathmoveto{\pgfqpoint{4.134547in}{0.843750in}}%
\pgfpathlineto{\pgfqpoint{4.245658in}{0.843750in}}%
\pgfpathlineto{\pgfqpoint{4.356769in}{0.843750in}}%
\pgfusepath{stroke}%
\end{pgfscope}%
\begin{pgfscope}%
\definecolor{textcolor}{rgb}{0.000000,0.000000,0.000000}%
\pgfsetstrokecolor{textcolor}%
\pgfsetfillcolor{textcolor}%
\pgftext[x=4.445658in,y=0.804861in,left,base]{\color{textcolor}{\rmfamily\fontsize{8.000000}{9.600000}\selectfont\catcode`\^=\active\def^{\ifmmode\sp\else\^{}\fi}\catcode`\%=\active\def%{\%}Baseline UDP}}%
\end{pgfscope}%
\begin{pgfscope}%
\pgfsetrectcap%
\pgfsetroundjoin%
\pgfsetlinewidth{1.806750pt}%
\definecolor{currentstroke}{rgb}{0.180392,0.490196,0.196078}%
\pgfsetstrokecolor{currentstroke}%
\pgfsetdash{}{0pt}%
\pgfpathmoveto{\pgfqpoint{4.134547in}{0.688812in}}%
\pgfpathlineto{\pgfqpoint{4.245658in}{0.688812in}}%
\pgfpathlineto{\pgfqpoint{4.356769in}{0.688812in}}%
\pgfusepath{stroke}%
\end{pgfscope}%
\begin{pgfscope}%
\definecolor{textcolor}{rgb}{0.000000,0.000000,0.000000}%
\pgfsetstrokecolor{textcolor}%
\pgfsetfillcolor{textcolor}%
\pgftext[x=4.445658in,y=0.649923in,left,base]{\color{textcolor}{\rmfamily\fontsize{8.000000}{9.600000}\selectfont\catcode`\^=\active\def^{\ifmmode\sp\else\^{}\fi}\catcode`\%=\active\def%{\%}CTaps}}%
\end{pgfscope}%
\end{pgfpicture}%
\makeatother%
\endgroup%

    \caption[UDP RTT experiment]{Overhead of CTaps when performing echoes over UDP on localhost.
        It shows that CTaps has a \qty{45}{\percent} increase in median RTT, but a \qty{73}{\percent} increase at p99,
        suggesting a disproportionate tail overhead. ($N = \num{100000}$) }
    \label{fig:udp-overhead}
\end{figure}

This experiment shows that while the overhead in absolute numbers
is small, the relative increase over a pure UDP client is quite large.
For real-world applications this effect would be dwarfed
by the latency of the network. The more interesting
finding is that CTaps has a longer tail, which is
discussed further in \cref{sec:disc-perf}.

\subsection{Candidate Racing Experiment}\label{sec:handshake-kde}
This experiment can be seen in~\cref{fig:handshake-kde}. It
demonstrates the benefits of candidate racing and how it helps find the
endpoint with the smallest delay.
This experiment was modeled after the theoretical
scenario displayed in~\cref{fig:racingSeqDiagram} in \cref{chap:background}.

Two possible paths were used with different Remote Endpoints.
The slow endpoint had a delay of \qty{150}{\milli\second}
while the fast endpoint had a delay of \qty{50}{\milli\second}.
Both the Picoquic
and CTaps client were configured to use the slow endpoint as an initial candidate.
The CTaps client was also given the fast endpoint as an alternative candidate to
perform candidate racing. This reflects a realistic post-DNS-resolution
scenario, where multiple endpoints are available, but their ordering
does not reflect actual performance. Simple sequential clients will always
use the first candidate, regardless of its performance.

In order to gather a distribution of handshake times, the test was run 200~times with a
jitter of \qty{7}{\milli\second} and a latency correlation of \qty{25}{\percent}.
The theoretical ideal handshake times for each endpoint are featured
in the graph. With a \qty{250}{\milli\second} stagger delay, 2~RTTs for a full handshake and delays of
\qtylist{50;150}{\milli\second},
the theoretical minimum handshake times work out to:
$\qty{50}{\milli\second} \cdot 4 + \qty{250}{\milli\second} = \qty{450}{\milli\second}$
and
$\qty{150}{\milli\second} \cdot 4 = \qty{600}{\milli\second}$ respectively. 

\begin{figure}[H]
    %% Creator: Matplotlib, PGF backend
%%
%% To include the figure in your LaTeX document, write
%%   \input{<filename>.pgf}
%%
%% Make sure the required packages are loaded in your preamble
%%   \usepackage{pgf}
%%
%% Also ensure that all the required font packages are loaded; for instance,
%% the lmodern package is sometimes necessary when using math font.
%%   \usepackage{lmodern}
%%
%% Figures using additional raster images can only be included by \input if
%% they are in the same directory as the main LaTeX file. For loading figures
%% from other directories you can use the `import` package
%%   \usepackage{import}
%%
%% and then include the figures with
%%   \import{<path to file>}{<filename>.pgf}
%%
%% Matplotlib used the following preamble
%%   \def\mathdefault#1{#1}
%%   \everymath=\expandafter{\the\everymath\displaystyle}
%%   \IfFileExists{scrextend.sty}{
%%     \usepackage[fontsize=10.000000pt]{scrextend}
%%   }{
%%     \renewcommand{\normalsize}{\fontsize{10.000000}{12.000000}\selectfont}
%%     \normalsize
%%   }
%%   
%%   \ifdefined\pdftexversion\else  % non-pdftex case.
%%     \usepackage{fontspec}
%%     \setmainfont{DejaVuSerif.ttf}[Path=\detokenize{/home/ikhovind/Documents/Skole/taps/benchmark/venv/lib/python3.11/site-packages/matplotlib/mpl-data/fonts/ttf/}]
%%     \setsansfont{DejaVuSans.ttf}[Path=\detokenize{/home/ikhovind/Documents/Skole/taps/benchmark/venv/lib/python3.11/site-packages/matplotlib/mpl-data/fonts/ttf/}]
%%     \setmonofont{DejaVuSansMono.ttf}[Path=\detokenize{/home/ikhovind/Documents/Skole/taps/benchmark/venv/lib/python3.11/site-packages/matplotlib/mpl-data/fonts/ttf/}]
%%   \fi
%%   \makeatletter\@ifpackageloaded{underscore}{}{\usepackage[strings]{underscore}}\makeatother
%%
\begingroup%
\makeatletter%
\begin{pgfpicture}%
\pgfpathrectangle{\pgfpointorigin}{\pgfqpoint{5.645393in}{3.255629in}}%
\pgfusepath{use as bounding box, clip}%
\begin{pgfscope}%
\pgfsetbuttcap%
\pgfsetmiterjoin%
\definecolor{currentfill}{rgb}{1.000000,1.000000,1.000000}%
\pgfsetfillcolor{currentfill}%
\pgfsetlinewidth{0.000000pt}%
\definecolor{currentstroke}{rgb}{1.000000,1.000000,1.000000}%
\pgfsetstrokecolor{currentstroke}%
\pgfsetdash{}{0pt}%
\pgfpathmoveto{\pgfqpoint{0.000000in}{0.000000in}}%
\pgfpathlineto{\pgfqpoint{5.645393in}{0.000000in}}%
\pgfpathlineto{\pgfqpoint{5.645393in}{3.255629in}}%
\pgfpathlineto{\pgfqpoint{0.000000in}{3.255629in}}%
\pgfpathlineto{\pgfqpoint{0.000000in}{0.000000in}}%
\pgfpathclose%
\pgfusepath{fill}%
\end{pgfscope}%
\begin{pgfscope}%
\pgfsetbuttcap%
\pgfsetmiterjoin%
\definecolor{currentfill}{rgb}{1.000000,1.000000,1.000000}%
\pgfsetfillcolor{currentfill}%
\pgfsetlinewidth{0.000000pt}%
\definecolor{currentstroke}{rgb}{0.000000,0.000000,0.000000}%
\pgfsetstrokecolor{currentstroke}%
\pgfsetstrokeopacity{0.000000}%
\pgfsetdash{}{0pt}%
\pgfpathmoveto{\pgfqpoint{0.784801in}{0.521603in}}%
\pgfpathlineto{\pgfqpoint{5.134273in}{0.521603in}}%
\pgfpathlineto{\pgfqpoint{5.134273in}{2.945669in}}%
\pgfpathlineto{\pgfqpoint{0.784801in}{2.945669in}}%
\pgfpathlineto{\pgfqpoint{0.784801in}{0.521603in}}%
\pgfpathclose%
\pgfusepath{fill}%
\end{pgfscope}%
\begin{pgfscope}%
\pgfpathrectangle{\pgfqpoint{0.784801in}{0.521603in}}{\pgfqpoint{4.349472in}{2.424065in}}%
\pgfusepath{clip}%
\pgfsetbuttcap%
\pgfsetroundjoin%
\definecolor{currentfill}{rgb}{0.564706,0.792157,0.976471}%
\pgfsetfillcolor{currentfill}%
\pgfsetfillopacity{0.400000}%
\pgfsetlinewidth{1.505625pt}%
\definecolor{currentstroke}{rgb}{0.564706,0.792157,0.976471}%
\pgfsetstrokecolor{currentstroke}%
\pgfsetdash{}{0pt}%
\pgfsys@defobject{currentmarker}{\pgfqpoint{2.654882in}{0.521603in}}{\pgfqpoint{3.484202in}{2.830237in}}{%
\pgfpathmoveto{\pgfqpoint{2.654882in}{0.522412in}}%
\pgfpathlineto{\pgfqpoint{2.654882in}{0.521603in}}%
\pgfpathlineto{\pgfqpoint{2.659050in}{0.521603in}}%
\pgfpathlineto{\pgfqpoint{2.663217in}{0.521603in}}%
\pgfpathlineto{\pgfqpoint{2.667385in}{0.521603in}}%
\pgfpathlineto{\pgfqpoint{2.671552in}{0.521603in}}%
\pgfpathlineto{\pgfqpoint{2.675720in}{0.521603in}}%
\pgfpathlineto{\pgfqpoint{2.679887in}{0.521603in}}%
\pgfpathlineto{\pgfqpoint{2.684054in}{0.521603in}}%
\pgfpathlineto{\pgfqpoint{2.688222in}{0.521603in}}%
\pgfpathlineto{\pgfqpoint{2.692389in}{0.521603in}}%
\pgfpathlineto{\pgfqpoint{2.696557in}{0.521603in}}%
\pgfpathlineto{\pgfqpoint{2.700724in}{0.521603in}}%
\pgfpathlineto{\pgfqpoint{2.704892in}{0.521603in}}%
\pgfpathlineto{\pgfqpoint{2.709059in}{0.521603in}}%
\pgfpathlineto{\pgfqpoint{2.713227in}{0.521603in}}%
\pgfpathlineto{\pgfqpoint{2.717394in}{0.521603in}}%
\pgfpathlineto{\pgfqpoint{2.721561in}{0.521603in}}%
\pgfpathlineto{\pgfqpoint{2.725729in}{0.521603in}}%
\pgfpathlineto{\pgfqpoint{2.729896in}{0.521603in}}%
\pgfpathlineto{\pgfqpoint{2.734064in}{0.521603in}}%
\pgfpathlineto{\pgfqpoint{2.738231in}{0.521603in}}%
\pgfpathlineto{\pgfqpoint{2.742399in}{0.521603in}}%
\pgfpathlineto{\pgfqpoint{2.746566in}{0.521603in}}%
\pgfpathlineto{\pgfqpoint{2.750733in}{0.521603in}}%
\pgfpathlineto{\pgfqpoint{2.754901in}{0.521603in}}%
\pgfpathlineto{\pgfqpoint{2.759068in}{0.521603in}}%
\pgfpathlineto{\pgfqpoint{2.763236in}{0.521603in}}%
\pgfpathlineto{\pgfqpoint{2.767403in}{0.521603in}}%
\pgfpathlineto{\pgfqpoint{2.771571in}{0.521603in}}%
\pgfpathlineto{\pgfqpoint{2.775738in}{0.521603in}}%
\pgfpathlineto{\pgfqpoint{2.779906in}{0.521603in}}%
\pgfpathlineto{\pgfqpoint{2.784073in}{0.521603in}}%
\pgfpathlineto{\pgfqpoint{2.788240in}{0.521603in}}%
\pgfpathlineto{\pgfqpoint{2.792408in}{0.521603in}}%
\pgfpathlineto{\pgfqpoint{2.796575in}{0.521603in}}%
\pgfpathlineto{\pgfqpoint{2.800743in}{0.521603in}}%
\pgfpathlineto{\pgfqpoint{2.804910in}{0.521603in}}%
\pgfpathlineto{\pgfqpoint{2.809078in}{0.521603in}}%
\pgfpathlineto{\pgfqpoint{2.813245in}{0.521603in}}%
\pgfpathlineto{\pgfqpoint{2.817412in}{0.521603in}}%
\pgfpathlineto{\pgfqpoint{2.821580in}{0.521603in}}%
\pgfpathlineto{\pgfqpoint{2.825747in}{0.521603in}}%
\pgfpathlineto{\pgfqpoint{2.829915in}{0.521603in}}%
\pgfpathlineto{\pgfqpoint{2.834082in}{0.521603in}}%
\pgfpathlineto{\pgfqpoint{2.838250in}{0.521603in}}%
\pgfpathlineto{\pgfqpoint{2.842417in}{0.521603in}}%
\pgfpathlineto{\pgfqpoint{2.846584in}{0.521603in}}%
\pgfpathlineto{\pgfqpoint{2.850752in}{0.521603in}}%
\pgfpathlineto{\pgfqpoint{2.854919in}{0.521603in}}%
\pgfpathlineto{\pgfqpoint{2.859087in}{0.521603in}}%
\pgfpathlineto{\pgfqpoint{2.863254in}{0.521603in}}%
\pgfpathlineto{\pgfqpoint{2.867422in}{0.521603in}}%
\pgfpathlineto{\pgfqpoint{2.871589in}{0.521603in}}%
\pgfpathlineto{\pgfqpoint{2.875757in}{0.521603in}}%
\pgfpathlineto{\pgfqpoint{2.879924in}{0.521603in}}%
\pgfpathlineto{\pgfqpoint{2.884091in}{0.521603in}}%
\pgfpathlineto{\pgfqpoint{2.888259in}{0.521603in}}%
\pgfpathlineto{\pgfqpoint{2.892426in}{0.521603in}}%
\pgfpathlineto{\pgfqpoint{2.896594in}{0.521603in}}%
\pgfpathlineto{\pgfqpoint{2.900761in}{0.521603in}}%
\pgfpathlineto{\pgfqpoint{2.904929in}{0.521603in}}%
\pgfpathlineto{\pgfqpoint{2.909096in}{0.521603in}}%
\pgfpathlineto{\pgfqpoint{2.913263in}{0.521603in}}%
\pgfpathlineto{\pgfqpoint{2.917431in}{0.521603in}}%
\pgfpathlineto{\pgfqpoint{2.921598in}{0.521603in}}%
\pgfpathlineto{\pgfqpoint{2.925766in}{0.521603in}}%
\pgfpathlineto{\pgfqpoint{2.929933in}{0.521603in}}%
\pgfpathlineto{\pgfqpoint{2.934101in}{0.521603in}}%
\pgfpathlineto{\pgfqpoint{2.938268in}{0.521603in}}%
\pgfpathlineto{\pgfqpoint{2.942436in}{0.521603in}}%
\pgfpathlineto{\pgfqpoint{2.946603in}{0.521603in}}%
\pgfpathlineto{\pgfqpoint{2.950770in}{0.521603in}}%
\pgfpathlineto{\pgfqpoint{2.954938in}{0.521603in}}%
\pgfpathlineto{\pgfqpoint{2.959105in}{0.521603in}}%
\pgfpathlineto{\pgfqpoint{2.963273in}{0.521603in}}%
\pgfpathlineto{\pgfqpoint{2.967440in}{0.521603in}}%
\pgfpathlineto{\pgfqpoint{2.971608in}{0.521603in}}%
\pgfpathlineto{\pgfqpoint{2.975775in}{0.521603in}}%
\pgfpathlineto{\pgfqpoint{2.979942in}{0.521603in}}%
\pgfpathlineto{\pgfqpoint{2.984110in}{0.521603in}}%
\pgfpathlineto{\pgfqpoint{2.988277in}{0.521603in}}%
\pgfpathlineto{\pgfqpoint{2.992445in}{0.521603in}}%
\pgfpathlineto{\pgfqpoint{2.996612in}{0.521603in}}%
\pgfpathlineto{\pgfqpoint{3.000780in}{0.521603in}}%
\pgfpathlineto{\pgfqpoint{3.004947in}{0.521603in}}%
\pgfpathlineto{\pgfqpoint{3.009115in}{0.521603in}}%
\pgfpathlineto{\pgfqpoint{3.013282in}{0.521603in}}%
\pgfpathlineto{\pgfqpoint{3.017449in}{0.521603in}}%
\pgfpathlineto{\pgfqpoint{3.021617in}{0.521603in}}%
\pgfpathlineto{\pgfqpoint{3.025784in}{0.521603in}}%
\pgfpathlineto{\pgfqpoint{3.029952in}{0.521603in}}%
\pgfpathlineto{\pgfqpoint{3.034119in}{0.521603in}}%
\pgfpathlineto{\pgfqpoint{3.038287in}{0.521603in}}%
\pgfpathlineto{\pgfqpoint{3.042454in}{0.521603in}}%
\pgfpathlineto{\pgfqpoint{3.046621in}{0.521603in}}%
\pgfpathlineto{\pgfqpoint{3.050789in}{0.521603in}}%
\pgfpathlineto{\pgfqpoint{3.054956in}{0.521603in}}%
\pgfpathlineto{\pgfqpoint{3.059124in}{0.521603in}}%
\pgfpathlineto{\pgfqpoint{3.063291in}{0.521603in}}%
\pgfpathlineto{\pgfqpoint{3.067459in}{0.521603in}}%
\pgfpathlineto{\pgfqpoint{3.071626in}{0.521603in}}%
\pgfpathlineto{\pgfqpoint{3.075793in}{0.521603in}}%
\pgfpathlineto{\pgfqpoint{3.079961in}{0.521603in}}%
\pgfpathlineto{\pgfqpoint{3.084128in}{0.521603in}}%
\pgfpathlineto{\pgfqpoint{3.088296in}{0.521603in}}%
\pgfpathlineto{\pgfqpoint{3.092463in}{0.521603in}}%
\pgfpathlineto{\pgfqpoint{3.096631in}{0.521603in}}%
\pgfpathlineto{\pgfqpoint{3.100798in}{0.521603in}}%
\pgfpathlineto{\pgfqpoint{3.104966in}{0.521603in}}%
\pgfpathlineto{\pgfqpoint{3.109133in}{0.521603in}}%
\pgfpathlineto{\pgfqpoint{3.113300in}{0.521603in}}%
\pgfpathlineto{\pgfqpoint{3.117468in}{0.521603in}}%
\pgfpathlineto{\pgfqpoint{3.121635in}{0.521603in}}%
\pgfpathlineto{\pgfqpoint{3.125803in}{0.521603in}}%
\pgfpathlineto{\pgfqpoint{3.129970in}{0.521603in}}%
\pgfpathlineto{\pgfqpoint{3.134138in}{0.521603in}}%
\pgfpathlineto{\pgfqpoint{3.138305in}{0.521603in}}%
\pgfpathlineto{\pgfqpoint{3.142472in}{0.521603in}}%
\pgfpathlineto{\pgfqpoint{3.146640in}{0.521603in}}%
\pgfpathlineto{\pgfqpoint{3.150807in}{0.521603in}}%
\pgfpathlineto{\pgfqpoint{3.154975in}{0.521603in}}%
\pgfpathlineto{\pgfqpoint{3.159142in}{0.521603in}}%
\pgfpathlineto{\pgfqpoint{3.163310in}{0.521603in}}%
\pgfpathlineto{\pgfqpoint{3.167477in}{0.521603in}}%
\pgfpathlineto{\pgfqpoint{3.171645in}{0.521603in}}%
\pgfpathlineto{\pgfqpoint{3.175812in}{0.521603in}}%
\pgfpathlineto{\pgfqpoint{3.179979in}{0.521603in}}%
\pgfpathlineto{\pgfqpoint{3.184147in}{0.521603in}}%
\pgfpathlineto{\pgfqpoint{3.188314in}{0.521603in}}%
\pgfpathlineto{\pgfqpoint{3.192482in}{0.521603in}}%
\pgfpathlineto{\pgfqpoint{3.196649in}{0.521603in}}%
\pgfpathlineto{\pgfqpoint{3.200817in}{0.521603in}}%
\pgfpathlineto{\pgfqpoint{3.204984in}{0.521603in}}%
\pgfpathlineto{\pgfqpoint{3.209151in}{0.521603in}}%
\pgfpathlineto{\pgfqpoint{3.213319in}{0.521603in}}%
\pgfpathlineto{\pgfqpoint{3.217486in}{0.521603in}}%
\pgfpathlineto{\pgfqpoint{3.221654in}{0.521603in}}%
\pgfpathlineto{\pgfqpoint{3.225821in}{0.521603in}}%
\pgfpathlineto{\pgfqpoint{3.229989in}{0.521603in}}%
\pgfpathlineto{\pgfqpoint{3.234156in}{0.521603in}}%
\pgfpathlineto{\pgfqpoint{3.238323in}{0.521603in}}%
\pgfpathlineto{\pgfqpoint{3.242491in}{0.521603in}}%
\pgfpathlineto{\pgfqpoint{3.246658in}{0.521603in}}%
\pgfpathlineto{\pgfqpoint{3.250826in}{0.521603in}}%
\pgfpathlineto{\pgfqpoint{3.254993in}{0.521603in}}%
\pgfpathlineto{\pgfqpoint{3.259161in}{0.521603in}}%
\pgfpathlineto{\pgfqpoint{3.263328in}{0.521603in}}%
\pgfpathlineto{\pgfqpoint{3.267496in}{0.521603in}}%
\pgfpathlineto{\pgfqpoint{3.271663in}{0.521603in}}%
\pgfpathlineto{\pgfqpoint{3.275830in}{0.521603in}}%
\pgfpathlineto{\pgfqpoint{3.279998in}{0.521603in}}%
\pgfpathlineto{\pgfqpoint{3.284165in}{0.521603in}}%
\pgfpathlineto{\pgfqpoint{3.288333in}{0.521603in}}%
\pgfpathlineto{\pgfqpoint{3.292500in}{0.521603in}}%
\pgfpathlineto{\pgfqpoint{3.296668in}{0.521603in}}%
\pgfpathlineto{\pgfqpoint{3.300835in}{0.521603in}}%
\pgfpathlineto{\pgfqpoint{3.305002in}{0.521603in}}%
\pgfpathlineto{\pgfqpoint{3.309170in}{0.521603in}}%
\pgfpathlineto{\pgfqpoint{3.313337in}{0.521603in}}%
\pgfpathlineto{\pgfqpoint{3.317505in}{0.521603in}}%
\pgfpathlineto{\pgfqpoint{3.321672in}{0.521603in}}%
\pgfpathlineto{\pgfqpoint{3.325840in}{0.521603in}}%
\pgfpathlineto{\pgfqpoint{3.330007in}{0.521603in}}%
\pgfpathlineto{\pgfqpoint{3.334175in}{0.521603in}}%
\pgfpathlineto{\pgfqpoint{3.338342in}{0.521603in}}%
\pgfpathlineto{\pgfqpoint{3.342509in}{0.521603in}}%
\pgfpathlineto{\pgfqpoint{3.346677in}{0.521603in}}%
\pgfpathlineto{\pgfqpoint{3.350844in}{0.521603in}}%
\pgfpathlineto{\pgfqpoint{3.355012in}{0.521603in}}%
\pgfpathlineto{\pgfqpoint{3.359179in}{0.521603in}}%
\pgfpathlineto{\pgfqpoint{3.363347in}{0.521603in}}%
\pgfpathlineto{\pgfqpoint{3.367514in}{0.521603in}}%
\pgfpathlineto{\pgfqpoint{3.371681in}{0.521603in}}%
\pgfpathlineto{\pgfqpoint{3.375849in}{0.521603in}}%
\pgfpathlineto{\pgfqpoint{3.380016in}{0.521603in}}%
\pgfpathlineto{\pgfqpoint{3.384184in}{0.521603in}}%
\pgfpathlineto{\pgfqpoint{3.388351in}{0.521603in}}%
\pgfpathlineto{\pgfqpoint{3.392519in}{0.521603in}}%
\pgfpathlineto{\pgfqpoint{3.396686in}{0.521603in}}%
\pgfpathlineto{\pgfqpoint{3.400853in}{0.521603in}}%
\pgfpathlineto{\pgfqpoint{3.405021in}{0.521603in}}%
\pgfpathlineto{\pgfqpoint{3.409188in}{0.521603in}}%
\pgfpathlineto{\pgfqpoint{3.413356in}{0.521603in}}%
\pgfpathlineto{\pgfqpoint{3.417523in}{0.521603in}}%
\pgfpathlineto{\pgfqpoint{3.421691in}{0.521603in}}%
\pgfpathlineto{\pgfqpoint{3.425858in}{0.521603in}}%
\pgfpathlineto{\pgfqpoint{3.430026in}{0.521603in}}%
\pgfpathlineto{\pgfqpoint{3.434193in}{0.521603in}}%
\pgfpathlineto{\pgfqpoint{3.438360in}{0.521603in}}%
\pgfpathlineto{\pgfqpoint{3.442528in}{0.521603in}}%
\pgfpathlineto{\pgfqpoint{3.446695in}{0.521603in}}%
\pgfpathlineto{\pgfqpoint{3.450863in}{0.521603in}}%
\pgfpathlineto{\pgfqpoint{3.455030in}{0.521603in}}%
\pgfpathlineto{\pgfqpoint{3.459198in}{0.521603in}}%
\pgfpathlineto{\pgfqpoint{3.463365in}{0.521603in}}%
\pgfpathlineto{\pgfqpoint{3.467532in}{0.521603in}}%
\pgfpathlineto{\pgfqpoint{3.471700in}{0.521603in}}%
\pgfpathlineto{\pgfqpoint{3.475867in}{0.521603in}}%
\pgfpathlineto{\pgfqpoint{3.480035in}{0.521603in}}%
\pgfpathlineto{\pgfqpoint{3.484202in}{0.521603in}}%
\pgfpathlineto{\pgfqpoint{3.484202in}{0.522589in}}%
\pgfpathlineto{\pgfqpoint{3.484202in}{0.522589in}}%
\pgfpathlineto{\pgfqpoint{3.480035in}{0.522961in}}%
\pgfpathlineto{\pgfqpoint{3.475867in}{0.523452in}}%
\pgfpathlineto{\pgfqpoint{3.471700in}{0.524092in}}%
\pgfpathlineto{\pgfqpoint{3.467532in}{0.524915in}}%
\pgfpathlineto{\pgfqpoint{3.463365in}{0.525962in}}%
\pgfpathlineto{\pgfqpoint{3.459198in}{0.527276in}}%
\pgfpathlineto{\pgfqpoint{3.455030in}{0.528904in}}%
\pgfpathlineto{\pgfqpoint{3.450863in}{0.530894in}}%
\pgfpathlineto{\pgfqpoint{3.446695in}{0.533295in}}%
\pgfpathlineto{\pgfqpoint{3.442528in}{0.536155in}}%
\pgfpathlineto{\pgfqpoint{3.438360in}{0.539515in}}%
\pgfpathlineto{\pgfqpoint{3.434193in}{0.543411in}}%
\pgfpathlineto{\pgfqpoint{3.430026in}{0.547865in}}%
\pgfpathlineto{\pgfqpoint{3.425858in}{0.552888in}}%
\pgfpathlineto{\pgfqpoint{3.421691in}{0.558476in}}%
\pgfpathlineto{\pgfqpoint{3.417523in}{0.564606in}}%
\pgfpathlineto{\pgfqpoint{3.413356in}{0.571235in}}%
\pgfpathlineto{\pgfqpoint{3.409188in}{0.578306in}}%
\pgfpathlineto{\pgfqpoint{3.405021in}{0.585743in}}%
\pgfpathlineto{\pgfqpoint{3.400853in}{0.593462in}}%
\pgfpathlineto{\pgfqpoint{3.396686in}{0.601367in}}%
\pgfpathlineto{\pgfqpoint{3.392519in}{0.609366in}}%
\pgfpathlineto{\pgfqpoint{3.388351in}{0.617373in}}%
\pgfpathlineto{\pgfqpoint{3.384184in}{0.625318in}}%
\pgfpathlineto{\pgfqpoint{3.380016in}{0.633157in}}%
\pgfpathlineto{\pgfqpoint{3.375849in}{0.640876in}}%
\pgfpathlineto{\pgfqpoint{3.371681in}{0.648503in}}%
\pgfpathlineto{\pgfqpoint{3.367514in}{0.656110in}}%
\pgfpathlineto{\pgfqpoint{3.363347in}{0.663815in}}%
\pgfpathlineto{\pgfqpoint{3.359179in}{0.671783in}}%
\pgfpathlineto{\pgfqpoint{3.355012in}{0.680222in}}%
\pgfpathlineto{\pgfqpoint{3.350844in}{0.689374in}}%
\pgfpathlineto{\pgfqpoint{3.346677in}{0.699510in}}%
\pgfpathlineto{\pgfqpoint{3.342509in}{0.710910in}}%
\pgfpathlineto{\pgfqpoint{3.338342in}{0.723857in}}%
\pgfpathlineto{\pgfqpoint{3.334175in}{0.738613in}}%
\pgfpathlineto{\pgfqpoint{3.330007in}{0.755407in}}%
\pgfpathlineto{\pgfqpoint{3.325840in}{0.774420in}}%
\pgfpathlineto{\pgfqpoint{3.321672in}{0.795768in}}%
\pgfpathlineto{\pgfqpoint{3.317505in}{0.819494in}}%
\pgfpathlineto{\pgfqpoint{3.313337in}{0.845560in}}%
\pgfpathlineto{\pgfqpoint{3.309170in}{0.873843in}}%
\pgfpathlineto{\pgfqpoint{3.305002in}{0.904136in}}%
\pgfpathlineto{\pgfqpoint{3.300835in}{0.936156in}}%
\pgfpathlineto{\pgfqpoint{3.296668in}{0.969553in}}%
\pgfpathlineto{\pgfqpoint{3.292500in}{1.003923in}}%
\pgfpathlineto{\pgfqpoint{3.288333in}{1.038829in}}%
\pgfpathlineto{\pgfqpoint{3.284165in}{1.073818in}}%
\pgfpathlineto{\pgfqpoint{3.279998in}{1.108446in}}%
\pgfpathlineto{\pgfqpoint{3.275830in}{1.142294in}}%
\pgfpathlineto{\pgfqpoint{3.271663in}{1.174988in}}%
\pgfpathlineto{\pgfqpoint{3.267496in}{1.206222in}}%
\pgfpathlineto{\pgfqpoint{3.263328in}{1.235762in}}%
\pgfpathlineto{\pgfqpoint{3.259161in}{1.263461in}}%
\pgfpathlineto{\pgfqpoint{3.254993in}{1.289261in}}%
\pgfpathlineto{\pgfqpoint{3.250826in}{1.313193in}}%
\pgfpathlineto{\pgfqpoint{3.246658in}{1.335369in}}%
\pgfpathlineto{\pgfqpoint{3.242491in}{1.355975in}}%
\pgfpathlineto{\pgfqpoint{3.238323in}{1.375257in}}%
\pgfpathlineto{\pgfqpoint{3.234156in}{1.393509in}}%
\pgfpathlineto{\pgfqpoint{3.229989in}{1.411054in}}%
\pgfpathlineto{\pgfqpoint{3.225821in}{1.428229in}}%
\pgfpathlineto{\pgfqpoint{3.221654in}{1.445375in}}%
\pgfpathlineto{\pgfqpoint{3.217486in}{1.462822in}}%
\pgfpathlineto{\pgfqpoint{3.213319in}{1.480882in}}%
\pgfpathlineto{\pgfqpoint{3.209151in}{1.499844in}}%
\pgfpathlineto{\pgfqpoint{3.204984in}{1.519972in}}%
\pgfpathlineto{\pgfqpoint{3.200817in}{1.541506in}}%
\pgfpathlineto{\pgfqpoint{3.196649in}{1.564665in}}%
\pgfpathlineto{\pgfqpoint{3.192482in}{1.589654in}}%
\pgfpathlineto{\pgfqpoint{3.188314in}{1.616663in}}%
\pgfpathlineto{\pgfqpoint{3.184147in}{1.645876in}}%
\pgfpathlineto{\pgfqpoint{3.179979in}{1.677471in}}%
\pgfpathlineto{\pgfqpoint{3.175812in}{1.711618in}}%
\pgfpathlineto{\pgfqpoint{3.171645in}{1.748476in}}%
\pgfpathlineto{\pgfqpoint{3.167477in}{1.788181in}}%
\pgfpathlineto{\pgfqpoint{3.163310in}{1.830838in}}%
\pgfpathlineto{\pgfqpoint{3.159142in}{1.876505in}}%
\pgfpathlineto{\pgfqpoint{3.154975in}{1.925173in}}%
\pgfpathlineto{\pgfqpoint{3.150807in}{1.976757in}}%
\pgfpathlineto{\pgfqpoint{3.146640in}{2.031071in}}%
\pgfpathlineto{\pgfqpoint{3.142472in}{2.087823in}}%
\pgfpathlineto{\pgfqpoint{3.138305in}{2.146605in}}%
\pgfpathlineto{\pgfqpoint{3.134138in}{2.206893in}}%
\pgfpathlineto{\pgfqpoint{3.129970in}{2.268054in}}%
\pgfpathlineto{\pgfqpoint{3.125803in}{2.329360in}}%
\pgfpathlineto{\pgfqpoint{3.121635in}{2.390012in}}%
\pgfpathlineto{\pgfqpoint{3.117468in}{2.449168in}}%
\pgfpathlineto{\pgfqpoint{3.113300in}{2.505982in}}%
\pgfpathlineto{\pgfqpoint{3.109133in}{2.559635in}}%
\pgfpathlineto{\pgfqpoint{3.104966in}{2.609383in}}%
\pgfpathlineto{\pgfqpoint{3.100798in}{2.654585in}}%
\pgfpathlineto{\pgfqpoint{3.096631in}{2.694743in}}%
\pgfpathlineto{\pgfqpoint{3.092463in}{2.729519in}}%
\pgfpathlineto{\pgfqpoint{3.088296in}{2.758756in}}%
\pgfpathlineto{\pgfqpoint{3.084128in}{2.782479in}}%
\pgfpathlineto{\pgfqpoint{3.079961in}{2.800886in}}%
\pgfpathlineto{\pgfqpoint{3.075793in}{2.814335in}}%
\pgfpathlineto{\pgfqpoint{3.071626in}{2.823314in}}%
\pgfpathlineto{\pgfqpoint{3.067459in}{2.828402in}}%
\pgfpathlineto{\pgfqpoint{3.063291in}{2.830237in}}%
\pgfpathlineto{\pgfqpoint{3.059124in}{2.829465in}}%
\pgfpathlineto{\pgfqpoint{3.054956in}{2.826703in}}%
\pgfpathlineto{\pgfqpoint{3.050789in}{2.822503in}}%
\pgfpathlineto{\pgfqpoint{3.046621in}{2.817317in}}%
\pgfpathlineto{\pgfqpoint{3.042454in}{2.811475in}}%
\pgfpathlineto{\pgfqpoint{3.038287in}{2.805171in}}%
\pgfpathlineto{\pgfqpoint{3.034119in}{2.798461in}}%
\pgfpathlineto{\pgfqpoint{3.029952in}{2.791264in}}%
\pgfpathlineto{\pgfqpoint{3.025784in}{2.783377in}}%
\pgfpathlineto{\pgfqpoint{3.021617in}{2.774496in}}%
\pgfpathlineto{\pgfqpoint{3.017449in}{2.764236in}}%
\pgfpathlineto{\pgfqpoint{3.013282in}{2.752160in}}%
\pgfpathlineto{\pgfqpoint{3.009115in}{2.737806in}}%
\pgfpathlineto{\pgfqpoint{3.004947in}{2.720711in}}%
\pgfpathlineto{\pgfqpoint{3.000780in}{2.700435in}}%
\pgfpathlineto{\pgfqpoint{2.996612in}{2.676579in}}%
\pgfpathlineto{\pgfqpoint{2.992445in}{2.648804in}}%
\pgfpathlineto{\pgfqpoint{2.988277in}{2.616838in}}%
\pgfpathlineto{\pgfqpoint{2.984110in}{2.580486in}}%
\pgfpathlineto{\pgfqpoint{2.979942in}{2.539634in}}%
\pgfpathlineto{\pgfqpoint{2.975775in}{2.494252in}}%
\pgfpathlineto{\pgfqpoint{2.971608in}{2.444392in}}%
\pgfpathlineto{\pgfqpoint{2.967440in}{2.390186in}}%
\pgfpathlineto{\pgfqpoint{2.963273in}{2.331846in}}%
\pgfpathlineto{\pgfqpoint{2.959105in}{2.269656in}}%
\pgfpathlineto{\pgfqpoint{2.954938in}{2.203972in}}%
\pgfpathlineto{\pgfqpoint{2.950770in}{2.135208in}}%
\pgfpathlineto{\pgfqpoint{2.946603in}{2.063838in}}%
\pgfpathlineto{\pgfqpoint{2.942436in}{1.990378in}}%
\pgfpathlineto{\pgfqpoint{2.938268in}{1.915383in}}%
\pgfpathlineto{\pgfqpoint{2.934101in}{1.839430in}}%
\pgfpathlineto{\pgfqpoint{2.929933in}{1.763110in}}%
\pgfpathlineto{\pgfqpoint{2.925766in}{1.687013in}}%
\pgfpathlineto{\pgfqpoint{2.921598in}{1.611711in}}%
\pgfpathlineto{\pgfqpoint{2.917431in}{1.537754in}}%
\pgfpathlineto{\pgfqpoint{2.913263in}{1.465649in}}%
\pgfpathlineto{\pgfqpoint{2.909096in}{1.395857in}}%
\pgfpathlineto{\pgfqpoint{2.904929in}{1.328780in}}%
\pgfpathlineto{\pgfqpoint{2.900761in}{1.264760in}}%
\pgfpathlineto{\pgfqpoint{2.896594in}{1.204071in}}%
\pgfpathlineto{\pgfqpoint{2.892426in}{1.146922in}}%
\pgfpathlineto{\pgfqpoint{2.888259in}{1.093460in}}%
\pgfpathlineto{\pgfqpoint{2.884091in}{1.043768in}}%
\pgfpathlineto{\pgfqpoint{2.879924in}{0.997877in}}%
\pgfpathlineto{\pgfqpoint{2.875757in}{0.955769in}}%
\pgfpathlineto{\pgfqpoint{2.871589in}{0.917383in}}%
\pgfpathlineto{\pgfqpoint{2.867422in}{0.882626in}}%
\pgfpathlineto{\pgfqpoint{2.863254in}{0.851375in}}%
\pgfpathlineto{\pgfqpoint{2.859087in}{0.823485in}}%
\pgfpathlineto{\pgfqpoint{2.854919in}{0.798795in}}%
\pgfpathlineto{\pgfqpoint{2.850752in}{0.777129in}}%
\pgfpathlineto{\pgfqpoint{2.846584in}{0.758297in}}%
\pgfpathlineto{\pgfqpoint{2.842417in}{0.742102in}}%
\pgfpathlineto{\pgfqpoint{2.838250in}{0.728334in}}%
\pgfpathlineto{\pgfqpoint{2.834082in}{0.716773in}}%
\pgfpathlineto{\pgfqpoint{2.829915in}{0.707193in}}%
\pgfpathlineto{\pgfqpoint{2.825747in}{0.699353in}}%
\pgfpathlineto{\pgfqpoint{2.821580in}{0.693008in}}%
\pgfpathlineto{\pgfqpoint{2.817412in}{0.687905in}}%
\pgfpathlineto{\pgfqpoint{2.813245in}{0.683788in}}%
\pgfpathlineto{\pgfqpoint{2.809078in}{0.680401in}}%
\pgfpathlineto{\pgfqpoint{2.804910in}{0.677496in}}%
\pgfpathlineto{\pgfqpoint{2.800743in}{0.674834in}}%
\pgfpathlineto{\pgfqpoint{2.796575in}{0.672195in}}%
\pgfpathlineto{\pgfqpoint{2.792408in}{0.669381in}}%
\pgfpathlineto{\pgfqpoint{2.788240in}{0.666225in}}%
\pgfpathlineto{\pgfqpoint{2.784073in}{0.662592in}}%
\pgfpathlineto{\pgfqpoint{2.779906in}{0.658384in}}%
\pgfpathlineto{\pgfqpoint{2.775738in}{0.653544in}}%
\pgfpathlineto{\pgfqpoint{2.771571in}{0.648051in}}%
\pgfpathlineto{\pgfqpoint{2.767403in}{0.641925in}}%
\pgfpathlineto{\pgfqpoint{2.763236in}{0.635219in}}%
\pgfpathlineto{\pgfqpoint{2.759068in}{0.628014in}}%
\pgfpathlineto{\pgfqpoint{2.754901in}{0.620415in}}%
\pgfpathlineto{\pgfqpoint{2.750733in}{0.612546in}}%
\pgfpathlineto{\pgfqpoint{2.746566in}{0.604537in}}%
\pgfpathlineto{\pgfqpoint{2.742399in}{0.596523in}}%
\pgfpathlineto{\pgfqpoint{2.738231in}{0.588633in}}%
\pgfpathlineto{\pgfqpoint{2.734064in}{0.580987in}}%
\pgfpathlineto{\pgfqpoint{2.729896in}{0.573691in}}%
\pgfpathlineto{\pgfqpoint{2.725729in}{0.566831in}}%
\pgfpathlineto{\pgfqpoint{2.721561in}{0.560475in}}%
\pgfpathlineto{\pgfqpoint{2.717394in}{0.554668in}}%
\pgfpathlineto{\pgfqpoint{2.713227in}{0.549437in}}%
\pgfpathlineto{\pgfqpoint{2.709059in}{0.544788in}}%
\pgfpathlineto{\pgfqpoint{2.704892in}{0.540713in}}%
\pgfpathlineto{\pgfqpoint{2.700724in}{0.537187in}}%
\pgfpathlineto{\pgfqpoint{2.696557in}{0.534177in}}%
\pgfpathlineto{\pgfqpoint{2.692389in}{0.531639in}}%
\pgfpathlineto{\pgfqpoint{2.688222in}{0.529528in}}%
\pgfpathlineto{\pgfqpoint{2.684054in}{0.527793in}}%
\pgfpathlineto{\pgfqpoint{2.679887in}{0.526386in}}%
\pgfpathlineto{\pgfqpoint{2.675720in}{0.525259in}}%
\pgfpathlineto{\pgfqpoint{2.671552in}{0.524367in}}%
\pgfpathlineto{\pgfqpoint{2.667385in}{0.523670in}}%
\pgfpathlineto{\pgfqpoint{2.663217in}{0.523131in}}%
\pgfpathlineto{\pgfqpoint{2.659050in}{0.522721in}}%
\pgfpathlineto{\pgfqpoint{2.654882in}{0.522412in}}%
\pgfpathlineto{\pgfqpoint{2.654882in}{0.522412in}}%
\pgfpathclose%
\pgfusepath{stroke,fill}%
}%
\begin{pgfscope}%
\pgfsys@transformshift{0.000000in}{0.000000in}%
\pgfsys@useobject{currentmarker}{}%
\end{pgfscope}%
\end{pgfscope}%
\begin{pgfscope}%
\pgfpathrectangle{\pgfqpoint{0.784801in}{0.521603in}}{\pgfqpoint{4.349472in}{2.424065in}}%
\pgfusepath{clip}%
\pgfsetbuttcap%
\pgfsetroundjoin%
\definecolor{currentfill}{rgb}{0.647059,0.839216,0.654902}%
\pgfsetfillcolor{currentfill}%
\pgfsetfillopacity{0.400000}%
\pgfsetlinewidth{1.505625pt}%
\definecolor{currentstroke}{rgb}{0.647059,0.839216,0.654902}%
\pgfsetstrokecolor{currentstroke}%
\pgfsetdash{}{0pt}%
\pgfsys@defobject{currentmarker}{\pgfqpoint{1.542129in}{0.521603in}}{\pgfqpoint{2.416058in}{2.790110in}}{%
\pgfpathmoveto{\pgfqpoint{1.542129in}{0.521998in}}%
\pgfpathlineto{\pgfqpoint{1.542129in}{0.521603in}}%
\pgfpathlineto{\pgfqpoint{1.546520in}{0.521603in}}%
\pgfpathlineto{\pgfqpoint{1.550912in}{0.521603in}}%
\pgfpathlineto{\pgfqpoint{1.555303in}{0.521603in}}%
\pgfpathlineto{\pgfqpoint{1.559695in}{0.521603in}}%
\pgfpathlineto{\pgfqpoint{1.564087in}{0.521603in}}%
\pgfpathlineto{\pgfqpoint{1.568478in}{0.521603in}}%
\pgfpathlineto{\pgfqpoint{1.572870in}{0.521603in}}%
\pgfpathlineto{\pgfqpoint{1.577261in}{0.521603in}}%
\pgfpathlineto{\pgfqpoint{1.581653in}{0.521603in}}%
\pgfpathlineto{\pgfqpoint{1.586045in}{0.521603in}}%
\pgfpathlineto{\pgfqpoint{1.590436in}{0.521603in}}%
\pgfpathlineto{\pgfqpoint{1.594828in}{0.521603in}}%
\pgfpathlineto{\pgfqpoint{1.599219in}{0.521603in}}%
\pgfpathlineto{\pgfqpoint{1.603611in}{0.521603in}}%
\pgfpathlineto{\pgfqpoint{1.608003in}{0.521603in}}%
\pgfpathlineto{\pgfqpoint{1.612394in}{0.521603in}}%
\pgfpathlineto{\pgfqpoint{1.616786in}{0.521603in}}%
\pgfpathlineto{\pgfqpoint{1.621177in}{0.521603in}}%
\pgfpathlineto{\pgfqpoint{1.625569in}{0.521603in}}%
\pgfpathlineto{\pgfqpoint{1.629961in}{0.521603in}}%
\pgfpathlineto{\pgfqpoint{1.634352in}{0.521603in}}%
\pgfpathlineto{\pgfqpoint{1.638744in}{0.521603in}}%
\pgfpathlineto{\pgfqpoint{1.643135in}{0.521603in}}%
\pgfpathlineto{\pgfqpoint{1.647527in}{0.521603in}}%
\pgfpathlineto{\pgfqpoint{1.651919in}{0.521603in}}%
\pgfpathlineto{\pgfqpoint{1.656310in}{0.521603in}}%
\pgfpathlineto{\pgfqpoint{1.660702in}{0.521603in}}%
\pgfpathlineto{\pgfqpoint{1.665094in}{0.521603in}}%
\pgfpathlineto{\pgfqpoint{1.669485in}{0.521603in}}%
\pgfpathlineto{\pgfqpoint{1.673877in}{0.521603in}}%
\pgfpathlineto{\pgfqpoint{1.678268in}{0.521603in}}%
\pgfpathlineto{\pgfqpoint{1.682660in}{0.521603in}}%
\pgfpathlineto{\pgfqpoint{1.687052in}{0.521603in}}%
\pgfpathlineto{\pgfqpoint{1.691443in}{0.521603in}}%
\pgfpathlineto{\pgfqpoint{1.695835in}{0.521603in}}%
\pgfpathlineto{\pgfqpoint{1.700226in}{0.521603in}}%
\pgfpathlineto{\pgfqpoint{1.704618in}{0.521603in}}%
\pgfpathlineto{\pgfqpoint{1.709010in}{0.521603in}}%
\pgfpathlineto{\pgfqpoint{1.713401in}{0.521603in}}%
\pgfpathlineto{\pgfqpoint{1.717793in}{0.521603in}}%
\pgfpathlineto{\pgfqpoint{1.722184in}{0.521603in}}%
\pgfpathlineto{\pgfqpoint{1.726576in}{0.521603in}}%
\pgfpathlineto{\pgfqpoint{1.730968in}{0.521603in}}%
\pgfpathlineto{\pgfqpoint{1.735359in}{0.521603in}}%
\pgfpathlineto{\pgfqpoint{1.739751in}{0.521603in}}%
\pgfpathlineto{\pgfqpoint{1.744142in}{0.521603in}}%
\pgfpathlineto{\pgfqpoint{1.748534in}{0.521603in}}%
\pgfpathlineto{\pgfqpoint{1.752926in}{0.521603in}}%
\pgfpathlineto{\pgfqpoint{1.757317in}{0.521603in}}%
\pgfpathlineto{\pgfqpoint{1.761709in}{0.521603in}}%
\pgfpathlineto{\pgfqpoint{1.766100in}{0.521603in}}%
\pgfpathlineto{\pgfqpoint{1.770492in}{0.521603in}}%
\pgfpathlineto{\pgfqpoint{1.774884in}{0.521603in}}%
\pgfpathlineto{\pgfqpoint{1.779275in}{0.521603in}}%
\pgfpathlineto{\pgfqpoint{1.783667in}{0.521603in}}%
\pgfpathlineto{\pgfqpoint{1.788058in}{0.521603in}}%
\pgfpathlineto{\pgfqpoint{1.792450in}{0.521603in}}%
\pgfpathlineto{\pgfqpoint{1.796842in}{0.521603in}}%
\pgfpathlineto{\pgfqpoint{1.801233in}{0.521603in}}%
\pgfpathlineto{\pgfqpoint{1.805625in}{0.521603in}}%
\pgfpathlineto{\pgfqpoint{1.810017in}{0.521603in}}%
\pgfpathlineto{\pgfqpoint{1.814408in}{0.521603in}}%
\pgfpathlineto{\pgfqpoint{1.818800in}{0.521603in}}%
\pgfpathlineto{\pgfqpoint{1.823191in}{0.521603in}}%
\pgfpathlineto{\pgfqpoint{1.827583in}{0.521603in}}%
\pgfpathlineto{\pgfqpoint{1.831975in}{0.521603in}}%
\pgfpathlineto{\pgfqpoint{1.836366in}{0.521603in}}%
\pgfpathlineto{\pgfqpoint{1.840758in}{0.521603in}}%
\pgfpathlineto{\pgfqpoint{1.845149in}{0.521603in}}%
\pgfpathlineto{\pgfqpoint{1.849541in}{0.521603in}}%
\pgfpathlineto{\pgfqpoint{1.853933in}{0.521603in}}%
\pgfpathlineto{\pgfqpoint{1.858324in}{0.521603in}}%
\pgfpathlineto{\pgfqpoint{1.862716in}{0.521603in}}%
\pgfpathlineto{\pgfqpoint{1.867107in}{0.521603in}}%
\pgfpathlineto{\pgfqpoint{1.871499in}{0.521603in}}%
\pgfpathlineto{\pgfqpoint{1.875891in}{0.521603in}}%
\pgfpathlineto{\pgfqpoint{1.880282in}{0.521603in}}%
\pgfpathlineto{\pgfqpoint{1.884674in}{0.521603in}}%
\pgfpathlineto{\pgfqpoint{1.889065in}{0.521603in}}%
\pgfpathlineto{\pgfqpoint{1.893457in}{0.521603in}}%
\pgfpathlineto{\pgfqpoint{1.897849in}{0.521603in}}%
\pgfpathlineto{\pgfqpoint{1.902240in}{0.521603in}}%
\pgfpathlineto{\pgfqpoint{1.906632in}{0.521603in}}%
\pgfpathlineto{\pgfqpoint{1.911023in}{0.521603in}}%
\pgfpathlineto{\pgfqpoint{1.915415in}{0.521603in}}%
\pgfpathlineto{\pgfqpoint{1.919807in}{0.521603in}}%
\pgfpathlineto{\pgfqpoint{1.924198in}{0.521603in}}%
\pgfpathlineto{\pgfqpoint{1.928590in}{0.521603in}}%
\pgfpathlineto{\pgfqpoint{1.932981in}{0.521603in}}%
\pgfpathlineto{\pgfqpoint{1.937373in}{0.521603in}}%
\pgfpathlineto{\pgfqpoint{1.941765in}{0.521603in}}%
\pgfpathlineto{\pgfqpoint{1.946156in}{0.521603in}}%
\pgfpathlineto{\pgfqpoint{1.950548in}{0.521603in}}%
\pgfpathlineto{\pgfqpoint{1.954940in}{0.521603in}}%
\pgfpathlineto{\pgfqpoint{1.959331in}{0.521603in}}%
\pgfpathlineto{\pgfqpoint{1.963723in}{0.521603in}}%
\pgfpathlineto{\pgfqpoint{1.968114in}{0.521603in}}%
\pgfpathlineto{\pgfqpoint{1.972506in}{0.521603in}}%
\pgfpathlineto{\pgfqpoint{1.976898in}{0.521603in}}%
\pgfpathlineto{\pgfqpoint{1.981289in}{0.521603in}}%
\pgfpathlineto{\pgfqpoint{1.985681in}{0.521603in}}%
\pgfpathlineto{\pgfqpoint{1.990072in}{0.521603in}}%
\pgfpathlineto{\pgfqpoint{1.994464in}{0.521603in}}%
\pgfpathlineto{\pgfqpoint{1.998856in}{0.521603in}}%
\pgfpathlineto{\pgfqpoint{2.003247in}{0.521603in}}%
\pgfpathlineto{\pgfqpoint{2.007639in}{0.521603in}}%
\pgfpathlineto{\pgfqpoint{2.012030in}{0.521603in}}%
\pgfpathlineto{\pgfqpoint{2.016422in}{0.521603in}}%
\pgfpathlineto{\pgfqpoint{2.020814in}{0.521603in}}%
\pgfpathlineto{\pgfqpoint{2.025205in}{0.521603in}}%
\pgfpathlineto{\pgfqpoint{2.029597in}{0.521603in}}%
\pgfpathlineto{\pgfqpoint{2.033988in}{0.521603in}}%
\pgfpathlineto{\pgfqpoint{2.038380in}{0.521603in}}%
\pgfpathlineto{\pgfqpoint{2.042772in}{0.521603in}}%
\pgfpathlineto{\pgfqpoint{2.047163in}{0.521603in}}%
\pgfpathlineto{\pgfqpoint{2.051555in}{0.521603in}}%
\pgfpathlineto{\pgfqpoint{2.055946in}{0.521603in}}%
\pgfpathlineto{\pgfqpoint{2.060338in}{0.521603in}}%
\pgfpathlineto{\pgfqpoint{2.064730in}{0.521603in}}%
\pgfpathlineto{\pgfqpoint{2.069121in}{0.521603in}}%
\pgfpathlineto{\pgfqpoint{2.073513in}{0.521603in}}%
\pgfpathlineto{\pgfqpoint{2.077904in}{0.521603in}}%
\pgfpathlineto{\pgfqpoint{2.082296in}{0.521603in}}%
\pgfpathlineto{\pgfqpoint{2.086688in}{0.521603in}}%
\pgfpathlineto{\pgfqpoint{2.091079in}{0.521603in}}%
\pgfpathlineto{\pgfqpoint{2.095471in}{0.521603in}}%
\pgfpathlineto{\pgfqpoint{2.099862in}{0.521603in}}%
\pgfpathlineto{\pgfqpoint{2.104254in}{0.521603in}}%
\pgfpathlineto{\pgfqpoint{2.108646in}{0.521603in}}%
\pgfpathlineto{\pgfqpoint{2.113037in}{0.521603in}}%
\pgfpathlineto{\pgfqpoint{2.117429in}{0.521603in}}%
\pgfpathlineto{\pgfqpoint{2.121821in}{0.521603in}}%
\pgfpathlineto{\pgfqpoint{2.126212in}{0.521603in}}%
\pgfpathlineto{\pgfqpoint{2.130604in}{0.521603in}}%
\pgfpathlineto{\pgfqpoint{2.134995in}{0.521603in}}%
\pgfpathlineto{\pgfqpoint{2.139387in}{0.521603in}}%
\pgfpathlineto{\pgfqpoint{2.143779in}{0.521603in}}%
\pgfpathlineto{\pgfqpoint{2.148170in}{0.521603in}}%
\pgfpathlineto{\pgfqpoint{2.152562in}{0.521603in}}%
\pgfpathlineto{\pgfqpoint{2.156953in}{0.521603in}}%
\pgfpathlineto{\pgfqpoint{2.161345in}{0.521603in}}%
\pgfpathlineto{\pgfqpoint{2.165737in}{0.521603in}}%
\pgfpathlineto{\pgfqpoint{2.170128in}{0.521603in}}%
\pgfpathlineto{\pgfqpoint{2.174520in}{0.521603in}}%
\pgfpathlineto{\pgfqpoint{2.178911in}{0.521603in}}%
\pgfpathlineto{\pgfqpoint{2.183303in}{0.521603in}}%
\pgfpathlineto{\pgfqpoint{2.187695in}{0.521603in}}%
\pgfpathlineto{\pgfqpoint{2.192086in}{0.521603in}}%
\pgfpathlineto{\pgfqpoint{2.196478in}{0.521603in}}%
\pgfpathlineto{\pgfqpoint{2.200869in}{0.521603in}}%
\pgfpathlineto{\pgfqpoint{2.205261in}{0.521603in}}%
\pgfpathlineto{\pgfqpoint{2.209653in}{0.521603in}}%
\pgfpathlineto{\pgfqpoint{2.214044in}{0.521603in}}%
\pgfpathlineto{\pgfqpoint{2.218436in}{0.521603in}}%
\pgfpathlineto{\pgfqpoint{2.222827in}{0.521603in}}%
\pgfpathlineto{\pgfqpoint{2.227219in}{0.521603in}}%
\pgfpathlineto{\pgfqpoint{2.231611in}{0.521603in}}%
\pgfpathlineto{\pgfqpoint{2.236002in}{0.521603in}}%
\pgfpathlineto{\pgfqpoint{2.240394in}{0.521603in}}%
\pgfpathlineto{\pgfqpoint{2.244785in}{0.521603in}}%
\pgfpathlineto{\pgfqpoint{2.249177in}{0.521603in}}%
\pgfpathlineto{\pgfqpoint{2.253569in}{0.521603in}}%
\pgfpathlineto{\pgfqpoint{2.257960in}{0.521603in}}%
\pgfpathlineto{\pgfqpoint{2.262352in}{0.521603in}}%
\pgfpathlineto{\pgfqpoint{2.266744in}{0.521603in}}%
\pgfpathlineto{\pgfqpoint{2.271135in}{0.521603in}}%
\pgfpathlineto{\pgfqpoint{2.275527in}{0.521603in}}%
\pgfpathlineto{\pgfqpoint{2.279918in}{0.521603in}}%
\pgfpathlineto{\pgfqpoint{2.284310in}{0.521603in}}%
\pgfpathlineto{\pgfqpoint{2.288702in}{0.521603in}}%
\pgfpathlineto{\pgfqpoint{2.293093in}{0.521603in}}%
\pgfpathlineto{\pgfqpoint{2.297485in}{0.521603in}}%
\pgfpathlineto{\pgfqpoint{2.301876in}{0.521603in}}%
\pgfpathlineto{\pgfqpoint{2.306268in}{0.521603in}}%
\pgfpathlineto{\pgfqpoint{2.310660in}{0.521603in}}%
\pgfpathlineto{\pgfqpoint{2.315051in}{0.521603in}}%
\pgfpathlineto{\pgfqpoint{2.319443in}{0.521603in}}%
\pgfpathlineto{\pgfqpoint{2.323834in}{0.521603in}}%
\pgfpathlineto{\pgfqpoint{2.328226in}{0.521603in}}%
\pgfpathlineto{\pgfqpoint{2.332618in}{0.521603in}}%
\pgfpathlineto{\pgfqpoint{2.337009in}{0.521603in}}%
\pgfpathlineto{\pgfqpoint{2.341401in}{0.521603in}}%
\pgfpathlineto{\pgfqpoint{2.345792in}{0.521603in}}%
\pgfpathlineto{\pgfqpoint{2.350184in}{0.521603in}}%
\pgfpathlineto{\pgfqpoint{2.354576in}{0.521603in}}%
\pgfpathlineto{\pgfqpoint{2.358967in}{0.521603in}}%
\pgfpathlineto{\pgfqpoint{2.363359in}{0.521603in}}%
\pgfpathlineto{\pgfqpoint{2.367750in}{0.521603in}}%
\pgfpathlineto{\pgfqpoint{2.372142in}{0.521603in}}%
\pgfpathlineto{\pgfqpoint{2.376534in}{0.521603in}}%
\pgfpathlineto{\pgfqpoint{2.380925in}{0.521603in}}%
\pgfpathlineto{\pgfqpoint{2.385317in}{0.521603in}}%
\pgfpathlineto{\pgfqpoint{2.389708in}{0.521603in}}%
\pgfpathlineto{\pgfqpoint{2.394100in}{0.521603in}}%
\pgfpathlineto{\pgfqpoint{2.398492in}{0.521603in}}%
\pgfpathlineto{\pgfqpoint{2.402883in}{0.521603in}}%
\pgfpathlineto{\pgfqpoint{2.407275in}{0.521603in}}%
\pgfpathlineto{\pgfqpoint{2.411667in}{0.521603in}}%
\pgfpathlineto{\pgfqpoint{2.416058in}{0.521603in}}%
\pgfpathlineto{\pgfqpoint{2.416058in}{0.522008in}}%
\pgfpathlineto{\pgfqpoint{2.416058in}{0.522008in}}%
\pgfpathlineto{\pgfqpoint{2.411667in}{0.522173in}}%
\pgfpathlineto{\pgfqpoint{2.407275in}{0.522396in}}%
\pgfpathlineto{\pgfqpoint{2.402883in}{0.522692in}}%
\pgfpathlineto{\pgfqpoint{2.398492in}{0.523080in}}%
\pgfpathlineto{\pgfqpoint{2.394100in}{0.523581in}}%
\pgfpathlineto{\pgfqpoint{2.389708in}{0.524219in}}%
\pgfpathlineto{\pgfqpoint{2.385317in}{0.525020in}}%
\pgfpathlineto{\pgfqpoint{2.380925in}{0.526011in}}%
\pgfpathlineto{\pgfqpoint{2.376534in}{0.527219in}}%
\pgfpathlineto{\pgfqpoint{2.372142in}{0.528670in}}%
\pgfpathlineto{\pgfqpoint{2.367750in}{0.530390in}}%
\pgfpathlineto{\pgfqpoint{2.363359in}{0.532397in}}%
\pgfpathlineto{\pgfqpoint{2.358967in}{0.534705in}}%
\pgfpathlineto{\pgfqpoint{2.354576in}{0.537320in}}%
\pgfpathlineto{\pgfqpoint{2.350184in}{0.540239in}}%
\pgfpathlineto{\pgfqpoint{2.345792in}{0.543449in}}%
\pgfpathlineto{\pgfqpoint{2.341401in}{0.546924in}}%
\pgfpathlineto{\pgfqpoint{2.337009in}{0.550629in}}%
\pgfpathlineto{\pgfqpoint{2.332618in}{0.554518in}}%
\pgfpathlineto{\pgfqpoint{2.328226in}{0.558537in}}%
\pgfpathlineto{\pgfqpoint{2.323834in}{0.562625in}}%
\pgfpathlineto{\pgfqpoint{2.319443in}{0.566719in}}%
\pgfpathlineto{\pgfqpoint{2.315051in}{0.570758in}}%
\pgfpathlineto{\pgfqpoint{2.310660in}{0.574685in}}%
\pgfpathlineto{\pgfqpoint{2.306268in}{0.578457in}}%
\pgfpathlineto{\pgfqpoint{2.301876in}{0.582043in}}%
\pgfpathlineto{\pgfqpoint{2.297485in}{0.585433in}}%
\pgfpathlineto{\pgfqpoint{2.293093in}{0.588643in}}%
\pgfpathlineto{\pgfqpoint{2.288702in}{0.591716in}}%
\pgfpathlineto{\pgfqpoint{2.284310in}{0.594721in}}%
\pgfpathlineto{\pgfqpoint{2.279918in}{0.597762in}}%
\pgfpathlineto{\pgfqpoint{2.275527in}{0.600969in}}%
\pgfpathlineto{\pgfqpoint{2.271135in}{0.604499in}}%
\pgfpathlineto{\pgfqpoint{2.266744in}{0.608531in}}%
\pgfpathlineto{\pgfqpoint{2.262352in}{0.613262in}}%
\pgfpathlineto{\pgfqpoint{2.257960in}{0.618898in}}%
\pgfpathlineto{\pgfqpoint{2.253569in}{0.625646in}}%
\pgfpathlineto{\pgfqpoint{2.249177in}{0.633711in}}%
\pgfpathlineto{\pgfqpoint{2.244785in}{0.643282in}}%
\pgfpathlineto{\pgfqpoint{2.240394in}{0.654532in}}%
\pgfpathlineto{\pgfqpoint{2.236002in}{0.667610in}}%
\pgfpathlineto{\pgfqpoint{2.231611in}{0.682639in}}%
\pgfpathlineto{\pgfqpoint{2.227219in}{0.699717in}}%
\pgfpathlineto{\pgfqpoint{2.222827in}{0.718918in}}%
\pgfpathlineto{\pgfqpoint{2.218436in}{0.740300in}}%
\pgfpathlineto{\pgfqpoint{2.214044in}{0.763907in}}%
\pgfpathlineto{\pgfqpoint{2.209653in}{0.789779in}}%
\pgfpathlineto{\pgfqpoint{2.205261in}{0.817960in}}%
\pgfpathlineto{\pgfqpoint{2.200869in}{0.848501in}}%
\pgfpathlineto{\pgfqpoint{2.196478in}{0.881468in}}%
\pgfpathlineto{\pgfqpoint{2.192086in}{0.916940in}}%
\pgfpathlineto{\pgfqpoint{2.187695in}{0.955010in}}%
\pgfpathlineto{\pgfqpoint{2.183303in}{0.995774in}}%
\pgfpathlineto{\pgfqpoint{2.178911in}{1.039327in}}%
\pgfpathlineto{\pgfqpoint{2.174520in}{1.085743in}}%
\pgfpathlineto{\pgfqpoint{2.170128in}{1.135066in}}%
\pgfpathlineto{\pgfqpoint{2.165737in}{1.187285in}}%
\pgfpathlineto{\pgfqpoint{2.161345in}{1.242326in}}%
\pgfpathlineto{\pgfqpoint{2.156953in}{1.300029in}}%
\pgfpathlineto{\pgfqpoint{2.152562in}{1.360141in}}%
\pgfpathlineto{\pgfqpoint{2.148170in}{1.422303in}}%
\pgfpathlineto{\pgfqpoint{2.143779in}{1.486057in}}%
\pgfpathlineto{\pgfqpoint{2.139387in}{1.550841in}}%
\pgfpathlineto{\pgfqpoint{2.134995in}{1.616010in}}%
\pgfpathlineto{\pgfqpoint{2.130604in}{1.680851in}}%
\pgfpathlineto{\pgfqpoint{2.126212in}{1.744613in}}%
\pgfpathlineto{\pgfqpoint{2.121821in}{1.806537in}}%
\pgfpathlineto{\pgfqpoint{2.117429in}{1.865893in}}%
\pgfpathlineto{\pgfqpoint{2.113037in}{1.922020in}}%
\pgfpathlineto{\pgfqpoint{2.108646in}{1.974361in}}%
\pgfpathlineto{\pgfqpoint{2.104254in}{2.022496in}}%
\pgfpathlineto{\pgfqpoint{2.099862in}{2.066168in}}%
\pgfpathlineto{\pgfqpoint{2.095471in}{2.105300in}}%
\pgfpathlineto{\pgfqpoint{2.091079in}{2.140001in}}%
\pgfpathlineto{\pgfqpoint{2.086688in}{2.170556in}}%
\pgfpathlineto{\pgfqpoint{2.082296in}{2.197403in}}%
\pgfpathlineto{\pgfqpoint{2.077904in}{2.221111in}}%
\pgfpathlineto{\pgfqpoint{2.073513in}{2.242328in}}%
\pgfpathlineto{\pgfqpoint{2.069121in}{2.261744in}}%
\pgfpathlineto{\pgfqpoint{2.064730in}{2.280041in}}%
\pgfpathlineto{\pgfqpoint{2.060338in}{2.297852in}}%
\pgfpathlineto{\pgfqpoint{2.055946in}{2.315722in}}%
\pgfpathlineto{\pgfqpoint{2.051555in}{2.334085in}}%
\pgfpathlineto{\pgfqpoint{2.047163in}{2.353250in}}%
\pgfpathlineto{\pgfqpoint{2.042772in}{2.373398in}}%
\pgfpathlineto{\pgfqpoint{2.038380in}{2.394597in}}%
\pgfpathlineto{\pgfqpoint{2.033988in}{2.416817in}}%
\pgfpathlineto{\pgfqpoint{2.029597in}{2.439960in}}%
\pgfpathlineto{\pgfqpoint{2.025205in}{2.463884in}}%
\pgfpathlineto{\pgfqpoint{2.020814in}{2.488428in}}%
\pgfpathlineto{\pgfqpoint{2.016422in}{2.513433in}}%
\pgfpathlineto{\pgfqpoint{2.012030in}{2.538749in}}%
\pgfpathlineto{\pgfqpoint{2.007639in}{2.564237in}}%
\pgfpathlineto{\pgfqpoint{2.003247in}{2.589762in}}%
\pgfpathlineto{\pgfqpoint{1.998856in}{2.615177in}}%
\pgfpathlineto{\pgfqpoint{1.994464in}{2.640299in}}%
\pgfpathlineto{\pgfqpoint{1.990072in}{2.664895in}}%
\pgfpathlineto{\pgfqpoint{1.985681in}{2.688659in}}%
\pgfpathlineto{\pgfqpoint{1.981289in}{2.711201in}}%
\pgfpathlineto{\pgfqpoint{1.976898in}{2.732046in}}%
\pgfpathlineto{\pgfqpoint{1.972506in}{2.750642in}}%
\pgfpathlineto{\pgfqpoint{1.968114in}{2.766378in}}%
\pgfpathlineto{\pgfqpoint{1.963723in}{2.778615in}}%
\pgfpathlineto{\pgfqpoint{1.959331in}{2.786720in}}%
\pgfpathlineto{\pgfqpoint{1.954940in}{2.790110in}}%
\pgfpathlineto{\pgfqpoint{1.950548in}{2.788288in}}%
\pgfpathlineto{\pgfqpoint{1.946156in}{2.780882in}}%
\pgfpathlineto{\pgfqpoint{1.941765in}{2.767674in}}%
\pgfpathlineto{\pgfqpoint{1.937373in}{2.748616in}}%
\pgfpathlineto{\pgfqpoint{1.932981in}{2.723836in}}%
\pgfpathlineto{\pgfqpoint{1.928590in}{2.693633in}}%
\pgfpathlineto{\pgfqpoint{1.924198in}{2.658455in}}%
\pgfpathlineto{\pgfqpoint{1.919807in}{2.618867in}}%
\pgfpathlineto{\pgfqpoint{1.915415in}{2.575518in}}%
\pgfpathlineto{\pgfqpoint{1.911023in}{2.529097in}}%
\pgfpathlineto{\pgfqpoint{1.906632in}{2.480293in}}%
\pgfpathlineto{\pgfqpoint{1.902240in}{2.429754in}}%
\pgfpathlineto{\pgfqpoint{1.897849in}{2.378058in}}%
\pgfpathlineto{\pgfqpoint{1.893457in}{2.325689in}}%
\pgfpathlineto{\pgfqpoint{1.889065in}{2.273023in}}%
\pgfpathlineto{\pgfqpoint{1.884674in}{2.220323in}}%
\pgfpathlineto{\pgfqpoint{1.880282in}{2.167745in}}%
\pgfpathlineto{\pgfqpoint{1.875891in}{2.115352in}}%
\pgfpathlineto{\pgfqpoint{1.871499in}{2.063125in}}%
\pgfpathlineto{\pgfqpoint{1.867107in}{2.010994in}}%
\pgfpathlineto{\pgfqpoint{1.862716in}{1.958850in}}%
\pgfpathlineto{\pgfqpoint{1.858324in}{1.906571in}}%
\pgfpathlineto{\pgfqpoint{1.853933in}{1.854040in}}%
\pgfpathlineto{\pgfqpoint{1.849541in}{1.801159in}}%
\pgfpathlineto{\pgfqpoint{1.845149in}{1.747863in}}%
\pgfpathlineto{\pgfqpoint{1.840758in}{1.694122in}}%
\pgfpathlineto{\pgfqpoint{1.836366in}{1.639953in}}%
\pgfpathlineto{\pgfqpoint{1.831975in}{1.585414in}}%
\pgfpathlineto{\pgfqpoint{1.827583in}{1.530610in}}%
\pgfpathlineto{\pgfqpoint{1.823191in}{1.475689in}}%
\pgfpathlineto{\pgfqpoint{1.818800in}{1.420837in}}%
\pgfpathlineto{\pgfqpoint{1.814408in}{1.366280in}}%
\pgfpathlineto{\pgfqpoint{1.810017in}{1.312274in}}%
\pgfpathlineto{\pgfqpoint{1.805625in}{1.259101in}}%
\pgfpathlineto{\pgfqpoint{1.801233in}{1.207067in}}%
\pgfpathlineto{\pgfqpoint{1.796842in}{1.156489in}}%
\pgfpathlineto{\pgfqpoint{1.792450in}{1.107687in}}%
\pgfpathlineto{\pgfqpoint{1.788058in}{1.060977in}}%
\pgfpathlineto{\pgfqpoint{1.783667in}{1.016656in}}%
\pgfpathlineto{\pgfqpoint{1.779275in}{0.974990in}}%
\pgfpathlineto{\pgfqpoint{1.774884in}{0.936204in}}%
\pgfpathlineto{\pgfqpoint{1.770492in}{0.900470in}}%
\pgfpathlineto{\pgfqpoint{1.766100in}{0.867895in}}%
\pgfpathlineto{\pgfqpoint{1.761709in}{0.838522in}}%
\pgfpathlineto{\pgfqpoint{1.757317in}{0.812318in}}%
\pgfpathlineto{\pgfqpoint{1.752926in}{0.789181in}}%
\pgfpathlineto{\pgfqpoint{1.748534in}{0.768941in}}%
\pgfpathlineto{\pgfqpoint{1.744142in}{0.751373in}}%
\pgfpathlineto{\pgfqpoint{1.739751in}{0.736200in}}%
\pgfpathlineto{\pgfqpoint{1.735359in}{0.723117in}}%
\pgfpathlineto{\pgfqpoint{1.730968in}{0.711799in}}%
\pgfpathlineto{\pgfqpoint{1.726576in}{0.701919in}}%
\pgfpathlineto{\pgfqpoint{1.722184in}{0.693164in}}%
\pgfpathlineto{\pgfqpoint{1.717793in}{0.685244in}}%
\pgfpathlineto{\pgfqpoint{1.713401in}{0.677906in}}%
\pgfpathlineto{\pgfqpoint{1.709010in}{0.670940in}}%
\pgfpathlineto{\pgfqpoint{1.704618in}{0.664180in}}%
\pgfpathlineto{\pgfqpoint{1.700226in}{0.657506in}}%
\pgfpathlineto{\pgfqpoint{1.695835in}{0.650842in}}%
\pgfpathlineto{\pgfqpoint{1.691443in}{0.644147in}}%
\pgfpathlineto{\pgfqpoint{1.687052in}{0.637411in}}%
\pgfpathlineto{\pgfqpoint{1.682660in}{0.630648in}}%
\pgfpathlineto{\pgfqpoint{1.678268in}{0.623885in}}%
\pgfpathlineto{\pgfqpoint{1.673877in}{0.617158in}}%
\pgfpathlineto{\pgfqpoint{1.669485in}{0.610507in}}%
\pgfpathlineto{\pgfqpoint{1.665094in}{0.603969in}}%
\pgfpathlineto{\pgfqpoint{1.660702in}{0.597577in}}%
\pgfpathlineto{\pgfqpoint{1.656310in}{0.591361in}}%
\pgfpathlineto{\pgfqpoint{1.651919in}{0.585343in}}%
\pgfpathlineto{\pgfqpoint{1.647527in}{0.579544in}}%
\pgfpathlineto{\pgfqpoint{1.643135in}{0.573980in}}%
\pgfpathlineto{\pgfqpoint{1.638744in}{0.568667in}}%
\pgfpathlineto{\pgfqpoint{1.634352in}{0.563619in}}%
\pgfpathlineto{\pgfqpoint{1.629961in}{0.558851in}}%
\pgfpathlineto{\pgfqpoint{1.625569in}{0.554375in}}%
\pgfpathlineto{\pgfqpoint{1.621177in}{0.550204in}}%
\pgfpathlineto{\pgfqpoint{1.616786in}{0.546350in}}%
\pgfpathlineto{\pgfqpoint{1.612394in}{0.542821in}}%
\pgfpathlineto{\pgfqpoint{1.608003in}{0.539622in}}%
\pgfpathlineto{\pgfqpoint{1.603611in}{0.536752in}}%
\pgfpathlineto{\pgfqpoint{1.599219in}{0.534206in}}%
\pgfpathlineto{\pgfqpoint{1.594828in}{0.531976in}}%
\pgfpathlineto{\pgfqpoint{1.590436in}{0.530046in}}%
\pgfpathlineto{\pgfqpoint{1.586045in}{0.528397in}}%
\pgfpathlineto{\pgfqpoint{1.581653in}{0.527006in}}%
\pgfpathlineto{\pgfqpoint{1.577261in}{0.525850in}}%
\pgfpathlineto{\pgfqpoint{1.572870in}{0.524901in}}%
\pgfpathlineto{\pgfqpoint{1.568478in}{0.524132in}}%
\pgfpathlineto{\pgfqpoint{1.564087in}{0.523519in}}%
\pgfpathlineto{\pgfqpoint{1.559695in}{0.523036in}}%
\pgfpathlineto{\pgfqpoint{1.555303in}{0.522662in}}%
\pgfpathlineto{\pgfqpoint{1.550912in}{0.522375in}}%
\pgfpathlineto{\pgfqpoint{1.546520in}{0.522159in}}%
\pgfpathlineto{\pgfqpoint{1.542129in}{0.521998in}}%
\pgfpathlineto{\pgfqpoint{1.542129in}{0.521998in}}%
\pgfpathclose%
\pgfusepath{stroke,fill}%
}%
\begin{pgfscope}%
\pgfsys@transformshift{0.000000in}{0.000000in}%
\pgfsys@useobject{currentmarker}{}%
\end{pgfscope}%
\end{pgfscope}%
\begin{pgfscope}%
\pgfpathrectangle{\pgfqpoint{0.784801in}{0.521603in}}{\pgfqpoint{4.349472in}{2.424065in}}%
\pgfusepath{clip}%
\pgfsetbuttcap%
\pgfsetroundjoin%
\pgfsetlinewidth{0.803000pt}%
\definecolor{currentstroke}{rgb}{0.690196,0.690196,0.690196}%
\pgfsetstrokecolor{currentstroke}%
\pgfsetstrokeopacity{0.400000}%
\pgfsetdash{{2.960000pt}{1.280000pt}}{0.000000pt}%
\pgfpathmoveto{\pgfqpoint{0.784801in}{0.521603in}}%
\pgfpathlineto{\pgfqpoint{0.784801in}{2.945669in}}%
\pgfusepath{stroke}%
\end{pgfscope}%
\begin{pgfscope}%
\pgfsetbuttcap%
\pgfsetroundjoin%
\definecolor{currentfill}{rgb}{0.000000,0.000000,0.000000}%
\pgfsetfillcolor{currentfill}%
\pgfsetlinewidth{0.803000pt}%
\definecolor{currentstroke}{rgb}{0.000000,0.000000,0.000000}%
\pgfsetstrokecolor{currentstroke}%
\pgfsetdash{}{0pt}%
\pgfsys@defobject{currentmarker}{\pgfqpoint{0.000000in}{-0.048611in}}{\pgfqpoint{0.000000in}{0.000000in}}{%
\pgfpathmoveto{\pgfqpoint{0.000000in}{0.000000in}}%
\pgfpathlineto{\pgfqpoint{0.000000in}{-0.048611in}}%
\pgfusepath{stroke,fill}%
}%
\begin{pgfscope}%
\pgfsys@transformshift{0.784801in}{0.521603in}%
\pgfsys@useobject{currentmarker}{}%
\end{pgfscope}%
\end{pgfscope}%
\begin{pgfscope}%
\definecolor{textcolor}{rgb}{0.000000,0.000000,0.000000}%
\pgfsetstrokecolor{textcolor}%
\pgfsetfillcolor{textcolor}%
\pgftext[x=0.784801in,y=0.424381in,,top]{\color{textcolor}{\sffamily\fontsize{10.000000}{12.000000}\selectfont\catcode`\^=\active\def^{\ifmmode\sp\else\^{}\fi}\catcode`\%=\active\def%{\%}300}}%
\end{pgfscope}%
\begin{pgfscope}%
\pgfpathrectangle{\pgfqpoint{0.784801in}{0.521603in}}{\pgfqpoint{4.349472in}{2.424065in}}%
\pgfusepath{clip}%
\pgfsetbuttcap%
\pgfsetroundjoin%
\pgfsetlinewidth{0.803000pt}%
\definecolor{currentstroke}{rgb}{0.690196,0.690196,0.690196}%
\pgfsetstrokecolor{currentstroke}%
\pgfsetstrokeopacity{0.400000}%
\pgfsetdash{{2.960000pt}{1.280000pt}}{0.000000pt}%
\pgfpathmoveto{\pgfqpoint{1.509713in}{0.521603in}}%
\pgfpathlineto{\pgfqpoint{1.509713in}{2.945669in}}%
\pgfusepath{stroke}%
\end{pgfscope}%
\begin{pgfscope}%
\pgfsetbuttcap%
\pgfsetroundjoin%
\definecolor{currentfill}{rgb}{0.000000,0.000000,0.000000}%
\pgfsetfillcolor{currentfill}%
\pgfsetlinewidth{0.803000pt}%
\definecolor{currentstroke}{rgb}{0.000000,0.000000,0.000000}%
\pgfsetstrokecolor{currentstroke}%
\pgfsetdash{}{0pt}%
\pgfsys@defobject{currentmarker}{\pgfqpoint{0.000000in}{-0.048611in}}{\pgfqpoint{0.000000in}{0.000000in}}{%
\pgfpathmoveto{\pgfqpoint{0.000000in}{0.000000in}}%
\pgfpathlineto{\pgfqpoint{0.000000in}{-0.048611in}}%
\pgfusepath{stroke,fill}%
}%
\begin{pgfscope}%
\pgfsys@transformshift{1.509713in}{0.521603in}%
\pgfsys@useobject{currentmarker}{}%
\end{pgfscope}%
\end{pgfscope}%
\begin{pgfscope}%
\definecolor{textcolor}{rgb}{0.000000,0.000000,0.000000}%
\pgfsetstrokecolor{textcolor}%
\pgfsetfillcolor{textcolor}%
\pgftext[x=1.509713in,y=0.424381in,,top]{\color{textcolor}{\sffamily\fontsize{10.000000}{12.000000}\selectfont\catcode`\^=\active\def^{\ifmmode\sp\else\^{}\fi}\catcode`\%=\active\def%{\%}400}}%
\end{pgfscope}%
\begin{pgfscope}%
\pgfpathrectangle{\pgfqpoint{0.784801in}{0.521603in}}{\pgfqpoint{4.349472in}{2.424065in}}%
\pgfusepath{clip}%
\pgfsetbuttcap%
\pgfsetroundjoin%
\pgfsetlinewidth{0.803000pt}%
\definecolor{currentstroke}{rgb}{0.690196,0.690196,0.690196}%
\pgfsetstrokecolor{currentstroke}%
\pgfsetstrokeopacity{0.400000}%
\pgfsetdash{{2.960000pt}{1.280000pt}}{0.000000pt}%
\pgfpathmoveto{\pgfqpoint{2.234625in}{0.521603in}}%
\pgfpathlineto{\pgfqpoint{2.234625in}{2.945669in}}%
\pgfusepath{stroke}%
\end{pgfscope}%
\begin{pgfscope}%
\pgfsetbuttcap%
\pgfsetroundjoin%
\definecolor{currentfill}{rgb}{0.000000,0.000000,0.000000}%
\pgfsetfillcolor{currentfill}%
\pgfsetlinewidth{0.803000pt}%
\definecolor{currentstroke}{rgb}{0.000000,0.000000,0.000000}%
\pgfsetstrokecolor{currentstroke}%
\pgfsetdash{}{0pt}%
\pgfsys@defobject{currentmarker}{\pgfqpoint{0.000000in}{-0.048611in}}{\pgfqpoint{0.000000in}{0.000000in}}{%
\pgfpathmoveto{\pgfqpoint{0.000000in}{0.000000in}}%
\pgfpathlineto{\pgfqpoint{0.000000in}{-0.048611in}}%
\pgfusepath{stroke,fill}%
}%
\begin{pgfscope}%
\pgfsys@transformshift{2.234625in}{0.521603in}%
\pgfsys@useobject{currentmarker}{}%
\end{pgfscope}%
\end{pgfscope}%
\begin{pgfscope}%
\definecolor{textcolor}{rgb}{0.000000,0.000000,0.000000}%
\pgfsetstrokecolor{textcolor}%
\pgfsetfillcolor{textcolor}%
\pgftext[x=2.234625in,y=0.424381in,,top]{\color{textcolor}{\sffamily\fontsize{10.000000}{12.000000}\selectfont\catcode`\^=\active\def^{\ifmmode\sp\else\^{}\fi}\catcode`\%=\active\def%{\%}500}}%
\end{pgfscope}%
\begin{pgfscope}%
\pgfpathrectangle{\pgfqpoint{0.784801in}{0.521603in}}{\pgfqpoint{4.349472in}{2.424065in}}%
\pgfusepath{clip}%
\pgfsetbuttcap%
\pgfsetroundjoin%
\pgfsetlinewidth{0.803000pt}%
\definecolor{currentstroke}{rgb}{0.690196,0.690196,0.690196}%
\pgfsetstrokecolor{currentstroke}%
\pgfsetstrokeopacity{0.400000}%
\pgfsetdash{{2.960000pt}{1.280000pt}}{0.000000pt}%
\pgfpathmoveto{\pgfqpoint{2.959537in}{0.521603in}}%
\pgfpathlineto{\pgfqpoint{2.959537in}{2.945669in}}%
\pgfusepath{stroke}%
\end{pgfscope}%
\begin{pgfscope}%
\pgfsetbuttcap%
\pgfsetroundjoin%
\definecolor{currentfill}{rgb}{0.000000,0.000000,0.000000}%
\pgfsetfillcolor{currentfill}%
\pgfsetlinewidth{0.803000pt}%
\definecolor{currentstroke}{rgb}{0.000000,0.000000,0.000000}%
\pgfsetstrokecolor{currentstroke}%
\pgfsetdash{}{0pt}%
\pgfsys@defobject{currentmarker}{\pgfqpoint{0.000000in}{-0.048611in}}{\pgfqpoint{0.000000in}{0.000000in}}{%
\pgfpathmoveto{\pgfqpoint{0.000000in}{0.000000in}}%
\pgfpathlineto{\pgfqpoint{0.000000in}{-0.048611in}}%
\pgfusepath{stroke,fill}%
}%
\begin{pgfscope}%
\pgfsys@transformshift{2.959537in}{0.521603in}%
\pgfsys@useobject{currentmarker}{}%
\end{pgfscope}%
\end{pgfscope}%
\begin{pgfscope}%
\definecolor{textcolor}{rgb}{0.000000,0.000000,0.000000}%
\pgfsetstrokecolor{textcolor}%
\pgfsetfillcolor{textcolor}%
\pgftext[x=2.959537in,y=0.424381in,,top]{\color{textcolor}{\sffamily\fontsize{10.000000}{12.000000}\selectfont\catcode`\^=\active\def^{\ifmmode\sp\else\^{}\fi}\catcode`\%=\active\def%{\%}600}}%
\end{pgfscope}%
\begin{pgfscope}%
\pgfpathrectangle{\pgfqpoint{0.784801in}{0.521603in}}{\pgfqpoint{4.349472in}{2.424065in}}%
\pgfusepath{clip}%
\pgfsetbuttcap%
\pgfsetroundjoin%
\pgfsetlinewidth{0.803000pt}%
\definecolor{currentstroke}{rgb}{0.690196,0.690196,0.690196}%
\pgfsetstrokecolor{currentstroke}%
\pgfsetstrokeopacity{0.400000}%
\pgfsetdash{{2.960000pt}{1.280000pt}}{0.000000pt}%
\pgfpathmoveto{\pgfqpoint{3.684449in}{0.521603in}}%
\pgfpathlineto{\pgfqpoint{3.684449in}{2.945669in}}%
\pgfusepath{stroke}%
\end{pgfscope}%
\begin{pgfscope}%
\pgfsetbuttcap%
\pgfsetroundjoin%
\definecolor{currentfill}{rgb}{0.000000,0.000000,0.000000}%
\pgfsetfillcolor{currentfill}%
\pgfsetlinewidth{0.803000pt}%
\definecolor{currentstroke}{rgb}{0.000000,0.000000,0.000000}%
\pgfsetstrokecolor{currentstroke}%
\pgfsetdash{}{0pt}%
\pgfsys@defobject{currentmarker}{\pgfqpoint{0.000000in}{-0.048611in}}{\pgfqpoint{0.000000in}{0.000000in}}{%
\pgfpathmoveto{\pgfqpoint{0.000000in}{0.000000in}}%
\pgfpathlineto{\pgfqpoint{0.000000in}{-0.048611in}}%
\pgfusepath{stroke,fill}%
}%
\begin{pgfscope}%
\pgfsys@transformshift{3.684449in}{0.521603in}%
\pgfsys@useobject{currentmarker}{}%
\end{pgfscope}%
\end{pgfscope}%
\begin{pgfscope}%
\definecolor{textcolor}{rgb}{0.000000,0.000000,0.000000}%
\pgfsetstrokecolor{textcolor}%
\pgfsetfillcolor{textcolor}%
\pgftext[x=3.684449in,y=0.424381in,,top]{\color{textcolor}{\sffamily\fontsize{10.000000}{12.000000}\selectfont\catcode`\^=\active\def^{\ifmmode\sp\else\^{}\fi}\catcode`\%=\active\def%{\%}700}}%
\end{pgfscope}%
\begin{pgfscope}%
\pgfpathrectangle{\pgfqpoint{0.784801in}{0.521603in}}{\pgfqpoint{4.349472in}{2.424065in}}%
\pgfusepath{clip}%
\pgfsetbuttcap%
\pgfsetroundjoin%
\pgfsetlinewidth{0.803000pt}%
\definecolor{currentstroke}{rgb}{0.690196,0.690196,0.690196}%
\pgfsetstrokecolor{currentstroke}%
\pgfsetstrokeopacity{0.400000}%
\pgfsetdash{{2.960000pt}{1.280000pt}}{0.000000pt}%
\pgfpathmoveto{\pgfqpoint{4.409361in}{0.521603in}}%
\pgfpathlineto{\pgfqpoint{4.409361in}{2.945669in}}%
\pgfusepath{stroke}%
\end{pgfscope}%
\begin{pgfscope}%
\pgfsetbuttcap%
\pgfsetroundjoin%
\definecolor{currentfill}{rgb}{0.000000,0.000000,0.000000}%
\pgfsetfillcolor{currentfill}%
\pgfsetlinewidth{0.803000pt}%
\definecolor{currentstroke}{rgb}{0.000000,0.000000,0.000000}%
\pgfsetstrokecolor{currentstroke}%
\pgfsetdash{}{0pt}%
\pgfsys@defobject{currentmarker}{\pgfqpoint{0.000000in}{-0.048611in}}{\pgfqpoint{0.000000in}{0.000000in}}{%
\pgfpathmoveto{\pgfqpoint{0.000000in}{0.000000in}}%
\pgfpathlineto{\pgfqpoint{0.000000in}{-0.048611in}}%
\pgfusepath{stroke,fill}%
}%
\begin{pgfscope}%
\pgfsys@transformshift{4.409361in}{0.521603in}%
\pgfsys@useobject{currentmarker}{}%
\end{pgfscope}%
\end{pgfscope}%
\begin{pgfscope}%
\definecolor{textcolor}{rgb}{0.000000,0.000000,0.000000}%
\pgfsetstrokecolor{textcolor}%
\pgfsetfillcolor{textcolor}%
\pgftext[x=4.409361in,y=0.424381in,,top]{\color{textcolor}{\sffamily\fontsize{10.000000}{12.000000}\selectfont\catcode`\^=\active\def^{\ifmmode\sp\else\^{}\fi}\catcode`\%=\active\def%{\%}800}}%
\end{pgfscope}%
\begin{pgfscope}%
\pgfpathrectangle{\pgfqpoint{0.784801in}{0.521603in}}{\pgfqpoint{4.349472in}{2.424065in}}%
\pgfusepath{clip}%
\pgfsetbuttcap%
\pgfsetroundjoin%
\pgfsetlinewidth{0.803000pt}%
\definecolor{currentstroke}{rgb}{0.690196,0.690196,0.690196}%
\pgfsetstrokecolor{currentstroke}%
\pgfsetstrokeopacity{0.400000}%
\pgfsetdash{{2.960000pt}{1.280000pt}}{0.000000pt}%
\pgfpathmoveto{\pgfqpoint{5.134273in}{0.521603in}}%
\pgfpathlineto{\pgfqpoint{5.134273in}{2.945669in}}%
\pgfusepath{stroke}%
\end{pgfscope}%
\begin{pgfscope}%
\pgfsetbuttcap%
\pgfsetroundjoin%
\definecolor{currentfill}{rgb}{0.000000,0.000000,0.000000}%
\pgfsetfillcolor{currentfill}%
\pgfsetlinewidth{0.803000pt}%
\definecolor{currentstroke}{rgb}{0.000000,0.000000,0.000000}%
\pgfsetstrokecolor{currentstroke}%
\pgfsetdash{}{0pt}%
\pgfsys@defobject{currentmarker}{\pgfqpoint{0.000000in}{-0.048611in}}{\pgfqpoint{0.000000in}{0.000000in}}{%
\pgfpathmoveto{\pgfqpoint{0.000000in}{0.000000in}}%
\pgfpathlineto{\pgfqpoint{0.000000in}{-0.048611in}}%
\pgfusepath{stroke,fill}%
}%
\begin{pgfscope}%
\pgfsys@transformshift{5.134273in}{0.521603in}%
\pgfsys@useobject{currentmarker}{}%
\end{pgfscope}%
\end{pgfscope}%
\begin{pgfscope}%
\definecolor{textcolor}{rgb}{0.000000,0.000000,0.000000}%
\pgfsetstrokecolor{textcolor}%
\pgfsetfillcolor{textcolor}%
\pgftext[x=5.134273in,y=0.424381in,,top]{\color{textcolor}{\sffamily\fontsize{10.000000}{12.000000}\selectfont\catcode`\^=\active\def^{\ifmmode\sp\else\^{}\fi}\catcode`\%=\active\def%{\%}900}}%
\end{pgfscope}%
\begin{pgfscope}%
\definecolor{textcolor}{rgb}{0.000000,0.000000,0.000000}%
\pgfsetstrokecolor{textcolor}%
\pgfsetfillcolor{textcolor}%
\pgftext[x=2.959537in,y=0.234413in,,top]{\color{textcolor}{\sffamily\fontsize{10.000000}{12.000000}\selectfont\catcode`\^=\active\def^{\ifmmode\sp\else\^{}\fi}\catcode`\%=\active\def%{\%}Handshake Time (ms)}}%
\end{pgfscope}%
\begin{pgfscope}%
\pgfpathrectangle{\pgfqpoint{0.784801in}{0.521603in}}{\pgfqpoint{4.349472in}{2.424065in}}%
\pgfusepath{clip}%
\pgfsetbuttcap%
\pgfsetroundjoin%
\pgfsetlinewidth{0.803000pt}%
\definecolor{currentstroke}{rgb}{0.690196,0.690196,0.690196}%
\pgfsetstrokecolor{currentstroke}%
\pgfsetstrokeopacity{0.400000}%
\pgfsetdash{{2.960000pt}{1.280000pt}}{0.000000pt}%
\pgfpathmoveto{\pgfqpoint{0.784801in}{0.521603in}}%
\pgfpathlineto{\pgfqpoint{5.134273in}{0.521603in}}%
\pgfusepath{stroke}%
\end{pgfscope}%
\begin{pgfscope}%
\pgfsetbuttcap%
\pgfsetroundjoin%
\definecolor{currentfill}{rgb}{0.000000,0.000000,0.000000}%
\pgfsetfillcolor{currentfill}%
\pgfsetlinewidth{0.803000pt}%
\definecolor{currentstroke}{rgb}{0.000000,0.000000,0.000000}%
\pgfsetstrokecolor{currentstroke}%
\pgfsetdash{}{0pt}%
\pgfsys@defobject{currentmarker}{\pgfqpoint{-0.048611in}{0.000000in}}{\pgfqpoint{-0.000000in}{0.000000in}}{%
\pgfpathmoveto{\pgfqpoint{-0.000000in}{0.000000in}}%
\pgfpathlineto{\pgfqpoint{-0.048611in}{0.000000in}}%
\pgfusepath{stroke,fill}%
}%
\begin{pgfscope}%
\pgfsys@transformshift{0.784801in}{0.521603in}%
\pgfsys@useobject{currentmarker}{}%
\end{pgfscope}%
\end{pgfscope}%
\begin{pgfscope}%
\definecolor{textcolor}{rgb}{0.000000,0.000000,0.000000}%
\pgfsetstrokecolor{textcolor}%
\pgfsetfillcolor{textcolor}%
\pgftext[x=0.289968in, y=0.468842in, left, base]{\color{textcolor}{\sffamily\fontsize{10.000000}{12.000000}\selectfont\catcode`\^=\active\def^{\ifmmode\sp\else\^{}\fi}\catcode`\%=\active\def%{\%}0.000}}%
\end{pgfscope}%
\begin{pgfscope}%
\pgfpathrectangle{\pgfqpoint{0.784801in}{0.521603in}}{\pgfqpoint{4.349472in}{2.424065in}}%
\pgfusepath{clip}%
\pgfsetbuttcap%
\pgfsetroundjoin%
\pgfsetlinewidth{0.803000pt}%
\definecolor{currentstroke}{rgb}{0.690196,0.690196,0.690196}%
\pgfsetstrokecolor{currentstroke}%
\pgfsetstrokeopacity{0.400000}%
\pgfsetdash{{2.960000pt}{1.280000pt}}{0.000000pt}%
\pgfpathmoveto{\pgfqpoint{0.784801in}{0.984355in}}%
\pgfpathlineto{\pgfqpoint{5.134273in}{0.984355in}}%
\pgfusepath{stroke}%
\end{pgfscope}%
\begin{pgfscope}%
\pgfsetbuttcap%
\pgfsetroundjoin%
\definecolor{currentfill}{rgb}{0.000000,0.000000,0.000000}%
\pgfsetfillcolor{currentfill}%
\pgfsetlinewidth{0.803000pt}%
\definecolor{currentstroke}{rgb}{0.000000,0.000000,0.000000}%
\pgfsetstrokecolor{currentstroke}%
\pgfsetdash{}{0pt}%
\pgfsys@defobject{currentmarker}{\pgfqpoint{-0.048611in}{0.000000in}}{\pgfqpoint{-0.000000in}{0.000000in}}{%
\pgfpathmoveto{\pgfqpoint{-0.000000in}{0.000000in}}%
\pgfpathlineto{\pgfqpoint{-0.048611in}{0.000000in}}%
\pgfusepath{stroke,fill}%
}%
\begin{pgfscope}%
\pgfsys@transformshift{0.784801in}{0.984355in}%
\pgfsys@useobject{currentmarker}{}%
\end{pgfscope}%
\end{pgfscope}%
\begin{pgfscope}%
\definecolor{textcolor}{rgb}{0.000000,0.000000,0.000000}%
\pgfsetstrokecolor{textcolor}%
\pgfsetfillcolor{textcolor}%
\pgftext[x=0.289968in, y=0.931594in, left, base]{\color{textcolor}{\sffamily\fontsize{10.000000}{12.000000}\selectfont\catcode`\^=\active\def^{\ifmmode\sp\else\^{}\fi}\catcode`\%=\active\def%{\%}0.005}}%
\end{pgfscope}%
\begin{pgfscope}%
\pgfpathrectangle{\pgfqpoint{0.784801in}{0.521603in}}{\pgfqpoint{4.349472in}{2.424065in}}%
\pgfusepath{clip}%
\pgfsetbuttcap%
\pgfsetroundjoin%
\pgfsetlinewidth{0.803000pt}%
\definecolor{currentstroke}{rgb}{0.690196,0.690196,0.690196}%
\pgfsetstrokecolor{currentstroke}%
\pgfsetstrokeopacity{0.400000}%
\pgfsetdash{{2.960000pt}{1.280000pt}}{0.000000pt}%
\pgfpathmoveto{\pgfqpoint{0.784801in}{1.447107in}}%
\pgfpathlineto{\pgfqpoint{5.134273in}{1.447107in}}%
\pgfusepath{stroke}%
\end{pgfscope}%
\begin{pgfscope}%
\pgfsetbuttcap%
\pgfsetroundjoin%
\definecolor{currentfill}{rgb}{0.000000,0.000000,0.000000}%
\pgfsetfillcolor{currentfill}%
\pgfsetlinewidth{0.803000pt}%
\definecolor{currentstroke}{rgb}{0.000000,0.000000,0.000000}%
\pgfsetstrokecolor{currentstroke}%
\pgfsetdash{}{0pt}%
\pgfsys@defobject{currentmarker}{\pgfqpoint{-0.048611in}{0.000000in}}{\pgfqpoint{-0.000000in}{0.000000in}}{%
\pgfpathmoveto{\pgfqpoint{-0.000000in}{0.000000in}}%
\pgfpathlineto{\pgfqpoint{-0.048611in}{0.000000in}}%
\pgfusepath{stroke,fill}%
}%
\begin{pgfscope}%
\pgfsys@transformshift{0.784801in}{1.447107in}%
\pgfsys@useobject{currentmarker}{}%
\end{pgfscope}%
\end{pgfscope}%
\begin{pgfscope}%
\definecolor{textcolor}{rgb}{0.000000,0.000000,0.000000}%
\pgfsetstrokecolor{textcolor}%
\pgfsetfillcolor{textcolor}%
\pgftext[x=0.289968in, y=1.394346in, left, base]{\color{textcolor}{\sffamily\fontsize{10.000000}{12.000000}\selectfont\catcode`\^=\active\def^{\ifmmode\sp\else\^{}\fi}\catcode`\%=\active\def%{\%}0.010}}%
\end{pgfscope}%
\begin{pgfscope}%
\pgfpathrectangle{\pgfqpoint{0.784801in}{0.521603in}}{\pgfqpoint{4.349472in}{2.424065in}}%
\pgfusepath{clip}%
\pgfsetbuttcap%
\pgfsetroundjoin%
\pgfsetlinewidth{0.803000pt}%
\definecolor{currentstroke}{rgb}{0.690196,0.690196,0.690196}%
\pgfsetstrokecolor{currentstroke}%
\pgfsetstrokeopacity{0.400000}%
\pgfsetdash{{2.960000pt}{1.280000pt}}{0.000000pt}%
\pgfpathmoveto{\pgfqpoint{0.784801in}{1.909859in}}%
\pgfpathlineto{\pgfqpoint{5.134273in}{1.909859in}}%
\pgfusepath{stroke}%
\end{pgfscope}%
\begin{pgfscope}%
\pgfsetbuttcap%
\pgfsetroundjoin%
\definecolor{currentfill}{rgb}{0.000000,0.000000,0.000000}%
\pgfsetfillcolor{currentfill}%
\pgfsetlinewidth{0.803000pt}%
\definecolor{currentstroke}{rgb}{0.000000,0.000000,0.000000}%
\pgfsetstrokecolor{currentstroke}%
\pgfsetdash{}{0pt}%
\pgfsys@defobject{currentmarker}{\pgfqpoint{-0.048611in}{0.000000in}}{\pgfqpoint{-0.000000in}{0.000000in}}{%
\pgfpathmoveto{\pgfqpoint{-0.000000in}{0.000000in}}%
\pgfpathlineto{\pgfqpoint{-0.048611in}{0.000000in}}%
\pgfusepath{stroke,fill}%
}%
\begin{pgfscope}%
\pgfsys@transformshift{0.784801in}{1.909859in}%
\pgfsys@useobject{currentmarker}{}%
\end{pgfscope}%
\end{pgfscope}%
\begin{pgfscope}%
\definecolor{textcolor}{rgb}{0.000000,0.000000,0.000000}%
\pgfsetstrokecolor{textcolor}%
\pgfsetfillcolor{textcolor}%
\pgftext[x=0.289968in, y=1.857098in, left, base]{\color{textcolor}{\sffamily\fontsize{10.000000}{12.000000}\selectfont\catcode`\^=\active\def^{\ifmmode\sp\else\^{}\fi}\catcode`\%=\active\def%{\%}0.015}}%
\end{pgfscope}%
\begin{pgfscope}%
\pgfpathrectangle{\pgfqpoint{0.784801in}{0.521603in}}{\pgfqpoint{4.349472in}{2.424065in}}%
\pgfusepath{clip}%
\pgfsetbuttcap%
\pgfsetroundjoin%
\pgfsetlinewidth{0.803000pt}%
\definecolor{currentstroke}{rgb}{0.690196,0.690196,0.690196}%
\pgfsetstrokecolor{currentstroke}%
\pgfsetstrokeopacity{0.400000}%
\pgfsetdash{{2.960000pt}{1.280000pt}}{0.000000pt}%
\pgfpathmoveto{\pgfqpoint{0.784801in}{2.372611in}}%
\pgfpathlineto{\pgfqpoint{5.134273in}{2.372611in}}%
\pgfusepath{stroke}%
\end{pgfscope}%
\begin{pgfscope}%
\pgfsetbuttcap%
\pgfsetroundjoin%
\definecolor{currentfill}{rgb}{0.000000,0.000000,0.000000}%
\pgfsetfillcolor{currentfill}%
\pgfsetlinewidth{0.803000pt}%
\definecolor{currentstroke}{rgb}{0.000000,0.000000,0.000000}%
\pgfsetstrokecolor{currentstroke}%
\pgfsetdash{}{0pt}%
\pgfsys@defobject{currentmarker}{\pgfqpoint{-0.048611in}{0.000000in}}{\pgfqpoint{-0.000000in}{0.000000in}}{%
\pgfpathmoveto{\pgfqpoint{-0.000000in}{0.000000in}}%
\pgfpathlineto{\pgfqpoint{-0.048611in}{0.000000in}}%
\pgfusepath{stroke,fill}%
}%
\begin{pgfscope}%
\pgfsys@transformshift{0.784801in}{2.372611in}%
\pgfsys@useobject{currentmarker}{}%
\end{pgfscope}%
\end{pgfscope}%
\begin{pgfscope}%
\definecolor{textcolor}{rgb}{0.000000,0.000000,0.000000}%
\pgfsetstrokecolor{textcolor}%
\pgfsetfillcolor{textcolor}%
\pgftext[x=0.289968in, y=2.319849in, left, base]{\color{textcolor}{\sffamily\fontsize{10.000000}{12.000000}\selectfont\catcode`\^=\active\def^{\ifmmode\sp\else\^{}\fi}\catcode`\%=\active\def%{\%}0.020}}%
\end{pgfscope}%
\begin{pgfscope}%
\pgfpathrectangle{\pgfqpoint{0.784801in}{0.521603in}}{\pgfqpoint{4.349472in}{2.424065in}}%
\pgfusepath{clip}%
\pgfsetbuttcap%
\pgfsetroundjoin%
\pgfsetlinewidth{0.803000pt}%
\definecolor{currentstroke}{rgb}{0.690196,0.690196,0.690196}%
\pgfsetstrokecolor{currentstroke}%
\pgfsetstrokeopacity{0.400000}%
\pgfsetdash{{2.960000pt}{1.280000pt}}{0.000000pt}%
\pgfpathmoveto{\pgfqpoint{0.784801in}{2.835363in}}%
\pgfpathlineto{\pgfqpoint{5.134273in}{2.835363in}}%
\pgfusepath{stroke}%
\end{pgfscope}%
\begin{pgfscope}%
\pgfsetbuttcap%
\pgfsetroundjoin%
\definecolor{currentfill}{rgb}{0.000000,0.000000,0.000000}%
\pgfsetfillcolor{currentfill}%
\pgfsetlinewidth{0.803000pt}%
\definecolor{currentstroke}{rgb}{0.000000,0.000000,0.000000}%
\pgfsetstrokecolor{currentstroke}%
\pgfsetdash{}{0pt}%
\pgfsys@defobject{currentmarker}{\pgfqpoint{-0.048611in}{0.000000in}}{\pgfqpoint{-0.000000in}{0.000000in}}{%
\pgfpathmoveto{\pgfqpoint{-0.000000in}{0.000000in}}%
\pgfpathlineto{\pgfqpoint{-0.048611in}{0.000000in}}%
\pgfusepath{stroke,fill}%
}%
\begin{pgfscope}%
\pgfsys@transformshift{0.784801in}{2.835363in}%
\pgfsys@useobject{currentmarker}{}%
\end{pgfscope}%
\end{pgfscope}%
\begin{pgfscope}%
\definecolor{textcolor}{rgb}{0.000000,0.000000,0.000000}%
\pgfsetstrokecolor{textcolor}%
\pgfsetfillcolor{textcolor}%
\pgftext[x=0.289968in, y=2.782601in, left, base]{\color{textcolor}{\sffamily\fontsize{10.000000}{12.000000}\selectfont\catcode`\^=\active\def^{\ifmmode\sp\else\^{}\fi}\catcode`\%=\active\def%{\%}0.025}}%
\end{pgfscope}%
\begin{pgfscope}%
\definecolor{textcolor}{rgb}{0.000000,0.000000,0.000000}%
\pgfsetstrokecolor{textcolor}%
\pgfsetfillcolor{textcolor}%
\pgftext[x=0.234413in,y=1.733636in,,bottom,rotate=90.000000]{\color{textcolor}{\sffamily\fontsize{10.000000}{12.000000}\selectfont\catcode`\^=\active\def^{\ifmmode\sp\else\^{}\fi}\catcode`\%=\active\def%{\%}Probability Density}}%
\end{pgfscope}%
\begin{pgfscope}%
\pgfpathrectangle{\pgfqpoint{0.784801in}{0.521603in}}{\pgfqpoint{4.349472in}{2.424065in}}%
\pgfusepath{clip}%
\pgfsetbuttcap%
\pgfsetroundjoin%
\pgfsetlinewidth{1.204500pt}%
\definecolor{currentstroke}{rgb}{1.000000,0.000000,0.000000}%
\pgfsetstrokecolor{currentstroke}%
\pgfsetstrokeopacity{0.800000}%
\pgfsetdash{{4.440000pt}{1.920000pt}}{0.000000pt}%
\pgfpathmoveto{\pgfqpoint{2.959537in}{0.521603in}}%
\pgfpathlineto{\pgfqpoint{2.959537in}{2.945669in}}%
\pgfusepath{stroke}%
\end{pgfscope}%
\begin{pgfscope}%
\pgfpathrectangle{\pgfqpoint{0.784801in}{0.521603in}}{\pgfqpoint{4.349472in}{2.424065in}}%
\pgfusepath{clip}%
\pgfsetbuttcap%
\pgfsetroundjoin%
\pgfsetlinewidth{1.204500pt}%
\definecolor{currentstroke}{rgb}{1.000000,0.000000,0.000000}%
\pgfsetstrokecolor{currentstroke}%
\pgfsetstrokeopacity{0.800000}%
\pgfsetdash{{4.440000pt}{1.920000pt}}{0.000000pt}%
\pgfpathmoveto{\pgfqpoint{1.872169in}{0.521603in}}%
\pgfpathlineto{\pgfqpoint{1.872169in}{2.945669in}}%
\pgfusepath{stroke}%
\end{pgfscope}%
\begin{pgfscope}%
\pgfsetrectcap%
\pgfsetmiterjoin%
\pgfsetlinewidth{0.803000pt}%
\definecolor{currentstroke}{rgb}{0.000000,0.000000,0.000000}%
\pgfsetstrokecolor{currentstroke}%
\pgfsetdash{}{0pt}%
\pgfpathmoveto{\pgfqpoint{0.784801in}{0.521603in}}%
\pgfpathlineto{\pgfqpoint{0.784801in}{2.945669in}}%
\pgfusepath{stroke}%
\end{pgfscope}%
\begin{pgfscope}%
\pgfsetrectcap%
\pgfsetmiterjoin%
\pgfsetlinewidth{0.803000pt}%
\definecolor{currentstroke}{rgb}{0.000000,0.000000,0.000000}%
\pgfsetstrokecolor{currentstroke}%
\pgfsetdash{}{0pt}%
\pgfpathmoveto{\pgfqpoint{0.784801in}{0.521603in}}%
\pgfpathlineto{\pgfqpoint{5.134273in}{0.521603in}}%
\pgfusepath{stroke}%
\end{pgfscope}%
\begin{pgfscope}%
\definecolor{textcolor}{rgb}{0.000000,0.000000,0.000000}%
\pgfsetstrokecolor{textcolor}%
\pgfsetfillcolor{textcolor}%
\pgftext[x=2.959537in,y=3.029002in,,base]{\color{textcolor}{\sffamily\fontsize{12.000000}{14.400000}\selectfont\catcode`\^=\active\def^{\ifmmode\sp\else\^{}\fi}\catcode`\%=\active\def%{\%}Effect of Candidate Racing on Connection Establishment Time}}%
\end{pgfscope}%
\begin{pgfscope}%
\pgfsetbuttcap%
\pgfsetmiterjoin%
\definecolor{currentfill}{rgb}{1.000000,1.000000,1.000000}%
\pgfsetfillcolor{currentfill}%
\pgfsetfillopacity{0.800000}%
\pgfsetlinewidth{1.003750pt}%
\definecolor{currentstroke}{rgb}{0.800000,0.800000,0.800000}%
\pgfsetstrokecolor{currentstroke}%
\pgfsetstrokeopacity{0.800000}%
\pgfsetdash{}{0pt}%
\pgfpathmoveto{\pgfqpoint{3.568974in}{2.246773in}}%
\pgfpathlineto{\pgfqpoint{5.175822in}{2.246773in}}%
\pgfpathquadraticcurveto{\pgfqpoint{5.200822in}{2.246773in}}{\pgfqpoint{5.200822in}{2.271773in}}%
\pgfpathlineto{\pgfqpoint{5.200822in}{2.809687in}}%
\pgfpathquadraticcurveto{\pgfqpoint{5.200822in}{2.834687in}}{\pgfqpoint{5.175822in}{2.834687in}}%
\pgfpathlineto{\pgfqpoint{3.568974in}{2.834687in}}%
\pgfpathquadraticcurveto{\pgfqpoint{3.543974in}{2.834687in}}{\pgfqpoint{3.543974in}{2.809687in}}%
\pgfpathlineto{\pgfqpoint{3.543974in}{2.271773in}}%
\pgfpathquadraticcurveto{\pgfqpoint{3.543974in}{2.246773in}}{\pgfqpoint{3.568974in}{2.246773in}}%
\pgfpathlineto{\pgfqpoint{3.568974in}{2.246773in}}%
\pgfpathclose%
\pgfusepath{stroke,fill}%
\end{pgfscope}%
\begin{pgfscope}%
\pgfsetbuttcap%
\pgfsetroundjoin%
\pgfsetlinewidth{1.204500pt}%
\definecolor{currentstroke}{rgb}{1.000000,0.000000,0.000000}%
\pgfsetstrokecolor{currentstroke}%
\pgfsetstrokeopacity{0.800000}%
\pgfsetdash{{4.440000pt}{1.920000pt}}{0.000000pt}%
\pgfpathmoveto{\pgfqpoint{3.593974in}{2.733467in}}%
\pgfpathlineto{\pgfqpoint{3.718974in}{2.733467in}}%
\pgfpathlineto{\pgfqpoint{3.843974in}{2.733467in}}%
\pgfusepath{stroke}%
\end{pgfscope}%
\begin{pgfscope}%
\definecolor{textcolor}{rgb}{0.000000,0.000000,0.000000}%
\pgfsetstrokecolor{textcolor}%
\pgfsetfillcolor{textcolor}%
\pgftext[x=3.943974in,y=2.689717in,left,base]{\color{textcolor}{\sffamily\fontsize{9.000000}{10.800000}\selectfont\catcode`\^=\active\def^{\ifmmode\sp\else\^{}\fi}\catcode`\%=\active\def%{\%}Theoretical minima}}%
\end{pgfscope}%
\begin{pgfscope}%
\pgfsetbuttcap%
\pgfsetmiterjoin%
\definecolor{currentfill}{rgb}{0.564706,0.792157,0.976471}%
\pgfsetfillcolor{currentfill}%
\pgfsetfillopacity{0.400000}%
\pgfsetlinewidth{1.505625pt}%
\definecolor{currentstroke}{rgb}{0.564706,0.792157,0.976471}%
\pgfsetstrokecolor{currentstroke}%
\pgfsetdash{}{0pt}%
\pgfpathmoveto{\pgfqpoint{3.593974in}{2.506245in}}%
\pgfpathlineto{\pgfqpoint{3.843974in}{2.506245in}}%
\pgfpathlineto{\pgfqpoint{3.843974in}{2.593745in}}%
\pgfpathlineto{\pgfqpoint{3.593974in}{2.593745in}}%
\pgfpathlineto{\pgfqpoint{3.593974in}{2.506245in}}%
\pgfpathclose%
\pgfusepath{stroke,fill}%
\end{pgfscope}%
\begin{pgfscope}%
\definecolor{textcolor}{rgb}{0.000000,0.000000,0.000000}%
\pgfsetstrokecolor{textcolor}%
\pgfsetfillcolor{textcolor}%
\pgftext[x=3.943974in,y=2.506245in,left,base]{\color{textcolor}{\sffamily\fontsize{9.000000}{10.800000}\selectfont\catcode`\^=\active\def^{\ifmmode\sp\else\^{}\fi}\catcode`\%=\active\def%{\%}Picoquic}}%
\end{pgfscope}%
\begin{pgfscope}%
\pgfsetbuttcap%
\pgfsetmiterjoin%
\definecolor{currentfill}{rgb}{0.647059,0.839216,0.654902}%
\pgfsetfillcolor{currentfill}%
\pgfsetfillopacity{0.400000}%
\pgfsetlinewidth{1.505625pt}%
\definecolor{currentstroke}{rgb}{0.647059,0.839216,0.654902}%
\pgfsetstrokecolor{currentstroke}%
\pgfsetdash{}{0pt}%
\pgfpathmoveto{\pgfqpoint{3.593974in}{2.322773in}}%
\pgfpathlineto{\pgfqpoint{3.843974in}{2.322773in}}%
\pgfpathlineto{\pgfqpoint{3.843974in}{2.410273in}}%
\pgfpathlineto{\pgfqpoint{3.593974in}{2.410273in}}%
\pgfpathlineto{\pgfqpoint{3.593974in}{2.322773in}}%
\pgfpathclose%
\pgfusepath{stroke,fill}%
\end{pgfscope}%
\begin{pgfscope}%
\definecolor{textcolor}{rgb}{0.000000,0.000000,0.000000}%
\pgfsetstrokecolor{textcolor}%
\pgfsetfillcolor{textcolor}%
\pgftext[x=3.943974in,y=2.322773in,left,base]{\color{textcolor}{\sffamily\fontsize{9.000000}{10.800000}\selectfont\catcode`\^=\active\def^{\ifmmode\sp\else\^{}\fi}\catcode`\%=\active\def%{\%}CTaps}}%
\end{pgfscope}%
\end{pgfpicture}%
\makeatother%
\endgroup%

    \caption[Candidate racing experiment]{Distribution of handshake times when a fast and a slow endpoint are available, but the slow
    endpoint is attempted first, showing the benefit of candidate racing. ($N = 200$)}
    \label{fig:handshake-kde}
\end{figure}

\subsection{Small File Transfer Experiment}\label{sec:smallfiletransfer}
This experiment can be seen in~\cref{fig:small-file-bar} and it
displays how an application can benefit
from QUIC multistreaming while using a generic CTaps client
compatible with both TCP and QUIC.

The measurement consists of the time taken to transfer a file with 70~packets'
worth of data, showing the lost capacity when performing cloning
over TCP. This lost capacity is also modeled in~\cref{fig:connection-clone-gains} in 
\cref{chap:background}.
In order to show that the performance gap between TCP and QUIC grows with increasing delay,
the experiment sweeps across four delays of:
\qtylist{25;50;75;100}{\milli\second}.

Before the small file transfer was started, a large file of \qty{4.6}{\mega\byte} was
transferred in order to increase the CWND to be greater than the size of the small file.
This ensures that the small file transfer could be performed in a single RTT if the
CWND is reused. The file is relatively large, since it is the same
file used in the migration experiment in \cref{sec:path-migration}.

To verify that the small file transfer can actually be completed in
a single RTT, two conditions must hold: the large file must be sufficient
to grow the CWND to at least 70~packets, and the bandwidth-delay product (BDP) must be large
enough to carry the small file in a single RTT.
The required number of RTTs for the CWND to grow to at least  70~packets
can be found by solving: $10 \cdot 2^{x-1} \geq 70$, which gives
a minimum of 4~RTTs.
Assuming an MTU of 1500 bytes and 1460 bytes of data in each packet after discounting
the \qty{20}{\Byte} IP and \qty{20}{\Byte} TCP headers, then
$\sum_{n=1}^{4} 1460 \cdot 10 \cdot 2^{n-1} = \qty{219}{\kilo\byte}$ of data would
be transferred during the minimum 4~RTTs.
This is well below the \qty{4.6}{\mega\byte} file used in the experiment.

To verify that the BDP is large enough to carry the entire small file,
we can calculate the minimum BDP for 70~packets of \qty{1500}{\Bytes}, which is:
$70 \cdot 1500 \cdot 8 = \qty{840}{\kilo\bit}$.

In our experiment, the smallest delay is \qty{25}{\milli\second},
which gives an RTT of \qty{50}{\milli\second}.
The bandwidth is constant at \qty{50}{\mega\bit\per\second}, which gives a minimum BDP of
$\qty{50}{\mega\bit\per\second} \cdot \qty{50}{\milli\second} = \qty{2.5}{\mega\bit}$.
Since the bandwidth is constant and BDP scales
with RTT, we can conclude that the BDP of the link is large enough
to carry the small file for all tested delay values.

\begin{figure}[H]
  %% Creator: Matplotlib, PGF backend
%%
%% To include the figure in your LaTeX document, write
%%   \input{<filename>.pgf}
%%
%% Make sure the required packages are loaded in your preamble
%%   \usepackage{pgf}
%%
%% Also ensure that all the required font packages are loaded; for instance,
%% the lmodern package is sometimes necessary when using math font.
%%   \usepackage{lmodern}
%%
%% Figures using additional raster images can only be included by \input if
%% they are in the same directory as the main LaTeX file. For loading figures
%% from other directories you can use the `import` package
%%   \usepackage{import}
%%
%% and then include the figures with
%%   \import{<path to file>}{<filename>.pgf}
%%
%% Matplotlib used the following preamble
%%   \def\mathdefault#1{#1}
%%   \everymath=\expandafter{\the\everymath\displaystyle}
%%   \IfFileExists{scrextend.sty}{
%%     \usepackage[fontsize=10.000000pt]{scrextend}
%%   }{
%%     \renewcommand{\normalsize}{\fontsize{10.000000}{12.000000}\selectfont}
%%     \normalsize
%%   }
%%   
%%   \makeatletter\@ifpackageloaded{underscore}{}{\usepackage[strings]{underscore}}\makeatother
%%
\begingroup%
\makeatletter%
\begin{pgfpicture}%
\pgfpathrectangle{\pgfpointorigin}{\pgfqpoint{5.369125in}{3.280176in}}%
\pgfusepath{use as bounding box, clip}%
\begin{pgfscope}%
\pgfsetbuttcap%
\pgfsetmiterjoin%
\definecolor{currentfill}{rgb}{1.000000,1.000000,1.000000}%
\pgfsetfillcolor{currentfill}%
\pgfsetlinewidth{0.000000pt}%
\definecolor{currentstroke}{rgb}{1.000000,1.000000,1.000000}%
\pgfsetstrokecolor{currentstroke}%
\pgfsetdash{}{0pt}%
\pgfpathmoveto{\pgfqpoint{-0.000000in}{0.000000in}}%
\pgfpathlineto{\pgfqpoint{5.369125in}{0.000000in}}%
\pgfpathlineto{\pgfqpoint{5.369125in}{3.280176in}}%
\pgfpathlineto{\pgfqpoint{-0.000000in}{3.280176in}}%
\pgfpathlineto{\pgfqpoint{-0.000000in}{0.000000in}}%
\pgfpathclose%
\pgfusepath{fill}%
\end{pgfscope}%
\begin{pgfscope}%
\pgfsetbuttcap%
\pgfsetmiterjoin%
\definecolor{currentfill}{rgb}{1.000000,1.000000,1.000000}%
\pgfsetfillcolor{currentfill}%
\pgfsetlinewidth{0.000000pt}%
\definecolor{currentstroke}{rgb}{0.000000,0.000000,0.000000}%
\pgfsetstrokecolor{currentstroke}%
\pgfsetstrokeopacity{0.000000}%
\pgfsetdash{}{0pt}%
\pgfpathmoveto{\pgfqpoint{0.613195in}{0.528318in}}%
\pgfpathlineto{\pgfqpoint{5.269125in}{0.528318in}}%
\pgfpathlineto{\pgfqpoint{5.269125in}{2.991230in}}%
\pgfpathlineto{\pgfqpoint{0.613195in}{2.991230in}}%
\pgfpathlineto{\pgfqpoint{0.613195in}{0.528318in}}%
\pgfpathclose%
\pgfusepath{fill}%
\end{pgfscope}%
\begin{pgfscope}%
\pgfpathrectangle{\pgfqpoint{0.613195in}{0.528318in}}{\pgfqpoint{4.655930in}{2.462912in}}%
\pgfusepath{clip}%
\pgfsetbuttcap%
\pgfsetmiterjoin%
\definecolor{currentfill}{rgb}{0.082353,0.396078,0.752941}%
\pgfsetfillcolor{currentfill}%
\pgfsetlinewidth{0.000000pt}%
\definecolor{currentstroke}{rgb}{0.000000,0.000000,0.000000}%
\pgfsetstrokecolor{currentstroke}%
\pgfsetstrokeopacity{0.000000}%
\pgfsetdash{}{0pt}%
\pgfpathmoveto{\pgfqpoint{0.824828in}{0.528318in}}%
\pgfpathlineto{\pgfqpoint{1.026384in}{0.528318in}}%
\pgfpathlineto{\pgfqpoint{1.026384in}{1.148496in}}%
\pgfpathlineto{\pgfqpoint{0.824828in}{1.148496in}}%
\pgfpathlineto{\pgfqpoint{0.824828in}{0.528318in}}%
\pgfpathclose%
\pgfusepath{fill}%
\end{pgfscope}%
\begin{pgfscope}%
\pgfpathrectangle{\pgfqpoint{0.613195in}{0.528318in}}{\pgfqpoint{4.655930in}{2.462912in}}%
\pgfusepath{clip}%
\pgfsetbuttcap%
\pgfsetmiterjoin%
\definecolor{currentfill}{rgb}{0.082353,0.396078,0.752941}%
\pgfsetfillcolor{currentfill}%
\pgfsetlinewidth{0.000000pt}%
\definecolor{currentstroke}{rgb}{0.000000,0.000000,0.000000}%
\pgfsetstrokecolor{currentstroke}%
\pgfsetstrokeopacity{0.000000}%
\pgfsetdash{}{0pt}%
\pgfpathmoveto{\pgfqpoint{1.944581in}{0.528318in}}%
\pgfpathlineto{\pgfqpoint{2.146136in}{0.528318in}}%
\pgfpathlineto{\pgfqpoint{2.146136in}{1.720348in}}%
\pgfpathlineto{\pgfqpoint{1.944581in}{1.720348in}}%
\pgfpathlineto{\pgfqpoint{1.944581in}{0.528318in}}%
\pgfpathclose%
\pgfusepath{fill}%
\end{pgfscope}%
\begin{pgfscope}%
\pgfpathrectangle{\pgfqpoint{0.613195in}{0.528318in}}{\pgfqpoint{4.655930in}{2.462912in}}%
\pgfusepath{clip}%
\pgfsetbuttcap%
\pgfsetmiterjoin%
\definecolor{currentfill}{rgb}{0.082353,0.396078,0.752941}%
\pgfsetfillcolor{currentfill}%
\pgfsetlinewidth{0.000000pt}%
\definecolor{currentstroke}{rgb}{0.000000,0.000000,0.000000}%
\pgfsetstrokecolor{currentstroke}%
\pgfsetstrokeopacity{0.000000}%
\pgfsetdash{}{0pt}%
\pgfpathmoveto{\pgfqpoint{3.064333in}{0.528318in}}%
\pgfpathlineto{\pgfqpoint{3.265888in}{0.528318in}}%
\pgfpathlineto{\pgfqpoint{3.265888in}{2.300178in}}%
\pgfpathlineto{\pgfqpoint{3.064333in}{2.300178in}}%
\pgfpathlineto{\pgfqpoint{3.064333in}{0.528318in}}%
\pgfpathclose%
\pgfusepath{fill}%
\end{pgfscope}%
\begin{pgfscope}%
\pgfpathrectangle{\pgfqpoint{0.613195in}{0.528318in}}{\pgfqpoint{4.655930in}{2.462912in}}%
\pgfusepath{clip}%
\pgfsetbuttcap%
\pgfsetmiterjoin%
\definecolor{currentfill}{rgb}{0.082353,0.396078,0.752941}%
\pgfsetfillcolor{currentfill}%
\pgfsetlinewidth{0.000000pt}%
\definecolor{currentstroke}{rgb}{0.000000,0.000000,0.000000}%
\pgfsetstrokecolor{currentstroke}%
\pgfsetstrokeopacity{0.000000}%
\pgfsetdash{}{0pt}%
\pgfpathmoveto{\pgfqpoint{4.184085in}{0.528318in}}%
\pgfpathlineto{\pgfqpoint{4.385641in}{0.528318in}}%
\pgfpathlineto{\pgfqpoint{4.385641in}{2.873948in}}%
\pgfpathlineto{\pgfqpoint{4.184085in}{2.873948in}}%
\pgfpathlineto{\pgfqpoint{4.184085in}{0.528318in}}%
\pgfpathclose%
\pgfusepath{fill}%
\end{pgfscope}%
\begin{pgfscope}%
\pgfpathrectangle{\pgfqpoint{0.613195in}{0.528318in}}{\pgfqpoint{4.655930in}{2.462912in}}%
\pgfusepath{clip}%
\pgfsetbuttcap%
\pgfsetmiterjoin%
\definecolor{currentfill}{rgb}{0.564706,0.792157,0.976471}%
\pgfsetfillcolor{currentfill}%
\pgfsetlinewidth{0.000000pt}%
\definecolor{currentstroke}{rgb}{0.000000,0.000000,0.000000}%
\pgfsetstrokecolor{currentstroke}%
\pgfsetstrokeopacity{0.000000}%
\pgfsetdash{}{0pt}%
\pgfpathmoveto{\pgfqpoint{1.048779in}{0.528318in}}%
\pgfpathlineto{\pgfqpoint{1.250334in}{0.528318in}}%
\pgfpathlineto{\pgfqpoint{1.250334in}{1.147345in}}%
\pgfpathlineto{\pgfqpoint{1.048779in}{1.147345in}}%
\pgfpathlineto{\pgfqpoint{1.048779in}{0.528318in}}%
\pgfpathclose%
\pgfusepath{fill}%
\end{pgfscope}%
\begin{pgfscope}%
\pgfpathrectangle{\pgfqpoint{0.613195in}{0.528318in}}{\pgfqpoint{4.655930in}{2.462912in}}%
\pgfusepath{clip}%
\pgfsetbuttcap%
\pgfsetmiterjoin%
\definecolor{currentfill}{rgb}{0.564706,0.792157,0.976471}%
\pgfsetfillcolor{currentfill}%
\pgfsetlinewidth{0.000000pt}%
\definecolor{currentstroke}{rgb}{0.000000,0.000000,0.000000}%
\pgfsetstrokecolor{currentstroke}%
\pgfsetstrokeopacity{0.000000}%
\pgfsetdash{}{0pt}%
\pgfpathmoveto{\pgfqpoint{2.168531in}{0.528318in}}%
\pgfpathlineto{\pgfqpoint{2.370086in}{0.528318in}}%
\pgfpathlineto{\pgfqpoint{2.370086in}{1.720386in}}%
\pgfpathlineto{\pgfqpoint{2.168531in}{1.720386in}}%
\pgfpathlineto{\pgfqpoint{2.168531in}{0.528318in}}%
\pgfpathclose%
\pgfusepath{fill}%
\end{pgfscope}%
\begin{pgfscope}%
\pgfpathrectangle{\pgfqpoint{0.613195in}{0.528318in}}{\pgfqpoint{4.655930in}{2.462912in}}%
\pgfusepath{clip}%
\pgfsetbuttcap%
\pgfsetmiterjoin%
\definecolor{currentfill}{rgb}{0.564706,0.792157,0.976471}%
\pgfsetfillcolor{currentfill}%
\pgfsetlinewidth{0.000000pt}%
\definecolor{currentstroke}{rgb}{0.000000,0.000000,0.000000}%
\pgfsetstrokecolor{currentstroke}%
\pgfsetstrokeopacity{0.000000}%
\pgfsetdash{}{0pt}%
\pgfpathmoveto{\pgfqpoint{3.288283in}{0.528318in}}%
\pgfpathlineto{\pgfqpoint{3.489839in}{0.528318in}}%
\pgfpathlineto{\pgfqpoint{3.489839in}{2.298836in}}%
\pgfpathlineto{\pgfqpoint{3.288283in}{2.298836in}}%
\pgfpathlineto{\pgfqpoint{3.288283in}{0.528318in}}%
\pgfpathclose%
\pgfusepath{fill}%
\end{pgfscope}%
\begin{pgfscope}%
\pgfpathrectangle{\pgfqpoint{0.613195in}{0.528318in}}{\pgfqpoint{4.655930in}{2.462912in}}%
\pgfusepath{clip}%
\pgfsetbuttcap%
\pgfsetmiterjoin%
\definecolor{currentfill}{rgb}{0.564706,0.792157,0.976471}%
\pgfsetfillcolor{currentfill}%
\pgfsetlinewidth{0.000000pt}%
\definecolor{currentstroke}{rgb}{0.000000,0.000000,0.000000}%
\pgfsetstrokecolor{currentstroke}%
\pgfsetstrokeopacity{0.000000}%
\pgfsetdash{}{0pt}%
\pgfpathmoveto{\pgfqpoint{4.408036in}{0.528318in}}%
\pgfpathlineto{\pgfqpoint{4.609591in}{0.528318in}}%
\pgfpathlineto{\pgfqpoint{4.609591in}{2.872644in}}%
\pgfpathlineto{\pgfqpoint{4.408036in}{2.872644in}}%
\pgfpathlineto{\pgfqpoint{4.408036in}{0.528318in}}%
\pgfpathclose%
\pgfusepath{fill}%
\end{pgfscope}%
\begin{pgfscope}%
\pgfpathrectangle{\pgfqpoint{0.613195in}{0.528318in}}{\pgfqpoint{4.655930in}{2.462912in}}%
\pgfusepath{clip}%
\pgfsetbuttcap%
\pgfsetmiterjoin%
\definecolor{currentfill}{rgb}{0.180392,0.490196,0.196078}%
\pgfsetfillcolor{currentfill}%
\pgfsetlinewidth{0.000000pt}%
\definecolor{currentstroke}{rgb}{0.000000,0.000000,0.000000}%
\pgfsetstrokecolor{currentstroke}%
\pgfsetstrokeopacity{0.000000}%
\pgfsetdash{}{0pt}%
\pgfpathmoveto{\pgfqpoint{1.272729in}{0.528318in}}%
\pgfpathlineto{\pgfqpoint{1.474285in}{0.528318in}}%
\pgfpathlineto{\pgfqpoint{1.474285in}{0.789659in}}%
\pgfpathlineto{\pgfqpoint{1.272729in}{0.789659in}}%
\pgfpathlineto{\pgfqpoint{1.272729in}{0.528318in}}%
\pgfpathclose%
\pgfusepath{fill}%
\end{pgfscope}%
\begin{pgfscope}%
\pgfpathrectangle{\pgfqpoint{0.613195in}{0.528318in}}{\pgfqpoint{4.655930in}{2.462912in}}%
\pgfusepath{clip}%
\pgfsetbuttcap%
\pgfsetmiterjoin%
\definecolor{currentfill}{rgb}{0.180392,0.490196,0.196078}%
\pgfsetfillcolor{currentfill}%
\pgfsetlinewidth{0.000000pt}%
\definecolor{currentstroke}{rgb}{0.000000,0.000000,0.000000}%
\pgfsetstrokecolor{currentstroke}%
\pgfsetstrokeopacity{0.000000}%
\pgfsetdash{}{0pt}%
\pgfpathmoveto{\pgfqpoint{2.392481in}{0.528318in}}%
\pgfpathlineto{\pgfqpoint{2.594037in}{0.528318in}}%
\pgfpathlineto{\pgfqpoint{2.594037in}{0.981389in}}%
\pgfpathlineto{\pgfqpoint{2.392481in}{0.981389in}}%
\pgfpathlineto{\pgfqpoint{2.392481in}{0.528318in}}%
\pgfpathclose%
\pgfusepath{fill}%
\end{pgfscope}%
\begin{pgfscope}%
\pgfpathrectangle{\pgfqpoint{0.613195in}{0.528318in}}{\pgfqpoint{4.655930in}{2.462912in}}%
\pgfusepath{clip}%
\pgfsetbuttcap%
\pgfsetmiterjoin%
\definecolor{currentfill}{rgb}{0.180392,0.490196,0.196078}%
\pgfsetfillcolor{currentfill}%
\pgfsetlinewidth{0.000000pt}%
\definecolor{currentstroke}{rgb}{0.000000,0.000000,0.000000}%
\pgfsetstrokecolor{currentstroke}%
\pgfsetstrokeopacity{0.000000}%
\pgfsetdash{}{0pt}%
\pgfpathmoveto{\pgfqpoint{3.512234in}{0.528318in}}%
\pgfpathlineto{\pgfqpoint{3.713789in}{0.528318in}}%
\pgfpathlineto{\pgfqpoint{3.713789in}{1.174116in}}%
\pgfpathlineto{\pgfqpoint{3.512234in}{1.174116in}}%
\pgfpathlineto{\pgfqpoint{3.512234in}{0.528318in}}%
\pgfpathclose%
\pgfusepath{fill}%
\end{pgfscope}%
\begin{pgfscope}%
\pgfpathrectangle{\pgfqpoint{0.613195in}{0.528318in}}{\pgfqpoint{4.655930in}{2.462912in}}%
\pgfusepath{clip}%
\pgfsetbuttcap%
\pgfsetmiterjoin%
\definecolor{currentfill}{rgb}{0.180392,0.490196,0.196078}%
\pgfsetfillcolor{currentfill}%
\pgfsetlinewidth{0.000000pt}%
\definecolor{currentstroke}{rgb}{0.000000,0.000000,0.000000}%
\pgfsetstrokecolor{currentstroke}%
\pgfsetstrokeopacity{0.000000}%
\pgfsetdash{}{0pt}%
\pgfpathmoveto{\pgfqpoint{4.631986in}{0.528318in}}%
\pgfpathlineto{\pgfqpoint{4.833542in}{0.528318in}}%
\pgfpathlineto{\pgfqpoint{4.833542in}{1.366191in}}%
\pgfpathlineto{\pgfqpoint{4.631986in}{1.366191in}}%
\pgfpathlineto{\pgfqpoint{4.631986in}{0.528318in}}%
\pgfpathclose%
\pgfusepath{fill}%
\end{pgfscope}%
\begin{pgfscope}%
\pgfpathrectangle{\pgfqpoint{0.613195in}{0.528318in}}{\pgfqpoint{4.655930in}{2.462912in}}%
\pgfusepath{clip}%
\pgfsetbuttcap%
\pgfsetmiterjoin%
\definecolor{currentfill}{rgb}{0.647059,0.839216,0.654902}%
\pgfsetfillcolor{currentfill}%
\pgfsetlinewidth{0.000000pt}%
\definecolor{currentstroke}{rgb}{0.000000,0.000000,0.000000}%
\pgfsetstrokecolor{currentstroke}%
\pgfsetstrokeopacity{0.000000}%
\pgfsetdash{}{0pt}%
\pgfpathmoveto{\pgfqpoint{1.496680in}{0.528318in}}%
\pgfpathlineto{\pgfqpoint{1.698235in}{0.528318in}}%
\pgfpathlineto{\pgfqpoint{1.698235in}{0.789161in}}%
\pgfpathlineto{\pgfqpoint{1.496680in}{0.789161in}}%
\pgfpathlineto{\pgfqpoint{1.496680in}{0.528318in}}%
\pgfpathclose%
\pgfusepath{fill}%
\end{pgfscope}%
\begin{pgfscope}%
\pgfpathrectangle{\pgfqpoint{0.613195in}{0.528318in}}{\pgfqpoint{4.655930in}{2.462912in}}%
\pgfusepath{clip}%
\pgfsetbuttcap%
\pgfsetmiterjoin%
\definecolor{currentfill}{rgb}{0.647059,0.839216,0.654902}%
\pgfsetfillcolor{currentfill}%
\pgfsetlinewidth{0.000000pt}%
\definecolor{currentstroke}{rgb}{0.000000,0.000000,0.000000}%
\pgfsetstrokecolor{currentstroke}%
\pgfsetstrokeopacity{0.000000}%
\pgfsetdash{}{0pt}%
\pgfpathmoveto{\pgfqpoint{2.616432in}{0.528318in}}%
\pgfpathlineto{\pgfqpoint{2.817987in}{0.528318in}}%
\pgfpathlineto{\pgfqpoint{2.817987in}{0.982923in}}%
\pgfpathlineto{\pgfqpoint{2.616432in}{0.982923in}}%
\pgfpathlineto{\pgfqpoint{2.616432in}{0.528318in}}%
\pgfpathclose%
\pgfusepath{fill}%
\end{pgfscope}%
\begin{pgfscope}%
\pgfpathrectangle{\pgfqpoint{0.613195in}{0.528318in}}{\pgfqpoint{4.655930in}{2.462912in}}%
\pgfusepath{clip}%
\pgfsetbuttcap%
\pgfsetmiterjoin%
\definecolor{currentfill}{rgb}{0.647059,0.839216,0.654902}%
\pgfsetfillcolor{currentfill}%
\pgfsetlinewidth{0.000000pt}%
\definecolor{currentstroke}{rgb}{0.000000,0.000000,0.000000}%
\pgfsetstrokecolor{currentstroke}%
\pgfsetstrokeopacity{0.000000}%
\pgfsetdash{}{0pt}%
\pgfpathmoveto{\pgfqpoint{3.736184in}{0.528318in}}%
\pgfpathlineto{\pgfqpoint{3.937740in}{0.528318in}}%
\pgfpathlineto{\pgfqpoint{3.937740in}{1.174193in}}%
\pgfpathlineto{\pgfqpoint{3.736184in}{1.174193in}}%
\pgfpathlineto{\pgfqpoint{3.736184in}{0.528318in}}%
\pgfpathclose%
\pgfusepath{fill}%
\end{pgfscope}%
\begin{pgfscope}%
\pgfpathrectangle{\pgfqpoint{0.613195in}{0.528318in}}{\pgfqpoint{4.655930in}{2.462912in}}%
\pgfusepath{clip}%
\pgfsetbuttcap%
\pgfsetmiterjoin%
\definecolor{currentfill}{rgb}{0.647059,0.839216,0.654902}%
\pgfsetfillcolor{currentfill}%
\pgfsetlinewidth{0.000000pt}%
\definecolor{currentstroke}{rgb}{0.000000,0.000000,0.000000}%
\pgfsetstrokecolor{currentstroke}%
\pgfsetstrokeopacity{0.000000}%
\pgfsetdash{}{0pt}%
\pgfpathmoveto{\pgfqpoint{4.855937in}{0.528318in}}%
\pgfpathlineto{\pgfqpoint{5.057492in}{0.528318in}}%
\pgfpathlineto{\pgfqpoint{5.057492in}{1.362317in}}%
\pgfpathlineto{\pgfqpoint{4.855937in}{1.362317in}}%
\pgfpathlineto{\pgfqpoint{4.855937in}{0.528318in}}%
\pgfpathclose%
\pgfusepath{fill}%
\end{pgfscope}%
\begin{pgfscope}%
\pgfpathrectangle{\pgfqpoint{0.613195in}{0.528318in}}{\pgfqpoint{4.655930in}{2.462912in}}%
\pgfusepath{clip}%
\pgfsetbuttcap%
\pgfsetroundjoin%
\pgfsetlinewidth{0.803000pt}%
\definecolor{currentstroke}{rgb}{0.690196,0.690196,0.690196}%
\pgfsetstrokecolor{currentstroke}%
\pgfsetstrokeopacity{0.400000}%
\pgfsetdash{{2.960000pt}{1.280000pt}}{0.000000pt}%
\pgfpathmoveto{\pgfqpoint{1.261532in}{0.528318in}}%
\pgfpathlineto{\pgfqpoint{1.261532in}{2.991230in}}%
\pgfusepath{stroke}%
\end{pgfscope}%
\begin{pgfscope}%
\pgfsetbuttcap%
\pgfsetroundjoin%
\definecolor{currentfill}{rgb}{0.000000,0.000000,0.000000}%
\pgfsetfillcolor{currentfill}%
\pgfsetlinewidth{0.803000pt}%
\definecolor{currentstroke}{rgb}{0.000000,0.000000,0.000000}%
\pgfsetstrokecolor{currentstroke}%
\pgfsetdash{}{0pt}%
\pgfsys@defobject{currentmarker}{\pgfqpoint{0.000000in}{-0.048611in}}{\pgfqpoint{0.000000in}{0.000000in}}{%
\pgfpathmoveto{\pgfqpoint{0.000000in}{0.000000in}}%
\pgfpathlineto{\pgfqpoint{0.000000in}{-0.048611in}}%
\pgfusepath{stroke,fill}%
}%
\begin{pgfscope}%
\pgfsys@transformshift{1.261532in}{0.528318in}%
\pgfsys@useobject{currentmarker}{}%
\end{pgfscope}%
\end{pgfscope}%
\begin{pgfscope}%
\definecolor{textcolor}{rgb}{0.000000,0.000000,0.000000}%
\pgfsetstrokecolor{textcolor}%
\pgfsetfillcolor{textcolor}%
\pgftext[x=1.261532in,y=0.431095in,,top]{\color{textcolor}{\rmfamily\fontsize{10.000000}{12.000000}\selectfont\catcode`\^=\active\def^{\ifmmode\sp\else\^{}\fi}\catcode`\%=\active\def%{\%}50}}%
\end{pgfscope}%
\begin{pgfscope}%
\pgfpathrectangle{\pgfqpoint{0.613195in}{0.528318in}}{\pgfqpoint{4.655930in}{2.462912in}}%
\pgfusepath{clip}%
\pgfsetbuttcap%
\pgfsetroundjoin%
\pgfsetlinewidth{0.803000pt}%
\definecolor{currentstroke}{rgb}{0.690196,0.690196,0.690196}%
\pgfsetstrokecolor{currentstroke}%
\pgfsetstrokeopacity{0.400000}%
\pgfsetdash{{2.960000pt}{1.280000pt}}{0.000000pt}%
\pgfpathmoveto{\pgfqpoint{2.381284in}{0.528318in}}%
\pgfpathlineto{\pgfqpoint{2.381284in}{2.991230in}}%
\pgfusepath{stroke}%
\end{pgfscope}%
\begin{pgfscope}%
\pgfsetbuttcap%
\pgfsetroundjoin%
\definecolor{currentfill}{rgb}{0.000000,0.000000,0.000000}%
\pgfsetfillcolor{currentfill}%
\pgfsetlinewidth{0.803000pt}%
\definecolor{currentstroke}{rgb}{0.000000,0.000000,0.000000}%
\pgfsetstrokecolor{currentstroke}%
\pgfsetdash{}{0pt}%
\pgfsys@defobject{currentmarker}{\pgfqpoint{0.000000in}{-0.048611in}}{\pgfqpoint{0.000000in}{0.000000in}}{%
\pgfpathmoveto{\pgfqpoint{0.000000in}{0.000000in}}%
\pgfpathlineto{\pgfqpoint{0.000000in}{-0.048611in}}%
\pgfusepath{stroke,fill}%
}%
\begin{pgfscope}%
\pgfsys@transformshift{2.381284in}{0.528318in}%
\pgfsys@useobject{currentmarker}{}%
\end{pgfscope}%
\end{pgfscope}%
\begin{pgfscope}%
\definecolor{textcolor}{rgb}{0.000000,0.000000,0.000000}%
\pgfsetstrokecolor{textcolor}%
\pgfsetfillcolor{textcolor}%
\pgftext[x=2.381284in,y=0.431095in,,top]{\color{textcolor}{\rmfamily\fontsize{10.000000}{12.000000}\selectfont\catcode`\^=\active\def^{\ifmmode\sp\else\^{}\fi}\catcode`\%=\active\def%{\%}100}}%
\end{pgfscope}%
\begin{pgfscope}%
\pgfpathrectangle{\pgfqpoint{0.613195in}{0.528318in}}{\pgfqpoint{4.655930in}{2.462912in}}%
\pgfusepath{clip}%
\pgfsetbuttcap%
\pgfsetroundjoin%
\pgfsetlinewidth{0.803000pt}%
\definecolor{currentstroke}{rgb}{0.690196,0.690196,0.690196}%
\pgfsetstrokecolor{currentstroke}%
\pgfsetstrokeopacity{0.400000}%
\pgfsetdash{{2.960000pt}{1.280000pt}}{0.000000pt}%
\pgfpathmoveto{\pgfqpoint{3.501036in}{0.528318in}}%
\pgfpathlineto{\pgfqpoint{3.501036in}{2.991230in}}%
\pgfusepath{stroke}%
\end{pgfscope}%
\begin{pgfscope}%
\pgfsetbuttcap%
\pgfsetroundjoin%
\definecolor{currentfill}{rgb}{0.000000,0.000000,0.000000}%
\pgfsetfillcolor{currentfill}%
\pgfsetlinewidth{0.803000pt}%
\definecolor{currentstroke}{rgb}{0.000000,0.000000,0.000000}%
\pgfsetstrokecolor{currentstroke}%
\pgfsetdash{}{0pt}%
\pgfsys@defobject{currentmarker}{\pgfqpoint{0.000000in}{-0.048611in}}{\pgfqpoint{0.000000in}{0.000000in}}{%
\pgfpathmoveto{\pgfqpoint{0.000000in}{0.000000in}}%
\pgfpathlineto{\pgfqpoint{0.000000in}{-0.048611in}}%
\pgfusepath{stroke,fill}%
}%
\begin{pgfscope}%
\pgfsys@transformshift{3.501036in}{0.528318in}%
\pgfsys@useobject{currentmarker}{}%
\end{pgfscope}%
\end{pgfscope}%
\begin{pgfscope}%
\definecolor{textcolor}{rgb}{0.000000,0.000000,0.000000}%
\pgfsetstrokecolor{textcolor}%
\pgfsetfillcolor{textcolor}%
\pgftext[x=3.501036in,y=0.431095in,,top]{\color{textcolor}{\rmfamily\fontsize{10.000000}{12.000000}\selectfont\catcode`\^=\active\def^{\ifmmode\sp\else\^{}\fi}\catcode`\%=\active\def%{\%}150}}%
\end{pgfscope}%
\begin{pgfscope}%
\pgfpathrectangle{\pgfqpoint{0.613195in}{0.528318in}}{\pgfqpoint{4.655930in}{2.462912in}}%
\pgfusepath{clip}%
\pgfsetbuttcap%
\pgfsetroundjoin%
\pgfsetlinewidth{0.803000pt}%
\definecolor{currentstroke}{rgb}{0.690196,0.690196,0.690196}%
\pgfsetstrokecolor{currentstroke}%
\pgfsetstrokeopacity{0.400000}%
\pgfsetdash{{2.960000pt}{1.280000pt}}{0.000000pt}%
\pgfpathmoveto{\pgfqpoint{4.620789in}{0.528318in}}%
\pgfpathlineto{\pgfqpoint{4.620789in}{2.991230in}}%
\pgfusepath{stroke}%
\end{pgfscope}%
\begin{pgfscope}%
\pgfsetbuttcap%
\pgfsetroundjoin%
\definecolor{currentfill}{rgb}{0.000000,0.000000,0.000000}%
\pgfsetfillcolor{currentfill}%
\pgfsetlinewidth{0.803000pt}%
\definecolor{currentstroke}{rgb}{0.000000,0.000000,0.000000}%
\pgfsetstrokecolor{currentstroke}%
\pgfsetdash{}{0pt}%
\pgfsys@defobject{currentmarker}{\pgfqpoint{0.000000in}{-0.048611in}}{\pgfqpoint{0.000000in}{0.000000in}}{%
\pgfpathmoveto{\pgfqpoint{0.000000in}{0.000000in}}%
\pgfpathlineto{\pgfqpoint{0.000000in}{-0.048611in}}%
\pgfusepath{stroke,fill}%
}%
\begin{pgfscope}%
\pgfsys@transformshift{4.620789in}{0.528318in}%
\pgfsys@useobject{currentmarker}{}%
\end{pgfscope}%
\end{pgfscope}%
\begin{pgfscope}%
\definecolor{textcolor}{rgb}{0.000000,0.000000,0.000000}%
\pgfsetstrokecolor{textcolor}%
\pgfsetfillcolor{textcolor}%
\pgftext[x=4.620789in,y=0.431095in,,top]{\color{textcolor}{\rmfamily\fontsize{10.000000}{12.000000}\selectfont\catcode`\^=\active\def^{\ifmmode\sp\else\^{}\fi}\catcode`\%=\active\def%{\%}200}}%
\end{pgfscope}%
\begin{pgfscope}%
\definecolor{textcolor}{rgb}{0.000000,0.000000,0.000000}%
\pgfsetstrokecolor{textcolor}%
\pgfsetfillcolor{textcolor}%
\pgftext[x=2.941160in,y=0.252083in,,top]{\color{textcolor}{\rmfamily\fontsize{11.000000}{13.200000}\selectfont\catcode`\^=\active\def^{\ifmmode\sp\else\^{}\fi}\catcode`\%=\active\def%{\%}RTT (ms)}}%
\end{pgfscope}%
\begin{pgfscope}%
\pgfpathrectangle{\pgfqpoint{0.613195in}{0.528318in}}{\pgfqpoint{4.655930in}{2.462912in}}%
\pgfusepath{clip}%
\pgfsetbuttcap%
\pgfsetroundjoin%
\pgfsetlinewidth{0.803000pt}%
\definecolor{currentstroke}{rgb}{0.690196,0.690196,0.690196}%
\pgfsetstrokecolor{currentstroke}%
\pgfsetstrokeopacity{0.400000}%
\pgfsetdash{{2.960000pt}{1.280000pt}}{0.000000pt}%
\pgfpathmoveto{\pgfqpoint{0.613195in}{0.528318in}}%
\pgfpathlineto{\pgfqpoint{5.269125in}{0.528318in}}%
\pgfusepath{stroke}%
\end{pgfscope}%
\begin{pgfscope}%
\pgfsetbuttcap%
\pgfsetroundjoin%
\definecolor{currentfill}{rgb}{0.000000,0.000000,0.000000}%
\pgfsetfillcolor{currentfill}%
\pgfsetlinewidth{0.803000pt}%
\definecolor{currentstroke}{rgb}{0.000000,0.000000,0.000000}%
\pgfsetstrokecolor{currentstroke}%
\pgfsetdash{}{0pt}%
\pgfsys@defobject{currentmarker}{\pgfqpoint{-0.048611in}{0.000000in}}{\pgfqpoint{-0.000000in}{0.000000in}}{%
\pgfpathmoveto{\pgfqpoint{-0.000000in}{0.000000in}}%
\pgfpathlineto{\pgfqpoint{-0.048611in}{0.000000in}}%
\pgfusepath{stroke,fill}%
}%
\begin{pgfscope}%
\pgfsys@transformshift{0.613195in}{0.528318in}%
\pgfsys@useobject{currentmarker}{}%
\end{pgfscope}%
\end{pgfscope}%
\begin{pgfscope}%
\definecolor{textcolor}{rgb}{0.000000,0.000000,0.000000}%
\pgfsetstrokecolor{textcolor}%
\pgfsetfillcolor{textcolor}%
\pgftext[x=0.446528in, y=0.480092in, left, base]{\color{textcolor}{\rmfamily\fontsize{10.000000}{12.000000}\selectfont\catcode`\^=\active\def^{\ifmmode\sp\else\^{}\fi}\catcode`\%=\active\def%{\%}$\mathdefault{0}$}}%
\end{pgfscope}%
\begin{pgfscope}%
\pgfpathrectangle{\pgfqpoint{0.613195in}{0.528318in}}{\pgfqpoint{4.655930in}{2.462912in}}%
\pgfusepath{clip}%
\pgfsetbuttcap%
\pgfsetroundjoin%
\pgfsetlinewidth{0.803000pt}%
\definecolor{currentstroke}{rgb}{0.690196,0.690196,0.690196}%
\pgfsetstrokecolor{currentstroke}%
\pgfsetstrokeopacity{0.400000}%
\pgfsetdash{{2.960000pt}{1.280000pt}}{0.000000pt}%
\pgfpathmoveto{\pgfqpoint{0.613195in}{0.911854in}}%
\pgfpathlineto{\pgfqpoint{5.269125in}{0.911854in}}%
\pgfusepath{stroke}%
\end{pgfscope}%
\begin{pgfscope}%
\pgfsetbuttcap%
\pgfsetroundjoin%
\definecolor{currentfill}{rgb}{0.000000,0.000000,0.000000}%
\pgfsetfillcolor{currentfill}%
\pgfsetlinewidth{0.803000pt}%
\definecolor{currentstroke}{rgb}{0.000000,0.000000,0.000000}%
\pgfsetstrokecolor{currentstroke}%
\pgfsetdash{}{0pt}%
\pgfsys@defobject{currentmarker}{\pgfqpoint{-0.048611in}{0.000000in}}{\pgfqpoint{-0.000000in}{0.000000in}}{%
\pgfpathmoveto{\pgfqpoint{-0.000000in}{0.000000in}}%
\pgfpathlineto{\pgfqpoint{-0.048611in}{0.000000in}}%
\pgfusepath{stroke,fill}%
}%
\begin{pgfscope}%
\pgfsys@transformshift{0.613195in}{0.911854in}%
\pgfsys@useobject{currentmarker}{}%
\end{pgfscope}%
\end{pgfscope}%
\begin{pgfscope}%
\definecolor{textcolor}{rgb}{0.000000,0.000000,0.000000}%
\pgfsetstrokecolor{textcolor}%
\pgfsetfillcolor{textcolor}%
\pgftext[x=0.307639in, y=0.863629in, left, base]{\color{textcolor}{\rmfamily\fontsize{10.000000}{12.000000}\selectfont\catcode`\^=\active\def^{\ifmmode\sp\else\^{}\fi}\catcode`\%=\active\def%{\%}$\mathdefault{100}$}}%
\end{pgfscope}%
\begin{pgfscope}%
\pgfpathrectangle{\pgfqpoint{0.613195in}{0.528318in}}{\pgfqpoint{4.655930in}{2.462912in}}%
\pgfusepath{clip}%
\pgfsetbuttcap%
\pgfsetroundjoin%
\pgfsetlinewidth{0.803000pt}%
\definecolor{currentstroke}{rgb}{0.690196,0.690196,0.690196}%
\pgfsetstrokecolor{currentstroke}%
\pgfsetstrokeopacity{0.400000}%
\pgfsetdash{{2.960000pt}{1.280000pt}}{0.000000pt}%
\pgfpathmoveto{\pgfqpoint{0.613195in}{1.295390in}}%
\pgfpathlineto{\pgfqpoint{5.269125in}{1.295390in}}%
\pgfusepath{stroke}%
\end{pgfscope}%
\begin{pgfscope}%
\pgfsetbuttcap%
\pgfsetroundjoin%
\definecolor{currentfill}{rgb}{0.000000,0.000000,0.000000}%
\pgfsetfillcolor{currentfill}%
\pgfsetlinewidth{0.803000pt}%
\definecolor{currentstroke}{rgb}{0.000000,0.000000,0.000000}%
\pgfsetstrokecolor{currentstroke}%
\pgfsetdash{}{0pt}%
\pgfsys@defobject{currentmarker}{\pgfqpoint{-0.048611in}{0.000000in}}{\pgfqpoint{-0.000000in}{0.000000in}}{%
\pgfpathmoveto{\pgfqpoint{-0.000000in}{0.000000in}}%
\pgfpathlineto{\pgfqpoint{-0.048611in}{0.000000in}}%
\pgfusepath{stroke,fill}%
}%
\begin{pgfscope}%
\pgfsys@transformshift{0.613195in}{1.295390in}%
\pgfsys@useobject{currentmarker}{}%
\end{pgfscope}%
\end{pgfscope}%
\begin{pgfscope}%
\definecolor{textcolor}{rgb}{0.000000,0.000000,0.000000}%
\pgfsetstrokecolor{textcolor}%
\pgfsetfillcolor{textcolor}%
\pgftext[x=0.307639in, y=1.247165in, left, base]{\color{textcolor}{\rmfamily\fontsize{10.000000}{12.000000}\selectfont\catcode`\^=\active\def^{\ifmmode\sp\else\^{}\fi}\catcode`\%=\active\def%{\%}$\mathdefault{200}$}}%
\end{pgfscope}%
\begin{pgfscope}%
\pgfpathrectangle{\pgfqpoint{0.613195in}{0.528318in}}{\pgfqpoint{4.655930in}{2.462912in}}%
\pgfusepath{clip}%
\pgfsetbuttcap%
\pgfsetroundjoin%
\pgfsetlinewidth{0.803000pt}%
\definecolor{currentstroke}{rgb}{0.690196,0.690196,0.690196}%
\pgfsetstrokecolor{currentstroke}%
\pgfsetstrokeopacity{0.400000}%
\pgfsetdash{{2.960000pt}{1.280000pt}}{0.000000pt}%
\pgfpathmoveto{\pgfqpoint{0.613195in}{1.678926in}}%
\pgfpathlineto{\pgfqpoint{5.269125in}{1.678926in}}%
\pgfusepath{stroke}%
\end{pgfscope}%
\begin{pgfscope}%
\pgfsetbuttcap%
\pgfsetroundjoin%
\definecolor{currentfill}{rgb}{0.000000,0.000000,0.000000}%
\pgfsetfillcolor{currentfill}%
\pgfsetlinewidth{0.803000pt}%
\definecolor{currentstroke}{rgb}{0.000000,0.000000,0.000000}%
\pgfsetstrokecolor{currentstroke}%
\pgfsetdash{}{0pt}%
\pgfsys@defobject{currentmarker}{\pgfqpoint{-0.048611in}{0.000000in}}{\pgfqpoint{-0.000000in}{0.000000in}}{%
\pgfpathmoveto{\pgfqpoint{-0.000000in}{0.000000in}}%
\pgfpathlineto{\pgfqpoint{-0.048611in}{0.000000in}}%
\pgfusepath{stroke,fill}%
}%
\begin{pgfscope}%
\pgfsys@transformshift{0.613195in}{1.678926in}%
\pgfsys@useobject{currentmarker}{}%
\end{pgfscope}%
\end{pgfscope}%
\begin{pgfscope}%
\definecolor{textcolor}{rgb}{0.000000,0.000000,0.000000}%
\pgfsetstrokecolor{textcolor}%
\pgfsetfillcolor{textcolor}%
\pgftext[x=0.307639in, y=1.630701in, left, base]{\color{textcolor}{\rmfamily\fontsize{10.000000}{12.000000}\selectfont\catcode`\^=\active\def^{\ifmmode\sp\else\^{}\fi}\catcode`\%=\active\def%{\%}$\mathdefault{300}$}}%
\end{pgfscope}%
\begin{pgfscope}%
\pgfpathrectangle{\pgfqpoint{0.613195in}{0.528318in}}{\pgfqpoint{4.655930in}{2.462912in}}%
\pgfusepath{clip}%
\pgfsetbuttcap%
\pgfsetroundjoin%
\pgfsetlinewidth{0.803000pt}%
\definecolor{currentstroke}{rgb}{0.690196,0.690196,0.690196}%
\pgfsetstrokecolor{currentstroke}%
\pgfsetstrokeopacity{0.400000}%
\pgfsetdash{{2.960000pt}{1.280000pt}}{0.000000pt}%
\pgfpathmoveto{\pgfqpoint{0.613195in}{2.062462in}}%
\pgfpathlineto{\pgfqpoint{5.269125in}{2.062462in}}%
\pgfusepath{stroke}%
\end{pgfscope}%
\begin{pgfscope}%
\pgfsetbuttcap%
\pgfsetroundjoin%
\definecolor{currentfill}{rgb}{0.000000,0.000000,0.000000}%
\pgfsetfillcolor{currentfill}%
\pgfsetlinewidth{0.803000pt}%
\definecolor{currentstroke}{rgb}{0.000000,0.000000,0.000000}%
\pgfsetstrokecolor{currentstroke}%
\pgfsetdash{}{0pt}%
\pgfsys@defobject{currentmarker}{\pgfqpoint{-0.048611in}{0.000000in}}{\pgfqpoint{-0.000000in}{0.000000in}}{%
\pgfpathmoveto{\pgfqpoint{-0.000000in}{0.000000in}}%
\pgfpathlineto{\pgfqpoint{-0.048611in}{0.000000in}}%
\pgfusepath{stroke,fill}%
}%
\begin{pgfscope}%
\pgfsys@transformshift{0.613195in}{2.062462in}%
\pgfsys@useobject{currentmarker}{}%
\end{pgfscope}%
\end{pgfscope}%
\begin{pgfscope}%
\definecolor{textcolor}{rgb}{0.000000,0.000000,0.000000}%
\pgfsetstrokecolor{textcolor}%
\pgfsetfillcolor{textcolor}%
\pgftext[x=0.307639in, y=2.014237in, left, base]{\color{textcolor}{\rmfamily\fontsize{10.000000}{12.000000}\selectfont\catcode`\^=\active\def^{\ifmmode\sp\else\^{}\fi}\catcode`\%=\active\def%{\%}$\mathdefault{400}$}}%
\end{pgfscope}%
\begin{pgfscope}%
\pgfpathrectangle{\pgfqpoint{0.613195in}{0.528318in}}{\pgfqpoint{4.655930in}{2.462912in}}%
\pgfusepath{clip}%
\pgfsetbuttcap%
\pgfsetroundjoin%
\pgfsetlinewidth{0.803000pt}%
\definecolor{currentstroke}{rgb}{0.690196,0.690196,0.690196}%
\pgfsetstrokecolor{currentstroke}%
\pgfsetstrokeopacity{0.400000}%
\pgfsetdash{{2.960000pt}{1.280000pt}}{0.000000pt}%
\pgfpathmoveto{\pgfqpoint{0.613195in}{2.445998in}}%
\pgfpathlineto{\pgfqpoint{5.269125in}{2.445998in}}%
\pgfusepath{stroke}%
\end{pgfscope}%
\begin{pgfscope}%
\pgfsetbuttcap%
\pgfsetroundjoin%
\definecolor{currentfill}{rgb}{0.000000,0.000000,0.000000}%
\pgfsetfillcolor{currentfill}%
\pgfsetlinewidth{0.803000pt}%
\definecolor{currentstroke}{rgb}{0.000000,0.000000,0.000000}%
\pgfsetstrokecolor{currentstroke}%
\pgfsetdash{}{0pt}%
\pgfsys@defobject{currentmarker}{\pgfqpoint{-0.048611in}{0.000000in}}{\pgfqpoint{-0.000000in}{0.000000in}}{%
\pgfpathmoveto{\pgfqpoint{-0.000000in}{0.000000in}}%
\pgfpathlineto{\pgfqpoint{-0.048611in}{0.000000in}}%
\pgfusepath{stroke,fill}%
}%
\begin{pgfscope}%
\pgfsys@transformshift{0.613195in}{2.445998in}%
\pgfsys@useobject{currentmarker}{}%
\end{pgfscope}%
\end{pgfscope}%
\begin{pgfscope}%
\definecolor{textcolor}{rgb}{0.000000,0.000000,0.000000}%
\pgfsetstrokecolor{textcolor}%
\pgfsetfillcolor{textcolor}%
\pgftext[x=0.307639in, y=2.397773in, left, base]{\color{textcolor}{\rmfamily\fontsize{10.000000}{12.000000}\selectfont\catcode`\^=\active\def^{\ifmmode\sp\else\^{}\fi}\catcode`\%=\active\def%{\%}$\mathdefault{500}$}}%
\end{pgfscope}%
\begin{pgfscope}%
\pgfpathrectangle{\pgfqpoint{0.613195in}{0.528318in}}{\pgfqpoint{4.655930in}{2.462912in}}%
\pgfusepath{clip}%
\pgfsetbuttcap%
\pgfsetroundjoin%
\pgfsetlinewidth{0.803000pt}%
\definecolor{currentstroke}{rgb}{0.690196,0.690196,0.690196}%
\pgfsetstrokecolor{currentstroke}%
\pgfsetstrokeopacity{0.400000}%
\pgfsetdash{{2.960000pt}{1.280000pt}}{0.000000pt}%
\pgfpathmoveto{\pgfqpoint{0.613195in}{2.829535in}}%
\pgfpathlineto{\pgfqpoint{5.269125in}{2.829535in}}%
\pgfusepath{stroke}%
\end{pgfscope}%
\begin{pgfscope}%
\pgfsetbuttcap%
\pgfsetroundjoin%
\definecolor{currentfill}{rgb}{0.000000,0.000000,0.000000}%
\pgfsetfillcolor{currentfill}%
\pgfsetlinewidth{0.803000pt}%
\definecolor{currentstroke}{rgb}{0.000000,0.000000,0.000000}%
\pgfsetstrokecolor{currentstroke}%
\pgfsetdash{}{0pt}%
\pgfsys@defobject{currentmarker}{\pgfqpoint{-0.048611in}{0.000000in}}{\pgfqpoint{-0.000000in}{0.000000in}}{%
\pgfpathmoveto{\pgfqpoint{-0.000000in}{0.000000in}}%
\pgfpathlineto{\pgfqpoint{-0.048611in}{0.000000in}}%
\pgfusepath{stroke,fill}%
}%
\begin{pgfscope}%
\pgfsys@transformshift{0.613195in}{2.829535in}%
\pgfsys@useobject{currentmarker}{}%
\end{pgfscope}%
\end{pgfscope}%
\begin{pgfscope}%
\definecolor{textcolor}{rgb}{0.000000,0.000000,0.000000}%
\pgfsetstrokecolor{textcolor}%
\pgfsetfillcolor{textcolor}%
\pgftext[x=0.307639in, y=2.781309in, left, base]{\color{textcolor}{\rmfamily\fontsize{10.000000}{12.000000}\selectfont\catcode`\^=\active\def^{\ifmmode\sp\else\^{}\fi}\catcode`\%=\active\def%{\%}$\mathdefault{600}$}}%
\end{pgfscope}%
\begin{pgfscope}%
\definecolor{textcolor}{rgb}{0.000000,0.000000,0.000000}%
\pgfsetstrokecolor{textcolor}%
\pgfsetfillcolor{textcolor}%
\pgftext[x=0.252083in,y=1.759774in,,bottom,rotate=90.000000]{\color{textcolor}{\rmfamily\fontsize{11.000000}{13.200000}\selectfont\catcode`\^=\active\def^{\ifmmode\sp\else\^{}\fi}\catcode`\%=\active\def%{\%}Small file transfer time (ms)}}%
\end{pgfscope}%
\begin{pgfscope}%
\pgfsetrectcap%
\pgfsetmiterjoin%
\pgfsetlinewidth{0.803000pt}%
\definecolor{currentstroke}{rgb}{0.000000,0.000000,0.000000}%
\pgfsetstrokecolor{currentstroke}%
\pgfsetdash{}{0pt}%
\pgfpathmoveto{\pgfqpoint{0.613195in}{0.528318in}}%
\pgfpathlineto{\pgfqpoint{0.613195in}{2.991230in}}%
\pgfusepath{stroke}%
\end{pgfscope}%
\begin{pgfscope}%
\pgfsetrectcap%
\pgfsetmiterjoin%
\pgfsetlinewidth{0.803000pt}%
\definecolor{currentstroke}{rgb}{0.000000,0.000000,0.000000}%
\pgfsetstrokecolor{currentstroke}%
\pgfsetdash{}{0pt}%
\pgfpathmoveto{\pgfqpoint{5.269125in}{0.528318in}}%
\pgfpathlineto{\pgfqpoint{5.269125in}{2.991230in}}%
\pgfusepath{stroke}%
\end{pgfscope}%
\begin{pgfscope}%
\pgfsetrectcap%
\pgfsetmiterjoin%
\pgfsetlinewidth{0.803000pt}%
\definecolor{currentstroke}{rgb}{0.000000,0.000000,0.000000}%
\pgfsetstrokecolor{currentstroke}%
\pgfsetdash{}{0pt}%
\pgfpathmoveto{\pgfqpoint{0.613195in}{0.528318in}}%
\pgfpathlineto{\pgfqpoint{5.269125in}{0.528318in}}%
\pgfusepath{stroke}%
\end{pgfscope}%
\begin{pgfscope}%
\pgfsetrectcap%
\pgfsetmiterjoin%
\pgfsetlinewidth{0.803000pt}%
\definecolor{currentstroke}{rgb}{0.000000,0.000000,0.000000}%
\pgfsetstrokecolor{currentstroke}%
\pgfsetdash{}{0pt}%
\pgfpathmoveto{\pgfqpoint{0.613195in}{2.991230in}}%
\pgfpathlineto{\pgfqpoint{5.269125in}{2.991230in}}%
\pgfusepath{stroke}%
\end{pgfscope}%
\begin{pgfscope}%
\definecolor{textcolor}{rgb}{0.000000,0.000000,0.000000}%
\pgfsetstrokecolor{textcolor}%
\pgfsetfillcolor{textcolor}%
\pgftext[x=2.941160in,y=3.074563in,,base]{\color{textcolor}{\rmfamily\fontsize{11.000000}{13.200000}\selectfont\catcode`\^=\active\def^{\ifmmode\sp\else\^{}\fi}\catcode`\%=\active\def%{\%}Small file transfer time by implementation and RTT}}%
\end{pgfscope}%
\begin{pgfscope}%
\pgfsetbuttcap%
\pgfsetmiterjoin%
\definecolor{currentfill}{rgb}{1.000000,1.000000,1.000000}%
\pgfsetfillcolor{currentfill}%
\pgfsetfillopacity{0.800000}%
\pgfsetlinewidth{1.003750pt}%
\definecolor{currentstroke}{rgb}{0.800000,0.800000,0.800000}%
\pgfsetstrokecolor{currentstroke}%
\pgfsetstrokeopacity{0.800000}%
\pgfsetdash{}{0pt}%
\pgfpathmoveto{\pgfqpoint{0.700695in}{2.516230in}}%
\pgfpathlineto{\pgfqpoint{3.442166in}{2.516230in}}%
\pgfpathquadraticcurveto{\pgfqpoint{3.467166in}{2.516230in}}{\pgfqpoint{3.467166in}{2.541230in}}%
\pgfpathlineto{\pgfqpoint{3.467166in}{2.903730in}}%
\pgfpathquadraticcurveto{\pgfqpoint{3.467166in}{2.928730in}}{\pgfqpoint{3.442166in}{2.928730in}}%
\pgfpathlineto{\pgfqpoint{0.700695in}{2.928730in}}%
\pgfpathquadraticcurveto{\pgfqpoint{0.675695in}{2.928730in}}{\pgfqpoint{0.675695in}{2.903730in}}%
\pgfpathlineto{\pgfqpoint{0.675695in}{2.541230in}}%
\pgfpathquadraticcurveto{\pgfqpoint{0.675695in}{2.516230in}}{\pgfqpoint{0.700695in}{2.516230in}}%
\pgfpathlineto{\pgfqpoint{0.700695in}{2.516230in}}%
\pgfpathclose%
\pgfusepath{stroke,fill}%
\end{pgfscope}%
\begin{pgfscope}%
\pgfsetbuttcap%
\pgfsetmiterjoin%
\definecolor{currentfill}{rgb}{0.082353,0.396078,0.752941}%
\pgfsetfillcolor{currentfill}%
\pgfsetlinewidth{0.000000pt}%
\definecolor{currentstroke}{rgb}{0.000000,0.000000,0.000000}%
\pgfsetstrokecolor{currentstroke}%
\pgfsetstrokeopacity{0.000000}%
\pgfsetdash{}{0pt}%
\pgfpathmoveto{\pgfqpoint{0.725695in}{2.784980in}}%
\pgfpathlineto{\pgfqpoint{0.975695in}{2.784980in}}%
\pgfpathlineto{\pgfqpoint{0.975695in}{2.872480in}}%
\pgfpathlineto{\pgfqpoint{0.725695in}{2.872480in}}%
\pgfpathlineto{\pgfqpoint{0.725695in}{2.784980in}}%
\pgfpathclose%
\pgfusepath{fill}%
\end{pgfscope}%
\begin{pgfscope}%
\definecolor{textcolor}{rgb}{0.000000,0.000000,0.000000}%
\pgfsetstrokecolor{textcolor}%
\pgfsetfillcolor{textcolor}%
\pgftext[x=1.075695in,y=2.784980in,left,base]{\color{textcolor}{\rmfamily\fontsize{9.000000}{10.800000}\selectfont\catcode`\^=\active\def^{\ifmmode\sp\else\^{}\fi}\catcode`\%=\active\def%{\%}TCP (Native)}}%
\end{pgfscope}%
\begin{pgfscope}%
\pgfsetbuttcap%
\pgfsetmiterjoin%
\definecolor{currentfill}{rgb}{0.564706,0.792157,0.976471}%
\pgfsetfillcolor{currentfill}%
\pgfsetlinewidth{0.000000pt}%
\definecolor{currentstroke}{rgb}{0.000000,0.000000,0.000000}%
\pgfsetstrokecolor{currentstroke}%
\pgfsetstrokeopacity{0.000000}%
\pgfsetdash{}{0pt}%
\pgfpathmoveto{\pgfqpoint{0.725695in}{2.597480in}}%
\pgfpathlineto{\pgfqpoint{0.975695in}{2.597480in}}%
\pgfpathlineto{\pgfqpoint{0.975695in}{2.684980in}}%
\pgfpathlineto{\pgfqpoint{0.725695in}{2.684980in}}%
\pgfpathlineto{\pgfqpoint{0.725695in}{2.597480in}}%
\pgfpathclose%
\pgfusepath{fill}%
\end{pgfscope}%
\begin{pgfscope}%
\definecolor{textcolor}{rgb}{0.000000,0.000000,0.000000}%
\pgfsetstrokecolor{textcolor}%
\pgfsetfillcolor{textcolor}%
\pgftext[x=1.075695in,y=2.597480in,left,base]{\color{textcolor}{\rmfamily\fontsize{9.000000}{10.800000}\selectfont\catcode`\^=\active\def^{\ifmmode\sp\else\^{}\fi}\catcode`\%=\active\def%{\%}TCP (CTaps)}}%
\end{pgfscope}%
\begin{pgfscope}%
\pgfsetbuttcap%
\pgfsetmiterjoin%
\definecolor{currentfill}{rgb}{0.180392,0.490196,0.196078}%
\pgfsetfillcolor{currentfill}%
\pgfsetlinewidth{0.000000pt}%
\definecolor{currentstroke}{rgb}{0.000000,0.000000,0.000000}%
\pgfsetstrokecolor{currentstroke}%
\pgfsetstrokeopacity{0.000000}%
\pgfsetdash{}{0pt}%
\pgfpathmoveto{\pgfqpoint{2.109032in}{2.784980in}}%
\pgfpathlineto{\pgfqpoint{2.359032in}{2.784980in}}%
\pgfpathlineto{\pgfqpoint{2.359032in}{2.872480in}}%
\pgfpathlineto{\pgfqpoint{2.109032in}{2.872480in}}%
\pgfpathlineto{\pgfqpoint{2.109032in}{2.784980in}}%
\pgfpathclose%
\pgfusepath{fill}%
\end{pgfscope}%
\begin{pgfscope}%
\definecolor{textcolor}{rgb}{0.000000,0.000000,0.000000}%
\pgfsetstrokecolor{textcolor}%
\pgfsetfillcolor{textcolor}%
\pgftext[x=2.459032in,y=2.784980in,left,base]{\color{textcolor}{\rmfamily\fontsize{9.000000}{10.800000}\selectfont\catcode`\^=\active\def^{\ifmmode\sp\else\^{}\fi}\catcode`\%=\active\def%{\%}QUIC (Picoquic)}}%
\end{pgfscope}%
\begin{pgfscope}%
\pgfsetbuttcap%
\pgfsetmiterjoin%
\definecolor{currentfill}{rgb}{0.647059,0.839216,0.654902}%
\pgfsetfillcolor{currentfill}%
\pgfsetlinewidth{0.000000pt}%
\definecolor{currentstroke}{rgb}{0.000000,0.000000,0.000000}%
\pgfsetstrokecolor{currentstroke}%
\pgfsetstrokeopacity{0.000000}%
\pgfsetdash{}{0pt}%
\pgfpathmoveto{\pgfqpoint{2.109032in}{2.597480in}}%
\pgfpathlineto{\pgfqpoint{2.359032in}{2.597480in}}%
\pgfpathlineto{\pgfqpoint{2.359032in}{2.684980in}}%
\pgfpathlineto{\pgfqpoint{2.109032in}{2.684980in}}%
\pgfpathlineto{\pgfqpoint{2.109032in}{2.597480in}}%
\pgfpathclose%
\pgfusepath{fill}%
\end{pgfscope}%
\begin{pgfscope}%
\definecolor{textcolor}{rgb}{0.000000,0.000000,0.000000}%
\pgfsetstrokecolor{textcolor}%
\pgfsetfillcolor{textcolor}%
\pgftext[x=2.459032in,y=2.597480in,left,base]{\color{textcolor}{\rmfamily\fontsize{9.000000}{10.800000}\selectfont\catcode`\^=\active\def^{\ifmmode\sp\else\^{}\fi}\catcode`\%=\active\def%{\%}QUIC (CTaps)}}%
\end{pgfscope}%
\end{pgfpicture}%
\makeatother%
\endgroup%

  \caption[Small file transfer experiment]{Transfer time of a \qty{100}{\kilo\byte}~file after first transferring a large file to grow the CWND, showing the benefit of multistreaming in QUIC.}
  \label{fig:small-file-bar}
\end{figure}

\subsection{Connection Migration Experiment}\label{sec:path-migration}
This experiment can be seen in~\cref{fig:migration-throughput}, and it
demonstrates the benefits of connection migration. A client connects
to a server from \mono{127.0.0.1} and starts downloading
a \qty{4.6}{\mega\byte} file over an emulated
\qty{5}{\mega\bit\per\second} link with a delay of \qty{25}{\milli\second}.
The transfer reaches its steady-state, but after five seconds the
local socket address is blocked through iptables, as
described in~\cref{sec:softwaretools}. The transfer is performed
three times, with a TCP, CTaps and Picoquic client.

The clients record the size and timestamp of each packet received.
This throughput is gathered into \qty{100}{\milli\second} buckets in the graph.

The CTaps and Picoquic clients both successfully migrate to \mono{127.0.0.2} after the local
socket address is blocked, while TCP times out.
The packet traces show an approximate \qty{75}{\milli\second} gap in received data
following the iptables block for both Picoquic and CTaps. After this gap, the throughput recovers to its prior level.

This experiment was performed with the same CTaps and QUIC clients used in the small
file transfer experiment. The complexity of these clients is described in \cref{sec:client-complexity}.

The file size and bandwidth were chosen to give the transfer enough time to reach
its steady-state, while still having data left to transfer after recovering from
path failure.

\begin{figure}[H]
    %% Creator: Matplotlib, PGF backend
%%
%% To include the figure in your LaTeX document, write
%%   \input{<filename>.pgf}
%%
%% Make sure the required packages are loaded in your preamble
%%   \usepackage{pgf}
%%
%% Also ensure that all the required font packages are loaded; for instance,
%% the lmodern package is sometimes necessary when using math font.
%%   \usepackage{lmodern}
%%
%% Figures using additional raster images can only be included by \input if
%% they are in the same directory as the main LaTeX file. For loading figures
%% from other directories you can use the `import` package
%%   \usepackage{import}
%%
%% and then include the figures with
%%   \import{<path to file>}{<filename>.pgf}
%%
%% Matplotlib used the following preamble
%%   \def\mathdefault#1{#1}
%%   \everymath=\expandafter{\the\everymath\displaystyle}
%%   \IfFileExists{scrextend.sty}{
%%     \usepackage[fontsize=10.000000pt]{scrextend}
%%   }{
%%     \renewcommand{\normalsize}{\fontsize{10.000000}{12.000000}\selectfont}
%%     \normalsize
%%   }
%%   
%%   \ifdefined\pdftexversion\else  % non-pdftex case.
%%     \usepackage{fontspec}
%%     \setmainfont{DejaVuSerif.ttf}[Path=\detokenize{/home/ikhovind/Documents/Skole/taps/venv/lib/python3.11/site-packages/matplotlib/mpl-data/fonts/ttf/}]
%%     \setsansfont{DejaVuSans.ttf}[Path=\detokenize{/home/ikhovind/Documents/Skole/taps/venv/lib/python3.11/site-packages/matplotlib/mpl-data/fonts/ttf/}]
%%     \setmonofont{DejaVuSansMono.ttf}[Path=\detokenize{/home/ikhovind/Documents/Skole/taps/venv/lib/python3.11/site-packages/matplotlib/mpl-data/fonts/ttf/}]
%%   \fi
%%   \makeatletter\@ifpackageloaded{underscore}{}{\usepackage[strings]{underscore}}\makeatother
%%
\begingroup%
\makeatletter%
\begin{pgfpicture}%
\pgfpathrectangle{\pgfpointorigin}{\pgfqpoint{5.366595in}{3.265587in}}%
\pgfusepath{use as bounding box, clip}%
\begin{pgfscope}%
\pgfsetbuttcap%
\pgfsetmiterjoin%
\definecolor{currentfill}{rgb}{1.000000,1.000000,1.000000}%
\pgfsetfillcolor{currentfill}%
\pgfsetlinewidth{0.000000pt}%
\definecolor{currentstroke}{rgb}{1.000000,1.000000,1.000000}%
\pgfsetstrokecolor{currentstroke}%
\pgfsetdash{}{0pt}%
\pgfpathmoveto{\pgfqpoint{-0.000000in}{0.000000in}}%
\pgfpathlineto{\pgfqpoint{5.366595in}{0.000000in}}%
\pgfpathlineto{\pgfqpoint{5.366595in}{3.265587in}}%
\pgfpathlineto{\pgfqpoint{-0.000000in}{3.265587in}}%
\pgfpathlineto{\pgfqpoint{-0.000000in}{0.000000in}}%
\pgfpathclose%
\pgfusepath{fill}%
\end{pgfscope}%
\begin{pgfscope}%
\pgfsetbuttcap%
\pgfsetmiterjoin%
\definecolor{currentfill}{rgb}{1.000000,1.000000,1.000000}%
\pgfsetfillcolor{currentfill}%
\pgfsetlinewidth{0.000000pt}%
\definecolor{currentstroke}{rgb}{0.000000,0.000000,0.000000}%
\pgfsetstrokecolor{currentstroke}%
\pgfsetstrokeopacity{0.000000}%
\pgfsetdash{}{0pt}%
\pgfpathmoveto{\pgfqpoint{0.488997in}{0.535045in}}%
\pgfpathlineto{\pgfqpoint{5.266595in}{0.535045in}}%
\pgfpathlineto{\pgfqpoint{5.266595in}{2.795110in}}%
\pgfpathlineto{\pgfqpoint{0.488997in}{2.795110in}}%
\pgfpathlineto{\pgfqpoint{0.488997in}{0.535045in}}%
\pgfpathclose%
\pgfusepath{fill}%
\end{pgfscope}%
\begin{pgfscope}%
\pgfpathrectangle{\pgfqpoint{0.488997in}{0.535045in}}{\pgfqpoint{4.777597in}{2.260065in}}%
\pgfusepath{clip}%
\pgfsetbuttcap%
\pgfsetroundjoin%
\definecolor{currentfill}{rgb}{0.082353,0.396078,0.752941}%
\pgfsetfillcolor{currentfill}%
\pgfsetfillopacity{0.150000}%
\pgfsetlinewidth{1.003750pt}%
\definecolor{currentstroke}{rgb}{0.082353,0.396078,0.752941}%
\pgfsetstrokecolor{currentstroke}%
\pgfsetstrokeopacity{0.150000}%
\pgfsetdash{}{0pt}%
\pgfsys@defobject{currentmarker}{\pgfqpoint{0.706161in}{0.535045in}}{\pgfqpoint{5.049431in}{1.909526in}}{%
\pgfpathmoveto{\pgfqpoint{0.706161in}{0.535045in}}%
\pgfpathlineto{\pgfqpoint{0.706161in}{0.535045in}}%
\pgfpathlineto{\pgfqpoint{0.758489in}{1.516817in}}%
\pgfpathlineto{\pgfqpoint{0.810818in}{1.647720in}}%
\pgfpathlineto{\pgfqpoint{0.863146in}{1.549543in}}%
\pgfpathlineto{\pgfqpoint{0.915475in}{1.614994in}}%
\pgfpathlineto{\pgfqpoint{0.967803in}{1.582268in}}%
\pgfpathlineto{\pgfqpoint{1.020132in}{1.582268in}}%
\pgfpathlineto{\pgfqpoint{1.072461in}{1.680446in}}%
\pgfpathlineto{\pgfqpoint{1.124789in}{1.745897in}}%
\pgfpathlineto{\pgfqpoint{1.177118in}{1.614994in}}%
\pgfpathlineto{\pgfqpoint{1.229446in}{1.582268in}}%
\pgfpathlineto{\pgfqpoint{1.281775in}{1.844074in}}%
\pgfpathlineto{\pgfqpoint{1.334103in}{1.582268in}}%
\pgfpathlineto{\pgfqpoint{1.386432in}{1.614994in}}%
\pgfpathlineto{\pgfqpoint{1.438760in}{1.844074in}}%
\pgfpathlineto{\pgfqpoint{1.491089in}{1.582268in}}%
\pgfpathlineto{\pgfqpoint{1.543418in}{1.582268in}}%
\pgfpathlineto{\pgfqpoint{1.595746in}{1.876800in}}%
\pgfpathlineto{\pgfqpoint{1.648075in}{1.614994in}}%
\pgfpathlineto{\pgfqpoint{1.700403in}{1.876800in}}%
\pgfpathlineto{\pgfqpoint{1.752732in}{1.582268in}}%
\pgfpathlineto{\pgfqpoint{1.805060in}{1.876800in}}%
\pgfpathlineto{\pgfqpoint{1.857389in}{1.614994in}}%
\pgfpathlineto{\pgfqpoint{1.909717in}{1.778623in}}%
\pgfpathlineto{\pgfqpoint{1.962046in}{1.647720in}}%
\pgfpathlineto{\pgfqpoint{2.014375in}{1.811349in}}%
\pgfpathlineto{\pgfqpoint{2.066703in}{1.680446in}}%
\pgfpathlineto{\pgfqpoint{2.119032in}{1.876800in}}%
\pgfpathlineto{\pgfqpoint{2.171360in}{1.582268in}}%
\pgfpathlineto{\pgfqpoint{2.223689in}{1.844074in}}%
\pgfpathlineto{\pgfqpoint{2.276017in}{1.582268in}}%
\pgfpathlineto{\pgfqpoint{2.328346in}{1.909526in}}%
\pgfpathlineto{\pgfqpoint{2.380674in}{1.582268in}}%
\pgfpathlineto{\pgfqpoint{2.433003in}{1.909526in}}%
\pgfpathlineto{\pgfqpoint{2.485332in}{1.844074in}}%
\pgfpathlineto{\pgfqpoint{2.537660in}{1.582268in}}%
\pgfpathlineto{\pgfqpoint{2.589989in}{1.909526in}}%
\pgfpathlineto{\pgfqpoint{2.642317in}{1.549543in}}%
\pgfpathlineto{\pgfqpoint{2.694646in}{1.876800in}}%
\pgfpathlineto{\pgfqpoint{2.746974in}{1.909526in}}%
\pgfpathlineto{\pgfqpoint{2.799303in}{1.582268in}}%
\pgfpathlineto{\pgfqpoint{2.851632in}{1.876800in}}%
\pgfpathlineto{\pgfqpoint{2.903960in}{1.745897in}}%
\pgfpathlineto{\pgfqpoint{2.956289in}{1.713171in}}%
\pgfpathlineto{\pgfqpoint{3.008617in}{1.876800in}}%
\pgfpathlineto{\pgfqpoint{3.060946in}{1.582268in}}%
\pgfpathlineto{\pgfqpoint{3.113274in}{1.909526in}}%
\pgfpathlineto{\pgfqpoint{3.165603in}{1.876800in}}%
\pgfpathlineto{\pgfqpoint{3.217931in}{1.582268in}}%
\pgfpathlineto{\pgfqpoint{3.270260in}{1.876800in}}%
\pgfpathlineto{\pgfqpoint{3.322589in}{1.876800in}}%
\pgfpathlineto{\pgfqpoint{3.374917in}{1.811349in}}%
\pgfpathlineto{\pgfqpoint{3.427246in}{0.535045in}}%
\pgfpathlineto{\pgfqpoint{3.479574in}{0.535045in}}%
\pgfpathlineto{\pgfqpoint{3.531903in}{0.567770in}}%
\pgfpathlineto{\pgfqpoint{3.584231in}{0.535045in}}%
\pgfpathlineto{\pgfqpoint{3.636560in}{0.535045in}}%
\pgfpathlineto{\pgfqpoint{3.688888in}{0.535045in}}%
\pgfpathlineto{\pgfqpoint{3.741217in}{0.535045in}}%
\pgfpathlineto{\pgfqpoint{3.793546in}{0.535045in}}%
\pgfpathlineto{\pgfqpoint{3.845874in}{0.535045in}}%
\pgfpathlineto{\pgfqpoint{3.898203in}{0.535045in}}%
\pgfpathlineto{\pgfqpoint{3.950531in}{0.535045in}}%
\pgfpathlineto{\pgfqpoint{4.002860in}{0.535045in}}%
\pgfpathlineto{\pgfqpoint{4.055188in}{0.535045in}}%
\pgfpathlineto{\pgfqpoint{4.107517in}{0.535045in}}%
\pgfpathlineto{\pgfqpoint{4.159846in}{0.535045in}}%
\pgfpathlineto{\pgfqpoint{4.212174in}{0.535045in}}%
\pgfpathlineto{\pgfqpoint{4.264503in}{0.535045in}}%
\pgfpathlineto{\pgfqpoint{4.316831in}{0.535045in}}%
\pgfpathlineto{\pgfqpoint{4.369160in}{0.535045in}}%
\pgfpathlineto{\pgfqpoint{4.421488in}{0.535045in}}%
\pgfpathlineto{\pgfqpoint{4.473817in}{0.535045in}}%
\pgfpathlineto{\pgfqpoint{4.526145in}{0.535045in}}%
\pgfpathlineto{\pgfqpoint{4.578474in}{0.535045in}}%
\pgfpathlineto{\pgfqpoint{4.630803in}{0.535045in}}%
\pgfpathlineto{\pgfqpoint{4.683131in}{0.535045in}}%
\pgfpathlineto{\pgfqpoint{4.735460in}{0.535045in}}%
\pgfpathlineto{\pgfqpoint{4.787788in}{0.535045in}}%
\pgfpathlineto{\pgfqpoint{4.840117in}{0.535045in}}%
\pgfpathlineto{\pgfqpoint{4.892445in}{0.535045in}}%
\pgfpathlineto{\pgfqpoint{4.944774in}{0.535045in}}%
\pgfpathlineto{\pgfqpoint{4.997102in}{0.535045in}}%
\pgfpathlineto{\pgfqpoint{5.049431in}{0.535045in}}%
\pgfpathlineto{\pgfqpoint{5.049431in}{0.535045in}}%
\pgfpathlineto{\pgfqpoint{5.049431in}{0.535045in}}%
\pgfpathlineto{\pgfqpoint{4.997102in}{0.535045in}}%
\pgfpathlineto{\pgfqpoint{4.944774in}{0.535045in}}%
\pgfpathlineto{\pgfqpoint{4.892445in}{0.535045in}}%
\pgfpathlineto{\pgfqpoint{4.840117in}{0.535045in}}%
\pgfpathlineto{\pgfqpoint{4.787788in}{0.535045in}}%
\pgfpathlineto{\pgfqpoint{4.735460in}{0.535045in}}%
\pgfpathlineto{\pgfqpoint{4.683131in}{0.535045in}}%
\pgfpathlineto{\pgfqpoint{4.630803in}{0.535045in}}%
\pgfpathlineto{\pgfqpoint{4.578474in}{0.535045in}}%
\pgfpathlineto{\pgfqpoint{4.526145in}{0.535045in}}%
\pgfpathlineto{\pgfqpoint{4.473817in}{0.535045in}}%
\pgfpathlineto{\pgfqpoint{4.421488in}{0.535045in}}%
\pgfpathlineto{\pgfqpoint{4.369160in}{0.535045in}}%
\pgfpathlineto{\pgfqpoint{4.316831in}{0.535045in}}%
\pgfpathlineto{\pgfqpoint{4.264503in}{0.535045in}}%
\pgfpathlineto{\pgfqpoint{4.212174in}{0.535045in}}%
\pgfpathlineto{\pgfqpoint{4.159846in}{0.535045in}}%
\pgfpathlineto{\pgfqpoint{4.107517in}{0.535045in}}%
\pgfpathlineto{\pgfqpoint{4.055188in}{0.535045in}}%
\pgfpathlineto{\pgfqpoint{4.002860in}{0.535045in}}%
\pgfpathlineto{\pgfqpoint{3.950531in}{0.535045in}}%
\pgfpathlineto{\pgfqpoint{3.898203in}{0.535045in}}%
\pgfpathlineto{\pgfqpoint{3.845874in}{0.535045in}}%
\pgfpathlineto{\pgfqpoint{3.793546in}{0.535045in}}%
\pgfpathlineto{\pgfqpoint{3.741217in}{0.535045in}}%
\pgfpathlineto{\pgfqpoint{3.688888in}{0.535045in}}%
\pgfpathlineto{\pgfqpoint{3.636560in}{0.535045in}}%
\pgfpathlineto{\pgfqpoint{3.584231in}{0.535045in}}%
\pgfpathlineto{\pgfqpoint{3.531903in}{0.567770in}}%
\pgfpathlineto{\pgfqpoint{3.479574in}{0.535045in}}%
\pgfpathlineto{\pgfqpoint{3.427246in}{0.535045in}}%
\pgfpathlineto{\pgfqpoint{3.374917in}{1.811349in}}%
\pgfpathlineto{\pgfqpoint{3.322589in}{1.876800in}}%
\pgfpathlineto{\pgfqpoint{3.270260in}{1.876800in}}%
\pgfpathlineto{\pgfqpoint{3.217931in}{1.582268in}}%
\pgfpathlineto{\pgfqpoint{3.165603in}{1.876800in}}%
\pgfpathlineto{\pgfqpoint{3.113274in}{1.909526in}}%
\pgfpathlineto{\pgfqpoint{3.060946in}{1.582268in}}%
\pgfpathlineto{\pgfqpoint{3.008617in}{1.876800in}}%
\pgfpathlineto{\pgfqpoint{2.956289in}{1.713171in}}%
\pgfpathlineto{\pgfqpoint{2.903960in}{1.745897in}}%
\pgfpathlineto{\pgfqpoint{2.851632in}{1.876800in}}%
\pgfpathlineto{\pgfqpoint{2.799303in}{1.582268in}}%
\pgfpathlineto{\pgfqpoint{2.746974in}{1.909526in}}%
\pgfpathlineto{\pgfqpoint{2.694646in}{1.876800in}}%
\pgfpathlineto{\pgfqpoint{2.642317in}{1.549543in}}%
\pgfpathlineto{\pgfqpoint{2.589989in}{1.909526in}}%
\pgfpathlineto{\pgfqpoint{2.537660in}{1.582268in}}%
\pgfpathlineto{\pgfqpoint{2.485332in}{1.844074in}}%
\pgfpathlineto{\pgfqpoint{2.433003in}{1.909526in}}%
\pgfpathlineto{\pgfqpoint{2.380674in}{1.582268in}}%
\pgfpathlineto{\pgfqpoint{2.328346in}{1.909526in}}%
\pgfpathlineto{\pgfqpoint{2.276017in}{1.582268in}}%
\pgfpathlineto{\pgfqpoint{2.223689in}{1.844074in}}%
\pgfpathlineto{\pgfqpoint{2.171360in}{1.582268in}}%
\pgfpathlineto{\pgfqpoint{2.119032in}{1.876800in}}%
\pgfpathlineto{\pgfqpoint{2.066703in}{1.680446in}}%
\pgfpathlineto{\pgfqpoint{2.014375in}{1.811349in}}%
\pgfpathlineto{\pgfqpoint{1.962046in}{1.647720in}}%
\pgfpathlineto{\pgfqpoint{1.909717in}{1.778623in}}%
\pgfpathlineto{\pgfqpoint{1.857389in}{1.614994in}}%
\pgfpathlineto{\pgfqpoint{1.805060in}{1.876800in}}%
\pgfpathlineto{\pgfqpoint{1.752732in}{1.582268in}}%
\pgfpathlineto{\pgfqpoint{1.700403in}{1.876800in}}%
\pgfpathlineto{\pgfqpoint{1.648075in}{1.614994in}}%
\pgfpathlineto{\pgfqpoint{1.595746in}{1.876800in}}%
\pgfpathlineto{\pgfqpoint{1.543418in}{1.582268in}}%
\pgfpathlineto{\pgfqpoint{1.491089in}{1.582268in}}%
\pgfpathlineto{\pgfqpoint{1.438760in}{1.844074in}}%
\pgfpathlineto{\pgfqpoint{1.386432in}{1.614994in}}%
\pgfpathlineto{\pgfqpoint{1.334103in}{1.582268in}}%
\pgfpathlineto{\pgfqpoint{1.281775in}{1.844074in}}%
\pgfpathlineto{\pgfqpoint{1.229446in}{1.582268in}}%
\pgfpathlineto{\pgfqpoint{1.177118in}{1.614994in}}%
\pgfpathlineto{\pgfqpoint{1.124789in}{1.745897in}}%
\pgfpathlineto{\pgfqpoint{1.072461in}{1.680446in}}%
\pgfpathlineto{\pgfqpoint{1.020132in}{1.582268in}}%
\pgfpathlineto{\pgfqpoint{0.967803in}{1.582268in}}%
\pgfpathlineto{\pgfqpoint{0.915475in}{1.614994in}}%
\pgfpathlineto{\pgfqpoint{0.863146in}{1.549543in}}%
\pgfpathlineto{\pgfqpoint{0.810818in}{1.647720in}}%
\pgfpathlineto{\pgfqpoint{0.758489in}{1.516817in}}%
\pgfpathlineto{\pgfqpoint{0.706161in}{0.535045in}}%
\pgfpathlineto{\pgfqpoint{0.706161in}{0.535045in}}%
\pgfpathclose%
\pgfusepath{stroke,fill}%
}%
\begin{pgfscope}%
\pgfsys@transformshift{0.000000in}{0.000000in}%
\pgfsys@useobject{currentmarker}{}%
\end{pgfscope}%
\end{pgfscope}%
\begin{pgfscope}%
\pgfpathrectangle{\pgfqpoint{0.488997in}{0.535045in}}{\pgfqpoint{4.777597in}{2.260065in}}%
\pgfusepath{clip}%
\pgfsetbuttcap%
\pgfsetroundjoin%
\definecolor{currentfill}{rgb}{0.180392,0.490196,0.196078}%
\pgfsetfillcolor{currentfill}%
\pgfsetfillopacity{0.150000}%
\pgfsetlinewidth{1.003750pt}%
\definecolor{currentstroke}{rgb}{0.180392,0.490196,0.196078}%
\pgfsetstrokecolor{currentstroke}%
\pgfsetstrokeopacity{0.150000}%
\pgfsetdash{}{0pt}%
\pgfsys@defobject{currentmarker}{\pgfqpoint{0.706161in}{0.535045in}}{\pgfqpoint{4.944774in}{1.901209in}}{%
\pgfpathmoveto{\pgfqpoint{0.706161in}{0.535045in}}%
\pgfpathlineto{\pgfqpoint{0.706161in}{0.535045in}}%
\pgfpathlineto{\pgfqpoint{0.758489in}{0.662060in}}%
\pgfpathlineto{\pgfqpoint{0.810818in}{1.012416in}}%
\pgfpathlineto{\pgfqpoint{0.863146in}{1.806648in}}%
\pgfpathlineto{\pgfqpoint{0.915475in}{1.838266in}}%
\pgfpathlineto{\pgfqpoint{0.967803in}{1.774171in}}%
\pgfpathlineto{\pgfqpoint{1.020132in}{1.870178in}}%
\pgfpathlineto{\pgfqpoint{1.072461in}{1.870359in}}%
\pgfpathlineto{\pgfqpoint{1.124789in}{1.870607in}}%
\pgfpathlineto{\pgfqpoint{1.177118in}{1.838334in}}%
\pgfpathlineto{\pgfqpoint{1.229446in}{1.870607in}}%
\pgfpathlineto{\pgfqpoint{1.281775in}{1.838289in}}%
\pgfpathlineto{\pgfqpoint{1.334103in}{1.870381in}}%
\pgfpathlineto{\pgfqpoint{1.386432in}{1.838582in}}%
\pgfpathlineto{\pgfqpoint{1.438760in}{1.870381in}}%
\pgfpathlineto{\pgfqpoint{1.491089in}{1.870314in}}%
\pgfpathlineto{\pgfqpoint{1.543418in}{1.838515in}}%
\pgfpathlineto{\pgfqpoint{1.595746in}{1.870381in}}%
\pgfpathlineto{\pgfqpoint{1.648075in}{1.870110in}}%
\pgfpathlineto{\pgfqpoint{1.700403in}{1.838560in}}%
\pgfpathlineto{\pgfqpoint{1.752732in}{1.870359in}}%
\pgfpathlineto{\pgfqpoint{1.805060in}{1.838582in}}%
\pgfpathlineto{\pgfqpoint{1.857389in}{1.870133in}}%
\pgfpathlineto{\pgfqpoint{1.909717in}{1.870607in}}%
\pgfpathlineto{\pgfqpoint{1.962046in}{1.838582in}}%
\pgfpathlineto{\pgfqpoint{2.014375in}{1.870110in}}%
\pgfpathlineto{\pgfqpoint{2.066703in}{1.870607in}}%
\pgfpathlineto{\pgfqpoint{2.119032in}{1.837724in}}%
\pgfpathlineto{\pgfqpoint{2.171360in}{1.869184in}}%
\pgfpathlineto{\pgfqpoint{2.223689in}{1.837656in}}%
\pgfpathlineto{\pgfqpoint{2.276017in}{1.869432in}}%
\pgfpathlineto{\pgfqpoint{2.328346in}{1.869432in}}%
\pgfpathlineto{\pgfqpoint{2.380674in}{1.837656in}}%
\pgfpathlineto{\pgfqpoint{2.433003in}{1.869432in}}%
\pgfpathlineto{\pgfqpoint{2.485332in}{1.837656in}}%
\pgfpathlineto{\pgfqpoint{2.537660in}{1.869410in}}%
\pgfpathlineto{\pgfqpoint{2.589989in}{1.869432in}}%
\pgfpathlineto{\pgfqpoint{2.642317in}{1.837656in}}%
\pgfpathlineto{\pgfqpoint{2.694646in}{1.869432in}}%
\pgfpathlineto{\pgfqpoint{2.746974in}{1.869184in}}%
\pgfpathlineto{\pgfqpoint{2.799303in}{1.837882in}}%
\pgfpathlineto{\pgfqpoint{2.851632in}{1.869184in}}%
\pgfpathlineto{\pgfqpoint{2.903960in}{1.837633in}}%
\pgfpathlineto{\pgfqpoint{2.956289in}{1.869432in}}%
\pgfpathlineto{\pgfqpoint{3.008617in}{1.869410in}}%
\pgfpathlineto{\pgfqpoint{3.060946in}{1.837633in}}%
\pgfpathlineto{\pgfqpoint{3.113274in}{1.869410in}}%
\pgfpathlineto{\pgfqpoint{3.165603in}{1.869432in}}%
\pgfpathlineto{\pgfqpoint{3.217931in}{1.837656in}}%
\pgfpathlineto{\pgfqpoint{3.270260in}{1.869432in}}%
\pgfpathlineto{\pgfqpoint{3.322589in}{1.551667in}}%
\pgfpathlineto{\pgfqpoint{3.374917in}{1.138798in}}%
\pgfpathlineto{\pgfqpoint{3.427246in}{1.901073in}}%
\pgfpathlineto{\pgfqpoint{3.479574in}{1.837588in}}%
\pgfpathlineto{\pgfqpoint{3.531903in}{1.869364in}}%
\pgfpathlineto{\pgfqpoint{3.584231in}{1.869410in}}%
\pgfpathlineto{\pgfqpoint{3.636560in}{1.837633in}}%
\pgfpathlineto{\pgfqpoint{3.688888in}{1.869387in}}%
\pgfpathlineto{\pgfqpoint{3.741217in}{1.869184in}}%
\pgfpathlineto{\pgfqpoint{3.793546in}{1.837407in}}%
\pgfpathlineto{\pgfqpoint{3.845874in}{1.869658in}}%
\pgfpathlineto{\pgfqpoint{3.898203in}{1.837430in}}%
\pgfpathlineto{\pgfqpoint{3.950531in}{1.869432in}}%
\pgfpathlineto{\pgfqpoint{4.002860in}{1.869410in}}%
\pgfpathlineto{\pgfqpoint{4.055188in}{1.837656in}}%
\pgfpathlineto{\pgfqpoint{4.107517in}{1.869410in}}%
\pgfpathlineto{\pgfqpoint{4.159846in}{1.837633in}}%
\pgfpathlineto{\pgfqpoint{4.212174in}{1.869432in}}%
\pgfpathlineto{\pgfqpoint{4.264503in}{1.869432in}}%
\pgfpathlineto{\pgfqpoint{4.316831in}{1.837656in}}%
\pgfpathlineto{\pgfqpoint{4.369160in}{1.869410in}}%
\pgfpathlineto{\pgfqpoint{4.421488in}{1.837000in}}%
\pgfpathlineto{\pgfqpoint{4.473817in}{1.869432in}}%
\pgfpathlineto{\pgfqpoint{4.526145in}{1.869432in}}%
\pgfpathlineto{\pgfqpoint{4.578474in}{1.837656in}}%
\pgfpathlineto{\pgfqpoint{4.630803in}{1.869432in}}%
\pgfpathlineto{\pgfqpoint{4.683131in}{1.901209in}}%
\pgfpathlineto{\pgfqpoint{4.735460in}{1.805879in}}%
\pgfpathlineto{\pgfqpoint{4.787788in}{1.869184in}}%
\pgfpathlineto{\pgfqpoint{4.840117in}{1.869410in}}%
\pgfpathlineto{\pgfqpoint{4.892445in}{1.837882in}}%
\pgfpathlineto{\pgfqpoint{4.944774in}{1.835893in}}%
\pgfpathlineto{\pgfqpoint{4.944774in}{1.835893in}}%
\pgfpathlineto{\pgfqpoint{4.944774in}{1.835893in}}%
\pgfpathlineto{\pgfqpoint{4.892445in}{1.837882in}}%
\pgfpathlineto{\pgfqpoint{4.840117in}{1.869410in}}%
\pgfpathlineto{\pgfqpoint{4.787788in}{1.869184in}}%
\pgfpathlineto{\pgfqpoint{4.735460in}{1.805879in}}%
\pgfpathlineto{\pgfqpoint{4.683131in}{1.901209in}}%
\pgfpathlineto{\pgfqpoint{4.630803in}{1.869432in}}%
\pgfpathlineto{\pgfqpoint{4.578474in}{1.837656in}}%
\pgfpathlineto{\pgfqpoint{4.526145in}{1.869432in}}%
\pgfpathlineto{\pgfqpoint{4.473817in}{1.869432in}}%
\pgfpathlineto{\pgfqpoint{4.421488in}{1.837000in}}%
\pgfpathlineto{\pgfqpoint{4.369160in}{1.869410in}}%
\pgfpathlineto{\pgfqpoint{4.316831in}{1.837656in}}%
\pgfpathlineto{\pgfqpoint{4.264503in}{1.869432in}}%
\pgfpathlineto{\pgfqpoint{4.212174in}{1.869432in}}%
\pgfpathlineto{\pgfqpoint{4.159846in}{1.837633in}}%
\pgfpathlineto{\pgfqpoint{4.107517in}{1.869410in}}%
\pgfpathlineto{\pgfqpoint{4.055188in}{1.837656in}}%
\pgfpathlineto{\pgfqpoint{4.002860in}{1.869410in}}%
\pgfpathlineto{\pgfqpoint{3.950531in}{1.869432in}}%
\pgfpathlineto{\pgfqpoint{3.898203in}{1.837430in}}%
\pgfpathlineto{\pgfqpoint{3.845874in}{1.869658in}}%
\pgfpathlineto{\pgfqpoint{3.793546in}{1.837407in}}%
\pgfpathlineto{\pgfqpoint{3.741217in}{1.869184in}}%
\pgfpathlineto{\pgfqpoint{3.688888in}{1.869387in}}%
\pgfpathlineto{\pgfqpoint{3.636560in}{1.837633in}}%
\pgfpathlineto{\pgfqpoint{3.584231in}{1.869410in}}%
\pgfpathlineto{\pgfqpoint{3.531903in}{1.869364in}}%
\pgfpathlineto{\pgfqpoint{3.479574in}{1.837588in}}%
\pgfpathlineto{\pgfqpoint{3.427246in}{1.901073in}}%
\pgfpathlineto{\pgfqpoint{3.374917in}{1.138798in}}%
\pgfpathlineto{\pgfqpoint{3.322589in}{1.551667in}}%
\pgfpathlineto{\pgfqpoint{3.270260in}{1.869432in}}%
\pgfpathlineto{\pgfqpoint{3.217931in}{1.837656in}}%
\pgfpathlineto{\pgfqpoint{3.165603in}{1.869432in}}%
\pgfpathlineto{\pgfqpoint{3.113274in}{1.869410in}}%
\pgfpathlineto{\pgfqpoint{3.060946in}{1.837633in}}%
\pgfpathlineto{\pgfqpoint{3.008617in}{1.869410in}}%
\pgfpathlineto{\pgfqpoint{2.956289in}{1.869432in}}%
\pgfpathlineto{\pgfqpoint{2.903960in}{1.837633in}}%
\pgfpathlineto{\pgfqpoint{2.851632in}{1.869184in}}%
\pgfpathlineto{\pgfqpoint{2.799303in}{1.837882in}}%
\pgfpathlineto{\pgfqpoint{2.746974in}{1.869184in}}%
\pgfpathlineto{\pgfqpoint{2.694646in}{1.869432in}}%
\pgfpathlineto{\pgfqpoint{2.642317in}{1.837656in}}%
\pgfpathlineto{\pgfqpoint{2.589989in}{1.869432in}}%
\pgfpathlineto{\pgfqpoint{2.537660in}{1.869410in}}%
\pgfpathlineto{\pgfqpoint{2.485332in}{1.837656in}}%
\pgfpathlineto{\pgfqpoint{2.433003in}{1.869432in}}%
\pgfpathlineto{\pgfqpoint{2.380674in}{1.837656in}}%
\pgfpathlineto{\pgfqpoint{2.328346in}{1.869432in}}%
\pgfpathlineto{\pgfqpoint{2.276017in}{1.869432in}}%
\pgfpathlineto{\pgfqpoint{2.223689in}{1.837656in}}%
\pgfpathlineto{\pgfqpoint{2.171360in}{1.869184in}}%
\pgfpathlineto{\pgfqpoint{2.119032in}{1.837724in}}%
\pgfpathlineto{\pgfqpoint{2.066703in}{1.870607in}}%
\pgfpathlineto{\pgfqpoint{2.014375in}{1.870110in}}%
\pgfpathlineto{\pgfqpoint{1.962046in}{1.838582in}}%
\pgfpathlineto{\pgfqpoint{1.909717in}{1.870607in}}%
\pgfpathlineto{\pgfqpoint{1.857389in}{1.870133in}}%
\pgfpathlineto{\pgfqpoint{1.805060in}{1.838582in}}%
\pgfpathlineto{\pgfqpoint{1.752732in}{1.870359in}}%
\pgfpathlineto{\pgfqpoint{1.700403in}{1.838560in}}%
\pgfpathlineto{\pgfqpoint{1.648075in}{1.870110in}}%
\pgfpathlineto{\pgfqpoint{1.595746in}{1.870381in}}%
\pgfpathlineto{\pgfqpoint{1.543418in}{1.838515in}}%
\pgfpathlineto{\pgfqpoint{1.491089in}{1.870314in}}%
\pgfpathlineto{\pgfqpoint{1.438760in}{1.870381in}}%
\pgfpathlineto{\pgfqpoint{1.386432in}{1.838582in}}%
\pgfpathlineto{\pgfqpoint{1.334103in}{1.870381in}}%
\pgfpathlineto{\pgfqpoint{1.281775in}{1.838289in}}%
\pgfpathlineto{\pgfqpoint{1.229446in}{1.870607in}}%
\pgfpathlineto{\pgfqpoint{1.177118in}{1.838334in}}%
\pgfpathlineto{\pgfqpoint{1.124789in}{1.870607in}}%
\pgfpathlineto{\pgfqpoint{1.072461in}{1.870359in}}%
\pgfpathlineto{\pgfqpoint{1.020132in}{1.870178in}}%
\pgfpathlineto{\pgfqpoint{0.967803in}{1.774171in}}%
\pgfpathlineto{\pgfqpoint{0.915475in}{1.838266in}}%
\pgfpathlineto{\pgfqpoint{0.863146in}{1.806648in}}%
\pgfpathlineto{\pgfqpoint{0.810818in}{1.012416in}}%
\pgfpathlineto{\pgfqpoint{0.758489in}{0.662060in}}%
\pgfpathlineto{\pgfqpoint{0.706161in}{0.535045in}}%
\pgfpathlineto{\pgfqpoint{0.706161in}{0.535045in}}%
\pgfpathclose%
\pgfusepath{stroke,fill}%
}%
\begin{pgfscope}%
\pgfsys@transformshift{0.000000in}{0.000000in}%
\pgfsys@useobject{currentmarker}{}%
\end{pgfscope}%
\end{pgfscope}%
\begin{pgfscope}%
\pgfpathrectangle{\pgfqpoint{0.488997in}{0.535045in}}{\pgfqpoint{4.777597in}{2.260065in}}%
\pgfusepath{clip}%
\pgfsetbuttcap%
\pgfsetroundjoin%
\definecolor{currentfill}{rgb}{0.647059,0.839216,0.654902}%
\pgfsetfillcolor{currentfill}%
\pgfsetfillopacity{0.150000}%
\pgfsetlinewidth{1.003750pt}%
\definecolor{currentstroke}{rgb}{0.647059,0.839216,0.654902}%
\pgfsetstrokecolor{currentstroke}%
\pgfsetstrokeopacity{0.150000}%
\pgfsetdash{}{0pt}%
\pgfsys@defobject{currentmarker}{\pgfqpoint{0.706161in}{0.535045in}}{\pgfqpoint{4.944774in}{1.901299in}}{%
\pgfpathmoveto{\pgfqpoint{0.706161in}{0.535045in}}%
\pgfpathlineto{\pgfqpoint{0.706161in}{0.535045in}}%
\pgfpathlineto{\pgfqpoint{0.758489in}{0.662422in}}%
\pgfpathlineto{\pgfqpoint{0.810818in}{1.076014in}}%
\pgfpathlineto{\pgfqpoint{0.863146in}{1.806602in}}%
\pgfpathlineto{\pgfqpoint{0.915475in}{1.806467in}}%
\pgfpathlineto{\pgfqpoint{0.967803in}{1.774510in}}%
\pgfpathlineto{\pgfqpoint{1.020132in}{1.870449in}}%
\pgfpathlineto{\pgfqpoint{1.072461in}{1.870472in}}%
\pgfpathlineto{\pgfqpoint{1.124789in}{1.870449in}}%
\pgfpathlineto{\pgfqpoint{1.177118in}{1.838605in}}%
\pgfpathlineto{\pgfqpoint{1.229446in}{1.870449in}}%
\pgfpathlineto{\pgfqpoint{1.281775in}{1.870449in}}%
\pgfpathlineto{\pgfqpoint{1.334103in}{1.838650in}}%
\pgfpathlineto{\pgfqpoint{1.386432in}{1.870449in}}%
\pgfpathlineto{\pgfqpoint{1.438760in}{1.838673in}}%
\pgfpathlineto{\pgfqpoint{1.491089in}{1.870449in}}%
\pgfpathlineto{\pgfqpoint{1.543418in}{1.870449in}}%
\pgfpathlineto{\pgfqpoint{1.595746in}{1.838650in}}%
\pgfpathlineto{\pgfqpoint{1.648075in}{1.870314in}}%
\pgfpathlineto{\pgfqpoint{1.700403in}{1.870449in}}%
\pgfpathlineto{\pgfqpoint{1.752732in}{1.838808in}}%
\pgfpathlineto{\pgfqpoint{1.805060in}{1.870314in}}%
\pgfpathlineto{\pgfqpoint{1.857389in}{1.838673in}}%
\pgfpathlineto{\pgfqpoint{1.909717in}{1.870472in}}%
\pgfpathlineto{\pgfqpoint{1.962046in}{1.870449in}}%
\pgfpathlineto{\pgfqpoint{2.014375in}{1.838673in}}%
\pgfpathlineto{\pgfqpoint{2.066703in}{1.870291in}}%
\pgfpathlineto{\pgfqpoint{2.119032in}{1.869703in}}%
\pgfpathlineto{\pgfqpoint{2.171360in}{1.837746in}}%
\pgfpathlineto{\pgfqpoint{2.223689in}{1.869500in}}%
\pgfpathlineto{\pgfqpoint{2.276017in}{1.837746in}}%
\pgfpathlineto{\pgfqpoint{2.328346in}{1.869500in}}%
\pgfpathlineto{\pgfqpoint{2.380674in}{1.869364in}}%
\pgfpathlineto{\pgfqpoint{2.433003in}{1.837746in}}%
\pgfpathlineto{\pgfqpoint{2.485332in}{1.869658in}}%
\pgfpathlineto{\pgfqpoint{2.537660in}{1.837588in}}%
\pgfpathlineto{\pgfqpoint{2.589989in}{1.869500in}}%
\pgfpathlineto{\pgfqpoint{2.642317in}{1.901299in}}%
\pgfpathlineto{\pgfqpoint{2.694646in}{1.805947in}}%
\pgfpathlineto{\pgfqpoint{2.746974in}{1.869523in}}%
\pgfpathlineto{\pgfqpoint{2.799303in}{1.901299in}}%
\pgfpathlineto{\pgfqpoint{2.851632in}{1.805970in}}%
\pgfpathlineto{\pgfqpoint{2.903960in}{1.869500in}}%
\pgfpathlineto{\pgfqpoint{2.956289in}{1.837724in}}%
\pgfpathlineto{\pgfqpoint{3.008617in}{1.869500in}}%
\pgfpathlineto{\pgfqpoint{3.060946in}{1.869500in}}%
\pgfpathlineto{\pgfqpoint{3.113274in}{1.837724in}}%
\pgfpathlineto{\pgfqpoint{3.165603in}{1.869523in}}%
\pgfpathlineto{\pgfqpoint{3.217931in}{1.869523in}}%
\pgfpathlineto{\pgfqpoint{3.270260in}{1.837746in}}%
\pgfpathlineto{\pgfqpoint{3.322589in}{1.329299in}}%
\pgfpathlineto{\pgfqpoint{3.374917in}{1.456225in}}%
\pgfpathlineto{\pgfqpoint{3.427246in}{1.869455in}}%
\pgfpathlineto{\pgfqpoint{3.479574in}{1.869455in}}%
\pgfpathlineto{\pgfqpoint{3.531903in}{1.837724in}}%
\pgfpathlineto{\pgfqpoint{3.584231in}{1.869500in}}%
\pgfpathlineto{\pgfqpoint{3.636560in}{1.837678in}}%
\pgfpathlineto{\pgfqpoint{3.688888in}{1.869364in}}%
\pgfpathlineto{\pgfqpoint{3.741217in}{1.869500in}}%
\pgfpathlineto{\pgfqpoint{3.793546in}{1.837724in}}%
\pgfpathlineto{\pgfqpoint{3.845874in}{1.869523in}}%
\pgfpathlineto{\pgfqpoint{3.898203in}{1.837724in}}%
\pgfpathlineto{\pgfqpoint{3.950531in}{1.869500in}}%
\pgfpathlineto{\pgfqpoint{4.002860in}{1.869500in}}%
\pgfpathlineto{\pgfqpoint{4.055188in}{1.837746in}}%
\pgfpathlineto{\pgfqpoint{4.107517in}{1.869500in}}%
\pgfpathlineto{\pgfqpoint{4.159846in}{1.869523in}}%
\pgfpathlineto{\pgfqpoint{4.212174in}{1.837746in}}%
\pgfpathlineto{\pgfqpoint{4.264503in}{1.869500in}}%
\pgfpathlineto{\pgfqpoint{4.316831in}{1.836978in}}%
\pgfpathlineto{\pgfqpoint{4.369160in}{1.869500in}}%
\pgfpathlineto{\pgfqpoint{4.421488in}{1.869523in}}%
\pgfpathlineto{\pgfqpoint{4.473817in}{1.837724in}}%
\pgfpathlineto{\pgfqpoint{4.526145in}{1.869342in}}%
\pgfpathlineto{\pgfqpoint{4.578474in}{1.869658in}}%
\pgfpathlineto{\pgfqpoint{4.630803in}{1.837724in}}%
\pgfpathlineto{\pgfqpoint{4.683131in}{1.869523in}}%
\pgfpathlineto{\pgfqpoint{4.735460in}{1.837724in}}%
\pgfpathlineto{\pgfqpoint{4.787788in}{1.869364in}}%
\pgfpathlineto{\pgfqpoint{4.840117in}{1.869500in}}%
\pgfpathlineto{\pgfqpoint{4.892445in}{1.837746in}}%
\pgfpathlineto{\pgfqpoint{4.944774in}{1.733037in}}%
\pgfpathlineto{\pgfqpoint{4.944774in}{1.733037in}}%
\pgfpathlineto{\pgfqpoint{4.944774in}{1.733037in}}%
\pgfpathlineto{\pgfqpoint{4.892445in}{1.837746in}}%
\pgfpathlineto{\pgfqpoint{4.840117in}{1.869500in}}%
\pgfpathlineto{\pgfqpoint{4.787788in}{1.869364in}}%
\pgfpathlineto{\pgfqpoint{4.735460in}{1.837724in}}%
\pgfpathlineto{\pgfqpoint{4.683131in}{1.869523in}}%
\pgfpathlineto{\pgfqpoint{4.630803in}{1.837724in}}%
\pgfpathlineto{\pgfqpoint{4.578474in}{1.869658in}}%
\pgfpathlineto{\pgfqpoint{4.526145in}{1.869342in}}%
\pgfpathlineto{\pgfqpoint{4.473817in}{1.837724in}}%
\pgfpathlineto{\pgfqpoint{4.421488in}{1.869523in}}%
\pgfpathlineto{\pgfqpoint{4.369160in}{1.869500in}}%
\pgfpathlineto{\pgfqpoint{4.316831in}{1.836978in}}%
\pgfpathlineto{\pgfqpoint{4.264503in}{1.869500in}}%
\pgfpathlineto{\pgfqpoint{4.212174in}{1.837746in}}%
\pgfpathlineto{\pgfqpoint{4.159846in}{1.869523in}}%
\pgfpathlineto{\pgfqpoint{4.107517in}{1.869500in}}%
\pgfpathlineto{\pgfqpoint{4.055188in}{1.837746in}}%
\pgfpathlineto{\pgfqpoint{4.002860in}{1.869500in}}%
\pgfpathlineto{\pgfqpoint{3.950531in}{1.869500in}}%
\pgfpathlineto{\pgfqpoint{3.898203in}{1.837724in}}%
\pgfpathlineto{\pgfqpoint{3.845874in}{1.869523in}}%
\pgfpathlineto{\pgfqpoint{3.793546in}{1.837724in}}%
\pgfpathlineto{\pgfqpoint{3.741217in}{1.869500in}}%
\pgfpathlineto{\pgfqpoint{3.688888in}{1.869364in}}%
\pgfpathlineto{\pgfqpoint{3.636560in}{1.837678in}}%
\pgfpathlineto{\pgfqpoint{3.584231in}{1.869500in}}%
\pgfpathlineto{\pgfqpoint{3.531903in}{1.837724in}}%
\pgfpathlineto{\pgfqpoint{3.479574in}{1.869455in}}%
\pgfpathlineto{\pgfqpoint{3.427246in}{1.869455in}}%
\pgfpathlineto{\pgfqpoint{3.374917in}{1.456225in}}%
\pgfpathlineto{\pgfqpoint{3.322589in}{1.329299in}}%
\pgfpathlineto{\pgfqpoint{3.270260in}{1.837746in}}%
\pgfpathlineto{\pgfqpoint{3.217931in}{1.869523in}}%
\pgfpathlineto{\pgfqpoint{3.165603in}{1.869523in}}%
\pgfpathlineto{\pgfqpoint{3.113274in}{1.837724in}}%
\pgfpathlineto{\pgfqpoint{3.060946in}{1.869500in}}%
\pgfpathlineto{\pgfqpoint{3.008617in}{1.869500in}}%
\pgfpathlineto{\pgfqpoint{2.956289in}{1.837724in}}%
\pgfpathlineto{\pgfqpoint{2.903960in}{1.869500in}}%
\pgfpathlineto{\pgfqpoint{2.851632in}{1.805970in}}%
\pgfpathlineto{\pgfqpoint{2.799303in}{1.901299in}}%
\pgfpathlineto{\pgfqpoint{2.746974in}{1.869523in}}%
\pgfpathlineto{\pgfqpoint{2.694646in}{1.805947in}}%
\pgfpathlineto{\pgfqpoint{2.642317in}{1.901299in}}%
\pgfpathlineto{\pgfqpoint{2.589989in}{1.869500in}}%
\pgfpathlineto{\pgfqpoint{2.537660in}{1.837588in}}%
\pgfpathlineto{\pgfqpoint{2.485332in}{1.869658in}}%
\pgfpathlineto{\pgfqpoint{2.433003in}{1.837746in}}%
\pgfpathlineto{\pgfqpoint{2.380674in}{1.869364in}}%
\pgfpathlineto{\pgfqpoint{2.328346in}{1.869500in}}%
\pgfpathlineto{\pgfqpoint{2.276017in}{1.837746in}}%
\pgfpathlineto{\pgfqpoint{2.223689in}{1.869500in}}%
\pgfpathlineto{\pgfqpoint{2.171360in}{1.837746in}}%
\pgfpathlineto{\pgfqpoint{2.119032in}{1.869703in}}%
\pgfpathlineto{\pgfqpoint{2.066703in}{1.870291in}}%
\pgfpathlineto{\pgfqpoint{2.014375in}{1.838673in}}%
\pgfpathlineto{\pgfqpoint{1.962046in}{1.870449in}}%
\pgfpathlineto{\pgfqpoint{1.909717in}{1.870472in}}%
\pgfpathlineto{\pgfqpoint{1.857389in}{1.838673in}}%
\pgfpathlineto{\pgfqpoint{1.805060in}{1.870314in}}%
\pgfpathlineto{\pgfqpoint{1.752732in}{1.838808in}}%
\pgfpathlineto{\pgfqpoint{1.700403in}{1.870449in}}%
\pgfpathlineto{\pgfqpoint{1.648075in}{1.870314in}}%
\pgfpathlineto{\pgfqpoint{1.595746in}{1.838650in}}%
\pgfpathlineto{\pgfqpoint{1.543418in}{1.870449in}}%
\pgfpathlineto{\pgfqpoint{1.491089in}{1.870449in}}%
\pgfpathlineto{\pgfqpoint{1.438760in}{1.838673in}}%
\pgfpathlineto{\pgfqpoint{1.386432in}{1.870449in}}%
\pgfpathlineto{\pgfqpoint{1.334103in}{1.838650in}}%
\pgfpathlineto{\pgfqpoint{1.281775in}{1.870449in}}%
\pgfpathlineto{\pgfqpoint{1.229446in}{1.870449in}}%
\pgfpathlineto{\pgfqpoint{1.177118in}{1.838605in}}%
\pgfpathlineto{\pgfqpoint{1.124789in}{1.870449in}}%
\pgfpathlineto{\pgfqpoint{1.072461in}{1.870472in}}%
\pgfpathlineto{\pgfqpoint{1.020132in}{1.870449in}}%
\pgfpathlineto{\pgfqpoint{0.967803in}{1.774510in}}%
\pgfpathlineto{\pgfqpoint{0.915475in}{1.806467in}}%
\pgfpathlineto{\pgfqpoint{0.863146in}{1.806602in}}%
\pgfpathlineto{\pgfqpoint{0.810818in}{1.076014in}}%
\pgfpathlineto{\pgfqpoint{0.758489in}{0.662422in}}%
\pgfpathlineto{\pgfqpoint{0.706161in}{0.535045in}}%
\pgfpathlineto{\pgfqpoint{0.706161in}{0.535045in}}%
\pgfpathclose%
\pgfusepath{stroke,fill}%
}%
\begin{pgfscope}%
\pgfsys@transformshift{0.000000in}{0.000000in}%
\pgfsys@useobject{currentmarker}{}%
\end{pgfscope}%
\end{pgfscope}%
\begin{pgfscope}%
\pgfpathrectangle{\pgfqpoint{0.488997in}{0.535045in}}{\pgfqpoint{4.777597in}{2.260065in}}%
\pgfusepath{clip}%
\pgfsetbuttcap%
\pgfsetroundjoin%
\pgfsetlinewidth{0.803000pt}%
\definecolor{currentstroke}{rgb}{0.690196,0.690196,0.690196}%
\pgfsetstrokecolor{currentstroke}%
\pgfsetstrokeopacity{0.400000}%
\pgfsetdash{{2.960000pt}{1.280000pt}}{0.000000pt}%
\pgfpathmoveto{\pgfqpoint{0.679996in}{0.535045in}}%
\pgfpathlineto{\pgfqpoint{0.679996in}{2.795110in}}%
\pgfusepath{stroke}%
\end{pgfscope}%
\begin{pgfscope}%
\pgfsetbuttcap%
\pgfsetroundjoin%
\definecolor{currentfill}{rgb}{0.000000,0.000000,0.000000}%
\pgfsetfillcolor{currentfill}%
\pgfsetlinewidth{0.803000pt}%
\definecolor{currentstroke}{rgb}{0.000000,0.000000,0.000000}%
\pgfsetstrokecolor{currentstroke}%
\pgfsetdash{}{0pt}%
\pgfsys@defobject{currentmarker}{\pgfqpoint{0.000000in}{-0.048611in}}{\pgfqpoint{0.000000in}{0.000000in}}{%
\pgfpathmoveto{\pgfqpoint{0.000000in}{0.000000in}}%
\pgfpathlineto{\pgfqpoint{0.000000in}{-0.048611in}}%
\pgfusepath{stroke,fill}%
}%
\begin{pgfscope}%
\pgfsys@transformshift{0.679996in}{0.535045in}%
\pgfsys@useobject{currentmarker}{}%
\end{pgfscope}%
\end{pgfscope}%
\begin{pgfscope}%
\definecolor{textcolor}{rgb}{0.000000,0.000000,0.000000}%
\pgfsetstrokecolor{textcolor}%
\pgfsetfillcolor{textcolor}%
\pgftext[x=0.679996in,y=0.437822in,,top]{\color{textcolor}{\sffamily\fontsize{10.000000}{12.000000}\selectfont\catcode`\^=\active\def^{\ifmmode\sp\else\^{}\fi}\catcode`\%=\active\def%{\%}0}}%
\end{pgfscope}%
\begin{pgfscope}%
\pgfpathrectangle{\pgfqpoint{0.488997in}{0.535045in}}{\pgfqpoint{4.777597in}{2.260065in}}%
\pgfusepath{clip}%
\pgfsetbuttcap%
\pgfsetroundjoin%
\pgfsetlinewidth{0.803000pt}%
\definecolor{currentstroke}{rgb}{0.690196,0.690196,0.690196}%
\pgfsetstrokecolor{currentstroke}%
\pgfsetstrokeopacity{0.400000}%
\pgfsetdash{{2.960000pt}{1.280000pt}}{0.000000pt}%
\pgfpathmoveto{\pgfqpoint{1.726568in}{0.535045in}}%
\pgfpathlineto{\pgfqpoint{1.726568in}{2.795110in}}%
\pgfusepath{stroke}%
\end{pgfscope}%
\begin{pgfscope}%
\pgfsetbuttcap%
\pgfsetroundjoin%
\definecolor{currentfill}{rgb}{0.000000,0.000000,0.000000}%
\pgfsetfillcolor{currentfill}%
\pgfsetlinewidth{0.803000pt}%
\definecolor{currentstroke}{rgb}{0.000000,0.000000,0.000000}%
\pgfsetstrokecolor{currentstroke}%
\pgfsetdash{}{0pt}%
\pgfsys@defobject{currentmarker}{\pgfqpoint{0.000000in}{-0.048611in}}{\pgfqpoint{0.000000in}{0.000000in}}{%
\pgfpathmoveto{\pgfqpoint{0.000000in}{0.000000in}}%
\pgfpathlineto{\pgfqpoint{0.000000in}{-0.048611in}}%
\pgfusepath{stroke,fill}%
}%
\begin{pgfscope}%
\pgfsys@transformshift{1.726568in}{0.535045in}%
\pgfsys@useobject{currentmarker}{}%
\end{pgfscope}%
\end{pgfscope}%
\begin{pgfscope}%
\definecolor{textcolor}{rgb}{0.000000,0.000000,0.000000}%
\pgfsetstrokecolor{textcolor}%
\pgfsetfillcolor{textcolor}%
\pgftext[x=1.726568in,y=0.437822in,,top]{\color{textcolor}{\sffamily\fontsize{10.000000}{12.000000}\selectfont\catcode`\^=\active\def^{\ifmmode\sp\else\^{}\fi}\catcode`\%=\active\def%{\%}2000}}%
\end{pgfscope}%
\begin{pgfscope}%
\pgfpathrectangle{\pgfqpoint{0.488997in}{0.535045in}}{\pgfqpoint{4.777597in}{2.260065in}}%
\pgfusepath{clip}%
\pgfsetbuttcap%
\pgfsetroundjoin%
\pgfsetlinewidth{0.803000pt}%
\definecolor{currentstroke}{rgb}{0.690196,0.690196,0.690196}%
\pgfsetstrokecolor{currentstroke}%
\pgfsetstrokeopacity{0.400000}%
\pgfsetdash{{2.960000pt}{1.280000pt}}{0.000000pt}%
\pgfpathmoveto{\pgfqpoint{2.773139in}{0.535045in}}%
\pgfpathlineto{\pgfqpoint{2.773139in}{2.795110in}}%
\pgfusepath{stroke}%
\end{pgfscope}%
\begin{pgfscope}%
\pgfsetbuttcap%
\pgfsetroundjoin%
\definecolor{currentfill}{rgb}{0.000000,0.000000,0.000000}%
\pgfsetfillcolor{currentfill}%
\pgfsetlinewidth{0.803000pt}%
\definecolor{currentstroke}{rgb}{0.000000,0.000000,0.000000}%
\pgfsetstrokecolor{currentstroke}%
\pgfsetdash{}{0pt}%
\pgfsys@defobject{currentmarker}{\pgfqpoint{0.000000in}{-0.048611in}}{\pgfqpoint{0.000000in}{0.000000in}}{%
\pgfpathmoveto{\pgfqpoint{0.000000in}{0.000000in}}%
\pgfpathlineto{\pgfqpoint{0.000000in}{-0.048611in}}%
\pgfusepath{stroke,fill}%
}%
\begin{pgfscope}%
\pgfsys@transformshift{2.773139in}{0.535045in}%
\pgfsys@useobject{currentmarker}{}%
\end{pgfscope}%
\end{pgfscope}%
\begin{pgfscope}%
\definecolor{textcolor}{rgb}{0.000000,0.000000,0.000000}%
\pgfsetstrokecolor{textcolor}%
\pgfsetfillcolor{textcolor}%
\pgftext[x=2.773139in,y=0.437822in,,top]{\color{textcolor}{\sffamily\fontsize{10.000000}{12.000000}\selectfont\catcode`\^=\active\def^{\ifmmode\sp\else\^{}\fi}\catcode`\%=\active\def%{\%}4000}}%
\end{pgfscope}%
\begin{pgfscope}%
\pgfpathrectangle{\pgfqpoint{0.488997in}{0.535045in}}{\pgfqpoint{4.777597in}{2.260065in}}%
\pgfusepath{clip}%
\pgfsetbuttcap%
\pgfsetroundjoin%
\pgfsetlinewidth{0.803000pt}%
\definecolor{currentstroke}{rgb}{0.690196,0.690196,0.690196}%
\pgfsetstrokecolor{currentstroke}%
\pgfsetstrokeopacity{0.400000}%
\pgfsetdash{{2.960000pt}{1.280000pt}}{0.000000pt}%
\pgfpathmoveto{\pgfqpoint{3.819710in}{0.535045in}}%
\pgfpathlineto{\pgfqpoint{3.819710in}{2.795110in}}%
\pgfusepath{stroke}%
\end{pgfscope}%
\begin{pgfscope}%
\pgfsetbuttcap%
\pgfsetroundjoin%
\definecolor{currentfill}{rgb}{0.000000,0.000000,0.000000}%
\pgfsetfillcolor{currentfill}%
\pgfsetlinewidth{0.803000pt}%
\definecolor{currentstroke}{rgb}{0.000000,0.000000,0.000000}%
\pgfsetstrokecolor{currentstroke}%
\pgfsetdash{}{0pt}%
\pgfsys@defobject{currentmarker}{\pgfqpoint{0.000000in}{-0.048611in}}{\pgfqpoint{0.000000in}{0.000000in}}{%
\pgfpathmoveto{\pgfqpoint{0.000000in}{0.000000in}}%
\pgfpathlineto{\pgfqpoint{0.000000in}{-0.048611in}}%
\pgfusepath{stroke,fill}%
}%
\begin{pgfscope}%
\pgfsys@transformshift{3.819710in}{0.535045in}%
\pgfsys@useobject{currentmarker}{}%
\end{pgfscope}%
\end{pgfscope}%
\begin{pgfscope}%
\definecolor{textcolor}{rgb}{0.000000,0.000000,0.000000}%
\pgfsetstrokecolor{textcolor}%
\pgfsetfillcolor{textcolor}%
\pgftext[x=3.819710in,y=0.437822in,,top]{\color{textcolor}{\sffamily\fontsize{10.000000}{12.000000}\selectfont\catcode`\^=\active\def^{\ifmmode\sp\else\^{}\fi}\catcode`\%=\active\def%{\%}6000}}%
\end{pgfscope}%
\begin{pgfscope}%
\pgfpathrectangle{\pgfqpoint{0.488997in}{0.535045in}}{\pgfqpoint{4.777597in}{2.260065in}}%
\pgfusepath{clip}%
\pgfsetbuttcap%
\pgfsetroundjoin%
\pgfsetlinewidth{0.803000pt}%
\definecolor{currentstroke}{rgb}{0.690196,0.690196,0.690196}%
\pgfsetstrokecolor{currentstroke}%
\pgfsetstrokeopacity{0.400000}%
\pgfsetdash{{2.960000pt}{1.280000pt}}{0.000000pt}%
\pgfpathmoveto{\pgfqpoint{4.866281in}{0.535045in}}%
\pgfpathlineto{\pgfqpoint{4.866281in}{2.795110in}}%
\pgfusepath{stroke}%
\end{pgfscope}%
\begin{pgfscope}%
\pgfsetbuttcap%
\pgfsetroundjoin%
\definecolor{currentfill}{rgb}{0.000000,0.000000,0.000000}%
\pgfsetfillcolor{currentfill}%
\pgfsetlinewidth{0.803000pt}%
\definecolor{currentstroke}{rgb}{0.000000,0.000000,0.000000}%
\pgfsetstrokecolor{currentstroke}%
\pgfsetdash{}{0pt}%
\pgfsys@defobject{currentmarker}{\pgfqpoint{0.000000in}{-0.048611in}}{\pgfqpoint{0.000000in}{0.000000in}}{%
\pgfpathmoveto{\pgfqpoint{0.000000in}{0.000000in}}%
\pgfpathlineto{\pgfqpoint{0.000000in}{-0.048611in}}%
\pgfusepath{stroke,fill}%
}%
\begin{pgfscope}%
\pgfsys@transformshift{4.866281in}{0.535045in}%
\pgfsys@useobject{currentmarker}{}%
\end{pgfscope}%
\end{pgfscope}%
\begin{pgfscope}%
\definecolor{textcolor}{rgb}{0.000000,0.000000,0.000000}%
\pgfsetstrokecolor{textcolor}%
\pgfsetfillcolor{textcolor}%
\pgftext[x=4.866281in,y=0.437822in,,top]{\color{textcolor}{\sffamily\fontsize{10.000000}{12.000000}\selectfont\catcode`\^=\active\def^{\ifmmode\sp\else\^{}\fi}\catcode`\%=\active\def%{\%}8000}}%
\end{pgfscope}%
\begin{pgfscope}%
\definecolor{textcolor}{rgb}{0.000000,0.000000,0.000000}%
\pgfsetstrokecolor{textcolor}%
\pgfsetfillcolor{textcolor}%
\pgftext[x=2.877796in,y=0.247854in,,top]{\color{textcolor}{\sffamily\fontsize{11.000000}{13.200000}\selectfont\catcode`\^=\active\def^{\ifmmode\sp\else\^{}\fi}\catcode`\%=\active\def%{\%}Time since transfer start (ms)}}%
\end{pgfscope}%
\begin{pgfscope}%
\pgfpathrectangle{\pgfqpoint{0.488997in}{0.535045in}}{\pgfqpoint{4.777597in}{2.260065in}}%
\pgfusepath{clip}%
\pgfsetbuttcap%
\pgfsetroundjoin%
\pgfsetlinewidth{0.803000pt}%
\definecolor{currentstroke}{rgb}{0.690196,0.690196,0.690196}%
\pgfsetstrokecolor{currentstroke}%
\pgfsetstrokeopacity{0.400000}%
\pgfsetdash{{2.960000pt}{1.280000pt}}{0.000000pt}%
\pgfpathmoveto{\pgfqpoint{0.488997in}{0.535045in}}%
\pgfpathlineto{\pgfqpoint{5.266595in}{0.535045in}}%
\pgfusepath{stroke}%
\end{pgfscope}%
\begin{pgfscope}%
\pgfsetbuttcap%
\pgfsetroundjoin%
\definecolor{currentfill}{rgb}{0.000000,0.000000,0.000000}%
\pgfsetfillcolor{currentfill}%
\pgfsetlinewidth{0.803000pt}%
\definecolor{currentstroke}{rgb}{0.000000,0.000000,0.000000}%
\pgfsetstrokecolor{currentstroke}%
\pgfsetdash{}{0pt}%
\pgfsys@defobject{currentmarker}{\pgfqpoint{-0.048611in}{0.000000in}}{\pgfqpoint{-0.000000in}{0.000000in}}{%
\pgfpathmoveto{\pgfqpoint{-0.000000in}{0.000000in}}%
\pgfpathlineto{\pgfqpoint{-0.048611in}{0.000000in}}%
\pgfusepath{stroke,fill}%
}%
\begin{pgfscope}%
\pgfsys@transformshift{0.488997in}{0.535045in}%
\pgfsys@useobject{currentmarker}{}%
\end{pgfscope}%
\end{pgfscope}%
\begin{pgfscope}%
\definecolor{textcolor}{rgb}{0.000000,0.000000,0.000000}%
\pgfsetstrokecolor{textcolor}%
\pgfsetfillcolor{textcolor}%
\pgftext[x=0.303410in, y=0.482283in, left, base]{\color{textcolor}{\sffamily\fontsize{10.000000}{12.000000}\selectfont\catcode`\^=\active\def^{\ifmmode\sp\else\^{}\fi}\catcode`\%=\active\def%{\%}0}}%
\end{pgfscope}%
\begin{pgfscope}%
\pgfpathrectangle{\pgfqpoint{0.488997in}{0.535045in}}{\pgfqpoint{4.777597in}{2.260065in}}%
\pgfusepath{clip}%
\pgfsetbuttcap%
\pgfsetroundjoin%
\pgfsetlinewidth{0.803000pt}%
\definecolor{currentstroke}{rgb}{0.690196,0.690196,0.690196}%
\pgfsetstrokecolor{currentstroke}%
\pgfsetstrokeopacity{0.400000}%
\pgfsetdash{{2.960000pt}{1.280000pt}}{0.000000pt}%
\pgfpathmoveto{\pgfqpoint{0.488997in}{0.817553in}}%
\pgfpathlineto{\pgfqpoint{5.266595in}{0.817553in}}%
\pgfusepath{stroke}%
\end{pgfscope}%
\begin{pgfscope}%
\pgfsetbuttcap%
\pgfsetroundjoin%
\definecolor{currentfill}{rgb}{0.000000,0.000000,0.000000}%
\pgfsetfillcolor{currentfill}%
\pgfsetlinewidth{0.803000pt}%
\definecolor{currentstroke}{rgb}{0.000000,0.000000,0.000000}%
\pgfsetstrokecolor{currentstroke}%
\pgfsetdash{}{0pt}%
\pgfsys@defobject{currentmarker}{\pgfqpoint{-0.048611in}{0.000000in}}{\pgfqpoint{-0.000000in}{0.000000in}}{%
\pgfpathmoveto{\pgfqpoint{-0.000000in}{0.000000in}}%
\pgfpathlineto{\pgfqpoint{-0.048611in}{0.000000in}}%
\pgfusepath{stroke,fill}%
}%
\begin{pgfscope}%
\pgfsys@transformshift{0.488997in}{0.817553in}%
\pgfsys@useobject{currentmarker}{}%
\end{pgfscope}%
\end{pgfscope}%
\begin{pgfscope}%
\definecolor{textcolor}{rgb}{0.000000,0.000000,0.000000}%
\pgfsetstrokecolor{textcolor}%
\pgfsetfillcolor{textcolor}%
\pgftext[x=0.303410in, y=0.764791in, left, base]{\color{textcolor}{\sffamily\fontsize{10.000000}{12.000000}\selectfont\catcode`\^=\active\def^{\ifmmode\sp\else\^{}\fi}\catcode`\%=\active\def%{\%}1}}%
\end{pgfscope}%
\begin{pgfscope}%
\pgfpathrectangle{\pgfqpoint{0.488997in}{0.535045in}}{\pgfqpoint{4.777597in}{2.260065in}}%
\pgfusepath{clip}%
\pgfsetbuttcap%
\pgfsetroundjoin%
\pgfsetlinewidth{0.803000pt}%
\definecolor{currentstroke}{rgb}{0.690196,0.690196,0.690196}%
\pgfsetstrokecolor{currentstroke}%
\pgfsetstrokeopacity{0.400000}%
\pgfsetdash{{2.960000pt}{1.280000pt}}{0.000000pt}%
\pgfpathmoveto{\pgfqpoint{0.488997in}{1.100061in}}%
\pgfpathlineto{\pgfqpoint{5.266595in}{1.100061in}}%
\pgfusepath{stroke}%
\end{pgfscope}%
\begin{pgfscope}%
\pgfsetbuttcap%
\pgfsetroundjoin%
\definecolor{currentfill}{rgb}{0.000000,0.000000,0.000000}%
\pgfsetfillcolor{currentfill}%
\pgfsetlinewidth{0.803000pt}%
\definecolor{currentstroke}{rgb}{0.000000,0.000000,0.000000}%
\pgfsetstrokecolor{currentstroke}%
\pgfsetdash{}{0pt}%
\pgfsys@defobject{currentmarker}{\pgfqpoint{-0.048611in}{0.000000in}}{\pgfqpoint{-0.000000in}{0.000000in}}{%
\pgfpathmoveto{\pgfqpoint{-0.000000in}{0.000000in}}%
\pgfpathlineto{\pgfqpoint{-0.048611in}{0.000000in}}%
\pgfusepath{stroke,fill}%
}%
\begin{pgfscope}%
\pgfsys@transformshift{0.488997in}{1.100061in}%
\pgfsys@useobject{currentmarker}{}%
\end{pgfscope}%
\end{pgfscope}%
\begin{pgfscope}%
\definecolor{textcolor}{rgb}{0.000000,0.000000,0.000000}%
\pgfsetstrokecolor{textcolor}%
\pgfsetfillcolor{textcolor}%
\pgftext[x=0.303410in, y=1.047299in, left, base]{\color{textcolor}{\sffamily\fontsize{10.000000}{12.000000}\selectfont\catcode`\^=\active\def^{\ifmmode\sp\else\^{}\fi}\catcode`\%=\active\def%{\%}2}}%
\end{pgfscope}%
\begin{pgfscope}%
\pgfpathrectangle{\pgfqpoint{0.488997in}{0.535045in}}{\pgfqpoint{4.777597in}{2.260065in}}%
\pgfusepath{clip}%
\pgfsetbuttcap%
\pgfsetroundjoin%
\pgfsetlinewidth{0.803000pt}%
\definecolor{currentstroke}{rgb}{0.690196,0.690196,0.690196}%
\pgfsetstrokecolor{currentstroke}%
\pgfsetstrokeopacity{0.400000}%
\pgfsetdash{{2.960000pt}{1.280000pt}}{0.000000pt}%
\pgfpathmoveto{\pgfqpoint{0.488997in}{1.382569in}}%
\pgfpathlineto{\pgfqpoint{5.266595in}{1.382569in}}%
\pgfusepath{stroke}%
\end{pgfscope}%
\begin{pgfscope}%
\pgfsetbuttcap%
\pgfsetroundjoin%
\definecolor{currentfill}{rgb}{0.000000,0.000000,0.000000}%
\pgfsetfillcolor{currentfill}%
\pgfsetlinewidth{0.803000pt}%
\definecolor{currentstroke}{rgb}{0.000000,0.000000,0.000000}%
\pgfsetstrokecolor{currentstroke}%
\pgfsetdash{}{0pt}%
\pgfsys@defobject{currentmarker}{\pgfqpoint{-0.048611in}{0.000000in}}{\pgfqpoint{-0.000000in}{0.000000in}}{%
\pgfpathmoveto{\pgfqpoint{-0.000000in}{0.000000in}}%
\pgfpathlineto{\pgfqpoint{-0.048611in}{0.000000in}}%
\pgfusepath{stroke,fill}%
}%
\begin{pgfscope}%
\pgfsys@transformshift{0.488997in}{1.382569in}%
\pgfsys@useobject{currentmarker}{}%
\end{pgfscope}%
\end{pgfscope}%
\begin{pgfscope}%
\definecolor{textcolor}{rgb}{0.000000,0.000000,0.000000}%
\pgfsetstrokecolor{textcolor}%
\pgfsetfillcolor{textcolor}%
\pgftext[x=0.303410in, y=1.329807in, left, base]{\color{textcolor}{\sffamily\fontsize{10.000000}{12.000000}\selectfont\catcode`\^=\active\def^{\ifmmode\sp\else\^{}\fi}\catcode`\%=\active\def%{\%}3}}%
\end{pgfscope}%
\begin{pgfscope}%
\pgfpathrectangle{\pgfqpoint{0.488997in}{0.535045in}}{\pgfqpoint{4.777597in}{2.260065in}}%
\pgfusepath{clip}%
\pgfsetbuttcap%
\pgfsetroundjoin%
\pgfsetlinewidth{0.803000pt}%
\definecolor{currentstroke}{rgb}{0.690196,0.690196,0.690196}%
\pgfsetstrokecolor{currentstroke}%
\pgfsetstrokeopacity{0.400000}%
\pgfsetdash{{2.960000pt}{1.280000pt}}{0.000000pt}%
\pgfpathmoveto{\pgfqpoint{0.488997in}{1.665077in}}%
\pgfpathlineto{\pgfqpoint{5.266595in}{1.665077in}}%
\pgfusepath{stroke}%
\end{pgfscope}%
\begin{pgfscope}%
\pgfsetbuttcap%
\pgfsetroundjoin%
\definecolor{currentfill}{rgb}{0.000000,0.000000,0.000000}%
\pgfsetfillcolor{currentfill}%
\pgfsetlinewidth{0.803000pt}%
\definecolor{currentstroke}{rgb}{0.000000,0.000000,0.000000}%
\pgfsetstrokecolor{currentstroke}%
\pgfsetdash{}{0pt}%
\pgfsys@defobject{currentmarker}{\pgfqpoint{-0.048611in}{0.000000in}}{\pgfqpoint{-0.000000in}{0.000000in}}{%
\pgfpathmoveto{\pgfqpoint{-0.000000in}{0.000000in}}%
\pgfpathlineto{\pgfqpoint{-0.048611in}{0.000000in}}%
\pgfusepath{stroke,fill}%
}%
\begin{pgfscope}%
\pgfsys@transformshift{0.488997in}{1.665077in}%
\pgfsys@useobject{currentmarker}{}%
\end{pgfscope}%
\end{pgfscope}%
\begin{pgfscope}%
\definecolor{textcolor}{rgb}{0.000000,0.000000,0.000000}%
\pgfsetstrokecolor{textcolor}%
\pgfsetfillcolor{textcolor}%
\pgftext[x=0.303410in, y=1.612316in, left, base]{\color{textcolor}{\sffamily\fontsize{10.000000}{12.000000}\selectfont\catcode`\^=\active\def^{\ifmmode\sp\else\^{}\fi}\catcode`\%=\active\def%{\%}4}}%
\end{pgfscope}%
\begin{pgfscope}%
\pgfpathrectangle{\pgfqpoint{0.488997in}{0.535045in}}{\pgfqpoint{4.777597in}{2.260065in}}%
\pgfusepath{clip}%
\pgfsetbuttcap%
\pgfsetroundjoin%
\pgfsetlinewidth{0.803000pt}%
\definecolor{currentstroke}{rgb}{0.690196,0.690196,0.690196}%
\pgfsetstrokecolor{currentstroke}%
\pgfsetstrokeopacity{0.400000}%
\pgfsetdash{{2.960000pt}{1.280000pt}}{0.000000pt}%
\pgfpathmoveto{\pgfqpoint{0.488997in}{1.947585in}}%
\pgfpathlineto{\pgfqpoint{5.266595in}{1.947585in}}%
\pgfusepath{stroke}%
\end{pgfscope}%
\begin{pgfscope}%
\pgfsetbuttcap%
\pgfsetroundjoin%
\definecolor{currentfill}{rgb}{0.000000,0.000000,0.000000}%
\pgfsetfillcolor{currentfill}%
\pgfsetlinewidth{0.803000pt}%
\definecolor{currentstroke}{rgb}{0.000000,0.000000,0.000000}%
\pgfsetstrokecolor{currentstroke}%
\pgfsetdash{}{0pt}%
\pgfsys@defobject{currentmarker}{\pgfqpoint{-0.048611in}{0.000000in}}{\pgfqpoint{-0.000000in}{0.000000in}}{%
\pgfpathmoveto{\pgfqpoint{-0.000000in}{0.000000in}}%
\pgfpathlineto{\pgfqpoint{-0.048611in}{0.000000in}}%
\pgfusepath{stroke,fill}%
}%
\begin{pgfscope}%
\pgfsys@transformshift{0.488997in}{1.947585in}%
\pgfsys@useobject{currentmarker}{}%
\end{pgfscope}%
\end{pgfscope}%
\begin{pgfscope}%
\definecolor{textcolor}{rgb}{0.000000,0.000000,0.000000}%
\pgfsetstrokecolor{textcolor}%
\pgfsetfillcolor{textcolor}%
\pgftext[x=0.303410in, y=1.894824in, left, base]{\color{textcolor}{\sffamily\fontsize{10.000000}{12.000000}\selectfont\catcode`\^=\active\def^{\ifmmode\sp\else\^{}\fi}\catcode`\%=\active\def%{\%}5}}%
\end{pgfscope}%
\begin{pgfscope}%
\pgfpathrectangle{\pgfqpoint{0.488997in}{0.535045in}}{\pgfqpoint{4.777597in}{2.260065in}}%
\pgfusepath{clip}%
\pgfsetbuttcap%
\pgfsetroundjoin%
\pgfsetlinewidth{0.803000pt}%
\definecolor{currentstroke}{rgb}{0.690196,0.690196,0.690196}%
\pgfsetstrokecolor{currentstroke}%
\pgfsetstrokeopacity{0.400000}%
\pgfsetdash{{2.960000pt}{1.280000pt}}{0.000000pt}%
\pgfpathmoveto{\pgfqpoint{0.488997in}{2.230093in}}%
\pgfpathlineto{\pgfqpoint{5.266595in}{2.230093in}}%
\pgfusepath{stroke}%
\end{pgfscope}%
\begin{pgfscope}%
\pgfsetbuttcap%
\pgfsetroundjoin%
\definecolor{currentfill}{rgb}{0.000000,0.000000,0.000000}%
\pgfsetfillcolor{currentfill}%
\pgfsetlinewidth{0.803000pt}%
\definecolor{currentstroke}{rgb}{0.000000,0.000000,0.000000}%
\pgfsetstrokecolor{currentstroke}%
\pgfsetdash{}{0pt}%
\pgfsys@defobject{currentmarker}{\pgfqpoint{-0.048611in}{0.000000in}}{\pgfqpoint{-0.000000in}{0.000000in}}{%
\pgfpathmoveto{\pgfqpoint{-0.000000in}{0.000000in}}%
\pgfpathlineto{\pgfqpoint{-0.048611in}{0.000000in}}%
\pgfusepath{stroke,fill}%
}%
\begin{pgfscope}%
\pgfsys@transformshift{0.488997in}{2.230093in}%
\pgfsys@useobject{currentmarker}{}%
\end{pgfscope}%
\end{pgfscope}%
\begin{pgfscope}%
\definecolor{textcolor}{rgb}{0.000000,0.000000,0.000000}%
\pgfsetstrokecolor{textcolor}%
\pgfsetfillcolor{textcolor}%
\pgftext[x=0.303410in, y=2.177332in, left, base]{\color{textcolor}{\sffamily\fontsize{10.000000}{12.000000}\selectfont\catcode`\^=\active\def^{\ifmmode\sp\else\^{}\fi}\catcode`\%=\active\def%{\%}6}}%
\end{pgfscope}%
\begin{pgfscope}%
\pgfpathrectangle{\pgfqpoint{0.488997in}{0.535045in}}{\pgfqpoint{4.777597in}{2.260065in}}%
\pgfusepath{clip}%
\pgfsetbuttcap%
\pgfsetroundjoin%
\pgfsetlinewidth{0.803000pt}%
\definecolor{currentstroke}{rgb}{0.690196,0.690196,0.690196}%
\pgfsetstrokecolor{currentstroke}%
\pgfsetstrokeopacity{0.400000}%
\pgfsetdash{{2.960000pt}{1.280000pt}}{0.000000pt}%
\pgfpathmoveto{\pgfqpoint{0.488997in}{2.512602in}}%
\pgfpathlineto{\pgfqpoint{5.266595in}{2.512602in}}%
\pgfusepath{stroke}%
\end{pgfscope}%
\begin{pgfscope}%
\pgfsetbuttcap%
\pgfsetroundjoin%
\definecolor{currentfill}{rgb}{0.000000,0.000000,0.000000}%
\pgfsetfillcolor{currentfill}%
\pgfsetlinewidth{0.803000pt}%
\definecolor{currentstroke}{rgb}{0.000000,0.000000,0.000000}%
\pgfsetstrokecolor{currentstroke}%
\pgfsetdash{}{0pt}%
\pgfsys@defobject{currentmarker}{\pgfqpoint{-0.048611in}{0.000000in}}{\pgfqpoint{-0.000000in}{0.000000in}}{%
\pgfpathmoveto{\pgfqpoint{-0.000000in}{0.000000in}}%
\pgfpathlineto{\pgfqpoint{-0.048611in}{0.000000in}}%
\pgfusepath{stroke,fill}%
}%
\begin{pgfscope}%
\pgfsys@transformshift{0.488997in}{2.512602in}%
\pgfsys@useobject{currentmarker}{}%
\end{pgfscope}%
\end{pgfscope}%
\begin{pgfscope}%
\definecolor{textcolor}{rgb}{0.000000,0.000000,0.000000}%
\pgfsetstrokecolor{textcolor}%
\pgfsetfillcolor{textcolor}%
\pgftext[x=0.303410in, y=2.459840in, left, base]{\color{textcolor}{\sffamily\fontsize{10.000000}{12.000000}\selectfont\catcode`\^=\active\def^{\ifmmode\sp\else\^{}\fi}\catcode`\%=\active\def%{\%}7}}%
\end{pgfscope}%
\begin{pgfscope}%
\pgfpathrectangle{\pgfqpoint{0.488997in}{0.535045in}}{\pgfqpoint{4.777597in}{2.260065in}}%
\pgfusepath{clip}%
\pgfsetbuttcap%
\pgfsetroundjoin%
\pgfsetlinewidth{0.803000pt}%
\definecolor{currentstroke}{rgb}{0.690196,0.690196,0.690196}%
\pgfsetstrokecolor{currentstroke}%
\pgfsetstrokeopacity{0.400000}%
\pgfsetdash{{2.960000pt}{1.280000pt}}{0.000000pt}%
\pgfpathmoveto{\pgfqpoint{0.488997in}{2.795110in}}%
\pgfpathlineto{\pgfqpoint{5.266595in}{2.795110in}}%
\pgfusepath{stroke}%
\end{pgfscope}%
\begin{pgfscope}%
\pgfsetbuttcap%
\pgfsetroundjoin%
\definecolor{currentfill}{rgb}{0.000000,0.000000,0.000000}%
\pgfsetfillcolor{currentfill}%
\pgfsetlinewidth{0.803000pt}%
\definecolor{currentstroke}{rgb}{0.000000,0.000000,0.000000}%
\pgfsetstrokecolor{currentstroke}%
\pgfsetdash{}{0pt}%
\pgfsys@defobject{currentmarker}{\pgfqpoint{-0.048611in}{0.000000in}}{\pgfqpoint{-0.000000in}{0.000000in}}{%
\pgfpathmoveto{\pgfqpoint{-0.000000in}{0.000000in}}%
\pgfpathlineto{\pgfqpoint{-0.048611in}{0.000000in}}%
\pgfusepath{stroke,fill}%
}%
\begin{pgfscope}%
\pgfsys@transformshift{0.488997in}{2.795110in}%
\pgfsys@useobject{currentmarker}{}%
\end{pgfscope}%
\end{pgfscope}%
\begin{pgfscope}%
\definecolor{textcolor}{rgb}{0.000000,0.000000,0.000000}%
\pgfsetstrokecolor{textcolor}%
\pgfsetfillcolor{textcolor}%
\pgftext[x=0.303410in, y=2.742348in, left, base]{\color{textcolor}{\sffamily\fontsize{10.000000}{12.000000}\selectfont\catcode`\^=\active\def^{\ifmmode\sp\else\^{}\fi}\catcode`\%=\active\def%{\%}8}}%
\end{pgfscope}%
\begin{pgfscope}%
\definecolor{textcolor}{rgb}{0.000000,0.000000,0.000000}%
\pgfsetstrokecolor{textcolor}%
\pgfsetfillcolor{textcolor}%
\pgftext[x=0.247854in,y=1.665077in,,bottom,rotate=90.000000]{\color{textcolor}{\sffamily\fontsize{11.000000}{13.200000}\selectfont\catcode`\^=\active\def^{\ifmmode\sp\else\^{}\fi}\catcode`\%=\active\def%{\%}Throughput (Mbps)}}%
\end{pgfscope}%
\begin{pgfscope}%
\pgfpathrectangle{\pgfqpoint{0.488997in}{0.535045in}}{\pgfqpoint{4.777597in}{2.260065in}}%
\pgfusepath{clip}%
\pgfsetrectcap%
\pgfsetroundjoin%
\pgfsetlinewidth{2.007500pt}%
\definecolor{currentstroke}{rgb}{0.082353,0.396078,0.752941}%
\pgfsetstrokecolor{currentstroke}%
\pgfsetdash{}{0pt}%
\pgfpathmoveto{\pgfqpoint{0.706161in}{0.535045in}}%
\pgfpathlineto{\pgfqpoint{0.758489in}{1.516817in}}%
\pgfpathlineto{\pgfqpoint{0.810818in}{1.647720in}}%
\pgfpathlineto{\pgfqpoint{0.863146in}{1.549543in}}%
\pgfpathlineto{\pgfqpoint{0.915475in}{1.614994in}}%
\pgfpathlineto{\pgfqpoint{0.967803in}{1.582268in}}%
\pgfpathlineto{\pgfqpoint{1.020132in}{1.582268in}}%
\pgfpathlineto{\pgfqpoint{1.072461in}{1.680446in}}%
\pgfpathlineto{\pgfqpoint{1.124789in}{1.745897in}}%
\pgfpathlineto{\pgfqpoint{1.177118in}{1.614994in}}%
\pgfpathlineto{\pgfqpoint{1.229446in}{1.582268in}}%
\pgfpathlineto{\pgfqpoint{1.281775in}{1.844074in}}%
\pgfpathlineto{\pgfqpoint{1.334103in}{1.582268in}}%
\pgfpathlineto{\pgfqpoint{1.386432in}{1.614994in}}%
\pgfpathlineto{\pgfqpoint{1.438760in}{1.844074in}}%
\pgfpathlineto{\pgfqpoint{1.491089in}{1.582268in}}%
\pgfpathlineto{\pgfqpoint{1.543418in}{1.582268in}}%
\pgfpathlineto{\pgfqpoint{1.595746in}{1.876800in}}%
\pgfpathlineto{\pgfqpoint{1.648075in}{1.614994in}}%
\pgfpathlineto{\pgfqpoint{1.700403in}{1.876800in}}%
\pgfpathlineto{\pgfqpoint{1.752732in}{1.582268in}}%
\pgfpathlineto{\pgfqpoint{1.805060in}{1.876800in}}%
\pgfpathlineto{\pgfqpoint{1.857389in}{1.614994in}}%
\pgfpathlineto{\pgfqpoint{1.909717in}{1.778623in}}%
\pgfpathlineto{\pgfqpoint{1.962046in}{1.647720in}}%
\pgfpathlineto{\pgfqpoint{2.014375in}{1.811349in}}%
\pgfpathlineto{\pgfqpoint{2.066703in}{1.680446in}}%
\pgfpathlineto{\pgfqpoint{2.119032in}{1.876800in}}%
\pgfpathlineto{\pgfqpoint{2.171360in}{1.582268in}}%
\pgfpathlineto{\pgfqpoint{2.223689in}{1.844074in}}%
\pgfpathlineto{\pgfqpoint{2.276017in}{1.582268in}}%
\pgfpathlineto{\pgfqpoint{2.328346in}{1.909526in}}%
\pgfpathlineto{\pgfqpoint{2.380674in}{1.582268in}}%
\pgfpathlineto{\pgfqpoint{2.433003in}{1.909526in}}%
\pgfpathlineto{\pgfqpoint{2.485332in}{1.844074in}}%
\pgfpathlineto{\pgfqpoint{2.537660in}{1.582268in}}%
\pgfpathlineto{\pgfqpoint{2.589989in}{1.909526in}}%
\pgfpathlineto{\pgfqpoint{2.642317in}{1.549543in}}%
\pgfpathlineto{\pgfqpoint{2.694646in}{1.876800in}}%
\pgfpathlineto{\pgfqpoint{2.746974in}{1.909526in}}%
\pgfpathlineto{\pgfqpoint{2.799303in}{1.582268in}}%
\pgfpathlineto{\pgfqpoint{2.851632in}{1.876800in}}%
\pgfpathlineto{\pgfqpoint{2.903960in}{1.745897in}}%
\pgfpathlineto{\pgfqpoint{2.956289in}{1.713171in}}%
\pgfpathlineto{\pgfqpoint{3.008617in}{1.876800in}}%
\pgfpathlineto{\pgfqpoint{3.060946in}{1.582268in}}%
\pgfpathlineto{\pgfqpoint{3.113274in}{1.909526in}}%
\pgfpathlineto{\pgfqpoint{3.165603in}{1.876800in}}%
\pgfpathlineto{\pgfqpoint{3.217931in}{1.582268in}}%
\pgfpathlineto{\pgfqpoint{3.270260in}{1.876800in}}%
\pgfpathlineto{\pgfqpoint{3.322589in}{1.876800in}}%
\pgfpathlineto{\pgfqpoint{3.374917in}{1.811349in}}%
\pgfpathlineto{\pgfqpoint{3.427246in}{0.535045in}}%
\pgfpathlineto{\pgfqpoint{3.479574in}{0.535045in}}%
\pgfpathlineto{\pgfqpoint{3.531903in}{0.567770in}}%
\pgfpathlineto{\pgfqpoint{3.584231in}{0.535045in}}%
\pgfpathlineto{\pgfqpoint{3.636560in}{0.535045in}}%
\pgfpathlineto{\pgfqpoint{3.688888in}{0.535045in}}%
\pgfpathlineto{\pgfqpoint{3.741217in}{0.535045in}}%
\pgfpathlineto{\pgfqpoint{3.793546in}{0.535045in}}%
\pgfpathlineto{\pgfqpoint{3.845874in}{0.535045in}}%
\pgfpathlineto{\pgfqpoint{3.898203in}{0.535045in}}%
\pgfpathlineto{\pgfqpoint{3.950531in}{0.535045in}}%
\pgfpathlineto{\pgfqpoint{4.002860in}{0.535045in}}%
\pgfpathlineto{\pgfqpoint{4.055188in}{0.535045in}}%
\pgfpathlineto{\pgfqpoint{4.107517in}{0.535045in}}%
\pgfpathlineto{\pgfqpoint{4.159846in}{0.535045in}}%
\pgfpathlineto{\pgfqpoint{4.212174in}{0.535045in}}%
\pgfpathlineto{\pgfqpoint{4.264503in}{0.535045in}}%
\pgfpathlineto{\pgfqpoint{4.316831in}{0.535045in}}%
\pgfpathlineto{\pgfqpoint{4.369160in}{0.535045in}}%
\pgfpathlineto{\pgfqpoint{4.421488in}{0.535045in}}%
\pgfpathlineto{\pgfqpoint{4.473817in}{0.535045in}}%
\pgfpathlineto{\pgfqpoint{4.526145in}{0.535045in}}%
\pgfpathlineto{\pgfqpoint{4.578474in}{0.535045in}}%
\pgfpathlineto{\pgfqpoint{4.630803in}{0.535045in}}%
\pgfpathlineto{\pgfqpoint{4.683131in}{0.535045in}}%
\pgfpathlineto{\pgfqpoint{4.735460in}{0.535045in}}%
\pgfpathlineto{\pgfqpoint{4.787788in}{0.535045in}}%
\pgfpathlineto{\pgfqpoint{4.840117in}{0.535045in}}%
\pgfpathlineto{\pgfqpoint{4.892445in}{0.535045in}}%
\pgfpathlineto{\pgfqpoint{4.944774in}{0.535045in}}%
\pgfpathlineto{\pgfqpoint{4.997102in}{0.535045in}}%
\pgfpathlineto{\pgfqpoint{5.049431in}{0.535045in}}%
\pgfusepath{stroke}%
\end{pgfscope}%
\begin{pgfscope}%
\pgfpathrectangle{\pgfqpoint{0.488997in}{0.535045in}}{\pgfqpoint{4.777597in}{2.260065in}}%
\pgfusepath{clip}%
\pgfsetbuttcap%
\pgfsetroundjoin%
\definecolor{currentfill}{rgb}{0.082353,0.396078,0.752941}%
\pgfsetfillcolor{currentfill}%
\pgfsetlinewidth{1.003750pt}%
\definecolor{currentstroke}{rgb}{0.082353,0.396078,0.752941}%
\pgfsetstrokecolor{currentstroke}%
\pgfsetdash{}{0pt}%
\pgfsys@defobject{currentmarker}{\pgfqpoint{-0.027778in}{-0.027778in}}{\pgfqpoint{0.027778in}{0.027778in}}{%
\pgfpathmoveto{\pgfqpoint{0.000000in}{-0.027778in}}%
\pgfpathcurveto{\pgfqpoint{0.007367in}{-0.027778in}}{\pgfqpoint{0.014433in}{-0.024851in}}{\pgfqpoint{0.019642in}{-0.019642in}}%
\pgfpathcurveto{\pgfqpoint{0.024851in}{-0.014433in}}{\pgfqpoint{0.027778in}{-0.007367in}}{\pgfqpoint{0.027778in}{0.000000in}}%
\pgfpathcurveto{\pgfqpoint{0.027778in}{0.007367in}}{\pgfqpoint{0.024851in}{0.014433in}}{\pgfqpoint{0.019642in}{0.019642in}}%
\pgfpathcurveto{\pgfqpoint{0.014433in}{0.024851in}}{\pgfqpoint{0.007367in}{0.027778in}}{\pgfqpoint{0.000000in}{0.027778in}}%
\pgfpathcurveto{\pgfqpoint{-0.007367in}{0.027778in}}{\pgfqpoint{-0.014433in}{0.024851in}}{\pgfqpoint{-0.019642in}{0.019642in}}%
\pgfpathcurveto{\pgfqpoint{-0.024851in}{0.014433in}}{\pgfqpoint{-0.027778in}{0.007367in}}{\pgfqpoint{-0.027778in}{0.000000in}}%
\pgfpathcurveto{\pgfqpoint{-0.027778in}{-0.007367in}}{\pgfqpoint{-0.024851in}{-0.014433in}}{\pgfqpoint{-0.019642in}{-0.019642in}}%
\pgfpathcurveto{\pgfqpoint{-0.014433in}{-0.024851in}}{\pgfqpoint{-0.007367in}{-0.027778in}}{\pgfqpoint{0.000000in}{-0.027778in}}%
\pgfpathlineto{\pgfqpoint{0.000000in}{-0.027778in}}%
\pgfpathclose%
\pgfusepath{stroke,fill}%
}%
\begin{pgfscope}%
\pgfsys@transformshift{0.706161in}{0.535045in}%
\pgfsys@useobject{currentmarker}{}%
\end{pgfscope}%
\begin{pgfscope}%
\pgfsys@transformshift{0.758489in}{1.516817in}%
\pgfsys@useobject{currentmarker}{}%
\end{pgfscope}%
\begin{pgfscope}%
\pgfsys@transformshift{0.810818in}{1.647720in}%
\pgfsys@useobject{currentmarker}{}%
\end{pgfscope}%
\begin{pgfscope}%
\pgfsys@transformshift{0.863146in}{1.549543in}%
\pgfsys@useobject{currentmarker}{}%
\end{pgfscope}%
\begin{pgfscope}%
\pgfsys@transformshift{0.915475in}{1.614994in}%
\pgfsys@useobject{currentmarker}{}%
\end{pgfscope}%
\begin{pgfscope}%
\pgfsys@transformshift{0.967803in}{1.582268in}%
\pgfsys@useobject{currentmarker}{}%
\end{pgfscope}%
\begin{pgfscope}%
\pgfsys@transformshift{1.020132in}{1.582268in}%
\pgfsys@useobject{currentmarker}{}%
\end{pgfscope}%
\begin{pgfscope}%
\pgfsys@transformshift{1.072461in}{1.680446in}%
\pgfsys@useobject{currentmarker}{}%
\end{pgfscope}%
\begin{pgfscope}%
\pgfsys@transformshift{1.124789in}{1.745897in}%
\pgfsys@useobject{currentmarker}{}%
\end{pgfscope}%
\begin{pgfscope}%
\pgfsys@transformshift{1.177118in}{1.614994in}%
\pgfsys@useobject{currentmarker}{}%
\end{pgfscope}%
\begin{pgfscope}%
\pgfsys@transformshift{1.229446in}{1.582268in}%
\pgfsys@useobject{currentmarker}{}%
\end{pgfscope}%
\begin{pgfscope}%
\pgfsys@transformshift{1.281775in}{1.844074in}%
\pgfsys@useobject{currentmarker}{}%
\end{pgfscope}%
\begin{pgfscope}%
\pgfsys@transformshift{1.334103in}{1.582268in}%
\pgfsys@useobject{currentmarker}{}%
\end{pgfscope}%
\begin{pgfscope}%
\pgfsys@transformshift{1.386432in}{1.614994in}%
\pgfsys@useobject{currentmarker}{}%
\end{pgfscope}%
\begin{pgfscope}%
\pgfsys@transformshift{1.438760in}{1.844074in}%
\pgfsys@useobject{currentmarker}{}%
\end{pgfscope}%
\begin{pgfscope}%
\pgfsys@transformshift{1.491089in}{1.582268in}%
\pgfsys@useobject{currentmarker}{}%
\end{pgfscope}%
\begin{pgfscope}%
\pgfsys@transformshift{1.543418in}{1.582268in}%
\pgfsys@useobject{currentmarker}{}%
\end{pgfscope}%
\begin{pgfscope}%
\pgfsys@transformshift{1.595746in}{1.876800in}%
\pgfsys@useobject{currentmarker}{}%
\end{pgfscope}%
\begin{pgfscope}%
\pgfsys@transformshift{1.648075in}{1.614994in}%
\pgfsys@useobject{currentmarker}{}%
\end{pgfscope}%
\begin{pgfscope}%
\pgfsys@transformshift{1.700403in}{1.876800in}%
\pgfsys@useobject{currentmarker}{}%
\end{pgfscope}%
\begin{pgfscope}%
\pgfsys@transformshift{1.752732in}{1.582268in}%
\pgfsys@useobject{currentmarker}{}%
\end{pgfscope}%
\begin{pgfscope}%
\pgfsys@transformshift{1.805060in}{1.876800in}%
\pgfsys@useobject{currentmarker}{}%
\end{pgfscope}%
\begin{pgfscope}%
\pgfsys@transformshift{1.857389in}{1.614994in}%
\pgfsys@useobject{currentmarker}{}%
\end{pgfscope}%
\begin{pgfscope}%
\pgfsys@transformshift{1.909717in}{1.778623in}%
\pgfsys@useobject{currentmarker}{}%
\end{pgfscope}%
\begin{pgfscope}%
\pgfsys@transformshift{1.962046in}{1.647720in}%
\pgfsys@useobject{currentmarker}{}%
\end{pgfscope}%
\begin{pgfscope}%
\pgfsys@transformshift{2.014375in}{1.811349in}%
\pgfsys@useobject{currentmarker}{}%
\end{pgfscope}%
\begin{pgfscope}%
\pgfsys@transformshift{2.066703in}{1.680446in}%
\pgfsys@useobject{currentmarker}{}%
\end{pgfscope}%
\begin{pgfscope}%
\pgfsys@transformshift{2.119032in}{1.876800in}%
\pgfsys@useobject{currentmarker}{}%
\end{pgfscope}%
\begin{pgfscope}%
\pgfsys@transformshift{2.171360in}{1.582268in}%
\pgfsys@useobject{currentmarker}{}%
\end{pgfscope}%
\begin{pgfscope}%
\pgfsys@transformshift{2.223689in}{1.844074in}%
\pgfsys@useobject{currentmarker}{}%
\end{pgfscope}%
\begin{pgfscope}%
\pgfsys@transformshift{2.276017in}{1.582268in}%
\pgfsys@useobject{currentmarker}{}%
\end{pgfscope}%
\begin{pgfscope}%
\pgfsys@transformshift{2.328346in}{1.909526in}%
\pgfsys@useobject{currentmarker}{}%
\end{pgfscope}%
\begin{pgfscope}%
\pgfsys@transformshift{2.380674in}{1.582268in}%
\pgfsys@useobject{currentmarker}{}%
\end{pgfscope}%
\begin{pgfscope}%
\pgfsys@transformshift{2.433003in}{1.909526in}%
\pgfsys@useobject{currentmarker}{}%
\end{pgfscope}%
\begin{pgfscope}%
\pgfsys@transformshift{2.485332in}{1.844074in}%
\pgfsys@useobject{currentmarker}{}%
\end{pgfscope}%
\begin{pgfscope}%
\pgfsys@transformshift{2.537660in}{1.582268in}%
\pgfsys@useobject{currentmarker}{}%
\end{pgfscope}%
\begin{pgfscope}%
\pgfsys@transformshift{2.589989in}{1.909526in}%
\pgfsys@useobject{currentmarker}{}%
\end{pgfscope}%
\begin{pgfscope}%
\pgfsys@transformshift{2.642317in}{1.549543in}%
\pgfsys@useobject{currentmarker}{}%
\end{pgfscope}%
\begin{pgfscope}%
\pgfsys@transformshift{2.694646in}{1.876800in}%
\pgfsys@useobject{currentmarker}{}%
\end{pgfscope}%
\begin{pgfscope}%
\pgfsys@transformshift{2.746974in}{1.909526in}%
\pgfsys@useobject{currentmarker}{}%
\end{pgfscope}%
\begin{pgfscope}%
\pgfsys@transformshift{2.799303in}{1.582268in}%
\pgfsys@useobject{currentmarker}{}%
\end{pgfscope}%
\begin{pgfscope}%
\pgfsys@transformshift{2.851632in}{1.876800in}%
\pgfsys@useobject{currentmarker}{}%
\end{pgfscope}%
\begin{pgfscope}%
\pgfsys@transformshift{2.903960in}{1.745897in}%
\pgfsys@useobject{currentmarker}{}%
\end{pgfscope}%
\begin{pgfscope}%
\pgfsys@transformshift{2.956289in}{1.713171in}%
\pgfsys@useobject{currentmarker}{}%
\end{pgfscope}%
\begin{pgfscope}%
\pgfsys@transformshift{3.008617in}{1.876800in}%
\pgfsys@useobject{currentmarker}{}%
\end{pgfscope}%
\begin{pgfscope}%
\pgfsys@transformshift{3.060946in}{1.582268in}%
\pgfsys@useobject{currentmarker}{}%
\end{pgfscope}%
\begin{pgfscope}%
\pgfsys@transformshift{3.113274in}{1.909526in}%
\pgfsys@useobject{currentmarker}{}%
\end{pgfscope}%
\begin{pgfscope}%
\pgfsys@transformshift{3.165603in}{1.876800in}%
\pgfsys@useobject{currentmarker}{}%
\end{pgfscope}%
\begin{pgfscope}%
\pgfsys@transformshift{3.217931in}{1.582268in}%
\pgfsys@useobject{currentmarker}{}%
\end{pgfscope}%
\begin{pgfscope}%
\pgfsys@transformshift{3.270260in}{1.876800in}%
\pgfsys@useobject{currentmarker}{}%
\end{pgfscope}%
\begin{pgfscope}%
\pgfsys@transformshift{3.322589in}{1.876800in}%
\pgfsys@useobject{currentmarker}{}%
\end{pgfscope}%
\begin{pgfscope}%
\pgfsys@transformshift{3.374917in}{1.811349in}%
\pgfsys@useobject{currentmarker}{}%
\end{pgfscope}%
\begin{pgfscope}%
\pgfsys@transformshift{3.427246in}{0.535045in}%
\pgfsys@useobject{currentmarker}{}%
\end{pgfscope}%
\begin{pgfscope}%
\pgfsys@transformshift{3.479574in}{0.535045in}%
\pgfsys@useobject{currentmarker}{}%
\end{pgfscope}%
\begin{pgfscope}%
\pgfsys@transformshift{3.531903in}{0.567770in}%
\pgfsys@useobject{currentmarker}{}%
\end{pgfscope}%
\begin{pgfscope}%
\pgfsys@transformshift{3.584231in}{0.535045in}%
\pgfsys@useobject{currentmarker}{}%
\end{pgfscope}%
\begin{pgfscope}%
\pgfsys@transformshift{3.636560in}{0.535045in}%
\pgfsys@useobject{currentmarker}{}%
\end{pgfscope}%
\begin{pgfscope}%
\pgfsys@transformshift{3.688888in}{0.535045in}%
\pgfsys@useobject{currentmarker}{}%
\end{pgfscope}%
\begin{pgfscope}%
\pgfsys@transformshift{3.741217in}{0.535045in}%
\pgfsys@useobject{currentmarker}{}%
\end{pgfscope}%
\begin{pgfscope}%
\pgfsys@transformshift{3.793546in}{0.535045in}%
\pgfsys@useobject{currentmarker}{}%
\end{pgfscope}%
\begin{pgfscope}%
\pgfsys@transformshift{3.845874in}{0.535045in}%
\pgfsys@useobject{currentmarker}{}%
\end{pgfscope}%
\begin{pgfscope}%
\pgfsys@transformshift{3.898203in}{0.535045in}%
\pgfsys@useobject{currentmarker}{}%
\end{pgfscope}%
\begin{pgfscope}%
\pgfsys@transformshift{3.950531in}{0.535045in}%
\pgfsys@useobject{currentmarker}{}%
\end{pgfscope}%
\begin{pgfscope}%
\pgfsys@transformshift{4.002860in}{0.535045in}%
\pgfsys@useobject{currentmarker}{}%
\end{pgfscope}%
\begin{pgfscope}%
\pgfsys@transformshift{4.055188in}{0.535045in}%
\pgfsys@useobject{currentmarker}{}%
\end{pgfscope}%
\begin{pgfscope}%
\pgfsys@transformshift{4.107517in}{0.535045in}%
\pgfsys@useobject{currentmarker}{}%
\end{pgfscope}%
\begin{pgfscope}%
\pgfsys@transformshift{4.159846in}{0.535045in}%
\pgfsys@useobject{currentmarker}{}%
\end{pgfscope}%
\begin{pgfscope}%
\pgfsys@transformshift{4.212174in}{0.535045in}%
\pgfsys@useobject{currentmarker}{}%
\end{pgfscope}%
\begin{pgfscope}%
\pgfsys@transformshift{4.264503in}{0.535045in}%
\pgfsys@useobject{currentmarker}{}%
\end{pgfscope}%
\begin{pgfscope}%
\pgfsys@transformshift{4.316831in}{0.535045in}%
\pgfsys@useobject{currentmarker}{}%
\end{pgfscope}%
\begin{pgfscope}%
\pgfsys@transformshift{4.369160in}{0.535045in}%
\pgfsys@useobject{currentmarker}{}%
\end{pgfscope}%
\begin{pgfscope}%
\pgfsys@transformshift{4.421488in}{0.535045in}%
\pgfsys@useobject{currentmarker}{}%
\end{pgfscope}%
\begin{pgfscope}%
\pgfsys@transformshift{4.473817in}{0.535045in}%
\pgfsys@useobject{currentmarker}{}%
\end{pgfscope}%
\begin{pgfscope}%
\pgfsys@transformshift{4.526145in}{0.535045in}%
\pgfsys@useobject{currentmarker}{}%
\end{pgfscope}%
\begin{pgfscope}%
\pgfsys@transformshift{4.578474in}{0.535045in}%
\pgfsys@useobject{currentmarker}{}%
\end{pgfscope}%
\begin{pgfscope}%
\pgfsys@transformshift{4.630803in}{0.535045in}%
\pgfsys@useobject{currentmarker}{}%
\end{pgfscope}%
\begin{pgfscope}%
\pgfsys@transformshift{4.683131in}{0.535045in}%
\pgfsys@useobject{currentmarker}{}%
\end{pgfscope}%
\begin{pgfscope}%
\pgfsys@transformshift{4.735460in}{0.535045in}%
\pgfsys@useobject{currentmarker}{}%
\end{pgfscope}%
\begin{pgfscope}%
\pgfsys@transformshift{4.787788in}{0.535045in}%
\pgfsys@useobject{currentmarker}{}%
\end{pgfscope}%
\begin{pgfscope}%
\pgfsys@transformshift{4.840117in}{0.535045in}%
\pgfsys@useobject{currentmarker}{}%
\end{pgfscope}%
\begin{pgfscope}%
\pgfsys@transformshift{4.892445in}{0.535045in}%
\pgfsys@useobject{currentmarker}{}%
\end{pgfscope}%
\begin{pgfscope}%
\pgfsys@transformshift{4.944774in}{0.535045in}%
\pgfsys@useobject{currentmarker}{}%
\end{pgfscope}%
\begin{pgfscope}%
\pgfsys@transformshift{4.997102in}{0.535045in}%
\pgfsys@useobject{currentmarker}{}%
\end{pgfscope}%
\begin{pgfscope}%
\pgfsys@transformshift{5.049431in}{0.535045in}%
\pgfsys@useobject{currentmarker}{}%
\end{pgfscope}%
\end{pgfscope}%
\begin{pgfscope}%
\pgfpathrectangle{\pgfqpoint{0.488997in}{0.535045in}}{\pgfqpoint{4.777597in}{2.260065in}}%
\pgfusepath{clip}%
\pgfsetrectcap%
\pgfsetroundjoin%
\pgfsetlinewidth{2.007500pt}%
\definecolor{currentstroke}{rgb}{0.180392,0.490196,0.196078}%
\pgfsetstrokecolor{currentstroke}%
\pgfsetdash{}{0pt}%
\pgfpathmoveto{\pgfqpoint{0.706161in}{0.535045in}}%
\pgfpathlineto{\pgfqpoint{0.758489in}{0.662060in}}%
\pgfpathlineto{\pgfqpoint{0.810818in}{1.012416in}}%
\pgfpathlineto{\pgfqpoint{0.863146in}{1.806648in}}%
\pgfpathlineto{\pgfqpoint{0.915475in}{1.838266in}}%
\pgfpathlineto{\pgfqpoint{0.967803in}{1.774171in}}%
\pgfpathlineto{\pgfqpoint{1.020132in}{1.870178in}}%
\pgfpathlineto{\pgfqpoint{1.072461in}{1.870359in}}%
\pgfpathlineto{\pgfqpoint{1.124789in}{1.870607in}}%
\pgfpathlineto{\pgfqpoint{1.177118in}{1.838334in}}%
\pgfpathlineto{\pgfqpoint{1.229446in}{1.870607in}}%
\pgfpathlineto{\pgfqpoint{1.281775in}{1.838289in}}%
\pgfpathlineto{\pgfqpoint{1.334103in}{1.870381in}}%
\pgfpathlineto{\pgfqpoint{1.386432in}{1.838582in}}%
\pgfpathlineto{\pgfqpoint{1.438760in}{1.870381in}}%
\pgfpathlineto{\pgfqpoint{1.491089in}{1.870314in}}%
\pgfpathlineto{\pgfqpoint{1.543418in}{1.838515in}}%
\pgfpathlineto{\pgfqpoint{1.595746in}{1.870381in}}%
\pgfpathlineto{\pgfqpoint{1.648075in}{1.870110in}}%
\pgfpathlineto{\pgfqpoint{1.700403in}{1.838560in}}%
\pgfpathlineto{\pgfqpoint{1.752732in}{1.870359in}}%
\pgfpathlineto{\pgfqpoint{1.805060in}{1.838582in}}%
\pgfpathlineto{\pgfqpoint{1.857389in}{1.870133in}}%
\pgfpathlineto{\pgfqpoint{1.909717in}{1.870607in}}%
\pgfpathlineto{\pgfqpoint{1.962046in}{1.838582in}}%
\pgfpathlineto{\pgfqpoint{2.014375in}{1.870110in}}%
\pgfpathlineto{\pgfqpoint{2.066703in}{1.870607in}}%
\pgfpathlineto{\pgfqpoint{2.119032in}{1.837724in}}%
\pgfpathlineto{\pgfqpoint{2.171360in}{1.869184in}}%
\pgfpathlineto{\pgfqpoint{2.223689in}{1.837656in}}%
\pgfpathlineto{\pgfqpoint{2.276017in}{1.869432in}}%
\pgfpathlineto{\pgfqpoint{2.328346in}{1.869432in}}%
\pgfpathlineto{\pgfqpoint{2.380674in}{1.837656in}}%
\pgfpathlineto{\pgfqpoint{2.433003in}{1.869432in}}%
\pgfpathlineto{\pgfqpoint{2.485332in}{1.837656in}}%
\pgfpathlineto{\pgfqpoint{2.537660in}{1.869410in}}%
\pgfpathlineto{\pgfqpoint{2.589989in}{1.869432in}}%
\pgfpathlineto{\pgfqpoint{2.642317in}{1.837656in}}%
\pgfpathlineto{\pgfqpoint{2.694646in}{1.869432in}}%
\pgfpathlineto{\pgfqpoint{2.746974in}{1.869184in}}%
\pgfpathlineto{\pgfqpoint{2.799303in}{1.837882in}}%
\pgfpathlineto{\pgfqpoint{2.851632in}{1.869184in}}%
\pgfpathlineto{\pgfqpoint{2.903960in}{1.837633in}}%
\pgfpathlineto{\pgfqpoint{2.956289in}{1.869432in}}%
\pgfpathlineto{\pgfqpoint{3.008617in}{1.869410in}}%
\pgfpathlineto{\pgfqpoint{3.060946in}{1.837633in}}%
\pgfpathlineto{\pgfqpoint{3.113274in}{1.869410in}}%
\pgfpathlineto{\pgfqpoint{3.165603in}{1.869432in}}%
\pgfpathlineto{\pgfqpoint{3.217931in}{1.837656in}}%
\pgfpathlineto{\pgfqpoint{3.270260in}{1.869432in}}%
\pgfpathlineto{\pgfqpoint{3.322589in}{1.551667in}}%
\pgfpathlineto{\pgfqpoint{3.374917in}{1.138798in}}%
\pgfpathlineto{\pgfqpoint{3.427246in}{1.901073in}}%
\pgfpathlineto{\pgfqpoint{3.479574in}{1.837588in}}%
\pgfpathlineto{\pgfqpoint{3.531903in}{1.869364in}}%
\pgfpathlineto{\pgfqpoint{3.584231in}{1.869410in}}%
\pgfpathlineto{\pgfqpoint{3.636560in}{1.837633in}}%
\pgfpathlineto{\pgfqpoint{3.688888in}{1.869387in}}%
\pgfpathlineto{\pgfqpoint{3.741217in}{1.869184in}}%
\pgfpathlineto{\pgfqpoint{3.793546in}{1.837407in}}%
\pgfpathlineto{\pgfqpoint{3.845874in}{1.869658in}}%
\pgfpathlineto{\pgfqpoint{3.898203in}{1.837430in}}%
\pgfpathlineto{\pgfqpoint{3.950531in}{1.869432in}}%
\pgfpathlineto{\pgfqpoint{4.002860in}{1.869410in}}%
\pgfpathlineto{\pgfqpoint{4.055188in}{1.837656in}}%
\pgfpathlineto{\pgfqpoint{4.107517in}{1.869410in}}%
\pgfpathlineto{\pgfqpoint{4.159846in}{1.837633in}}%
\pgfpathlineto{\pgfqpoint{4.212174in}{1.869432in}}%
\pgfpathlineto{\pgfqpoint{4.264503in}{1.869432in}}%
\pgfpathlineto{\pgfqpoint{4.316831in}{1.837656in}}%
\pgfpathlineto{\pgfqpoint{4.369160in}{1.869410in}}%
\pgfpathlineto{\pgfqpoint{4.421488in}{1.837000in}}%
\pgfpathlineto{\pgfqpoint{4.473817in}{1.869432in}}%
\pgfpathlineto{\pgfqpoint{4.526145in}{1.869432in}}%
\pgfpathlineto{\pgfqpoint{4.578474in}{1.837656in}}%
\pgfpathlineto{\pgfqpoint{4.630803in}{1.869432in}}%
\pgfpathlineto{\pgfqpoint{4.683131in}{1.901209in}}%
\pgfpathlineto{\pgfqpoint{4.735460in}{1.805879in}}%
\pgfpathlineto{\pgfqpoint{4.787788in}{1.869184in}}%
\pgfpathlineto{\pgfqpoint{4.840117in}{1.869410in}}%
\pgfpathlineto{\pgfqpoint{4.892445in}{1.837882in}}%
\pgfpathlineto{\pgfqpoint{4.944774in}{1.835893in}}%
\pgfusepath{stroke}%
\end{pgfscope}%
\begin{pgfscope}%
\pgfpathrectangle{\pgfqpoint{0.488997in}{0.535045in}}{\pgfqpoint{4.777597in}{2.260065in}}%
\pgfusepath{clip}%
\pgfsetbuttcap%
\pgfsetroundjoin%
\definecolor{currentfill}{rgb}{0.180392,0.490196,0.196078}%
\pgfsetfillcolor{currentfill}%
\pgfsetlinewidth{1.003750pt}%
\definecolor{currentstroke}{rgb}{0.180392,0.490196,0.196078}%
\pgfsetstrokecolor{currentstroke}%
\pgfsetdash{}{0pt}%
\pgfsys@defobject{currentmarker}{\pgfqpoint{-0.027778in}{-0.027778in}}{\pgfqpoint{0.027778in}{0.027778in}}{%
\pgfpathmoveto{\pgfqpoint{0.000000in}{-0.027778in}}%
\pgfpathcurveto{\pgfqpoint{0.007367in}{-0.027778in}}{\pgfqpoint{0.014433in}{-0.024851in}}{\pgfqpoint{0.019642in}{-0.019642in}}%
\pgfpathcurveto{\pgfqpoint{0.024851in}{-0.014433in}}{\pgfqpoint{0.027778in}{-0.007367in}}{\pgfqpoint{0.027778in}{0.000000in}}%
\pgfpathcurveto{\pgfqpoint{0.027778in}{0.007367in}}{\pgfqpoint{0.024851in}{0.014433in}}{\pgfqpoint{0.019642in}{0.019642in}}%
\pgfpathcurveto{\pgfqpoint{0.014433in}{0.024851in}}{\pgfqpoint{0.007367in}{0.027778in}}{\pgfqpoint{0.000000in}{0.027778in}}%
\pgfpathcurveto{\pgfqpoint{-0.007367in}{0.027778in}}{\pgfqpoint{-0.014433in}{0.024851in}}{\pgfqpoint{-0.019642in}{0.019642in}}%
\pgfpathcurveto{\pgfqpoint{-0.024851in}{0.014433in}}{\pgfqpoint{-0.027778in}{0.007367in}}{\pgfqpoint{-0.027778in}{0.000000in}}%
\pgfpathcurveto{\pgfqpoint{-0.027778in}{-0.007367in}}{\pgfqpoint{-0.024851in}{-0.014433in}}{\pgfqpoint{-0.019642in}{-0.019642in}}%
\pgfpathcurveto{\pgfqpoint{-0.014433in}{-0.024851in}}{\pgfqpoint{-0.007367in}{-0.027778in}}{\pgfqpoint{0.000000in}{-0.027778in}}%
\pgfpathlineto{\pgfqpoint{0.000000in}{-0.027778in}}%
\pgfpathclose%
\pgfusepath{stroke,fill}%
}%
\begin{pgfscope}%
\pgfsys@transformshift{0.706161in}{0.535045in}%
\pgfsys@useobject{currentmarker}{}%
\end{pgfscope}%
\begin{pgfscope}%
\pgfsys@transformshift{0.758489in}{0.662060in}%
\pgfsys@useobject{currentmarker}{}%
\end{pgfscope}%
\begin{pgfscope}%
\pgfsys@transformshift{0.810818in}{1.012416in}%
\pgfsys@useobject{currentmarker}{}%
\end{pgfscope}%
\begin{pgfscope}%
\pgfsys@transformshift{0.863146in}{1.806648in}%
\pgfsys@useobject{currentmarker}{}%
\end{pgfscope}%
\begin{pgfscope}%
\pgfsys@transformshift{0.915475in}{1.838266in}%
\pgfsys@useobject{currentmarker}{}%
\end{pgfscope}%
\begin{pgfscope}%
\pgfsys@transformshift{0.967803in}{1.774171in}%
\pgfsys@useobject{currentmarker}{}%
\end{pgfscope}%
\begin{pgfscope}%
\pgfsys@transformshift{1.020132in}{1.870178in}%
\pgfsys@useobject{currentmarker}{}%
\end{pgfscope}%
\begin{pgfscope}%
\pgfsys@transformshift{1.072461in}{1.870359in}%
\pgfsys@useobject{currentmarker}{}%
\end{pgfscope}%
\begin{pgfscope}%
\pgfsys@transformshift{1.124789in}{1.870607in}%
\pgfsys@useobject{currentmarker}{}%
\end{pgfscope}%
\begin{pgfscope}%
\pgfsys@transformshift{1.177118in}{1.838334in}%
\pgfsys@useobject{currentmarker}{}%
\end{pgfscope}%
\begin{pgfscope}%
\pgfsys@transformshift{1.229446in}{1.870607in}%
\pgfsys@useobject{currentmarker}{}%
\end{pgfscope}%
\begin{pgfscope}%
\pgfsys@transformshift{1.281775in}{1.838289in}%
\pgfsys@useobject{currentmarker}{}%
\end{pgfscope}%
\begin{pgfscope}%
\pgfsys@transformshift{1.334103in}{1.870381in}%
\pgfsys@useobject{currentmarker}{}%
\end{pgfscope}%
\begin{pgfscope}%
\pgfsys@transformshift{1.386432in}{1.838582in}%
\pgfsys@useobject{currentmarker}{}%
\end{pgfscope}%
\begin{pgfscope}%
\pgfsys@transformshift{1.438760in}{1.870381in}%
\pgfsys@useobject{currentmarker}{}%
\end{pgfscope}%
\begin{pgfscope}%
\pgfsys@transformshift{1.491089in}{1.870314in}%
\pgfsys@useobject{currentmarker}{}%
\end{pgfscope}%
\begin{pgfscope}%
\pgfsys@transformshift{1.543418in}{1.838515in}%
\pgfsys@useobject{currentmarker}{}%
\end{pgfscope}%
\begin{pgfscope}%
\pgfsys@transformshift{1.595746in}{1.870381in}%
\pgfsys@useobject{currentmarker}{}%
\end{pgfscope}%
\begin{pgfscope}%
\pgfsys@transformshift{1.648075in}{1.870110in}%
\pgfsys@useobject{currentmarker}{}%
\end{pgfscope}%
\begin{pgfscope}%
\pgfsys@transformshift{1.700403in}{1.838560in}%
\pgfsys@useobject{currentmarker}{}%
\end{pgfscope}%
\begin{pgfscope}%
\pgfsys@transformshift{1.752732in}{1.870359in}%
\pgfsys@useobject{currentmarker}{}%
\end{pgfscope}%
\begin{pgfscope}%
\pgfsys@transformshift{1.805060in}{1.838582in}%
\pgfsys@useobject{currentmarker}{}%
\end{pgfscope}%
\begin{pgfscope}%
\pgfsys@transformshift{1.857389in}{1.870133in}%
\pgfsys@useobject{currentmarker}{}%
\end{pgfscope}%
\begin{pgfscope}%
\pgfsys@transformshift{1.909717in}{1.870607in}%
\pgfsys@useobject{currentmarker}{}%
\end{pgfscope}%
\begin{pgfscope}%
\pgfsys@transformshift{1.962046in}{1.838582in}%
\pgfsys@useobject{currentmarker}{}%
\end{pgfscope}%
\begin{pgfscope}%
\pgfsys@transformshift{2.014375in}{1.870110in}%
\pgfsys@useobject{currentmarker}{}%
\end{pgfscope}%
\begin{pgfscope}%
\pgfsys@transformshift{2.066703in}{1.870607in}%
\pgfsys@useobject{currentmarker}{}%
\end{pgfscope}%
\begin{pgfscope}%
\pgfsys@transformshift{2.119032in}{1.837724in}%
\pgfsys@useobject{currentmarker}{}%
\end{pgfscope}%
\begin{pgfscope}%
\pgfsys@transformshift{2.171360in}{1.869184in}%
\pgfsys@useobject{currentmarker}{}%
\end{pgfscope}%
\begin{pgfscope}%
\pgfsys@transformshift{2.223689in}{1.837656in}%
\pgfsys@useobject{currentmarker}{}%
\end{pgfscope}%
\begin{pgfscope}%
\pgfsys@transformshift{2.276017in}{1.869432in}%
\pgfsys@useobject{currentmarker}{}%
\end{pgfscope}%
\begin{pgfscope}%
\pgfsys@transformshift{2.328346in}{1.869432in}%
\pgfsys@useobject{currentmarker}{}%
\end{pgfscope}%
\begin{pgfscope}%
\pgfsys@transformshift{2.380674in}{1.837656in}%
\pgfsys@useobject{currentmarker}{}%
\end{pgfscope}%
\begin{pgfscope}%
\pgfsys@transformshift{2.433003in}{1.869432in}%
\pgfsys@useobject{currentmarker}{}%
\end{pgfscope}%
\begin{pgfscope}%
\pgfsys@transformshift{2.485332in}{1.837656in}%
\pgfsys@useobject{currentmarker}{}%
\end{pgfscope}%
\begin{pgfscope}%
\pgfsys@transformshift{2.537660in}{1.869410in}%
\pgfsys@useobject{currentmarker}{}%
\end{pgfscope}%
\begin{pgfscope}%
\pgfsys@transformshift{2.589989in}{1.869432in}%
\pgfsys@useobject{currentmarker}{}%
\end{pgfscope}%
\begin{pgfscope}%
\pgfsys@transformshift{2.642317in}{1.837656in}%
\pgfsys@useobject{currentmarker}{}%
\end{pgfscope}%
\begin{pgfscope}%
\pgfsys@transformshift{2.694646in}{1.869432in}%
\pgfsys@useobject{currentmarker}{}%
\end{pgfscope}%
\begin{pgfscope}%
\pgfsys@transformshift{2.746974in}{1.869184in}%
\pgfsys@useobject{currentmarker}{}%
\end{pgfscope}%
\begin{pgfscope}%
\pgfsys@transformshift{2.799303in}{1.837882in}%
\pgfsys@useobject{currentmarker}{}%
\end{pgfscope}%
\begin{pgfscope}%
\pgfsys@transformshift{2.851632in}{1.869184in}%
\pgfsys@useobject{currentmarker}{}%
\end{pgfscope}%
\begin{pgfscope}%
\pgfsys@transformshift{2.903960in}{1.837633in}%
\pgfsys@useobject{currentmarker}{}%
\end{pgfscope}%
\begin{pgfscope}%
\pgfsys@transformshift{2.956289in}{1.869432in}%
\pgfsys@useobject{currentmarker}{}%
\end{pgfscope}%
\begin{pgfscope}%
\pgfsys@transformshift{3.008617in}{1.869410in}%
\pgfsys@useobject{currentmarker}{}%
\end{pgfscope}%
\begin{pgfscope}%
\pgfsys@transformshift{3.060946in}{1.837633in}%
\pgfsys@useobject{currentmarker}{}%
\end{pgfscope}%
\begin{pgfscope}%
\pgfsys@transformshift{3.113274in}{1.869410in}%
\pgfsys@useobject{currentmarker}{}%
\end{pgfscope}%
\begin{pgfscope}%
\pgfsys@transformshift{3.165603in}{1.869432in}%
\pgfsys@useobject{currentmarker}{}%
\end{pgfscope}%
\begin{pgfscope}%
\pgfsys@transformshift{3.217931in}{1.837656in}%
\pgfsys@useobject{currentmarker}{}%
\end{pgfscope}%
\begin{pgfscope}%
\pgfsys@transformshift{3.270260in}{1.869432in}%
\pgfsys@useobject{currentmarker}{}%
\end{pgfscope}%
\begin{pgfscope}%
\pgfsys@transformshift{3.322589in}{1.551667in}%
\pgfsys@useobject{currentmarker}{}%
\end{pgfscope}%
\begin{pgfscope}%
\pgfsys@transformshift{3.374917in}{1.138798in}%
\pgfsys@useobject{currentmarker}{}%
\end{pgfscope}%
\begin{pgfscope}%
\pgfsys@transformshift{3.427246in}{1.901073in}%
\pgfsys@useobject{currentmarker}{}%
\end{pgfscope}%
\begin{pgfscope}%
\pgfsys@transformshift{3.479574in}{1.837588in}%
\pgfsys@useobject{currentmarker}{}%
\end{pgfscope}%
\begin{pgfscope}%
\pgfsys@transformshift{3.531903in}{1.869364in}%
\pgfsys@useobject{currentmarker}{}%
\end{pgfscope}%
\begin{pgfscope}%
\pgfsys@transformshift{3.584231in}{1.869410in}%
\pgfsys@useobject{currentmarker}{}%
\end{pgfscope}%
\begin{pgfscope}%
\pgfsys@transformshift{3.636560in}{1.837633in}%
\pgfsys@useobject{currentmarker}{}%
\end{pgfscope}%
\begin{pgfscope}%
\pgfsys@transformshift{3.688888in}{1.869387in}%
\pgfsys@useobject{currentmarker}{}%
\end{pgfscope}%
\begin{pgfscope}%
\pgfsys@transformshift{3.741217in}{1.869184in}%
\pgfsys@useobject{currentmarker}{}%
\end{pgfscope}%
\begin{pgfscope}%
\pgfsys@transformshift{3.793546in}{1.837407in}%
\pgfsys@useobject{currentmarker}{}%
\end{pgfscope}%
\begin{pgfscope}%
\pgfsys@transformshift{3.845874in}{1.869658in}%
\pgfsys@useobject{currentmarker}{}%
\end{pgfscope}%
\begin{pgfscope}%
\pgfsys@transformshift{3.898203in}{1.837430in}%
\pgfsys@useobject{currentmarker}{}%
\end{pgfscope}%
\begin{pgfscope}%
\pgfsys@transformshift{3.950531in}{1.869432in}%
\pgfsys@useobject{currentmarker}{}%
\end{pgfscope}%
\begin{pgfscope}%
\pgfsys@transformshift{4.002860in}{1.869410in}%
\pgfsys@useobject{currentmarker}{}%
\end{pgfscope}%
\begin{pgfscope}%
\pgfsys@transformshift{4.055188in}{1.837656in}%
\pgfsys@useobject{currentmarker}{}%
\end{pgfscope}%
\begin{pgfscope}%
\pgfsys@transformshift{4.107517in}{1.869410in}%
\pgfsys@useobject{currentmarker}{}%
\end{pgfscope}%
\begin{pgfscope}%
\pgfsys@transformshift{4.159846in}{1.837633in}%
\pgfsys@useobject{currentmarker}{}%
\end{pgfscope}%
\begin{pgfscope}%
\pgfsys@transformshift{4.212174in}{1.869432in}%
\pgfsys@useobject{currentmarker}{}%
\end{pgfscope}%
\begin{pgfscope}%
\pgfsys@transformshift{4.264503in}{1.869432in}%
\pgfsys@useobject{currentmarker}{}%
\end{pgfscope}%
\begin{pgfscope}%
\pgfsys@transformshift{4.316831in}{1.837656in}%
\pgfsys@useobject{currentmarker}{}%
\end{pgfscope}%
\begin{pgfscope}%
\pgfsys@transformshift{4.369160in}{1.869410in}%
\pgfsys@useobject{currentmarker}{}%
\end{pgfscope}%
\begin{pgfscope}%
\pgfsys@transformshift{4.421488in}{1.837000in}%
\pgfsys@useobject{currentmarker}{}%
\end{pgfscope}%
\begin{pgfscope}%
\pgfsys@transformshift{4.473817in}{1.869432in}%
\pgfsys@useobject{currentmarker}{}%
\end{pgfscope}%
\begin{pgfscope}%
\pgfsys@transformshift{4.526145in}{1.869432in}%
\pgfsys@useobject{currentmarker}{}%
\end{pgfscope}%
\begin{pgfscope}%
\pgfsys@transformshift{4.578474in}{1.837656in}%
\pgfsys@useobject{currentmarker}{}%
\end{pgfscope}%
\begin{pgfscope}%
\pgfsys@transformshift{4.630803in}{1.869432in}%
\pgfsys@useobject{currentmarker}{}%
\end{pgfscope}%
\begin{pgfscope}%
\pgfsys@transformshift{4.683131in}{1.901209in}%
\pgfsys@useobject{currentmarker}{}%
\end{pgfscope}%
\begin{pgfscope}%
\pgfsys@transformshift{4.735460in}{1.805879in}%
\pgfsys@useobject{currentmarker}{}%
\end{pgfscope}%
\begin{pgfscope}%
\pgfsys@transformshift{4.787788in}{1.869184in}%
\pgfsys@useobject{currentmarker}{}%
\end{pgfscope}%
\begin{pgfscope}%
\pgfsys@transformshift{4.840117in}{1.869410in}%
\pgfsys@useobject{currentmarker}{}%
\end{pgfscope}%
\begin{pgfscope}%
\pgfsys@transformshift{4.892445in}{1.837882in}%
\pgfsys@useobject{currentmarker}{}%
\end{pgfscope}%
\begin{pgfscope}%
\pgfsys@transformshift{4.944774in}{1.835893in}%
\pgfsys@useobject{currentmarker}{}%
\end{pgfscope}%
\end{pgfscope}%
\begin{pgfscope}%
\pgfpathrectangle{\pgfqpoint{0.488997in}{0.535045in}}{\pgfqpoint{4.777597in}{2.260065in}}%
\pgfusepath{clip}%
\pgfsetrectcap%
\pgfsetroundjoin%
\pgfsetlinewidth{2.007500pt}%
\definecolor{currentstroke}{rgb}{0.647059,0.839216,0.654902}%
\pgfsetstrokecolor{currentstroke}%
\pgfsetdash{}{0pt}%
\pgfpathmoveto{\pgfqpoint{0.706161in}{0.535045in}}%
\pgfpathlineto{\pgfqpoint{0.758489in}{0.662422in}}%
\pgfpathlineto{\pgfqpoint{0.810818in}{1.076014in}}%
\pgfpathlineto{\pgfqpoint{0.863146in}{1.806602in}}%
\pgfpathlineto{\pgfqpoint{0.915475in}{1.806467in}}%
\pgfpathlineto{\pgfqpoint{0.967803in}{1.774510in}}%
\pgfpathlineto{\pgfqpoint{1.020132in}{1.870449in}}%
\pgfpathlineto{\pgfqpoint{1.072461in}{1.870472in}}%
\pgfpathlineto{\pgfqpoint{1.124789in}{1.870449in}}%
\pgfpathlineto{\pgfqpoint{1.177118in}{1.838605in}}%
\pgfpathlineto{\pgfqpoint{1.229446in}{1.870449in}}%
\pgfpathlineto{\pgfqpoint{1.281775in}{1.870449in}}%
\pgfpathlineto{\pgfqpoint{1.334103in}{1.838650in}}%
\pgfpathlineto{\pgfqpoint{1.386432in}{1.870449in}}%
\pgfpathlineto{\pgfqpoint{1.438760in}{1.838673in}}%
\pgfpathlineto{\pgfqpoint{1.491089in}{1.870449in}}%
\pgfpathlineto{\pgfqpoint{1.543418in}{1.870449in}}%
\pgfpathlineto{\pgfqpoint{1.595746in}{1.838650in}}%
\pgfpathlineto{\pgfqpoint{1.648075in}{1.870314in}}%
\pgfpathlineto{\pgfqpoint{1.700403in}{1.870449in}}%
\pgfpathlineto{\pgfqpoint{1.752732in}{1.838808in}}%
\pgfpathlineto{\pgfqpoint{1.805060in}{1.870314in}}%
\pgfpathlineto{\pgfqpoint{1.857389in}{1.838673in}}%
\pgfpathlineto{\pgfqpoint{1.909717in}{1.870472in}}%
\pgfpathlineto{\pgfqpoint{1.962046in}{1.870449in}}%
\pgfpathlineto{\pgfqpoint{2.014375in}{1.838673in}}%
\pgfpathlineto{\pgfqpoint{2.066703in}{1.870291in}}%
\pgfpathlineto{\pgfqpoint{2.119032in}{1.869703in}}%
\pgfpathlineto{\pgfqpoint{2.171360in}{1.837746in}}%
\pgfpathlineto{\pgfqpoint{2.223689in}{1.869500in}}%
\pgfpathlineto{\pgfqpoint{2.276017in}{1.837746in}}%
\pgfpathlineto{\pgfqpoint{2.328346in}{1.869500in}}%
\pgfpathlineto{\pgfqpoint{2.380674in}{1.869364in}}%
\pgfpathlineto{\pgfqpoint{2.433003in}{1.837746in}}%
\pgfpathlineto{\pgfqpoint{2.485332in}{1.869658in}}%
\pgfpathlineto{\pgfqpoint{2.537660in}{1.837588in}}%
\pgfpathlineto{\pgfqpoint{2.589989in}{1.869500in}}%
\pgfpathlineto{\pgfqpoint{2.642317in}{1.901299in}}%
\pgfpathlineto{\pgfqpoint{2.694646in}{1.805947in}}%
\pgfpathlineto{\pgfqpoint{2.746974in}{1.869523in}}%
\pgfpathlineto{\pgfqpoint{2.799303in}{1.901299in}}%
\pgfpathlineto{\pgfqpoint{2.851632in}{1.805970in}}%
\pgfpathlineto{\pgfqpoint{2.903960in}{1.869500in}}%
\pgfpathlineto{\pgfqpoint{2.956289in}{1.837724in}}%
\pgfpathlineto{\pgfqpoint{3.008617in}{1.869500in}}%
\pgfpathlineto{\pgfqpoint{3.060946in}{1.869500in}}%
\pgfpathlineto{\pgfqpoint{3.113274in}{1.837724in}}%
\pgfpathlineto{\pgfqpoint{3.165603in}{1.869523in}}%
\pgfpathlineto{\pgfqpoint{3.217931in}{1.869523in}}%
\pgfpathlineto{\pgfqpoint{3.270260in}{1.837746in}}%
\pgfpathlineto{\pgfqpoint{3.322589in}{1.329299in}}%
\pgfpathlineto{\pgfqpoint{3.374917in}{1.456225in}}%
\pgfpathlineto{\pgfqpoint{3.427246in}{1.869455in}}%
\pgfpathlineto{\pgfqpoint{3.479574in}{1.869455in}}%
\pgfpathlineto{\pgfqpoint{3.531903in}{1.837724in}}%
\pgfpathlineto{\pgfqpoint{3.584231in}{1.869500in}}%
\pgfpathlineto{\pgfqpoint{3.636560in}{1.837678in}}%
\pgfpathlineto{\pgfqpoint{3.688888in}{1.869364in}}%
\pgfpathlineto{\pgfqpoint{3.741217in}{1.869500in}}%
\pgfpathlineto{\pgfqpoint{3.793546in}{1.837724in}}%
\pgfpathlineto{\pgfqpoint{3.845874in}{1.869523in}}%
\pgfpathlineto{\pgfqpoint{3.898203in}{1.837724in}}%
\pgfpathlineto{\pgfqpoint{3.950531in}{1.869500in}}%
\pgfpathlineto{\pgfqpoint{4.002860in}{1.869500in}}%
\pgfpathlineto{\pgfqpoint{4.055188in}{1.837746in}}%
\pgfpathlineto{\pgfqpoint{4.107517in}{1.869500in}}%
\pgfpathlineto{\pgfqpoint{4.159846in}{1.869523in}}%
\pgfpathlineto{\pgfqpoint{4.212174in}{1.837746in}}%
\pgfpathlineto{\pgfqpoint{4.264503in}{1.869500in}}%
\pgfpathlineto{\pgfqpoint{4.316831in}{1.836978in}}%
\pgfpathlineto{\pgfqpoint{4.369160in}{1.869500in}}%
\pgfpathlineto{\pgfqpoint{4.421488in}{1.869523in}}%
\pgfpathlineto{\pgfqpoint{4.473817in}{1.837724in}}%
\pgfpathlineto{\pgfqpoint{4.526145in}{1.869342in}}%
\pgfpathlineto{\pgfqpoint{4.578474in}{1.869658in}}%
\pgfpathlineto{\pgfqpoint{4.630803in}{1.837724in}}%
\pgfpathlineto{\pgfqpoint{4.683131in}{1.869523in}}%
\pgfpathlineto{\pgfqpoint{4.735460in}{1.837724in}}%
\pgfpathlineto{\pgfqpoint{4.787788in}{1.869364in}}%
\pgfpathlineto{\pgfqpoint{4.840117in}{1.869500in}}%
\pgfpathlineto{\pgfqpoint{4.892445in}{1.837746in}}%
\pgfpathlineto{\pgfqpoint{4.944774in}{1.733037in}}%
\pgfusepath{stroke}%
\end{pgfscope}%
\begin{pgfscope}%
\pgfpathrectangle{\pgfqpoint{0.488997in}{0.535045in}}{\pgfqpoint{4.777597in}{2.260065in}}%
\pgfusepath{clip}%
\pgfsetbuttcap%
\pgfsetroundjoin%
\definecolor{currentfill}{rgb}{0.647059,0.839216,0.654902}%
\pgfsetfillcolor{currentfill}%
\pgfsetlinewidth{1.003750pt}%
\definecolor{currentstroke}{rgb}{0.647059,0.839216,0.654902}%
\pgfsetstrokecolor{currentstroke}%
\pgfsetdash{}{0pt}%
\pgfsys@defobject{currentmarker}{\pgfqpoint{-0.027778in}{-0.027778in}}{\pgfqpoint{0.027778in}{0.027778in}}{%
\pgfpathmoveto{\pgfqpoint{0.000000in}{-0.027778in}}%
\pgfpathcurveto{\pgfqpoint{0.007367in}{-0.027778in}}{\pgfqpoint{0.014433in}{-0.024851in}}{\pgfqpoint{0.019642in}{-0.019642in}}%
\pgfpathcurveto{\pgfqpoint{0.024851in}{-0.014433in}}{\pgfqpoint{0.027778in}{-0.007367in}}{\pgfqpoint{0.027778in}{0.000000in}}%
\pgfpathcurveto{\pgfqpoint{0.027778in}{0.007367in}}{\pgfqpoint{0.024851in}{0.014433in}}{\pgfqpoint{0.019642in}{0.019642in}}%
\pgfpathcurveto{\pgfqpoint{0.014433in}{0.024851in}}{\pgfqpoint{0.007367in}{0.027778in}}{\pgfqpoint{0.000000in}{0.027778in}}%
\pgfpathcurveto{\pgfqpoint{-0.007367in}{0.027778in}}{\pgfqpoint{-0.014433in}{0.024851in}}{\pgfqpoint{-0.019642in}{0.019642in}}%
\pgfpathcurveto{\pgfqpoint{-0.024851in}{0.014433in}}{\pgfqpoint{-0.027778in}{0.007367in}}{\pgfqpoint{-0.027778in}{0.000000in}}%
\pgfpathcurveto{\pgfqpoint{-0.027778in}{-0.007367in}}{\pgfqpoint{-0.024851in}{-0.014433in}}{\pgfqpoint{-0.019642in}{-0.019642in}}%
\pgfpathcurveto{\pgfqpoint{-0.014433in}{-0.024851in}}{\pgfqpoint{-0.007367in}{-0.027778in}}{\pgfqpoint{0.000000in}{-0.027778in}}%
\pgfpathlineto{\pgfqpoint{0.000000in}{-0.027778in}}%
\pgfpathclose%
\pgfusepath{stroke,fill}%
}%
\begin{pgfscope}%
\pgfsys@transformshift{0.706161in}{0.535045in}%
\pgfsys@useobject{currentmarker}{}%
\end{pgfscope}%
\begin{pgfscope}%
\pgfsys@transformshift{0.758489in}{0.662422in}%
\pgfsys@useobject{currentmarker}{}%
\end{pgfscope}%
\begin{pgfscope}%
\pgfsys@transformshift{0.810818in}{1.076014in}%
\pgfsys@useobject{currentmarker}{}%
\end{pgfscope}%
\begin{pgfscope}%
\pgfsys@transformshift{0.863146in}{1.806602in}%
\pgfsys@useobject{currentmarker}{}%
\end{pgfscope}%
\begin{pgfscope}%
\pgfsys@transformshift{0.915475in}{1.806467in}%
\pgfsys@useobject{currentmarker}{}%
\end{pgfscope}%
\begin{pgfscope}%
\pgfsys@transformshift{0.967803in}{1.774510in}%
\pgfsys@useobject{currentmarker}{}%
\end{pgfscope}%
\begin{pgfscope}%
\pgfsys@transformshift{1.020132in}{1.870449in}%
\pgfsys@useobject{currentmarker}{}%
\end{pgfscope}%
\begin{pgfscope}%
\pgfsys@transformshift{1.072461in}{1.870472in}%
\pgfsys@useobject{currentmarker}{}%
\end{pgfscope}%
\begin{pgfscope}%
\pgfsys@transformshift{1.124789in}{1.870449in}%
\pgfsys@useobject{currentmarker}{}%
\end{pgfscope}%
\begin{pgfscope}%
\pgfsys@transformshift{1.177118in}{1.838605in}%
\pgfsys@useobject{currentmarker}{}%
\end{pgfscope}%
\begin{pgfscope}%
\pgfsys@transformshift{1.229446in}{1.870449in}%
\pgfsys@useobject{currentmarker}{}%
\end{pgfscope}%
\begin{pgfscope}%
\pgfsys@transformshift{1.281775in}{1.870449in}%
\pgfsys@useobject{currentmarker}{}%
\end{pgfscope}%
\begin{pgfscope}%
\pgfsys@transformshift{1.334103in}{1.838650in}%
\pgfsys@useobject{currentmarker}{}%
\end{pgfscope}%
\begin{pgfscope}%
\pgfsys@transformshift{1.386432in}{1.870449in}%
\pgfsys@useobject{currentmarker}{}%
\end{pgfscope}%
\begin{pgfscope}%
\pgfsys@transformshift{1.438760in}{1.838673in}%
\pgfsys@useobject{currentmarker}{}%
\end{pgfscope}%
\begin{pgfscope}%
\pgfsys@transformshift{1.491089in}{1.870449in}%
\pgfsys@useobject{currentmarker}{}%
\end{pgfscope}%
\begin{pgfscope}%
\pgfsys@transformshift{1.543418in}{1.870449in}%
\pgfsys@useobject{currentmarker}{}%
\end{pgfscope}%
\begin{pgfscope}%
\pgfsys@transformshift{1.595746in}{1.838650in}%
\pgfsys@useobject{currentmarker}{}%
\end{pgfscope}%
\begin{pgfscope}%
\pgfsys@transformshift{1.648075in}{1.870314in}%
\pgfsys@useobject{currentmarker}{}%
\end{pgfscope}%
\begin{pgfscope}%
\pgfsys@transformshift{1.700403in}{1.870449in}%
\pgfsys@useobject{currentmarker}{}%
\end{pgfscope}%
\begin{pgfscope}%
\pgfsys@transformshift{1.752732in}{1.838808in}%
\pgfsys@useobject{currentmarker}{}%
\end{pgfscope}%
\begin{pgfscope}%
\pgfsys@transformshift{1.805060in}{1.870314in}%
\pgfsys@useobject{currentmarker}{}%
\end{pgfscope}%
\begin{pgfscope}%
\pgfsys@transformshift{1.857389in}{1.838673in}%
\pgfsys@useobject{currentmarker}{}%
\end{pgfscope}%
\begin{pgfscope}%
\pgfsys@transformshift{1.909717in}{1.870472in}%
\pgfsys@useobject{currentmarker}{}%
\end{pgfscope}%
\begin{pgfscope}%
\pgfsys@transformshift{1.962046in}{1.870449in}%
\pgfsys@useobject{currentmarker}{}%
\end{pgfscope}%
\begin{pgfscope}%
\pgfsys@transformshift{2.014375in}{1.838673in}%
\pgfsys@useobject{currentmarker}{}%
\end{pgfscope}%
\begin{pgfscope}%
\pgfsys@transformshift{2.066703in}{1.870291in}%
\pgfsys@useobject{currentmarker}{}%
\end{pgfscope}%
\begin{pgfscope}%
\pgfsys@transformshift{2.119032in}{1.869703in}%
\pgfsys@useobject{currentmarker}{}%
\end{pgfscope}%
\begin{pgfscope}%
\pgfsys@transformshift{2.171360in}{1.837746in}%
\pgfsys@useobject{currentmarker}{}%
\end{pgfscope}%
\begin{pgfscope}%
\pgfsys@transformshift{2.223689in}{1.869500in}%
\pgfsys@useobject{currentmarker}{}%
\end{pgfscope}%
\begin{pgfscope}%
\pgfsys@transformshift{2.276017in}{1.837746in}%
\pgfsys@useobject{currentmarker}{}%
\end{pgfscope}%
\begin{pgfscope}%
\pgfsys@transformshift{2.328346in}{1.869500in}%
\pgfsys@useobject{currentmarker}{}%
\end{pgfscope}%
\begin{pgfscope}%
\pgfsys@transformshift{2.380674in}{1.869364in}%
\pgfsys@useobject{currentmarker}{}%
\end{pgfscope}%
\begin{pgfscope}%
\pgfsys@transformshift{2.433003in}{1.837746in}%
\pgfsys@useobject{currentmarker}{}%
\end{pgfscope}%
\begin{pgfscope}%
\pgfsys@transformshift{2.485332in}{1.869658in}%
\pgfsys@useobject{currentmarker}{}%
\end{pgfscope}%
\begin{pgfscope}%
\pgfsys@transformshift{2.537660in}{1.837588in}%
\pgfsys@useobject{currentmarker}{}%
\end{pgfscope}%
\begin{pgfscope}%
\pgfsys@transformshift{2.589989in}{1.869500in}%
\pgfsys@useobject{currentmarker}{}%
\end{pgfscope}%
\begin{pgfscope}%
\pgfsys@transformshift{2.642317in}{1.901299in}%
\pgfsys@useobject{currentmarker}{}%
\end{pgfscope}%
\begin{pgfscope}%
\pgfsys@transformshift{2.694646in}{1.805947in}%
\pgfsys@useobject{currentmarker}{}%
\end{pgfscope}%
\begin{pgfscope}%
\pgfsys@transformshift{2.746974in}{1.869523in}%
\pgfsys@useobject{currentmarker}{}%
\end{pgfscope}%
\begin{pgfscope}%
\pgfsys@transformshift{2.799303in}{1.901299in}%
\pgfsys@useobject{currentmarker}{}%
\end{pgfscope}%
\begin{pgfscope}%
\pgfsys@transformshift{2.851632in}{1.805970in}%
\pgfsys@useobject{currentmarker}{}%
\end{pgfscope}%
\begin{pgfscope}%
\pgfsys@transformshift{2.903960in}{1.869500in}%
\pgfsys@useobject{currentmarker}{}%
\end{pgfscope}%
\begin{pgfscope}%
\pgfsys@transformshift{2.956289in}{1.837724in}%
\pgfsys@useobject{currentmarker}{}%
\end{pgfscope}%
\begin{pgfscope}%
\pgfsys@transformshift{3.008617in}{1.869500in}%
\pgfsys@useobject{currentmarker}{}%
\end{pgfscope}%
\begin{pgfscope}%
\pgfsys@transformshift{3.060946in}{1.869500in}%
\pgfsys@useobject{currentmarker}{}%
\end{pgfscope}%
\begin{pgfscope}%
\pgfsys@transformshift{3.113274in}{1.837724in}%
\pgfsys@useobject{currentmarker}{}%
\end{pgfscope}%
\begin{pgfscope}%
\pgfsys@transformshift{3.165603in}{1.869523in}%
\pgfsys@useobject{currentmarker}{}%
\end{pgfscope}%
\begin{pgfscope}%
\pgfsys@transformshift{3.217931in}{1.869523in}%
\pgfsys@useobject{currentmarker}{}%
\end{pgfscope}%
\begin{pgfscope}%
\pgfsys@transformshift{3.270260in}{1.837746in}%
\pgfsys@useobject{currentmarker}{}%
\end{pgfscope}%
\begin{pgfscope}%
\pgfsys@transformshift{3.322589in}{1.329299in}%
\pgfsys@useobject{currentmarker}{}%
\end{pgfscope}%
\begin{pgfscope}%
\pgfsys@transformshift{3.374917in}{1.456225in}%
\pgfsys@useobject{currentmarker}{}%
\end{pgfscope}%
\begin{pgfscope}%
\pgfsys@transformshift{3.427246in}{1.869455in}%
\pgfsys@useobject{currentmarker}{}%
\end{pgfscope}%
\begin{pgfscope}%
\pgfsys@transformshift{3.479574in}{1.869455in}%
\pgfsys@useobject{currentmarker}{}%
\end{pgfscope}%
\begin{pgfscope}%
\pgfsys@transformshift{3.531903in}{1.837724in}%
\pgfsys@useobject{currentmarker}{}%
\end{pgfscope}%
\begin{pgfscope}%
\pgfsys@transformshift{3.584231in}{1.869500in}%
\pgfsys@useobject{currentmarker}{}%
\end{pgfscope}%
\begin{pgfscope}%
\pgfsys@transformshift{3.636560in}{1.837678in}%
\pgfsys@useobject{currentmarker}{}%
\end{pgfscope}%
\begin{pgfscope}%
\pgfsys@transformshift{3.688888in}{1.869364in}%
\pgfsys@useobject{currentmarker}{}%
\end{pgfscope}%
\begin{pgfscope}%
\pgfsys@transformshift{3.741217in}{1.869500in}%
\pgfsys@useobject{currentmarker}{}%
\end{pgfscope}%
\begin{pgfscope}%
\pgfsys@transformshift{3.793546in}{1.837724in}%
\pgfsys@useobject{currentmarker}{}%
\end{pgfscope}%
\begin{pgfscope}%
\pgfsys@transformshift{3.845874in}{1.869523in}%
\pgfsys@useobject{currentmarker}{}%
\end{pgfscope}%
\begin{pgfscope}%
\pgfsys@transformshift{3.898203in}{1.837724in}%
\pgfsys@useobject{currentmarker}{}%
\end{pgfscope}%
\begin{pgfscope}%
\pgfsys@transformshift{3.950531in}{1.869500in}%
\pgfsys@useobject{currentmarker}{}%
\end{pgfscope}%
\begin{pgfscope}%
\pgfsys@transformshift{4.002860in}{1.869500in}%
\pgfsys@useobject{currentmarker}{}%
\end{pgfscope}%
\begin{pgfscope}%
\pgfsys@transformshift{4.055188in}{1.837746in}%
\pgfsys@useobject{currentmarker}{}%
\end{pgfscope}%
\begin{pgfscope}%
\pgfsys@transformshift{4.107517in}{1.869500in}%
\pgfsys@useobject{currentmarker}{}%
\end{pgfscope}%
\begin{pgfscope}%
\pgfsys@transformshift{4.159846in}{1.869523in}%
\pgfsys@useobject{currentmarker}{}%
\end{pgfscope}%
\begin{pgfscope}%
\pgfsys@transformshift{4.212174in}{1.837746in}%
\pgfsys@useobject{currentmarker}{}%
\end{pgfscope}%
\begin{pgfscope}%
\pgfsys@transformshift{4.264503in}{1.869500in}%
\pgfsys@useobject{currentmarker}{}%
\end{pgfscope}%
\begin{pgfscope}%
\pgfsys@transformshift{4.316831in}{1.836978in}%
\pgfsys@useobject{currentmarker}{}%
\end{pgfscope}%
\begin{pgfscope}%
\pgfsys@transformshift{4.369160in}{1.869500in}%
\pgfsys@useobject{currentmarker}{}%
\end{pgfscope}%
\begin{pgfscope}%
\pgfsys@transformshift{4.421488in}{1.869523in}%
\pgfsys@useobject{currentmarker}{}%
\end{pgfscope}%
\begin{pgfscope}%
\pgfsys@transformshift{4.473817in}{1.837724in}%
\pgfsys@useobject{currentmarker}{}%
\end{pgfscope}%
\begin{pgfscope}%
\pgfsys@transformshift{4.526145in}{1.869342in}%
\pgfsys@useobject{currentmarker}{}%
\end{pgfscope}%
\begin{pgfscope}%
\pgfsys@transformshift{4.578474in}{1.869658in}%
\pgfsys@useobject{currentmarker}{}%
\end{pgfscope}%
\begin{pgfscope}%
\pgfsys@transformshift{4.630803in}{1.837724in}%
\pgfsys@useobject{currentmarker}{}%
\end{pgfscope}%
\begin{pgfscope}%
\pgfsys@transformshift{4.683131in}{1.869523in}%
\pgfsys@useobject{currentmarker}{}%
\end{pgfscope}%
\begin{pgfscope}%
\pgfsys@transformshift{4.735460in}{1.837724in}%
\pgfsys@useobject{currentmarker}{}%
\end{pgfscope}%
\begin{pgfscope}%
\pgfsys@transformshift{4.787788in}{1.869364in}%
\pgfsys@useobject{currentmarker}{}%
\end{pgfscope}%
\begin{pgfscope}%
\pgfsys@transformshift{4.840117in}{1.869500in}%
\pgfsys@useobject{currentmarker}{}%
\end{pgfscope}%
\begin{pgfscope}%
\pgfsys@transformshift{4.892445in}{1.837746in}%
\pgfsys@useobject{currentmarker}{}%
\end{pgfscope}%
\begin{pgfscope}%
\pgfsys@transformshift{4.944774in}{1.733037in}%
\pgfsys@useobject{currentmarker}{}%
\end{pgfscope}%
\end{pgfscope}%
\begin{pgfscope}%
\pgfpathrectangle{\pgfqpoint{0.488997in}{0.535045in}}{\pgfqpoint{4.777597in}{2.260065in}}%
\pgfusepath{clip}%
\pgfsetbuttcap%
\pgfsetroundjoin%
\pgfsetlinewidth{1.505625pt}%
\definecolor{currentstroke}{rgb}{1.000000,0.000000,0.000000}%
\pgfsetstrokecolor{currentstroke}%
\pgfsetdash{{5.550000pt}{2.400000pt}}{0.000000pt}%
\pgfpathmoveto{\pgfqpoint{3.296424in}{0.535045in}}%
\pgfpathlineto{\pgfqpoint{3.296424in}{2.795110in}}%
\pgfusepath{stroke}%
\end{pgfscope}%
\begin{pgfscope}%
\pgfsetrectcap%
\pgfsetmiterjoin%
\pgfsetlinewidth{0.803000pt}%
\definecolor{currentstroke}{rgb}{0.000000,0.000000,0.000000}%
\pgfsetstrokecolor{currentstroke}%
\pgfsetdash{}{0pt}%
\pgfpathmoveto{\pgfqpoint{0.488997in}{0.535045in}}%
\pgfpathlineto{\pgfqpoint{0.488997in}{2.795110in}}%
\pgfusepath{stroke}%
\end{pgfscope}%
\begin{pgfscope}%
\pgfsetrectcap%
\pgfsetmiterjoin%
\pgfsetlinewidth{0.803000pt}%
\definecolor{currentstroke}{rgb}{0.000000,0.000000,0.000000}%
\pgfsetstrokecolor{currentstroke}%
\pgfsetdash{}{0pt}%
\pgfpathmoveto{\pgfqpoint{5.266595in}{0.535045in}}%
\pgfpathlineto{\pgfqpoint{5.266595in}{2.795110in}}%
\pgfusepath{stroke}%
\end{pgfscope}%
\begin{pgfscope}%
\pgfsetrectcap%
\pgfsetmiterjoin%
\pgfsetlinewidth{0.803000pt}%
\definecolor{currentstroke}{rgb}{0.000000,0.000000,0.000000}%
\pgfsetstrokecolor{currentstroke}%
\pgfsetdash{}{0pt}%
\pgfpathmoveto{\pgfqpoint{0.488997in}{0.535045in}}%
\pgfpathlineto{\pgfqpoint{5.266595in}{0.535045in}}%
\pgfusepath{stroke}%
\end{pgfscope}%
\begin{pgfscope}%
\pgfsetrectcap%
\pgfsetmiterjoin%
\pgfsetlinewidth{0.803000pt}%
\definecolor{currentstroke}{rgb}{0.000000,0.000000,0.000000}%
\pgfsetstrokecolor{currentstroke}%
\pgfsetdash{}{0pt}%
\pgfpathmoveto{\pgfqpoint{0.488997in}{2.795110in}}%
\pgfpathlineto{\pgfqpoint{5.266595in}{2.795110in}}%
\pgfusepath{stroke}%
\end{pgfscope}%
\begin{pgfscope}%
\definecolor{textcolor}{rgb}{0.000000,0.000000,0.000000}%
\pgfsetstrokecolor{textcolor}%
\pgfsetfillcolor{textcolor}%
\pgftext[x=1.284222in, y=3.049512in, left, base]{\color{textcolor}{\sffamily\fontsize{11.000000}{13.200000}\selectfont\catcode`\^=\active\def^{\ifmmode\sp\else\^{}\fi}\catcode`\%=\active\def%{\%}QUIC recovery after connection migration}}%
\end{pgfscope}%
\begin{pgfscope}%
\definecolor{textcolor}{rgb}{0.000000,0.000000,0.000000}%
\pgfsetstrokecolor{textcolor}%
\pgfsetfillcolor{textcolor}%
\pgftext[x=2.877796in, y=2.878443in, left, base]{\color{textcolor}{\sffamily\fontsize{11.000000}{13.200000}\selectfont\catcode`\^=\active\def^{\ifmmode\sp\else\^{}\fi}\catcode`\%=\active\def%{\%}}}%
\end{pgfscope}%
\begin{pgfscope}%
\pgfsetbuttcap%
\pgfsetmiterjoin%
\definecolor{currentfill}{rgb}{1.000000,1.000000,1.000000}%
\pgfsetfillcolor{currentfill}%
\pgfsetfillopacity{0.800000}%
\pgfsetlinewidth{1.003750pt}%
\definecolor{currentstroke}{rgb}{0.800000,0.800000,0.800000}%
\pgfsetstrokecolor{currentstroke}%
\pgfsetstrokeopacity{0.800000}%
\pgfsetdash{}{0pt}%
\pgfpathmoveto{\pgfqpoint{2.313763in}{1.961224in}}%
\pgfpathlineto{\pgfqpoint{5.179095in}{1.961224in}}%
\pgfpathquadraticcurveto{\pgfqpoint{5.204095in}{1.961224in}}{\pgfqpoint{5.204095in}{1.986224in}}%
\pgfpathlineto{\pgfqpoint{5.204095in}{2.707610in}}%
\pgfpathquadraticcurveto{\pgfqpoint{5.204095in}{2.732610in}}{\pgfqpoint{5.179095in}{2.732610in}}%
\pgfpathlineto{\pgfqpoint{2.313763in}{2.732610in}}%
\pgfpathquadraticcurveto{\pgfqpoint{2.288763in}{2.732610in}}{\pgfqpoint{2.288763in}{2.707610in}}%
\pgfpathlineto{\pgfqpoint{2.288763in}{1.986224in}}%
\pgfpathquadraticcurveto{\pgfqpoint{2.288763in}{1.961224in}}{\pgfqpoint{2.313763in}{1.961224in}}%
\pgfpathlineto{\pgfqpoint{2.313763in}{1.961224in}}%
\pgfpathclose%
\pgfusepath{stroke,fill}%
\end{pgfscope}%
\begin{pgfscope}%
\pgfsetrectcap%
\pgfsetroundjoin%
\pgfsetlinewidth{2.007500pt}%
\definecolor{currentstroke}{rgb}{0.082353,0.396078,0.752941}%
\pgfsetstrokecolor{currentstroke}%
\pgfsetdash{}{0pt}%
\pgfpathmoveto{\pgfqpoint{2.338763in}{2.631389in}}%
\pgfpathlineto{\pgfqpoint{2.463763in}{2.631389in}}%
\pgfpathlineto{\pgfqpoint{2.588763in}{2.631389in}}%
\pgfusepath{stroke}%
\end{pgfscope}%
\begin{pgfscope}%
\pgfsetbuttcap%
\pgfsetroundjoin%
\definecolor{currentfill}{rgb}{0.082353,0.396078,0.752941}%
\pgfsetfillcolor{currentfill}%
\pgfsetlinewidth{1.003750pt}%
\definecolor{currentstroke}{rgb}{0.082353,0.396078,0.752941}%
\pgfsetstrokecolor{currentstroke}%
\pgfsetdash{}{0pt}%
\pgfsys@defobject{currentmarker}{\pgfqpoint{-0.027778in}{-0.027778in}}{\pgfqpoint{0.027778in}{0.027778in}}{%
\pgfpathmoveto{\pgfqpoint{0.000000in}{-0.027778in}}%
\pgfpathcurveto{\pgfqpoint{0.007367in}{-0.027778in}}{\pgfqpoint{0.014433in}{-0.024851in}}{\pgfqpoint{0.019642in}{-0.019642in}}%
\pgfpathcurveto{\pgfqpoint{0.024851in}{-0.014433in}}{\pgfqpoint{0.027778in}{-0.007367in}}{\pgfqpoint{0.027778in}{0.000000in}}%
\pgfpathcurveto{\pgfqpoint{0.027778in}{0.007367in}}{\pgfqpoint{0.024851in}{0.014433in}}{\pgfqpoint{0.019642in}{0.019642in}}%
\pgfpathcurveto{\pgfqpoint{0.014433in}{0.024851in}}{\pgfqpoint{0.007367in}{0.027778in}}{\pgfqpoint{0.000000in}{0.027778in}}%
\pgfpathcurveto{\pgfqpoint{-0.007367in}{0.027778in}}{\pgfqpoint{-0.014433in}{0.024851in}}{\pgfqpoint{-0.019642in}{0.019642in}}%
\pgfpathcurveto{\pgfqpoint{-0.024851in}{0.014433in}}{\pgfqpoint{-0.027778in}{0.007367in}}{\pgfqpoint{-0.027778in}{0.000000in}}%
\pgfpathcurveto{\pgfqpoint{-0.027778in}{-0.007367in}}{\pgfqpoint{-0.024851in}{-0.014433in}}{\pgfqpoint{-0.019642in}{-0.019642in}}%
\pgfpathcurveto{\pgfqpoint{-0.014433in}{-0.024851in}}{\pgfqpoint{-0.007367in}{-0.027778in}}{\pgfqpoint{0.000000in}{-0.027778in}}%
\pgfpathlineto{\pgfqpoint{0.000000in}{-0.027778in}}%
\pgfpathclose%
\pgfusepath{stroke,fill}%
}%
\begin{pgfscope}%
\pgfsys@transformshift{2.463763in}{2.631389in}%
\pgfsys@useobject{currentmarker}{}%
\end{pgfscope}%
\end{pgfscope}%
\begin{pgfscope}%
\definecolor{textcolor}{rgb}{0.000000,0.000000,0.000000}%
\pgfsetstrokecolor{textcolor}%
\pgfsetfillcolor{textcolor}%
\pgftext[x=2.688763in,y=2.587639in,left,base]{\color{textcolor}{\sffamily\fontsize{9.000000}{10.800000}\selectfont\catcode`\^=\active\def^{\ifmmode\sp\else\^{}\fi}\catcode`\%=\active\def%{\%}TCP (Native)}}%
\end{pgfscope}%
\begin{pgfscope}%
\pgfsetrectcap%
\pgfsetroundjoin%
\pgfsetlinewidth{2.007500pt}%
\definecolor{currentstroke}{rgb}{0.180392,0.490196,0.196078}%
\pgfsetstrokecolor{currentstroke}%
\pgfsetdash{}{0pt}%
\pgfpathmoveto{\pgfqpoint{2.338763in}{2.447918in}}%
\pgfpathlineto{\pgfqpoint{2.463763in}{2.447918in}}%
\pgfpathlineto{\pgfqpoint{2.588763in}{2.447918in}}%
\pgfusepath{stroke}%
\end{pgfscope}%
\begin{pgfscope}%
\pgfsetbuttcap%
\pgfsetroundjoin%
\definecolor{currentfill}{rgb}{0.180392,0.490196,0.196078}%
\pgfsetfillcolor{currentfill}%
\pgfsetlinewidth{1.003750pt}%
\definecolor{currentstroke}{rgb}{0.180392,0.490196,0.196078}%
\pgfsetstrokecolor{currentstroke}%
\pgfsetdash{}{0pt}%
\pgfsys@defobject{currentmarker}{\pgfqpoint{-0.027778in}{-0.027778in}}{\pgfqpoint{0.027778in}{0.027778in}}{%
\pgfpathmoveto{\pgfqpoint{0.000000in}{-0.027778in}}%
\pgfpathcurveto{\pgfqpoint{0.007367in}{-0.027778in}}{\pgfqpoint{0.014433in}{-0.024851in}}{\pgfqpoint{0.019642in}{-0.019642in}}%
\pgfpathcurveto{\pgfqpoint{0.024851in}{-0.014433in}}{\pgfqpoint{0.027778in}{-0.007367in}}{\pgfqpoint{0.027778in}{0.000000in}}%
\pgfpathcurveto{\pgfqpoint{0.027778in}{0.007367in}}{\pgfqpoint{0.024851in}{0.014433in}}{\pgfqpoint{0.019642in}{0.019642in}}%
\pgfpathcurveto{\pgfqpoint{0.014433in}{0.024851in}}{\pgfqpoint{0.007367in}{0.027778in}}{\pgfqpoint{0.000000in}{0.027778in}}%
\pgfpathcurveto{\pgfqpoint{-0.007367in}{0.027778in}}{\pgfqpoint{-0.014433in}{0.024851in}}{\pgfqpoint{-0.019642in}{0.019642in}}%
\pgfpathcurveto{\pgfqpoint{-0.024851in}{0.014433in}}{\pgfqpoint{-0.027778in}{0.007367in}}{\pgfqpoint{-0.027778in}{0.000000in}}%
\pgfpathcurveto{\pgfqpoint{-0.027778in}{-0.007367in}}{\pgfqpoint{-0.024851in}{-0.014433in}}{\pgfqpoint{-0.019642in}{-0.019642in}}%
\pgfpathcurveto{\pgfqpoint{-0.014433in}{-0.024851in}}{\pgfqpoint{-0.007367in}{-0.027778in}}{\pgfqpoint{0.000000in}{-0.027778in}}%
\pgfpathlineto{\pgfqpoint{0.000000in}{-0.027778in}}%
\pgfpathclose%
\pgfusepath{stroke,fill}%
}%
\begin{pgfscope}%
\pgfsys@transformshift{2.463763in}{2.447918in}%
\pgfsys@useobject{currentmarker}{}%
\end{pgfscope}%
\end{pgfscope}%
\begin{pgfscope}%
\definecolor{textcolor}{rgb}{0.000000,0.000000,0.000000}%
\pgfsetstrokecolor{textcolor}%
\pgfsetfillcolor{textcolor}%
\pgftext[x=2.688763in,y=2.404168in,left,base]{\color{textcolor}{\sffamily\fontsize{9.000000}{10.800000}\selectfont\catcode`\^=\active\def^{\ifmmode\sp\else\^{}\fi}\catcode`\%=\active\def%{\%}QUIC (Picoquic)}}%
\end{pgfscope}%
\begin{pgfscope}%
\pgfsetrectcap%
\pgfsetroundjoin%
\pgfsetlinewidth{2.007500pt}%
\definecolor{currentstroke}{rgb}{0.647059,0.839216,0.654902}%
\pgfsetstrokecolor{currentstroke}%
\pgfsetdash{}{0pt}%
\pgfpathmoveto{\pgfqpoint{2.338763in}{2.264446in}}%
\pgfpathlineto{\pgfqpoint{2.463763in}{2.264446in}}%
\pgfpathlineto{\pgfqpoint{2.588763in}{2.264446in}}%
\pgfusepath{stroke}%
\end{pgfscope}%
\begin{pgfscope}%
\pgfsetbuttcap%
\pgfsetroundjoin%
\definecolor{currentfill}{rgb}{0.647059,0.839216,0.654902}%
\pgfsetfillcolor{currentfill}%
\pgfsetlinewidth{1.003750pt}%
\definecolor{currentstroke}{rgb}{0.647059,0.839216,0.654902}%
\pgfsetstrokecolor{currentstroke}%
\pgfsetdash{}{0pt}%
\pgfsys@defobject{currentmarker}{\pgfqpoint{-0.027778in}{-0.027778in}}{\pgfqpoint{0.027778in}{0.027778in}}{%
\pgfpathmoveto{\pgfqpoint{0.000000in}{-0.027778in}}%
\pgfpathcurveto{\pgfqpoint{0.007367in}{-0.027778in}}{\pgfqpoint{0.014433in}{-0.024851in}}{\pgfqpoint{0.019642in}{-0.019642in}}%
\pgfpathcurveto{\pgfqpoint{0.024851in}{-0.014433in}}{\pgfqpoint{0.027778in}{-0.007367in}}{\pgfqpoint{0.027778in}{0.000000in}}%
\pgfpathcurveto{\pgfqpoint{0.027778in}{0.007367in}}{\pgfqpoint{0.024851in}{0.014433in}}{\pgfqpoint{0.019642in}{0.019642in}}%
\pgfpathcurveto{\pgfqpoint{0.014433in}{0.024851in}}{\pgfqpoint{0.007367in}{0.027778in}}{\pgfqpoint{0.000000in}{0.027778in}}%
\pgfpathcurveto{\pgfqpoint{-0.007367in}{0.027778in}}{\pgfqpoint{-0.014433in}{0.024851in}}{\pgfqpoint{-0.019642in}{0.019642in}}%
\pgfpathcurveto{\pgfqpoint{-0.024851in}{0.014433in}}{\pgfqpoint{-0.027778in}{0.007367in}}{\pgfqpoint{-0.027778in}{0.000000in}}%
\pgfpathcurveto{\pgfqpoint{-0.027778in}{-0.007367in}}{\pgfqpoint{-0.024851in}{-0.014433in}}{\pgfqpoint{-0.019642in}{-0.019642in}}%
\pgfpathcurveto{\pgfqpoint{-0.014433in}{-0.024851in}}{\pgfqpoint{-0.007367in}{-0.027778in}}{\pgfqpoint{0.000000in}{-0.027778in}}%
\pgfpathlineto{\pgfqpoint{0.000000in}{-0.027778in}}%
\pgfpathclose%
\pgfusepath{stroke,fill}%
}%
\begin{pgfscope}%
\pgfsys@transformshift{2.463763in}{2.264446in}%
\pgfsys@useobject{currentmarker}{}%
\end{pgfscope}%
\end{pgfscope}%
\begin{pgfscope}%
\definecolor{textcolor}{rgb}{0.000000,0.000000,0.000000}%
\pgfsetstrokecolor{textcolor}%
\pgfsetfillcolor{textcolor}%
\pgftext[x=2.688763in,y=2.220696in,left,base]{\color{textcolor}{\sffamily\fontsize{9.000000}{10.800000}\selectfont\catcode`\^=\active\def^{\ifmmode\sp\else\^{}\fi}\catcode`\%=\active\def%{\%}QUIC (CTaps)}}%
\end{pgfscope}%
\begin{pgfscope}%
\pgfsetbuttcap%
\pgfsetroundjoin%
\pgfsetlinewidth{1.505625pt}%
\definecolor{currentstroke}{rgb}{1.000000,0.000000,0.000000}%
\pgfsetstrokecolor{currentstroke}%
\pgfsetdash{{5.550000pt}{2.400000pt}}{0.000000pt}%
\pgfpathmoveto{\pgfqpoint{2.338763in}{2.080974in}}%
\pgfpathlineto{\pgfqpoint{2.463763in}{2.080974in}}%
\pgfpathlineto{\pgfqpoint{2.588763in}{2.080974in}}%
\pgfusepath{stroke}%
\end{pgfscope}%
\begin{pgfscope}%
\definecolor{textcolor}{rgb}{0.000000,0.000000,0.000000}%
\pgfsetstrokecolor{textcolor}%
\pgfsetfillcolor{textcolor}%
\pgftext[x=2.688763in,y=2.037224in,left,base]{\color{textcolor}{\sffamily\fontsize{9.000000}{10.800000}\selectfont\catcode`\^=\active\def^{\ifmmode\sp\else\^{}\fi}\catcode`\%=\active\def%{\%}Path becomes unavailable (T=5000ms)}}%
\end{pgfscope}%
\end{pgfpicture}%
\makeatother%
\endgroup%

    \caption[Connection migration experiment]{Transfer of a \qty{4.6}{\mega\byte}~file over a \qty{5}{\mega\bit\per\second}~link with CTaps, Picoquic and TCP.
    The first path is blocked after 5~seconds.}
    \label{fig:migration-throughput}
\end{figure}

\subsection{Client Complexity}\label{sec:client-complexity}
While the advantages of QUIC are clear, they
are also available through using a QUIC library directly
and do not require a TAPS System.
The largest benefit of TAPS is therefore not the features
of the underlying protocol, but that it makes these
features available with minimal complexity for application developers.

When collecting these experiments, the same Picoquic, CTaps and TCP clients
were used for the small file transfer and the connection migration experiment.
Their relative complexity measured in lines of code
can be seen in~\cref{tab:loc_complexity}. This
shows that while a CTaps client is more complex than a QUIC or TCP client
individually, it is less complex than each of them combined, while
still supporting both protocols.

The two handshake experiments used custom clients which immediately
terminated the connection after the handshake had completed, in order
to optimize the experiment execution time.

\begin{table}[htbp]
\centering
\begin{threeparttable}
    \caption[Lines of code for experiment clients]{Lines of code of experiment clients used in the small file transfer and migration experiments.\tnote{1}}
\label{tab:loc_complexity}
\begin{tabular}{lrrrr}
\toprule
\rowcolor{gray!25}
\textbf{Component} & \textbf{TCP} & \textbf{QUIC} & \textbf{TCP \& QUIC} & \textbf{CTaps} \\
\midrule
\multicolumn{1}{l||}{Initialization \& context setup}  & 33  &  69 & 102  & 73  \\
\multicolumn{1}{l||}{Connection establishment}         & 16  &  37 &  53  & 44  \\
\multicolumn{1}{l||}{Data transfer}                    & 22  &  29 &  51  & 57  \\
\multicolumn{1}{l||}{Error handling}                   & 63  &  26 &  89  & 57  \\
\multicolumn{1}{l||}{Teardown}                         & 15  &  27 &  42  & 30  \\
\multicolumn{1}{l||}{Migration}                        &  0  &  24 &  24  &  4  \\
\midrule
\textbf{Total}                   & 149 & 212 & 361 & 265 \\
\bottomrule
\end{tabular}
\begin{tablenotes}
      \footnotesize
      \item[1] The counts were obtained by classifying each non-blank,
non-comment line into one of six categories
for each of the three clients used for the experiments.
Annotated source files used for counting are available at \mono{benchmark/annotations/} in the CTaps repository.
    \end{tablenotes}
\end{threeparttable}
\end{table}

\section{Discussion} \label{sec:discussion}
\subsection{Handshake Time Experiment}\label{sec:single-endpoint-handshake-discussion}
The results in~\cref{fig:racing-handshake} 
demonstrate how a single CTaps client can successfully connect to multiple
servers utilizing different protocols with negligible overhead in an I/O-bound experiment.

The experiment also shows that candidate racing
is not implemented optimally in CTaps.
Currently, CTaps uses a static stagger delay
of 250~milliseconds between candidates as recommended by Happy Eyeballs V2~\cite{rfc8305}, but it does
not implement interleaving as described in the same document.
This means that all combinations involving the most desirable protocol are attempted
before a new protocol is attempted.

This becomes a significant delay
when racing both path combinations and protocol options.
If an application is faced with a large number of Endpoints,
for example after a DNS resolution,
then it could face a multi-second delay
if the most desirable protocol is
not compatible with the server.

This lack of protocol interleaving is made worse by CTaps treating
each individual ALPN value as a separate protocol option, as this
inflates the number of candidates for racing.
It is possible to utilize multiple ALPN values for a
single connection when performing ALPN, but not when performing
session resumption, since this relies on using the same ALPN as
the previous session. It is possible to generate candidates
more efficiently when session resumption is not expected,
but this was deemed to be out of scope for CTaps.

Potential mitigations for the high cost of the stagger delay as
it exists currently are discussed in~\cref{sec:future-work}.

\subsection{UDP RTT Experiment}
While the experiment in \cref{sec:udp-rtt-ex} shows that the overhead of CTaps is measurable,
it would be negligible for most networking applications in real-world
settings.

Part of the overhead shown likely stems from CTaps using libuv as
its underlying event loop library. It has been assumed that since libuv
is a mature library and is used in large software projects, the performance gaps to other
event loop libraries are relatively small. Benchmarking libuv
is out of scope for CTaps, so more profiling should be performed
if performance becomes a concern.

The performed experiment measures overhead in a fairly ideal case and
does not show the overhead of CTaps
in more complicated situations. This experiment therefore
shows the absolute minimum added overhead of CTaps and does not show
how CTaps would react for cases with multiple clients transferring
concurrently.

The long tail is of interest, especially if it scales poorly under
load, although this cannot be inferred from the data. If CTaps were to
be utilized in a latency-sensitive environment, more testing should be
performed in order to establish the behavior under load.

\subsection{Candidate Racing Experiment}
The high cost of the stagger delay is present not only in the handshake experiment,
but also in the candidate racing experiment as the experiment itself becomes
more contrived than needed. While it still serves to
show the benefit gained on a suboptimal ordering from a DNS resolution, the difference in the two endpoints has to be
excessively large, so that candidate racing in CTaps can overcome the high cost of the stagger delay.

\subsection{Small File Transfer Experiment}
The small file experiment shows that CTaps can utilize Connection
Cloning to make use of QUIC multistreaming, with transfer times
dominated by network conditions rather than CTaps overhead.
While this shows clear benefit in
this specific scenario, downloading large files as required for
this experiment is uncommon. In 2025, Cloudflare
measured that the median number of packets sent to clients of their service
was 12~\cite{sulemanahmadMeasuringCharacteristicsTCP2025}.
This suggests that most connections never grow a large CWND,
which limits the benefit of CWND reuse in practice. However, for applications
that open multiple connections to the same endpoint, multistreaming could
consolidate these connections into streams on a single connection,
potentially making CWND reuse relevant.

This experiment also shows that a single CTaps client can be used with both a
TCP and a QUIC server, and that application developers using CTaps can write networking
code independently of the underlying protocol.

% Discuss connection migration
\subsection{Connection Migration Experiment}\label{sec:connection-mig-benchmark-discssion}
The connection migration experiment shows that connection migration is available through
CTaps with minimal complexity for the CTaps client. While this is a promising result,
the scenario in the experiment is also very specific since it requires a synchronous
error from \mono{sendmsg()}. Since the error is synchronous, recovery is also very fast.
CTaps has not been tested with timeout-based migration, and verifying this is left for future work.

The biggest takeaway from the experiment is that, like the multistreaming
usage in the small file transfer experiment, it shows
that CTaps can make use of a QUIC-specific advantage
from a client that also supports TCP.

The overall connection migration functionality in CTaps is limited by the inability to bind
to multiple sockets as discussed in~\cref{sec:connection-migration-impl}.
This naturally limits the number of situations in which the client is actually able to migrate.
Adding support for multiple sockets
within CTaps is discussed further in \cref{sec:future-conn-mig}.

% Discuss CTaps client complexity
\subsection{Client Complexity}
The complexity of each client as seen in~\cref{tab:loc_complexity}
shows that CTaps successfully abstracts away complexity for application programmers.
CTaps achieves this not only when accessing advanced features of QUIC such as connection migration, but also in the protocol
selection and the error handling itself.
This results in simpler application code and easier migration to future protocols.

Lines of code is a coarse proxy for complexity, but it provides a
concrete comparison across these implementations.
While CTaps is the most complex
individual client by lines of code, it provides TCP and QUIC support, connection racing,
and several QUIC benefits at less than the combined complexity of the QUIC
and TCP clients. A combined TCP and QUIC client would likely feature
some shared logic, so the naive sum in~\cref{tab:loc_complexity} can
be viewed as an upper bound for the non-CTaps clients. 
In order to develop a replacement for CTaps, the TCP and QUIC
clients would also have to implement
fallback logic, which is not accounted for in the table.

CTaps abstracts away significant complexity, but this comes at the cost of making it a heavy
dependency. Applications relying on CTaps are therefore exposed to any correctness issues within it.

Additionally, CTaps forces applications into an asynchronous, callback-driven model
through its built-in event loop.
Applications must be designed around this model from the start,
which creates a barrier for adoption in existing codebases.

% Discuss performance in general
\subsection{Performance}\label{sec:disc-perf}
Overall, the experiments show that for simple, I/O-bound applications,
the overhead of CTaps is negligible.
In addition, the UDP RTT experiment
shows that CTaps features a longer tail in worst-case
performance, but it is not possible to say if this will remain
relatively constant, or scale up with load.

CTaps has not been written with performance in mind and
there are several potential causes of the overhead seen.
The duplicate memory allocations for every passed parameter are likely a
part of the delay seen for the UDP RTT experiment, and
could be replaced with reference counting if performance is
essential.

Before any steps are taken to optimize CTaps, a
software profiler should be utilized to determine exactly which
parts are the source of the overhead.

\subsection{Summary}
To summarize, the experiments show that while CTaps does have a measurable overhead,
it is negligible when compared to a realistic level of network latency.
More testing is required to establish CTaps' behavior under load.

Beyond performance, CTaps successfully exposes QUIC-specific features such as connection migration and multistreaming,
as well as TAPS features such as candidate racing, through a single generic client.
However, the current static stagger delay implementation imposes a significant cost
on connection establishment time. The client complexity comparison in~\cref{sec:client-complexity}
further shows that CTaps provides these features at a lower complexity than a combined TCP and QUIC
implementation. The broader implications of these results are discussed in~\cref{chap:conc}.
